% %%%%%%%%%%%%%%%%%%%%%%%%%%%%%%%%%%%%%%%%%%%%%%%%%%%%%%%%%%%%%%%%%%%%%%% %
%                                                                         %
% The Project Gutenberg EBook of The First Steps in Algebra, by           %
% G. A. (George Albert) Wentworth                                         %
%                                                                         %
% This eBook is for the use of anyone anywhere at no cost and with        %
% almost no restrictions whatsoever.  You may copy it, give it away or    %
% re-use it under the terms of the Project Gutenberg License included     %
% with this eBook or online at www.gutenberg.net                          %
%                                                                         %
%                                                                         %
% Title: The First Steps in Algebra                                       %
%                                                                         %
% Author: G. A. (George Albert) Wentworth                                 %
%                                                                         %
% Release Date: July 9, 2011 [EBook #36670]                               %
% Most recently updated: June 11, 2021                         %
%                                                                         %
% Language: English                                                       %
%                                                                         %
% Character set encoding: UTF-8                                           %
%                                                                         %
% *** START OF THIS PROJECT GUTENBERG EBOOK THE FIRST STEPS IN ALGEBRA ***%
%                                                                         %
% %%%%%%%%%%%%%%%%%%%%%%%%%%%%%%%%%%%%%%%%%%%%%%%%%%%%%%%%%%%%%%%%%%%%%%% %

\def\ebook{36670}
%%%%%%%%%%%%%%%%%%%%%%%%%%%%%%%%%%%%%%%%%%%%%%%%%%%%%%%%%%%%%%%%%%%%%%
%%                                                                  %%
%% Packages and substitutions:                                      %%
%%                                                                  %%
%% book:     Required.                                              %%
%% inputenc: Standard DP encoding. Required.                        %%
%%                                                                  %%
%% fix-cm:   For larger title page fonts. Optional.                 %%
%% ifthen:   Logical conditionals. Required.                        %%
%%                                                                  %%
%% amsmath:  AMS mathematics enhancements. Required.                %%
%% amssymb:  Additional mathematical symbols. Required.             %%
%%                                                                  %%
%% alltt:    Fixed-width font environment. Required.                %%
%% array:    Enhanced tabular features. Required.                   %%
%%                                                                  %%
%% indentfirst: Indent first word of each sectional unit. Optional. %%
%% textcase: Apply \MakeUppercase (et al.) only to text, not math.  %%
%%           Required.                                              %%
%%                                                                  %%
%% calc:     Length calculations. Required.                         %%
%% soul:     Spaced text. Optional.                                 %%
%%                                                                  %%
%% fancyhdr: Enhanced running headers and footers. Required.        %%
%%                                                                  %%
%% graphicx: Standard interface for graphics inclusion. Required.   %%
%% wrapfig:  Illustrations surrounded by text. Required.            %%
%%                                                                  %%
%% geometry: Enhanced page layout package. Required.                %%
%% hyperref: Hypertext embellishments for pdf output. Required.     %%
%%                                                                  %%
%%                                                                  %%
%% Producer's Comments:                                             %%
%%                                                                  %%
%%   Changes are noted in this file in two ways.                    %%
%%   1. \DPtypo{}{} for typographical corrections, showing original %%
%%      and replacement text side-by-side.                          %%
%%   2. [** TN: Note]s for lengthier or stylistic comments.         %%
%%                                                                  %%
%% Compilation Flags:                                               %%
%%                                                                  %%
%%   The following behavior may be controlled by boolean flags.     %%
%%                                                                  %%
%%   ForPrinting (false by default):                                %%
%%   Compile a screen-optimized PDF file. Set to true for print-    %%
%%   optimized file (large text block, two-sided layout, black      %%
%%   hyperlinks).                                                   %%
%%                                                                  %%
%%   Both print and screen layout are relatively loose, and contain %%
%%   hard-coded page breaks (\PrintBreak, \ScreenBreak, \newpage).  %%
%%                                                                  %%
%%                                                                  %%
%% PDF pages: 269 (if ForPrinting set to false)                     %%
%% PDF page size: 4.75 x 7" (non-standard)                          %%
%% PDF bookmarks: created, point to ToC entries                     %%
%% PDF document info: filled in                                     %%
%% Images: 3 pdf diagrams                                           %%
%%                                                                  %%
%% Summary of log file:                                             %%
%% * Large numbers of visually-harmless over-full hboxes and vboxes %%
%%   from DPalign* and DPgather* environments.                      %%
%%                                                                  %%
%%                                                                  %%
%% Compile History:                                                 %%
%%                                                                  %%
%% July, 2011: adhere (Andrew D. Hwang)                             %%
%%             texlive2007, GNU/Linux                               %%
%%                                                                  %%
%% Command block:                                                   %%
%%                                                                  %%
%%     pdflatex x2                                                  %%
%%                                                                  %%
%%                                                                  %%
%% July 2011: pglatex.                                              %%
%%   Compile this project with:                                     %%
%%   pdflatex 36670-t.tex ..... TWO times                           %%
%%                                                                  %%
%%   pdfTeXk, Version 3.141592-1.40.3 (Web2C 7.5.6)                 %%
%%                                                                  %%
%%%%%%%%%%%%%%%%%%%%%%%%%%%%%%%%%%%%%%%%%%%%%%%%%%%%%%%%%%%%%%%%%%%%%%
\listfiles
\documentclass[12pt]{book}[2005/09/16]

%%%%%%%%%%%%%%%%%%%%%%%%%%%%% PACKAGES %%%%%%%%%%%%%%%%%%%%%%%%%%%%%%%
\usepackage[utf8]{inputenc}[2006/05/05]

\usepackage{ifthen}[2001/05/26]  %% Logical conditionals

\usepackage{amsmath}[2000/07/18] %% Displayed equations
\usepackage{amssymb}[2002/01/22] %% and additional symbols

\usepackage{alltt}[1997/06/16]   %% boilerplate, credits, license

\usepackage{array}[2005/08/23]   %% extended array/tabular features

\usepackage{multicol}

\usepackage{graphicx}[1999/02/16]%% For diagrams

\usepackage{indentfirst}[1995/11/23]
\usepackage{textcase}[2004/10/07]

\usepackage{calc}[2005/08/06]

% for running heads
\usepackage{fancyhdr}

%%%%%%%%%%%%%%%%%%%%%%%%%%%%%%%%%%%%%%%%%%%%%%%%%%%%%%%%%%%%%%%%%
%%%% Interlude:  Set up PRINTING (default) or SCREEN VIEWING %%%%
%%%%%%%%%%%%%%%%%%%%%%%%%%%%%%%%%%%%%%%%%%%%%%%%%%%%%%%%%%%%%%%%%

% ForPrinting=true (default)           false
% Asymmetric margins                   Symmetric margins
% Black hyperlinks                     Blue hyperlinks
% Start Preface, ToC, etc. recto       No blank verso pages
%
% Chapter-like ``Sections'' start both recto and verso in the scanned
% book. This behavior has been retained.
\newboolean{ForPrinting}

%% UNCOMMENT the next line for a PRINT-OPTIMIZED VERSION of the text %%
%\setboolean{ForPrinting}{true}

%% Initialize values to ForPrinting=false
\newcommand{\Margins}{hmarginratio=1:1}     % Symmetric margins
\newcommand{\HLinkColor}{blue}              % Hyperlink color
\newcommand{\PDFPageLayout}{SinglePage}
\newcommand{\TransNote}{Transcriber's Note}
\newcommand{\TransNoteCommon}{%
  Minor typographical corrections and presentational changes have
  been made without comment.
  \bigskip
}

\newcommand{\TransNoteText}{%
  \TransNoteCommon

  This PDF file is optimized for screen viewing, but may easily be
  recompiled for printing. Please see the preamble of the \LaTeX\
  source file for instructions.
}
%% Re-set if ForPrinting=true
\ifthenelse{\boolean{ForPrinting}}{%
  \renewcommand{\Margins}{hmarginratio=2:3} % Asymmetric margins
  \renewcommand{\HLinkColor}{black}         % Hyperlink color
  \renewcommand{\PDFPageLayout}{TwoPageRight}
  \renewcommand{\TransNote}{Transcriber's Note}
  \renewcommand{\TransNoteText}{%
    \TransNoteCommon

    This PDF file is optimized for printing, but may easily be
    recompiled for screen viewing. Please see the preamble of the
    \LaTeX\ source file for instructions.
  }
}{% If ForPrinting=false, don't skip to recto
  \renewcommand{\cleardoublepage}{\clearpage}
}
%%%%%%%%%%%%%%%%%%%%%%%%%%%%%%%%%%%%%%%%%%%%%%%%%%%%%%%%%%%%%%%%%
%%%%  End of PRINTING/SCREEN VIEWING code; back to packages  %%%%
%%%%%%%%%%%%%%%%%%%%%%%%%%%%%%%%%%%%%%%%%%%%%%%%%%%%%%%%%%%%%%%%%

\ifthenelse{\boolean{ForPrinting}}{%
  \setlength{\paperwidth}{8.5in}%
  \setlength{\paperheight}{11in}%
  \usepackage[body={5in,8in},\Margins]{geometry}[2002/07/08]
}{%
  \setlength{\paperwidth}{4.75in}%
  \setlength{\paperheight}{7in}%
  \raggedbottom
  \usepackage[body={4.5in,6in},\Margins,includeheadfoot]{geometry}[2002/07/08]
}

\providecommand{\ebook}{00000}    % Overridden during white-washing
\usepackage[pdftex,
  hyperref,
  hyperfootnotes=false,
  pdftitle={The Project Gutenberg eBook \#\ebook: The First Steps in Algebra},
  pdfauthor={George Albert Wentworth},
  pdfkeywords={Peter Vachuska, Chuck Greif, Andrew D. Hwang
               Project Gutenberg Online Distributed Proofreading Team},
  pdfstartview=Fit,    % default value
  pdfstartpage=1,      % default value
  pdfpagemode=UseNone, % default value
  bookmarks=true,      % default value
  linktocpage=false,   % default value
  pdfpagelayout=\PDFPageLayout,
  pdfdisplaydoctitle,
  pdfpagelabels=true,
  bookmarksopen=true,
  bookmarksopenlevel=1,
  colorlinks=true,
  linkcolor=\HLinkColor]{hyperref}[2007/02/07]

%%%% Fixed-width environment to format PG boilerplate %%%%
\newenvironment{PGtext}{%
\begin{alltt}
\fontsize{8.1}{9}\ttfamily\selectfont}%
{\end{alltt}}

%% No hrule in page header
\renewcommand{\headrulewidth}{0pt}
\setlength{\extrarowheight}{2pt}
\setlength{\arraycolsep}{2pt}

% Tighten up spacing surrounding multicols environments
\setlength{\multicolsep}{8pt plus 4pt minus 8pt}
\setlength{\parskip}{6pt plus 6pt minus 4pt}

% Check for specified remaining height on current page
\newcommand{\Require}[1][2\baselineskip]{%
  \setlength{\TmpLen}{\pagegoal-\pagetotal}%
  \ifthenelse{\lengthtest{0pt < \TmpLen}\and\lengthtest{\TmpLen < #1}}{%
    \vfil\newpage%
  }{}%
}

% Hard-coded page breaks for print- and screen-formatted versions
\newcommand{\PrintBreak}{\ifthenelse{\boolean{ForPrinting}}{\newpage}{}}
\newcommand{\ScreenBreak}{\ifthenelse{\not\boolean{ForPrinting}}{\newpage}{}}

\newcommand{\PrintStretch}[1]{%
  \ifthenelse{\boolean{ForPrinting}}{\enlargethispage{#1}}{}%
}

% Running heads
\newcommand{\FlushRunningHeads}{\clearpage\fancyhf{}\cleardoublepage}

\newcommand{\RHead}[1]{\footnotesize\MakeUppercase{#1}}
\newcommand{\SetRunningHeads}[2][First Steps in Algebra.]{%
  \pagestyle{fancy}
  \fancyhf{}
  \setlength{\headheight}{15pt}
  \thispagestyle{plain}
  \fancyhead[CE]{\RHead{#1}}
  \fancyhead[CO]{\RHead{#2}}

  \ifthenelse{\boolean{ForPrinting}}
             {\fancyhead[RO,LE]{\thepage}}
             {\fancyhead[R]{\thepage}}
}

\newcommand{\BookMark}[2]{\phantomsection\pdfbookmark[#1]{#2}{#2}}

%\Chapter[ToC entry]{Number.}{Heading title}
\newcommand{\Chapter}[3][]{%
  \cleardoublepage
  \phantomsection
  \label{chap:#2}

  \ifthenelse{\equal{#2}{I.}}{%
    \begin{center}
      \textbf{\Large FIRST STEPS IN ALGEBRA.}

      \tb
    \end{center}
  }{}%
  \section*{\normalfont\centering\large CHAPTER #2}
  \subsection*{\normalfont\centering\small \MakeUppercase{#3}}
  \SetRunningHeads{#3}%
  \ifthenelse{\equal{#1}{}}{%
    \addtocontents{toc}{\protect\ToCLine{#2}{#3}}%
    \BookMark{0}{Chapter #2 #3}
  }{%
    \addtocontents{toc}{\protect\ToCLine{#2}{#1}}%
    \BookMark{0}{Chapter #2 #1}
  }
}

\newcommand{\Section}[1]{\subsection*{\centering\normalfont\small\scshape #1}}

% Numbered units
\newcommand{\Paragraph}[1]{\medskip\par\textbf{#1}}

\newcommand{\Exercise}[2][]{%
  \Require[3\baselineskip]
  \subsection*{\centering\small Exercise #2}
  \ifthenelse{\not\equal{#1}{}}{%
    \subsubsection*{\centering\normalfont\small\scshape #1}
  }{}
  \phantomsection\label{exer:#2}%
}

\newcommand{\Preface}{%
  \FlushRunningHeads
  \pagestyle{fancy}
  \section*{\centering\Large PREFACE.}
  \SetRunningHeads[Preface.]{Preface.}
  \BookMark{0}{Preface.}
}

\newcommand{\Signature}[2]{%
  \medskip%
  \null\hfill#1 \\
  \hspace*{\parindent}#2%
}

% N.B. Re-defines the \Item command
\newcommand{\Answers}{
  \FlushRunningHeads
  \pagestyle{fancy}
  \phantomsection
  \label{answers}
  \section*{\centering\Large ANSWERS.}
  \SetRunningHeads[Answers.]{Answers.}
  \addtocontents{toc}{\protect\ToCLine{}{Answers.}}%
  \BookMark{0}{Answers.}
  \ifthenelse{\not\boolean{ForPrinting}}{\footnotesize}{\small}%

  \renewcommand{\Item}[1]{
    \vspace{4pt plus 16pt}\makebox[1em][r]{\textbf{##1}\ }\hangindent 2.5em%
  }
}

% For internal use to keep track of current number of columns
\newcommand{\CurrCols}{1}

\newcommand{\ResetCols}[1][1]{%
  \ifthenelse{\not\equal{\CurrCols}{1}}{%
    \end{multicols}
  }{}
  \renewcommand{\CurrCols}{#1}
  \ifthenelse{\not\equal{#1}{1}}{%
    \begin{multicols}{#1}%
    }{}
}

% For internal use
\newcommand{\AnsSectHead}[2]{
  \subsection*{%
    \centering\small #1 #2 \hyperref[exer:#2]{Page \pageref{exer:#2}}
  }
}

% \AnsTo[cols]{Exercise}{42.}
\newcommand{\AnsTo}[3][3]{%
  \ifthenelse{\not\equal{#3}{1.}}{% Not first group
    \ifthenelse{\not\equal{\CurrCols}{1}}{%
    \end{multicols}
    }
  }{}

  \renewcommand{\CurrCols}{#1}%
  \ifthenelse{\not\equal{#1}{1}}{%
    \begin{multicols}{#1}[\AnsSectHead{#2}{#3}]
    }{%
      \AnsSectHead{#2}{#3}
    }
  }

% Semantic boldface
\newcommand{\Defn}[1]{\textbf{#1}}
\newcommand{\Dictum}[1]{\Require\textbf{#1}}

% Semantic small type
\newenvironment{Remark}[1][]{\par\small\textsc{#1}}{\par}
\newenvironment{Soln}[1][]{\par\small\textsc{#1}}{\par}

% Semantic italics
\newenvironment{Theorem}[1][]{%
  \ifthenelse{\not\equal{#1}{nopar}}{\smallskip\par#1}{}%
  \itshape}{\medskip\upshape\par}


\newcommand{\Graphic}[1]{%
\begin{center}
  \includegraphics[width=\textwidth]{./images/fig#1.pdf}%
\end{center}
}

\newcommand{\Item}[1]{\vspace{4pt plus 16pt}\makebox[1em][r]{\textbf{#1}\ }}
\newcommand{\Ax}[1]{\vspace{4pt plus 16pt}\makebox[3.5em][l]{Ax.\ #1}\ignorespaces}
\newcommand{\Ans}[1]{%
  \hfil\allowbreak\null\nobreak\hfill\nobreak\mbox{#1 \textit{Ans.}}%
}
% Trailing words in displays
\newcommand{\EqText}[1]{\rlap{\quad\text{#1}}}

\newcommand{\First}[1]{\textsc{#1}}

% For corrections.
\newcommand{\DPtypo}[2]{#2}
\newcommand{\Add}[1]{\DPtypo{}{#1}}

% \PadTo[#1]{#2}{#3} sets #3 in a box of width #2, aligned at #1 (default [r])
% e.g., \PadTo{3x^{2} - 5x}{x}
\newlength{\TmpLen}
\newcommand{\PadTo}[3][r]{%
  \settowidth{\TmpLen}{$#2$}%
  \makebox[\TmpLen][#1]{$#3$}%
}

% Decorative rule
\newcommand{\tb}[1][0.75in]{%
  \rule{#1}{0.5pt}\\
  \medskip
}

% Common abbreviations
\newcommand{\HCF}{H.\,C.\,F.}
\newcommand{\LCM}{L.\,C.\,M.}
\newcommand{\LCD}{L.\,C.\,D.}

% For rectangular arrays illustrating multiplication
\newcommand{\DOT}{\bullet}
\newcommand{\NOT}{\bullet\llap{$/$}}

% Miscellaneous alignment conveniences
\newcommand{\Z}{\phantom{0}}
\newcommand{\Neg}{\phantom{-}}

% In-line "vertical sum" and vinculum
\newcommand{\VSum}[2]{\begin{array}{r}#1\\[-4pt]+#2\end{array}\Big|}
\newcommand{\Vinc}[1]{\overline{#1\strut}}

% Table entries providing their own vertical bar (or not)
\newcommand{\NoBar}[1]{\multicolumn{1}{@{}r}{#1}}
\newcommand{\TbBar}[1]{\multicolumn{1}{@{}r|}{#1}}

% Big brace, for table on 168.png
\newcommand{\threelines}{\rule[-18pt]{0pt}{36pt}}
\newcommand{\BB}{\makebox[2pt][r]{\smash{$\left.\threelines\right\}$}}}

\DeclareUnicodeCharacter{00A3}{\pounds}
\DeclareUnicodeCharacter{00B0}{{}^\circ}
\DeclareUnicodeCharacter{00B1}{\pm}
\DeclareUnicodeCharacter{00B7}{\cdot}
\DeclareUnicodeCharacter{00D7}{\times}
\DeclareUnicodeCharacter{00F7}{\div}

% ToC formatting
\AtBeginDocument{\renewcommand{\contentsname}%
  {\protect\centering\protect\Large CONTENTS.}}

\newcommand{\TableofContents}{%
  \tableofcontents
  \SetRunningHeads{Contents.}
  \BookMark{0}{Contents.}
}

% Project-dependent logic
\newcommand{\ToCLine}[2]{%
  \settowidth{\TmpLen}{\textsc{\footnotesize Chapter99}}%
  \ifthenelse{\equal{#1}{I.}}{%
    \noindent\textsc{\footnotesize Chapter\hfill Page}\\
  }{}%
  \ifthenelse{\not\equal{#1}{}}{%
    \noindent\parbox[b]{\textwidth}{\strut\raggedright\hangindent\TmpLen%
      \makebox[\TmpLen][r]{#1\quad}\textsc{#2}\dotfill\pageref{chap:#1}%
    }%
  }{%
    \noindent\parbox[b]{\textwidth}{\strut\raggedright\hangindent\TmpLen%
      \makebox[\TmpLen][r]{#1\quad}\textsc{#2}\dotfill\pageref{answers}%
    }%
  }
}

%% Major document divisions %%
\newcommand{\PGBoilerPlate}{%
  \pagenumbering{Alph}
  \pagestyle{empty}
  \BookMark{0}{PG Boilerplate.}
}
\newcommand{\FrontMatter}{%
  \cleardoublepage % pagestyle still empty; \Preface calls \pagestyle{fancy}
  \frontmatter
  \BookMark{-1}{Front Matter.}
}

\newcommand{\MainMatter}{%
  \FlushRunningHeads
  \mainmatter
  \BookMark{-1}{Main Matter.}
}

\newcommand{\BackMatter}{%
  \FlushRunningHeads
  \pagenumbering{Roman}
  \backmatter
  \BookMark{-1}{Back Matter.}
  \BookMark{0}{PG License.}
  \SetRunningHeads[License.]{License.}
}

\newcommand{\Tag}[1]{\tag*{#1}}

% DPalign, DPgather
\makeatletter
\providecommand\shortintertext\intertext
\newcount\DP@lign@no
\newtoks\DP@lignb@dy
\newif\ifDP@cr
\newif\ifbr@ce
\def\f@@zl@bar{\null}
\def\addto@DPbody#1{\global\DP@lignb@dy\@xp{\the\DP@lignb@dy#1}}
\def\parseb@dy#1{\ifx\f@@zl@bar#1\f@@zl@bar
    \addto@DPbody{{}}\let\@next\parseb@dy
  \else\ifx\end#1
    \let\@next\process@DPb@dy
    \ifDP@cr\else\addto@DPbody{\DPh@@kr&\DP@rint}\@xp\addto@DPbody\@xp{\@xp{\the\DP@lign@no}&}\fi
    \addto@DPbody{\end}
  \else\ifx\intertext#1
    \def\@next{\eat@command0}%
  \else\ifx\shortintertext#1
    \def\@next{\eat@command1}%
  \else\ifDP@cr\addto@DPbody{&\DP@lint}\@xp\addto@DPbody\@xp{\@xp{\the\DP@lign@no}&\DPh@@kl}
          \DP@crfalse\fi
    \ifx\begin#1\def\begin@stack{b}
      \let\@next\eat@environment
  \else\ifx\lintertext#1
    \let\@next\linter@text
  \else\ifx\rintertext#1
    \let\@next\rinter@text
  \else\ifx\\#1
    \addto@DPbody{\DPh@@kr&\DP@rint}\@xp\addto@DPbody\@xp{\@xp{\the\DP@lign@no}&\\}\DP@crtrue
    \global\advance\DP@lign@no\@ne
    \let\@next\parse@cr
  \else\check@braces#1!Q!Q!Q!\ifbr@ce\addto@DPbody{{#1}}\else
    \addto@DPbody{#1}\fi
    \let\@next\parseb@dy
  \fi\fi\fi\fi\fi\fi\fi\fi\@next}
\def\process@DPb@dy{\let\lintertext\@gobble\let\rintertext\@gobble
  \@xp\start@align\@xp\tw@\@xp\st@rredtrue\@xp\m@ne\the\DP@lignb@dy}
\def\linter@text#1{\@xp\DPlint\@xp{\the\DP@lign@no}{#1}\parseb@dy}
\def\rinter@text#1{\@xp\DPrint\@xp{\the\DP@lign@no}{#1}\parseb@dy}
\def\DPlint#1#2{\@xp\def\csname DP@lint:#1\endcsname{\text{#2}}}
\def\DPrint#1#2{\@xp\def\csname DP@rint:#1\endcsname{\text{#2}}}
\def\DP@lint#1{\ifbalancedlrint\@xp\ifx\csname
DP@lint:#1\endcsname\relax\phantom
  {\csname DP@rint:#1\endcsname}\else\csname DP@lint:#1\endcsname\fi
  \else\csname DP@lint:#1\endcsname\fi}
\def\DP@rint#1{\ifbalancedlrint\@xp\ifx\csname
DP@rint:#1\endcsname\relax\phantom
  {\csname DP@lint:#1\endcsname}\else\csname DP@rint:#1\endcsname\fi
  \else\csname DP@rint:#1\endcsname\fi}
\def\eat@command#1#2{\ifcase#1\addto@DPbody{\intertext{#2}}\or
  \addto@DPbody{\shortintertext{#2}}\fi\DP@crtrue
  \global\advance\DP@lign@no\@ne\parseb@dy}
\def\parse@cr{\new@ifnextchar*{\parse@crst}{\parse@crst{}}}
\def\parse@crst#1{\addto@DPbody{#1}\new@ifnextchar[{\parse@crb}{\parseb@dy}}
\def\parse@crb[#1]{\addto@DPbody{[#1]}\parseb@dy}
\def\check@braces#1#2!Q!Q!Q!{\def\dp@lignt@stm@cro{#2}\ifx
  \empty\dp@lignt@stm@cro\br@cefalse\else\br@cetrue\fi}
\def\eat@environment#1{\addto@DPbody{\begin{#1}}\begingroup
  \def\@currenvir{#1}\let\@next\digest@env\@next}
\def\digest@env#1\end#2{%
  \edef\begin@stack{\push@begins#1\begin\end \@xp\@gobble\begin@stack}%
  \ifx\@empty\begin@stack
    \@checkend{#2}
    \endgroup\let\@next\parseb@dy\fi
    \addto@DPbody{#1\end{#2}}
    \@next}
\def\lintertext{lint}\def\rintertext{rint}
\newif\ifbalancedlrint
\let\DPh@@kl\empty\let\DPh@@kr\empty
\def\DPg@therl{&\omit\hfil$\displaystyle}
\def\DPg@therr{$\hfil}

\newenvironment{DPalign*}[1][a]{%
  \setlength{\abovedisplayskip}{8pt plus 4pt minus 6pt}
  \setlength{\belowdisplayskip}{8pt plus 4pt minus 6pt}
  \if m#1\balancedlrintfalse\else\balancedlrinttrue\fi
  \global\DP@lign@no\z@\DP@crfalse
  \DP@lignb@dy{&\DP@lint0&}\parseb@dy
}{%
  \endalign
}
\newenvironment{DPgather*}[1][a]{%
  \setlength{\abovedisplayskip}{8pt plus 4pt minus 6pt}
  \setlength{\belowdisplayskip}{8pt plus 4pt minus 6pt}
  \if m#1\balancedlrintfalse\else\balancedlrinttrue\fi
  \global\DP@lign@no\z@\DP@crfalse
  \let\DPh@@kl\DPg@therl
  \let\DPh@@kr\DPg@therr
  \DP@lignb@dy{&\DP@lint0&\DPh@@kl}\parseb@dy
}{%
  \endalign
}
\makeatother
%%%%%%%%%%%%%%%%%%%%%%%% START OF DOCUMENT %%%%%%%%%%%%%%%%%%%%%%%%%%

\begin{document}
%%%% PG BOILERPLATE %%%%
\PGBoilerPlate

\begin{center}
\begin{minipage}{\textwidth}
\small
\begin{PGtext}
The Project Gutenberg EBook of The First Steps in Algebra, by 
G. A. (George Albert) Wentworth

This eBook is for the use of anyone anywhere at no cost and with
almost no restrictions whatsoever.  You may copy it, give it away or
re-use it under the terms of the Project Gutenberg License included
with this eBook or online at www.gutenberg.net


Title: The First Steps in Algebra

Author: G. A. (George Albert) Wentworth

Release Date: July 9, 2011 [EBook #36670]
Most recently updated: June 11, 2021

Language: English

Character set encoding: UTF-8     

*** START OF THIS PROJECT GUTENBERG EBOOK THE FIRST STEPS IN ALGEBRA ***
\end{PGtext}
\end{minipage}
\end{center}
\clearpage
%%%% Credits and transcriber's note %%%%
\begin{center}
\begin{minipage}{\textwidth}
\begin{PGtext}
Produced by Andrew D. Hwang, Peter Vachuska, Chuck Greif
and the Online Distributed Proofreading Team at
http://www.pgdp.net.
\end{PGtext}
\end{minipage}
\end{center}
\vfill

\begin{minipage}{0.85\textwidth}
\small
\BookMark{0}{Transcriber's Note}
\subsection*{\centering\normalfont\scshape
\normalsize\MakeLowercase{\TransNote}}

\raggedright
\TransNoteText
\end{minipage}
%%%%%%%%%%%%%%%%%%%%%%%%%%% FRONT MATTER %%%%%%%%%%%%%%%%%%%%%%%%%%
\FrontMatter
%% -----File: 001.png---Folio i-------
\begin{center}
\Large THE
\vfill

\textbf{\LARGE FIRST STEPS IN ALGEBRA.}
\vfill

\normalsize BY \\[12pt]
\Large G. A. WENTWORTH, A.M. \\[12pt]
\footnotesize AUTHOR OF A SERIES OF TEXT-BOOKS IN MATHEMATICS.
\vfill\vfill

\large BOSTON, U.S.A.: \\
PUBLISHED BY GINN \& COMPANY. \\
1904.
\end{center}
%% -----File: 002.png---Folio ii-------
\newpage
\null\vfill
\begin{center}
\footnotesize
Entered according to Act of Congress, in the year 1894, by \\[4pt]
G. A. WENTWORTH, \\[4pt]
in the Office of the Librarian of Congress, at Washington. \\[4pt]
\tb
\textsc{All Rights Reserved.}
\vfill

\textsc{Typography by J. S. Cushing \& Co., Boston, U.S.A.} \\
\tb[2.5in]
\textsc{Presswork by Ginn \& Co., Boston, U.S.A.}
\end{center}
%% -----File: 003.png---Folio iii-------


\Preface

\First{This} book is written for pupils in the upper grades of
grammar schools and the lower grades of high schools.
The introduction of the simple elements of Algebra into
these grades will, it is thought, so stimulate the mental
activity of the pupils, that they will make considerable
progress in Algebra without detriment to their progress
in Arithmetic, even if no more time is allowed for the
two studies than is usually given to Arithmetic alone.

The great danger in preparing an Algebra for very
young pupils is that the author, in endeavoring to smooth
the path of the learner, will sacrifice much of the educational
value of the study.  To avoid this real and serious
danger, and at the same time to gain the required simplicity,
great care has been given to the explanations of
the fundamental operations and rules, the arrangement
of topics, the model solutions of examples, and the making
of easy examples for the pupils to solve.

Nearly all the examples throughout the book are new,
and made expressly for beginners.

The first chapter clears the way for quite a full treatment
of simple integral equations with one unknown number.
In the first two chapters only \emph{positive} numbers are
%% -----File: 004.png---Folio iv-------
involved, and the learner is led to see the practical advantages
of Algebra in its most interesting applications before
he faces the difficulties of negative numbers.

The third chapter contains a simple explanation of negative
numbers.  The recognition of the facts that the real
nature of subtraction is counting backwards, and that the
real nature of multiplication is forming the product from
the multiplicand precisely as the multiplier is formed from
unity, makes an easy road to the laws of addition and subtraction
of algebraic numbers, and to the law of signs in
multiplication and division.  All the principles and rules
of this chapter are illustrated and enforced by numerous
examples involving \emph{simple} algebraic expressions only.

The ordinary processes with \emph{compound} expressions, including
simple cases of resolution into factors, and the
treatment of fractions, naturally follow the third chapter.
The immediate succession of topics that require similar
work is of the highest importance to the beginner, and it
is hoped that the half-dozen chapters on algebraic expressions
will prove interesting, and give sufficient readiness
in the use of symbols.

A chapter on fractional equations with one unknown
number, a chapter on simultaneous equations with two
unknown numbers, and a chapter on quadratics follow in
order.  Only one method of elimination is given in simultaneous
equations and one method of completing the
square in quadratics.  Moreover, the solution of the examples
in quadratics requires the square roots of only small
numbers such as every pupil knows who has learned the
%% -----File: 005.png---Folio v-------
multiplication table.  In each of these three chapters a
considerable number of problems is given to \emph{state} and solve.
By this means the learner is led to exercise his reasoning
faculty, and to realize that the methods of Algebra require
a strictly logical process.  These problems, however, are
divided into classes, and a model solution of an example
of each class is given as a guide to the solution of other
examples of that class.

The course may end with the chapter on quadratics, but
the simple questions of arithmetical progression and of
geometrical progression are so interesting in themselves,
and show so clearly the power of Algebra, that it will
be a great loss not to take the short chapters on these
series.

The last chapter is on square and cube roots.  It is
expected that pupils who use this book will learn how to
extract the square and cube roots by the simple formulas
of Algebra, and be spared the necessity of committing to
memory the long and tedious rules given in Arithmetic,
rules that are generally forgotten in less time than they
are learned.

Any corrections or suggestions will be thankfully received
by the author.

A teachers' edition is in press, containing solutions of
examples, and such suggestions as experience with beginners
has shown to be valuable.

\Signature{G\Add{.} A. WENTWORTH.}
{\textsc{Exeter}, NH, April, 1894}
%% -----File: 006.png---Folio vi-------
\TableofContents
\iffalse
CONTENTS.

Chapter Page

I.     Introduction..............       1

II.    Simple Equations.............     19

III.   Positive and Negative Numbers.......     33

IV.    Addition and Subtraction.........     46

V.     Multiplication and Division........     53

VI.    Special Rules in Multiplication and Division .     64

VII.   Factors...............     71

VIII.  Common Factors and Multiples.......     84

IX.    Fractions................     89

X.     Fractional Equations...........   103

XI.    Simultaneous Equations of the First Degree     .    122

XII.   Quadratic Equations...........    132

XIII.  Arithmetical Progression.........    142

XIV.   Geometrical Progression..........   148

XV.    Square and Cube Roots..........    152

       Answers................    165
\fi
%% -----File: 007.png---Folio 1-------
\MainMatter
% FIRST STEPS IN  ALGEBRA.
% [** TN: Chapter macro prints preceding line]

\Chapter{I.}{Introduction.}

\begin{Remark}[\First{Note}\Add{.}]
The principal definitions are put at the beginning of the
book for convenient reference.  They are not to be committed to
memory.  It is a good plan to have definitions and explanations
read aloud in the class, and to encourage pupils to make comments
upon them, and ask questions about them.
\end{Remark}

\Paragraph{1.  Algebra.}  Algebra, like Arithmetic, treats of numbers.

\Paragraph{2.  Units.}  In counting separate objects or in measuring
magnitudes, the \emph{standards} by which we count or measure
are called \Defn{units}.

\begin{Remark}
Thus, in counting the boys in a school, the unit is a boy; in selling
eggs by the dozen, the unit is a dozen eggs; in selling bricks by
the thousand, the unit is a thousand bricks; in measuring short distances,
the unit is an inch, a foot, or a yard; in measuring long
distances, the unit is a rod or a mile.
\end{Remark}

\Paragraph{3.  Numbers.}  \emph{Repetitions of the unit} are expressed by
numbers.

\Paragraph{4.  Quantities.}  A number of specified units of any kind
is called a quantity; as, $4$~pounds, $5$~oranges.

\Paragraph{5.  Number-Symbols in Arithmetic.}  Arithmetic employs
the arbitrary symbols, $1$,~$2$, $3$, $4$, $5$, $6$, $7$, $8$, $9$,~$0$, called
\Defn{figures}, to represent numbers.
%% -----File: 008.png---Folio 2-------

\Paragraph{6.  Number-Symbols in Algebra.}  Algebra employs \emph{the
letters of the alphabet} in addition to the figures of Arithmetic
to represent numbers.  Letters are used as \emph{general}
symbols of numbers to which \emph{any particular values} may be
assigned.


\Section{PRINCIPAL SIGNS OF OPERATIONS.}

\Paragraph{7.}  The signs of the fundamental operations are the same
in Algebra as in Arithmetic.

\Paragraph{8.  The Sign of Addition,~$+$.} The sign~$+$ is read \emph{plus}.

\begin{Remark}
Thus, $4 + 3$, read $4$~plus~$3$, indicates that the number~$3$ is to be
added to the number~$4$, $a + b$, read $a$~plus~$b$, indicates that the number~$b$
is to be added to the number~$a$.
\end{Remark}

\Paragraph{9.  The Sign of Subtraction,~$-$\Add{.}} The sign~$-$ is read \emph{minus}.

\begin{Remark}
Thus, $4 - 3$, read $4$~minus~$3$, indicates that the number~$3$ is to be
subtracted from the number~$4$, $a - b$, read $a$~minus~$b$, indicates that
the number~$b$ is to be subtracted from the number~$a$.
\end{Remark}

\Paragraph{10.  The Sign of Multiplication,~$×$\Add{.}} The sign~$×$ is read
\emph{times}.

\begin{Remark}
Thus, $4 × 3$, read $4$~times~$3$, indicates that the number~$3$ is to be
multiplied by~$4$, $a × b$, read $a$~times~$b$, indicates that the number~$b$
is to be multiplied by the number~$a$.
\end{Remark}

A dot is sometimes used for the sign of multiplication.
Thus $2 · 3 · 4 · 5$ means the same as $2 × 3 × 4 × 5$.  Either
sign is read \emph{multiplied by} when followed by the multiplier.
\$$a × b$, or \$$a · b$, is read $a$~dollars multiplied by~$b$.

\Paragraph{11.  The Sign of Division,~$÷$.} The sign~$÷$ is read \emph{divided by}.

\begin{Remark}
Thus, $4 ÷ 2$, read $4$~divided by~$2$, indicates that the number~$4$ is
to be divided by~$2$, $a ÷ b$, read $a$~divided by~$b$, indicates that the
number~$a$ is to be divided by the number~$b$.
\end{Remark}
%% -----File: 009.png---Folio 3-------

Division is also indicated by writing the dividend above
the divisor with a horizontal line between them.

\begin{Remark}
Thus, $\dfrac{4}{2}$ means the same as~$4 ÷ 2$; $\dfrac{a}{b}$~means the same as~$a ÷ b$.
\end{Remark}


\Section{OTHER SIGNS USED IN ALGEBRA.}

\Paragraph{12.  The Sign of Equality,~$=$.} The sign~$=$ is read \emph{is equal
to}, when placed between two numbers and indicates that
these two numbers are equal.

\begin{Remark}
Thus, $8 + 4 = 12$ means that $8 + 4$ and~$12$ stand for \emph{equal} numbers;
$x + y = 20$ means that $x + y$ and~$20$ stand for equal numbers.
\end{Remark}

\Paragraph{13.  The Sign of Inequality, $>$~or~$<$.} The sign $>$~or~$<$ is
read \emph{is greater than} and \emph{is less than} respectively, and when
placed between two numbers indicates that these two numbers
are unequal, and that the number toward which the
sign opens is the greater.

\begin{Remark}
Thus, $9 + 6 > 12$ means that $9 + 6$ is  greater than~$12$; and
$9 + 6 < 16$ means that $9 + 6$ is less than~$16$.
\end{Remark}

%[** TN: [sic] no punctuation after \therefore]
\Paragraph{14.  The Sign of Deduction,~$\therefore$}\quad   The sign~$\therefore$ is read \emph{hence}
or \emph{therefore}.

\Paragraph{15.  The Sign of Continuation,~$\dots$.}    The sign~$\dots$ is read
\emph{and so on}.

\Paragraph{16.  The Signs of Aggregation.}  The signs of aggregation
are the bar~$|$, the vinculum~$\overline{\phantom{a+x}}$, the parenthesis~$(\ )$, the
bracket~$[\ ]$, and the brace~$\{\ \}$.

\begin{Remark}
Thus, each of the expressions $\VSum{a}{b}$, $\Vinc{a + b}$, $(a + b)$, $[a + b]$, $\{a + b\}$,
signifies that $a + b$ is to be treated as a single number.
\end{Remark}
%% -----File: 010.png---Folio 4-------


\Section{FACTORS. COEFFICIENTS. POWERS.}

\Paragraph{17.  Factors.}  When a number consists of the product of
two or more numbers, each of these numbers is called a
\Defn{factor} of the product.

The sign~$×$ is generally omitted between a figure and a
letter, or between letters; thus, instead of $63 × a × b$, we
write~$63ab$; instead of $a × b × c$, we write~$abc$.

The expression~$abc$ must not be confounded with $a + b + c$.
$abc$~is a product; $a + b + c$ is a sum.
\begin{DPalign*}
\lintertext{\indent If}
a = 2,\quad b &= 3,\quad c = 4, \\
\lintertext{then}
abc &= 2 × 3 × 4 = 24; \\
\lintertext{but}
a + b + c &= 2 + 3 + 4 = 9.
\end{DPalign*}

\begin{Remark}[Note.]
When a sign of operation is omitted in the notation of
Arithmetic, it is always the \emph{sign of addition}; but when a sign of
operation is omitted in the notation of Algebra, it is always the
\emph{sign of multiplication}.  Thus, $456$~means $400 + 50 + 6$, but $4ab$
means $4 × a × b$.
\end{Remark}

\Paragraph{18.}  Factors expressed by letters are called \Defn{literal} factors;
factors expressed by figures are called \Defn{numerical} factors.

\Paragraph{19.}  If one factor of a product is equal to~$0$, the product
is equal to~$0$, whatever the values of the other factors.
Such a factor is called a \Defn{zero factor}.

\Paragraph{20.  Coefficients.}  A known factor of a product which is
prefixed to another factor, to show the number of times that
factor is taken, is called a \Defn{coefficient}.

\begin{Remark}
Thus, in~$7c$, $7$~is the coefficient of~$c$; in~$7ax$, $7$~is the
coefficient of~$ax$,
or, if $a$~is known, $7a$~is the coefficient of~$x$.
\end{Remark}
%% -----File: 011.png---Folio 5-------

By coefficient, we generally mean the \textbf{numerical coefficient
with its sign}.  If no numerical coefficient is written, $1$~is
understood.  Thus, $ax$~means the same as~$1ax$.

\Paragraph{21.  Powers and Roots.}  A product consisting of two or
more \textbf{equal factors} is called a \Defn{power} of that factor, and one
of the equal factors is called a \Defn{root} of the number.

\begin{Remark}
Thus, $9 = 3 × 3$; that is, $9$~is a power of~$3$, and $3$~is a root of~$9$.
\end{Remark}

\Paragraph{22.  Indices or Exponents.}  An index or  exponent is a
number-symbol written at the right of, and a little above,
a number.

If the index is a \emph{whole number}, it shows the number
of times the given number is taken as a factor.

\begin{Remark}
Thus, $a^{1}$, or simply~$a$, denotes that $a$~is taken \emph{once} as a
factor; $a^{2}$~denotes
that $a$~is taken \emph{twice} as a factor; $a^{3}$~denotes that $a$~is taken
\emph{three times} as a factor; and $a^{4}$~denotes that $a$~is taken \emph{four times} as a
factor; and so on.  These are read: the first power of~$a$; the second
power of~$a$; the third power of~$a$; the fourth power of~$a$; and so on.

$a^{3}$~is written instead of~$aaa$.

$a^{4}$~is written instead of~$aaaa$.
\end{Remark}

\Paragraph{23.}  The meaning of coefficient and exponent must be
carefully distinguished.  Thus,
\begin{DPalign*}
4a &= a + a + a + a; \\
a^{4} &= a× a× a× a. \displaybreak[1] \\
\lintertext{\indent If $a = 3$,}
4a &= 3 + 3 + 3 + 3 = 12. \\
a^{4} &= 3 × 3 × 3 × 3 = 81.
\end{DPalign*}

\begin{Remark}
The second power of a number is generally called the \emph{square} of
that number; thus, $a^{2}$~is called the \emph{square} of~$a$, because if $a$~denotes
the number of units of length in the side of a square, $a^{2}$~denotes the
number of units of surface in the square.  The third power of a number
is generally called the \emph{cube} of that number; thus, $a^{3}$~is called the
\emph{cube} of~$a$, because if $a$~denotes the number of units of length in the
edge of a cube, $a^{3}$~denotes the number of units of volume in the
cube.
\end{Remark}
%% -----File: 012.png---Folio 6-------


\Section{ALGEBRAIC EXPRESSIONS.}

\Paragraph{24.  An Algebraic Expression.}  An algebraic expression is
a number written with algebraic symbols.  An algebraic
expression may consist of one symbol, or of several symbols
connected by signs.

\begin{Remark}
Thus, $a$, $3abc$, $5a + 2b - 3c$, are algebraic expressions.
\end{Remark}

\Paragraph{25.  Terms.}  A \Defn{term} is an algebraic expression, the parts
of which are not separated by the sign $+$~or~$-$.

\begin{Remark}
Thus, $a$, $5xy$, $2ab × 4cd$, $\dfrac{3ab}{4cd}$ are algebraic expressions of one
term each.  A term may be separated into parts by the sign $×$~or~$÷$.
\end{Remark}

\Paragraph{26.  Simple Expressions.}  An algebraic expression of \emph{one
term} is called a \Defn{simple expression} or \Defn{monomial}.

\begin{Remark}
Thus, $5xy$, $7a × 2b$, $7a ÷ 2b$, are simple expressions.
\end{Remark}

\Paragraph{27.  Compound Expressions.}  An algebraic expression of
\emph{two or more terms} is called a \Defn{compound expression} or \Defn{polynomial}.

\begin{Remark}
Thus, $5xy + 7a$, $2x - y - 3z$, $4a - 3b + 2c - 3d$ are compound
expressions.
\end{Remark}

\Paragraph{28.}  A polynomial of two terms is called a \Defn{binomial}; of
three terms, a \Defn{trinomial}.

\begin{Remark}
Thus, $3a - b$ is a binomial; and $3a - b + c$ is a trinomial.
\end{Remark}

\Paragraph{29.  Positive and Negative Terms.}  The terms of a compound
expression preceded by the sign~$+$ are called \Defn{positive
terms}, and the terms preceded by the sign~$-$ are called
\Defn{negative terms}.  The sign~$+$ before the first term is omitted.

\Paragraph{30.}  A positive and a negative term of the same numerical
value cancel each other when combined.
%% -----File: 013.png---Folio 7-------

\Paragraph{31.  Like Terms.}  Terms which have the same combination
of \emph{letters} are called \Defn{like} or \Defn{similar} terms; terms which
do not have the same combination of letters are called
\Defn{unlike} or \Defn{dissimilar} terms.

\begin{Remark}
Thus, $5a^{2}bc$, $-7a^{2}bc$, $a^{2}bc$, are like terms; but $5a^{2}bc$, $5ab^{2}c$,
$5abc^{2}$, are unlike terms.
\end{Remark}

\Paragraph{32.  Degree of a Term.}  A term that is the product of
three letters is said to be of the \emph{third degree}; a term of
four letters is of the \emph{fourth degree}; and so on.

\begin{Remark}
Thus, $5abc$~is of the third degree; $2a^{2}b^{2}c^{2}$, that is, $2aabbcc$, is of
the sixth degree.
\end{Remark}

\Paragraph{33.  Degree of a Compound Expression.}  The degree of a
compound expression is the degree of that term of the
expression which is of the \emph{highest degree}.

\begin{Remark}
Thus, $a^{2}x^{2} + bx + c$ is of the fourth degree, since $a^{2}x^{2}$~is of the
fourth degree.
\end{Remark}

\Paragraph{34.  Dominant Letter.}  It often happens that there is one
letter in an expression of more importance than the rest,
and this is, therefore, called the \Defn{dominant letter}.  In such
cases the degree of the expression is generally called by
the degree of the \emph{dominant letter}.

\begin{Remark}
Thus, $a^{2}x^{2} + bx + c$ is of the \emph{second degree in~$x$}.
\end{Remark}

\Paragraph{35.  Arrangement of a Compound Expression.}  A compound
expression is said to be \emph{arranged} according to the powers
of some letter when the exponents of that letter, reckoning
from left to right, either descend or ascend in \emph{the order of
magnitude}.

\begin{Remark}
Thus, $3ax^{3} - 4bx^{2} - 6ax + 8b$ is arranged according to the descending
powers of~$x$, and $8b - 6ax - 4bx^{2} + 3ax^{3}$ is arranged
according to the ascending powers of~$x$.
\end{Remark}
%% -----File: 014.png---Folio 8-------


\Section{PARENTHESES.}

\Paragraph{36.}  If a compound expression is to be treated as a whole,
it is enclosed in a parenthesis.

\begin{Remark}
Thus, $2 × (10 + 5)$ means that we are to add $5$~to~$10$ and multiply
the result by~$2$; if we were to omit the parenthesis and write
$2 × 10 + 5$, the meaning would be that we were to multiply $10$~by~$2$
and add~$5$ to the result.
\end{Remark}

Like the parenthesis, we use with the same meaning any
other sign of aggregation.

\begin{Remark}
Thus, $(5 + 2)$, $[5 + 2]$, $\{5 + 2\}$, $\Vinc{5 + 2}$, $\VSum{5}{2}$, all mean that the
expression $5 + 2$ is to be treated as the single symbol~$7$.
\end{Remark}

\Paragraph{37.  Parentheses preceded by~$+$.} If a man has $10$~dollars
and afterwards collects $3$~dollars and then $2$~dollars, it
makes no difference whether he adds the $3$~dollars to his
$10$~dollars, and then the $2$~dollars, or puts the $3$~and~$2$
dollars together and adds their sum to his $10$~dollars.

The first process is represented by $10 + 3 + 2$.

The second process is represented by $10 + (3 + 2)$.

\begin{DPgather*}
\lintertext{\indent Hence,}
10 + (3 + 2) = 10 + 3 + 2.
\Tag{(1)}
\end{DPgather*}

If a man has $10$~dollars and afterwards collects $3$~dollars
and then pays a bill of $2$~dollars, it makes no difference
whether he adds the $3$~dollars collected to his $10$~dollars
and pays out of this sum his bill of $2$~dollars, or pays the
$2$~dollars from the $3$~dollars collected and adds the remainder
to his $10$~dollars.

The first process is represented by $10 + 3 - 2$.

The second process is represented by $10 + (3 - 2)$.

\begin{DPgather*}
\lintertext{\indent Hence,}
10 + (3 - 2) = 10 + 3 - 2.
\Tag{(2)}
\end{DPgather*}
%% -----File: 015.png---Folio 9-------

From (1)~and~(2) it follows that

\begin{Theorem}
If an expression within a parenthesis is preceded by the
sign~$+$, the parenthesis can be removed without making any
change in the signs of the expression.
\end{Theorem}

\begin{Theorem}[\textsc{Conversely.}] Any part of an expression can lie enclosed
within a parenthesis and the sign~$+$ prefixed, without making
any change in the signs of the terms thus enclosed.
\end{Theorem}

\Paragraph{38.  Parentheses preceded by~$-$.} If a man has $10$~dollars
and has to pay two bills, one of $3$~dollars and one of $2$~dollars,
it makes no difference whether he takes $3$~dollars
and $2$~dollars in succession, or takes the $3$~and~$2$ dollars at
one time, from his $10$~dollars.

The first process is represented by $10 - 3 - 2$.

The second process is represented by $10 - (3 + 2)$.

\begin{DPgather*}
\lintertext{\indent Hence,}
10 - (3 + 2) = 10 - 3 - 2.
\Tag{(3)}
\end{DPgather*}

If a man has $10$~dollars consisting of $2$~five-dollar bills,
and has a debt of $3$~dollars to pay, he can pay his debt by
giving a five-dollar bill and receiving $2$~dollars.

This process is represented by $10 - 5 + 2$.

Since the debt paid is $3$~dollars, that is, $(5 - 2)$~dollars,
the number of dollars he has left can evidently be
expressed by
\begin{DPalign*}
10 &- (5 - 2). \\
\lintertext{\indent Hence,}
10 &- (5 - 2) = 10 - 5 + 2.
\Tag{(4)}
\end{DPalign*}

From (3)~and~(4) it follows that

\begin{Theorem}
If an expression within a parenthesis is preceded by the
sign~$-$, the parenthesis can be removed, provided the sign
before each term within the parenthesis is changed, the
sign~$+$ to~$-$, and the sign~$-$ to~$+$.
\end{Theorem}
%% -----File: 016.png---Folio 10-------

\begin{Theorem}[\textsc{Conversely.}] Any part of an expression can be enclosed
within a parenthesis and the sign~$-$ prefixed, provided the
sign of each term enclosed is changed, the sign~$+$ to~$-$, and
the sign~$-$ to~$+$.
\end{Theorem}


\Exercise{1.}

Remove the parentheses, and combine:
\begin{multicols}{2}
\Item{1.}  $9 + (3 + 2)$.

\Item{2.}  $9 + (3 - 2)$.

\Item{3.}  $7 + (5 + 1)$.

\Item{4.}  $7 + (5 - 1)$.

\Item{5.}  $6 + (4 + 3)$.

\Item{6.}  $6 + (4 - 3)$.

\Item{7.}  $3 + (8 - 2)$.

\Item{8.}  $9 - (8 - 6)$.

\Item{9.}  $10 - (9 - 5)$.

\Item{10.}  $9 - (6 + 1)$.

\Item{11.}  $8 - (3 + 2)$.

\Item{12.}  $7 - (3 - 2)$.

\Item{13.}  $9 - (4 + 3)$.

\Item{14.}  $9 - (4 - 3)$.

\Item{15.}  $7 - (5 - 2)$.

\Item{16.}  $7 - (7 - 3)$.

\Item{17.}  $(8 - 6) - 1$.

\Item{18.}  $(3 - 2) - (1 - 1)$.

\Item{19.}  $(7 - 3) - (3 - 2)$.

\Item{20.}  $(8 - 2) - (5 - 3)$.

\Item{21.}  $15 - (10 - 3 - 2)$.
\end{multicols}

\Paragraph{39.  Multiplying a Compound Expression.}  The expression
$4(5 + 3)$ means that we are to take the sum of the numbers
$5$~and~$3$ four times.  The process can be represented by
placing five dots in a line, and a little to the right three
more dots in the same line, and then placing a second,
third, and fourth line of dots underneath the first line and
exactly similar to it.
\[
\begin{array}{*{9}{>{\ }r}}
\DOT & \DOT & \DOT & \DOT & \DOT & \quad & \DOT & \DOT & \DOT \\[-4pt]
\DOT & \DOT & \DOT & \DOT & \DOT &       & \DOT & \DOT & \DOT \\[-4pt]
\DOT & \DOT & \DOT & \DOT & \DOT &       & \DOT & \DOT & \DOT \\[-4pt]
\DOT & \DOT & \DOT & \DOT & \DOT &       & \DOT & \DOT & \DOT \\[-4pt]
\end{array}
\]

There are $(5 + 3)$ dots in each line, and $4$~lines.  The
total number of dots, therefore, is $4 × (5 + 3)$.

We see that in the left-hand group there are $4 × 5$ dots,
and in the right-hand group $4 × 3$ dots.  The sum of these
%% -----File: 017.png---Folio 11-------
two numbers $(4 × 5) + (4 × 3)$ must be equal to the total
number; that is,
\begin{align*}
4(5 + 3) &= (4 × 5) + (4 × 3) \\
         &= 20 + 12.
\end{align*}

Again, the expression $4(8 - 3)$ means that we are to
take the difference of the numbers $8$~and~$3$ four times.
The process can be represented by placing eight dots in a
line and crossing the last three, and then placing a second,
third, and fourth line of dots underneath the first line and
exactly similar to it.
\[
%[** TN: Added gap between dot groups.]
\begin{array}{*{9}{>{\ }r}}
\DOT & \DOT & \DOT & \DOT & \DOT & \quad & \NOT & \NOT & \NOT \\[-4pt]
\DOT & \DOT & \DOT & \DOT & \DOT &       & \NOT & \NOT & \NOT \\[-4pt]
\DOT & \DOT & \DOT & \DOT & \DOT &       & \NOT & \NOT & \NOT \\[-4pt]
\DOT & \DOT & \DOT & \DOT & \DOT &       & \NOT & \NOT & \NOT \\[-4pt]
\end{array}
\]

The whole number of dots not crossed in each line is
evidently $(8 - 3)$, and the whole number of lines is~$4$.
Therefore the total number of dots not crossed is
\[
4 × (8 - 3).
\]

The total number of dots (crossed and not crossed) is
$(4 × 8)$, and the total number of dots crossed is~$(4 × 3)$.
Therefore the total number of dots not crossed is
\begin{DPalign*}
(4 × 8) &- (4 × 3); \\
\lintertext{that is,}
4(8 - 3) &= (4 × 8) - (4 × 3) \\
         &= 32 - 12. \displaybreak[1] \\
\intertext{\indent If $a$, $b$, and~$c$ stand for any three numbers, we have}
a (b + c) &= ab + ac, \\
\lintertext{and}
a(b - c) &= ab - ac.
\EqText{Therefore,}
\end{DPalign*}

\Dictum{To multiply a compound expression by a simple one},
\begin{Theorem}
Multiply each term by the multiplier, and write the successive
products with the same signs as those of the original
terms.
\end{Theorem}
%% -----File: 018.png---Folio 12-------

\Exercise{2.}

Multiply and remove parentheses:
\begin{multicols}{3}
\Item{1.}  $7(8 + 5)$.

\Item{2.}  $7(8 - 5)$.

\Item{3.}  $6(7 + 3)$.

\Item{4.}  $6(7 - 3)$.

\Item{5.}  $8(7 + 5)$.

\Item{6.}  $8(7 - 5)$.

\Item{7.}  $9(6 - 2)$.

\Item{8.}  $4(a + b)$.

\Item{9.}  $4(a - b)$.

\Item{10.}  $2(a^{2} + b^{2})$.

\Item{11.}  $2(a^{2} - b^{2})$.

\Item{12.}  $3(ab + c)$.

\Item{13.}  $3(ab - c)$.

\Item{14.}  $3(c - ab)$.

\Item{15.}  $a(b + c)$.

\Item{16.}  $a(b - c)$.

\Item{17.}  $3a(b + c)$.

\Item{18.}  $3a(b - c)$.

\Item{19.}  $5a(b^{2} + c)$.

\Item{20.}  $5a(b^{2} - c^{2})$.

\Item{21.}  $5a^{2}(b^{2} - c)$.
\end{multicols}

\Paragraph{40.}  The numerical value of an algebraic expression is the
number obtained by putting for the letters involved the
numbers for which these letters stand, and then performing
the operations required by the signs.

\Item{1.}  If $b = 4$, find the value of~$3b^{2}$.
\begin{Soln}
\begin{DPgather*}
\lintertext{\indent Here}
3b^{2} = 3 × 4^{2} = 3 × 16 = 48.
\end{DPgather*}
\end{Soln}

\Item{2.}  If $a = 7$, $b = 2$, $c = 3$, find the value of~$5ab^{2}c^{3}$.
\begin{Soln}
\begin{DPgather*}
\lintertext{\indent Here}
5ab^{2}c^{3} = 5 × 7 × 2^{2} × 3^{3} = 3780.
\end{DPgather*}
\end{Soln}

\Exercise{3.}

If $a = 7$, $b = 5$, $c = 3$, find the value of
\begin{multicols}{3}
\Item{1.}  $9a$.

\Item{2.}  $8ab$.

\Item{3.}  $4b^{2}c$.

\Item{4.}  $2a^{2}$.

\Item{5.}  $3c^{3}$.

\Item{6.}  $2b^{4}$.

\Item{7.}  $5ac$.

\Item{8.}  $abc$.

\Item{9.}  $abc^{2}$.

\Item{10.}  $\frac{1}{3}abc$.

\Item{11.}  $\frac{1}{5}ab^{2}c$.

\Item{12.}  $\frac{1}{7}a^{2}bc$.
\end{multicols}

If $a = 5$, $b = 2$, $c = 0$, $x = 1$, $y = 3$, find the value of
\begin{multicols}{3}
\Item{13.}  $4acy^{2}$.

\Item{14.}  $3ax^{5}y^{2}$.

\Item{15.}  $2ab^{2}y$.

\Item{16.}  $2a^{2}b^{2}c^{2}y^{2}$.

\Item{17.}  $2a^{2}b^{2}x^{2}y^{2}$.

\Item{18.}  $2abx^{3}y^{3}$.

\Item{19.}  $3abcxy$.

\Item{20.}  $3abx^{3}y^{2}$.

\Item{21.}  $3ab^{2}xy^{2}$.
\end{multicols}
%% -----File: 019.png---Folio 13-------

\Paragraph{41.  The Numerical Value of a Compound Expression.}

If $a$~stands for~$10$, $b$~for~$4$, and $c$~for~$3$, find the value of
the expression $5ab - 10c^{2} - 5b^{2}$.

Find the value of each term, and combine the results.
\begin{Soln}
\begin{gather*}
\begin{alignedat}{3}
    5ab &\text{ stands for } &5 × 10 &× 4    &&= 200; \\
10c^{2} &\text{ stands for } &    10 &× 3^{2} &&= \Z90; \\
5b^{2} &\text{ stands for }  &     5 &× 4^{2} &&= \Z80.
\end{alignedat} \\
\begin{aligned}
&\therefore 5ab - 10c^{2} - 5b^{2} \\
&= 200 - 90 - 80 \\
&=  30.
\end{aligned}
\end{gather*}
\end{Soln}

\Paragraph{42.}  In finding the value of a compound expression the
operations indicated \emph{for each term} must be performed \emph{before}
the operation indicated by the sign prefixed to the term.

When there is no sign expressed between single symbols
or between \emph{simple} and \emph{compound expressions}, it must be
remembered that the sign understood is the \emph{sign of multiplication}.
Thus $2(a - b)$ has the same meaning as $2 × (a - b)$.

\Exercise{4.}

If $a = 5$, $b = 4$, $c = 3$, find the value of
\begin{multicols}{2}
\Item{1.}  $9a - 2bc$.

\Item{2.}  $ab + 2c$.

\Item{3.}  $abc + bc$.

\Item{4.}  $5ac + 2a$.

\Item{5.}  $2abc - 2ac^{2}$.

\Item{6.}  $ab + bc - ac$.

\Item{7.}  $ac - (b + c)$.

\Item{8.}  $a^{2} + (b^{2} + c^{2})$\Add{.}

\Item{9.}  $2a + (2b + 2c)$.

\Item{10.}  $a^{2} - b^{2} - c^{2}$.

\Item{11.}  $3(a - b + c)$.

\Item{12.}  $6ab - (bc + 8)$.

\Item{13.}  $7bc - c^{2} + a$.

\Item{14.}  $5ac - b^{2} + 3b$.

\Item{15.}  $4b^{2}c - 5c^{2} - 2b$.

\Item{16.}  $2a + (b + c)$.

\Item{17.}  $b + 2(a - c)$.

\Item{18.}  $c + 2(a - b)$.

\Item{19.}  $2a - (b + c)$.

\Item{20.}  $2b - (a - c)$.

\Item{21.}  $2c - (a - b)$.

\Item{22.}  $2c - 5(a - b)$.

\Item{23.}  $2b - 3(a - c)$.

\Item{24.}  $2c - b(a - b)$.
\end{multicols}
%% -----File: 020.png---Folio 14-------


\Section{ALGEBRAIC NOTATION.}

\Exercise{5.}

\Item{1.}  Read $a + b$; $a - b$; $ab$; $a÷b$.

\Item{2.}  Write six increased by four. \Ans{$6 + 4$.}

\Item{3.}  Write $a$ increased by~$b$.

\Item{4.}  Write six diminished by four. \Ans{$6 - 4$.}

\Item{5.}  Write $a$ diminished by~$b$.

\Item{6.}  By how much does twenty-five exceed sixteen?
\Ans{$25 - 16$.}

\Item{7.}  By how much does $x$ exceed~$y$?

\Item{8.}  Write four times three; the fourth power of three.
\Ans{$4 × 3$; $3^{4}$.}

\Item{9.}  Write four times~$x$; the fourth power of~$x$.

\Item{10.}  If one part of twenty-five is  fifteen, what is the
other part?
\Ans{$25 - 15$.}

\Item{11.}  If one part of~$35$ is~$x$, what is the other part?

\Item{12.}  If one part of~$x$ is~$a$, what is the other part?

\Item{13.}  How much does ten lack of being twelve?
\Ans{$12 - 10$.}

\Item{14.}  How much does $x$ lack of being fourteen?

\Item{15.}  How much does $x$ lack of being~$a$?

\Item{16.}  If a man walks four miles an hour, how many miles
will he walk in three hours?
\Ans{$3 × 4$.}

\Item{17.}  If a man walks $y$~miles an hour, how many miles
will he walk in $x$~hours?

\Item{18.}  If a man walks $y$~miles an hour, how many hours
will it take him to walk $x$~miles?
%% -----File: 021.png---Folio 15-------

\Exercise{6.}

\Item{1.}  If the dividend is twenty and the quotient five,
what is the divisor?
\Ans{$\frac{20}{5}$.}

\Item{2.}  If the dividend is~$a$ and the quotient~$b$, what is the
divisor?

\Item{3.}  If John is twenty years old to-day, how old was he
four years ago?   How old will he be five years hence?
\Ans{$20 - 4$; $20 + 5$.}

\Item{4.}  If James is $x$~years old to-day, how old was he three
years ago?    How old will he be seven years hence?

\Item{5.}  Write four times the expression seven minus five.
\Ans{$4(7 - 5)$.}

\Item{6.}  Write seven times the expression $2x$~minus~$y$.

\Item{7.}  Write the next integral number above four.
\Ans{$4 + 1$.}

\Item{8.}  If $x$~is an integral number, write the next integral
number above it; the next integral number below it.

\Item{9.}  What number is less than $20$ by~$d$?

\Item{10.}  If the difference of two numbers is five, and the
smaller number is fifteen, what is the greater number?
\Ans{$15 + 5$.}

\Item{11.}  If the difference of two numbers is eight, and the
smaller number is~$x$, what is the greater number?

\Item{12.}  If the sum of two numbers is~$30$, and one of them
is~$20$, what is the other?
\Ans{$30 - 20$.}

\Item{13.}  If the sum of two numbers is~$x$, and one of them is~$10$,
what is the other?

\Item{14.}  If $100$ contains~$x$ ten times, what is the value of~$x$?
%% -----File: 022.png---Folio 16-------

\Exercise{7.}

\Item{1.}  In $x$~years a man will be $40$~years old; what is his
present age?

\Item{2.}  How old will a man be in $y$~years, if his present age
is $a$~years?

\Item{3.}  What is the value of~$x$ if $7x$~equals~$28$?

\Item{4.}  If it takes $3$~men $4$~days to reap a field, how many
days will it take one man to reap it?
\Ans{$3 × 4$.}

\Item{5.}  If it takes $a$~men $b$~days to reap a field, how many
days will it take one man to reap it?

\Item{6.}  What is the excess of~$5x$ over~$3x$?

\Item{7.}  By how much does $20 - 3$ exceed $(10 + 1)$?
\Ans{$20 - 3 - (10 + 1)$.}

\Item{8.}  By how much does $2x - 3$ exceed $(x + 1)$?

\Item{9.}  If $x$~stands for~$10$, find the value of~$4(3x - 20)$.

\Item{10.}  If $a$~stands for~$10$, and $b$~for~$2$, find the value of
$2(a - 2b)$.

\Item{11.}  How many cents in $a$~dollars, $b$~quarters, and $c$~dimes?

\Item{12.}  A book-shelf contains French, Latin, and Greek
books.  There are $100$~books in all, and there are $x$~Latin
and $y$~Greek books.  How many French books are there?

\Item{13.}  A regiment of men is drawn up in $10$~ranks of $80$~men
each, and there are $15$~men over.  How many men
are there in the regiment?
\Ans{$10 × 80 + 15$.}

\Item{14.}  A regiment of men is drawn up in $x$~ranks of $y$~men
each, and there are $c$~men over.  How many men are there
in the regiment?
%% -----File: 023.png---Folio 17-------

\Exercise{8.}

\Item{1.}  A room is $10$~yards long and $8$~yards wide.  In the
middle there is  a carpet $6$~yards square.  How many
square yards of oilcloth will be required to cover the rest
of the floor?
\Ans{$10 × 8 - 6^{2}$.}

\Item{2.}  A room is $x$~yards long and $y$~yards wide.  In the
middle there is  a carpet $a$~yards square.  How many
square yards of oilcloth will be required to cover the rest
of the floor?

\Item{3.}  How many rolls of paper $g$~feet long and $k$~feet
wide will be required to paper a room, the perimeter of
which, after proper allowance is made for doors and windows,
is $p$~feet and the height $h$~feet?

\Item{4.}  Write six times the square of~$m$, plus five~$c$ times
the expression $d$~plus $b$~minus~$a$.

\Item{5.}  Write five times the expression two~$n$ plus one,
diminished by six times the expression $c$~minus $a$~plus~$b$.

\Item{6.}  A lady bought a dress for $a$~dollars, a cloak for $b$~dollars,
two pairs of gloves for $c$~dollars a pair.  She gave
a hundred-dollar bill in  payment.  How much money
should be returned to her?

\Item{7.}  If a man can perform a piece of work in $4$~days, how
much of it can he do in one day?
\Ans{$\frac{1}{4}$.}

\Item{8.}  If a man can perform a piece of work in $x$~days,
how much of it can he do in one day?

\Item{9.}  If A~can do a piece of work in $x$~days, B~in $y$~days,
C~in $z$~days, how much of it can they all do in one day,
working together?

\Item{10.}  Write an expression for the sum, and also for the
product, of three consecutive numbers of which the least is~$n$.
%% -----File: 024.png---Folio 18-------

\Item{11.}  The product of two factors is~$36$; if one of the
factors is~$x$, what is the other factor?

\Item{12.}  If $d$~is the divisor and $q$~the quotient, what is the
dividend?

\Item{13.}  If $d$~is the divisor, $q$~the quotient, and $r$~the remainder,
what is the dividend?

\Item{14.}  If $x$~oranges can be bought for $50$~cents, how many
oranges can be bought for $100$~cents?

\Item{15.}  What is the price in cents of $x$~apples, if they are
ten cents a dozen?

\Item{16.}  If $b$~oranges cost $6$~cents, what will $a$~oranges cost?

\Item{17.}  How many miles between two places, if a train
travelling $m$~miles an hour requires $4$~hours to make the
journey?

\Item{18.}  If a man was $x$~years old $10$~years ago, how many
years old will he be $7$~years hence?

\Item{19.}  If a man was $x$~years old $y$~years ago, how many
years old will he be $c$~years hence?

\Item{20.}  If a floor is $3x$~yards long and $12$~yards wide, how
many square yards does the floor contain?

\Item{21.}  How many hours will it take to walk $c$~miles, at
the rate of one mile in $15$~minutes?

\Item{22.}  Write three consecutive numbers of which $x$~is the
middle number.

\Item{23.}  If an odd number is represented by~$2n + 1$, what
will represent the next odd number?
%% -----File: 025.png---Folio 19-------


\Chapter{II.}{Simple Equations.}

\Paragraph{43.  Equations.}  An equation is a statement in symbols
that two expressions stand for the same number.

\begin{Remark}
Thus, the equation $3x + 2 = 8$ states that $3x + 2$ and~$8$ stand for
the same number.
\end{Remark}

\Paragraph{44.}  That part of the equation which precedes the sign
of equality is called the \Defn{first member}, or \Defn{left side}, and that
which follows the sign of equality is called the \Defn{second member},
or \Defn{right side}.

\Paragraph{45.}  The statement of equality between two algebraic
expressions, if true for all values of the letters involved, is
called an \Defn{identical equation};  but if true only for certain
particular values of the letters involved, it is called an
\Defn{equation of condition}.

\begin{Remark}
Thus, $a + b = b + a$, which is true for \emph{all values} of $a$~and~$b$, is an
\emph{identical equation}, and $3x + 2 = 8$, which is true only when $x$~stands
for~$2$, is an \emph{equation of condition}\Add{.}
\end{Remark}

For brevity, an identical equation is called an \Defn{identity},
and an equation of condition is called simply an \Defn{equation}.

\Paragraph{46.}  We often employ an equation to discover an \emph{unknown
number} from its relation to known numbers.  We usually
represent the unknown number by one of the \emph{last} letters
of the alphabet, as $x$,~$y$,~$z$; and by way of distinction, we
use the \emph{first} letters, $a$,~$b$, $c$,~etc., to represent numbers that
%% -----File: 026.png---Folio 20-------
are supposed to be known, though not expressed in the
number-symbols of Arithmetic.

\begin{Remark}
Thus, in the equation $ax + b = c$, $x$~is supposed to represent an
unknown number, and $a$,~$b$, and~$c$ are supposed to represent known
numbers.
\end{Remark}

\Paragraph{47.  Simple Equations.}  An equation which contains the
first power of~$x$, the symbol for the unknown number, and
no higher power, is called a \Defn{simple equation}, or an \Defn{equation
of the first degree}.

\begin{Remark}
Thus, $ax + b = c$ is a simple equation, or an equation of the first
degree \emph{in~$x$}.
\end{Remark}

\Paragraph{48.  Solution of an Equation.}  To solve an equation is to
find the unknown number; that is, the number which, when
substituted for its symbol in the given equation, renders the
equation an identity.  This number is said to \emph{satisfy} the
equation, and is called the \Defn{root} of the equation.

\Paragraph{49.  Axioms.}  In solving an equation, we make use of the
following axioms:

\Ax{1.} If equal numbers be added to equal numbers,
the sums will be equal.

\Ax{2.} If equal numbers be subtracted from equal numbers,
the remainders will be equal.

\Ax{3.} If equal numbers be multiplied by equal numbers,
the products will be equal.

\Ax{4.} If equal numbers be divided by equal numbers,
the quotients will be equal.

\begin{Theorem}
If, therefore, the two sides of an equation be increased by,
diminished by, multiplied by, or divided by equal numbers,
the results will be equal.
\end{Theorem}

%[** TN: Next paragraph set in normal-size type in the original]
\begin{Remark}
Thus, if $8x = 24$, then $8x + 4 = 24 + 4$, $8x - 4 = 24 - 4$,
$4 × 8x = 4 × 24$, and $8x ÷ 4 = 24 ÷ 4$.
\end{Remark}
%% -----File: 027.png---Folio 21-------

\Paragraph{50.  Transposition of Terms.}  It becomes necessary in solving
an equation to bring all the terms that contain the
symbol for the unknown number to one side of the equation,
and all the other terms to the other side.  This is
called \Defn{transposing the terms}. We will illustrate by examples:

\Item{1.}  Find the number for which $x$~stands when
\[
14x - 11 = 5x + 70.
\]

The first object to be attained is to get all the terms
which contain~$x$ on the left side of the equation, and all the
other terms on the right side.  This can be done by first
subtracting~$5x$ from both sides (Ax.~2), which gives
\[
9x - 11 = 70,
\]
and then adding~$11$ to these equals (Ax.~1), which gives
\begin{DPalign*}
9x + 11 - 11 &= 70 + 11. \\
\lintertext{\indent Combine,}
9x &= 81. \\
\lintertext{\indent Divide by~$9$,}
x &= 9.
\end{DPalign*}

\Item{2.}  Find the number for which $x$~stands when $x + b = a$.
\begin{DPalign*}[m]
\lintertext{\indent The equation is}
x + b &= a. \\
\lintertext{\indent Subtract~$b$ from each side,}
x + b - b &= a - b.
\rintertext{(Ax.~2)}
\end{DPalign*}

Since $+b$~and~$-b$ in the left side cancel each other
(§~30), we have $x = a - b$.

\Item{3.}  Find the number for which $x$~stands when $x - b = a$.
\begin{DPalign*}
\lintertext{\indent The equation is}
x - b &= a. \\
\lintertext{Add $+b$ to each side,}
x + b - b &= a + b.
\rintertext{(Ax.~1)}
\end{DPalign*}

Since $+b$~and~$-b$ in the left side cancel each other
(§~30), we have $x = a + b$.
%% -----File: 028.png---Folio 22-------

\Paragraph{51.}  The effect of the operation in the preceding equations,
when Axioms (1)~and~(2) are used, is to take a term
from one side and put it on the other side with its sign
changed.  We can proceed in a like manner in any other
case.  Hence the general rule:

\Paragraph{52.} \begin{Theorem}[nopar] Any term may be transposed from one side of an
equation to the other, provided its sign is changed.
\end{Theorem}

\Paragraph{53.}  Any term, therefore, which occurs on both sides
with \emph{the same sign} may be removed from both without
affecting the equality; and the sign of every term of an
equation may be changed without affecting the equality.

\Paragraph{54.  Verification.}  When the root is substituted for its
symbol in the given equation, and the equation reduces to
an \emph{identity}, the root is said to be \Defn{verified}.  We will illustrate
by examples:

\Item{1.}  What number added to twice itself gives~$24$?

Let $x$~stand for the number; \\
then $2x$~will stand for twice the number, \\
and the number added to twice itself will be $x + 2x$.

But the number added to twice itself is~$24$.
\begin{DPalign*}
\therefore  x + 2x &= 24. \\
\lintertext{\indent Combine $x$~and~$2x$,}
3x &= 24. \\
\intertext{\indent Divide by~$3$, the coefficient of~$x$,}
x &= 8.
\rintertext{(Ax.~4)}
\end{DPalign*}

Therefore the required number is~$8$.

\begin{DPalign*}
\lintertext{\indent\textsc{Verification.}}
x + 2x &= 24, \\
8 + 2 × 8 &= 24, \\
8 + 16 &= 24, \\
24 &= 24.
\end{DPalign*}
%% -----File: 029.png---Folio 23-------

\ScreenBreak
\Item{2.}  If $4x - 5$ stands for~$19$, for what number does $x$~stand?

We have the equation
\begin{DPalign*}
4x - 5 &= 19. \\
\lintertext{\indent Transpose $-5$,}
4x &= 19 + 5. \\
\lintertext{\indent Combine,}
4x &= 24. \\
\lintertext{\indent Divide by~$4$,}
x &= 6.
\rintertext{(Ax.~4)} \displaybreak[1] \\
\lintertext{\indent \textsc{Verification.}}
4x - 5 &= 19, \\
4 × 6 - 5 &= 19, \\
24 - 5 &= 19, \\
19 &= 19.
\end{DPalign*}

\Item{3.}  If $3x - 7$ stands for the same number as $14 - 4x$,
what number does $x$~stand for?

We have the equation
\[
3x - 7 = 14 - 4x.
\]

Transpose $-4x$ to the left side, and $-7$ to the right side,
\begin{DPalign*}
3x + 4x &= 14 + 7. \\
\lintertext{\indent Combine,}
7x &= 21. \\
\lintertext{\indent Divide by~$7$,}
x &= 3. \displaybreak[1] \\
\lintertext{\indent \textsc{Verification.}}
3x - 7 &= 14 - 4x, \\
3 × 3 - 7 &= 14 - 4 × 3, \\
2 &= 2.
\end{DPalign*}

\Item{4.}  Solve the equation
\[
7(x - 1) - 30 = 4(x - 4).
\]
We have the equation
\[
7(x - 1) - 30 = 4(x - 4).
\]
%% -----File: 030.png---Folio 24-------
%[** TN: Equation repeated at page break in the original]
% 7(x - 1) - 30 = 4(x - 4).

Remove the parentheses,
\begin{DPalign*}
7x - 7 - 30 &= 4x - 16. \\
\lintertext{\indent Then}
7x - 4x &= 7 + 30 - 16. \\
\lintertext{\indent Combine,}
3x &= 21. \\
\lintertext{\indent Divide by~$3$,}
x &= 7. \displaybreak[1] \\
\lintertext{\indent \textsc{Verification.}}
7(7 - 1) - 30 &= 4(7 - 4), \\
7 × 6 - 30 &= 4 x 3, \\
42 - 30 &= 12, \\
12 &= 12.
\end{DPalign*}

\Exercise{9.}

Find the number that $x$~stands for, if:
\begin{multicols}{2}
\Item{1.}  $3x = x + 8$.

\Item{2.}  $3x = 2x + 5$.

\Item{3.}  $3x + 4 = x + 10$.

\Item{4.}  $4x + 6 = x + 9$.

\Item{5.}  $7x - 19 = 5x + 7$.

\Item{6.}  $3(x - 2) = 2(x - 3)$.

\Item{7.}  $8x + 7 = 4x + 27$.

\Item{8.}  $3x + 10 = x + 20$.

\Item{9.}  $5(x - 2) = 3x + 4$.

\Item{10.}  $3(x - 2) = 2(x - 1)$.

\Item{11.}  $2x + 3 = 16 - (2x - 3)$.

\Item{12.}  $19x - 3 = 2(7 + x)$.

\Item{13.}  $7x - 70 = 5x - 20$.

\Item{14.}  $2x - 22 = 108 - 2x$.

\Item{15.}  $2(x + 5) + 5(x - 4) = 32$.

\Item{16.}  $2(3x - 25) = 10$.

\Item{17.}  $33x - 70 = 3x + 20$.

\Item{18.}  $4(1 + x) + 3(2 + x) = 17$.

\Item{19.}  $8x - (x + 2) = 47$.

\Item{20.}  $3(x - 2) = 50 - (2x - 9)$.
\end{multicols}
%% -----File: 031.png---Folio 25-------

\Item{21.}  $2x - (3 + 4x - 3x + 5) = 4$.

\Item{22.}  $5(2 - x) + 7x - 21 = x + 3$.

\Item{23.}  $3(x - 2) + 2(x - 3) + (x - 4) = 3x + 5$.

\Item{24.}  $x + 1 + x + 2 + x + 4 = 2x + 12$.

\Item{25.}  $(2x - 5) - (x - 4) + (x - 3) = x - 4$.

\Item{26.}  $4 - 5x - (1 - 8x) = 63 - x$.

\Item{27.}  $3x - (x + 10) - (x - 3) = 14 - x$.

\Item{28.}  $x^{2} - 2x - 3 = x^{2} - 3x + 1$.

\Item{29.}  $(x^{2} - 9) - (x^{2} - 16) + x = 10$.

\Item{30.}  $x^{2} + 8x - (x^{2} - x - 2) = 5(x + 3) + 3$.

\Item{31.}  $x^{2} + x - 2 + x^{2} + 2x - 3 = 2x^{2} - 7x - 1$.

\Item{32.}  $10x - (x - 5) = 2x + 47$.

\Item{33.}  $7x - 5 - (6 - 8x) + 2 = 3x - 7 + 106$.

\Item{34.}  $6x + 3 - (3x + 2) = (2x - 1) + 9$.

\Item{35.}  $3(x + 10) + 4(x + 20) + 5x - 170 = 15 - 3x$.

\Item{36.}  $20 - x + 4(x - 1) - (x - 2) = 30$.

\Item{37.}  $5x + 3 - (2x - 2) + (1 - x) = 6(9 - x)$.


\Paragraph{55.  Statement and Solution of Problems.}  The difficulties
which the beginner usually meets in stating problems will
be quickly overcome if he will observe the following directions:

Study the problem until you clearly understand its meaning
and just what is required to be found.

Remember that $x$~must not be put for money, length,
time, weight,~etc., but for the \textbf{required number} of \emph{specified
units} of money, length, time, weight,~etc.

Express each statement carefully in algebraic language,
and write out in full just what each expression stands for.
%% -----File: 032.png---Folio 26-------

Do not attempt to form the equation until all the statements
are made in symbols.

We will illustrate by examples:

\Item{1.}  John has three times as many oranges as James, and
they together have~$32$.  How many has each?
\begin{Soln}
\begin{DPgather*}
\lintertext{\indent Let}
\text{$x$~stand for the \emph{number} of oranges James has;} \\
\lintertext{then}
\text{$3x$~is the number of oranges John has;} \\
\lintertext{and}
\text{$x + 3x$ is the number of oranges they together have.}
\end{DPgather*}

But $32$~is the number of oranges they together have.
\begin{DPalign*}
\therefore x + 3x &= 32; \\
\lintertext{or,}
4x &= 32, \\
\lintertext{and}
x &= 8. \\
\lintertext{\indent Since $x = 8$,}
3x &= 24.
\end{DPalign*}
\end{Soln}

Therefore James has $8$~oranges, and John has $24$~oranges.

\begin{Remark}[Note.] Beginners in stating the preceding problem generally write:
\[
\text{Let $x = {}$\emph{what} James had.}
\]

Now, we know \emph{what} James had.  He had oranges, and we are to
discover simply the \emph{number} of oranges he had.
\end{Remark}

\Item{2.}  James and John together have~\$$24$, and James has
\$$8$~more than John.  How many dollars has each?
\begin{Soln}
\begin{DPgather*}
\lintertext{\indent Let}
\text{$x$~stand for the number of dollars John has;} \\
\lintertext{then}
\text{$x + 8$ is the number of dollars James has;} \\
\lintertext{and}
\text{$x + (x + 8)$ is the number of dollars they both have.}
\end{DPgather*}

But $24$~is the number of dollars they both have.
\[
\therefore x + (x + 8) = 24.
\]

Removing the parenthesis,
\[
x + x + 8 = 24.
\]
\begin{DPalign*}
\therefore 2x &= 16. \\
\lintertext{\indent Dividing by~$2$,}
x &= 8. \\
\lintertext{\indent Since $x = 8$,}
x + 8 &= 16.
\end{DPalign*}
\end{Soln}

Therefore John has~\$$8$, and James has~\$$16$.
%% -----File: 033.png---Folio 27-------

\begin{Remark}[Note\Add{.}] The beginner must avoid the mistake of writing
\[
\text{Let $x = {}$John's money\Add{.}}
\]

We are required to find the \emph{number} of dollars John has, and therefore
$x$~must represent this required number.
\end{Remark}

\Item{3.}  The sum of two numbers is~$18$, and three times the
greater number exceeds four times the less by~$5$.  Find the
numbers.

\begin{Soln}
Let $x = {}$the greater number.

Then, since $18$~is the sum and $x$~is one of the numbers, the other
number must be the sum minus~$x$.  Hence
\[
18 - x = \text{the smaller number}\Add{.}
\]

Now, three times the greater number is~$3x$, and four times the less
number is~$4(18 - x)$\Add{.}
\begin{DPalign*}
\lintertext{\indent Hence,}
&3x - 4(18 - x) = \text{the excess}\Add{.} \\
\lintertext{\indent But}
&5 = \text{the excess}\Add{,} \\
\therefore\ &3x - 4(18 - x) = 5 \\
\therefore\ &3x - (72 - 4x) = 5, \\
\lintertext{or}
&3x - 72 + 4x = 5. \\
\therefore\ &7x = 77, \\
\lintertext{and}
&\Z x = 11\Add{.}
\end{DPalign*}
\end{Soln}

Therefore the numbers are $11$~and~$7$.

\Exercise{10.}

\Item{1.}  If a number is multiplied by~$9$, the product is~$270$.
Find the number.

\Item{2.}  If the sum of the ages of a father and son is $60$~years,
and the father is $5$~times as old as the son, what is the
age of each?

\Item{3.}  The sum of two numbers is~$91$, and the greater is $6$~times
the less.  Find the numbers.
%% -----File: 034.png---Folio 28-------

\Item{4.}  A tree $90$~feet high was broken so that the part
broken off was $8$~times the length of the part left standing.
Find the length of each part.

\Item{5.}  The difference of two numbers is~$7$, and their sum is~$53$.
Find the numbers.

\Item{6.}  The difference of two numbers is~$12$, and their sum is~$84$.
Find the numbers.

\Item{7.}  Divide $35$ into two parts so that one part shall be
greater by~$5$ than the other part.

\Item{8.}  Three times a given number is equal to the number
increased by~$40$.  Find the number.

\Item{9.}  Three times a given number diminished by~$24$ is
equal to the given number.  Find the number.

\Item{10.}  One number is $4$~times another, and their difference
is~$30$.  Find the numbers.

\Item{11.}  The sum of two numbers is~$36$, and one of them
exceeds twice the other by~$6$.  Find the numbers.

\begin{Remark}[Hint.]  Let $x$~equal the greater number: then $36 - x$ will equal the
smaller.
\end{Remark}

\Item{12.}  The sum of two numbers is~$40$, and $5$~times the
smaller exceeds $2$~times the greater by~$25$.  Find the
numbers.

\Item{13.}  The number $30$ is divided into two parts such that
$4$~times the greater part exceeds $5$~times the smaller part
by~$30$.  Find the parts.

\Item{14.}  The sum of two numbers is~$27$, and twice the greater
number increased by $3$~times the less is~$61$.  Find the
numbers.

\Item{15.}  The sum of two numbers is~$32$, and five times the
smaller is $3$~times the greater number.  Find the numbers.
%% -----File: 035.png---Folio 29-------

\Exercise{11.}

\Item{1.}  A farmer sold a horse and a cow for~\$$210$.  He
sold the horse for four times as much as the cow.  How
much did he get for each?

\Item{2.}  Three times the excess of a certain number over~$6$
is equal to the number plus~$144$.  Find the number.

\Item{3.}  Thirty-one times a certain number is as much above~$40$
as nine times the number is below~$40$.  Find the number.

\Item{4.}  Two numbers differ by~$10$, and their sum is equal to
seven times their difference.  Find the numbers.

\Item{5.}  Find three consecutive numbers, $x$, $x + 1$, and~$x + 2$,
whose sum is~$78$.

\Item{6.}  Find five consecutive numbers whose sum is~$35$.

\Item{7.}  The sum of the ages of A~and~B is $40$~years, and $10$~years
hence A~will be twice as old as~B\@. Find their
present ages.

\Item{8.}  A father is four times as old as his son, and in $5$~years
he will be only three times as old.  Find their present ages.

\Item{9.}  One man is $60$~years old, and another man is $50$~years.
How many years ago was the first man twice as
old as the second?

\Item{10.}  A man $50$~years old has a son $10$~years old.  In
how many years will the father be three times as old as
the son?

\Item{11.}  A~has~\$$100$, and B~has~\$$20$.  How much must A
give B in order that they may each have the same sum?

\Item{12.}  A banker paid \$$63$ in $5$-dollar bills and $2$-dollar
bills.  He paid just as many $5$-dollar bills as $2$-dollar bills.
How many bills of each kind did he pay?
%% -----File: 036.png---Folio 30-------

\Exercise{12.}

\Item{1.}  In a  company of $90$~persons, composed of men,
women, and children, there are three times as many children
as men, and twice as many women as men.  How many
are there of each?

\Item{2.}  Find the number whose double exceeds~$70$ by as
much as the number itself is less than~$80$.

\Item{3.}  A farmer employed two men to build $112$~rods of
wall.  One of them built on the average $4$~rods a day, and
the other $3$~rods a day.  How many days did they work?

\Item{4.}  Two men travel in \emph{opposite} directions, one $30$~miles
a day, and the other $20$~miles a day.  In how many days
will they be $350$~miles apart?

\Item{5.}  Two men travel in the same direction, one $30$~miles
a day, and the other $20$~miles a day.  In how many days
will they be $350$~miles apart?

\Item{6.}  A man bought $3$~equal lots of hay for~\$$408$.  For
the first lot he gave \$$17$~a ton, for the second~\$$16$, for the
third~\$$18$.  How many tons did he buy in all?

\Item{7.}  A farmer sold a quantity of wood for~\$$84$, one half
of it at \$$3$~a cord, and the other half at \$$4$~a cord.  How
many cords did he sell?

\begin{Remark}[Hint.]
Let $2x$~equal the number of cords.
\end{Remark}

\Item{8.}  If $2x - 3$ stands for~$29$, for what number will
$4 + x$ stand?

\Item{9.}  At an election two opposing candidates received
together $2044$~votes, and one received $104$~more votes than
the other.  How many votes did each candidate receive?
%% -----File: 037.png---Folio 31-------

\Exercise{13.}

\Item{1.}  A man walks $4$~miles an  hour for $x$~hours, and
another man walks $3$~miles an hour for $x + 2$~hours.  If
they each walk the same distance, how many miles does
each walk?

\Item{2.}  A has twice as much money as~B; but if A gives~B
\$$30$, it will take twice as much as A has left to equal~B's.
How much money has each?

\Item{3.}  I have \$$12.75$ in two-dollar bills and twenty-five
cent pieces, and I have twice as many bills as twenty-five
cent pieces.  How many have I of each?

\Item{4.}  I have in mind a certain number.  If this number
is diminished by~$8$ and the remainder multiplied by~$8$, the
result is the same as if the number was diminished by~$6$
and the remainder multiplied by~$6$.  What is the number?

\Item{5.}  I have five times as many half-dollars as quarters, and
the half-dollars and quarters amount to~\$$11$.  How many
of each have I?

\Item{6.}  A man pays a debt of~\$$91$ with ten-dollar bills and
one-dollar bills, paying three times as many one-dollar bills
as ten-dollar bills.  How many bills of each kind does he
pay?

\Item{7.}  A father is four times as old as his son, but $4$~years
hence he will be only three times as old as his son.  How
old is each?

\Item{8.}  A workman was employed for $24$~days.  For every
day he worked he was to receive~\$$1.50$, and for every day
he was idle he was to pay \$$0.50$ for his board.  At the end
of the time he received~\$$28$.  How many days did he
work?
%% -----File: 038.png---Folio 32-------

\Exercise{14.}

\Item{1.}  A boy has $4$~hours at his disposal.  How far can he
ride into the country at the rate of $9$~miles an hour and
walk back at the rate of $3$~miles an hour, if he returns just
on time?

\begin{Remark}[Hint.]  Let $x = {}$the number of hours he rides.

Then $4 - x = {}$the number he walks.
\end{Remark}

\Item{2.}  A has~\$$180$, and B has~\$$80$.  How much must A
give B in order that six times B's money shall be equal to
$7$~times~A's?

\Item{3.}  A grocer has two kinds of tea, one kind worth $45$~cents
a pound, and the other worth $65$~cents a pound.
How many pounds of each kind must he take to make $80$~pounds,
worth $50$~cents a pound?

\Item{4.}  A tank holding $1200$~gallons has three pipes.  The
first lets in $8$~gallons a minute, the second $10$~gallons, and
the third $12$~gallons a minute.  In how many minutes will
the tank be filled?

\Item{5.}  The fore and hind wheels of a carriage are $10$~feet
and $12$~feet respectively in  circumference.  How many
feet will the carriage have passed over when the fore wheel
has made $250$~revolutions more than the hind wheel?

\Item{6.}  Divide a yard of tape into two parts so that one part
shall be $6$~inches longer than the other part.

\Item{7.}  A boy bought $7$~dozen oranges for~\$$1.50$.  For a part
he paid $20$~cents a dozen; and for the remainder, $25$~cents
a dozen.  How many dozen of each kind did he buy?

\Item{8.}  How can a bill of~\$$3.30$ be paid in quarters and ten-cent
pieces so as to pay three times as many ten-cent
pieces as quarters?
%% -----File: 039.png---Folio 33-------


\Chapter{III.}{Positive and Negative Numbers.}

\Paragraph{56. Quantities Opposite in Kind.}  If a person is engaged
in trade, his capital will be \emph{increased} by his \emph{gains}, and
\emph{diminished} by his \emph{losses}.

\emph{Increase} in temperature is measured by the number of
degrees the mercury rises in a thermometer, and \emph{decrease}
in temperature by the number of degrees the mercury \emph{falls}.

In considering any quantity whatever, a quantity that
\emph{increases} the quantity considered is called a \emph{positive quantity};
and a quantity that \emph{decreases} the quantity considered
is called a \emph{negative quantity}.

\Paragraph{57.  Positive and Negative Numbers.}  If from a  given
point, marked~$0$, we draw a straight line to the right, and
beginning from the \emph{zero} point lay off units of length on this
line, the successive repetitions of the unit will be expressed
by the \emph{natural series of numbers}, $1$,~$2$, $3$, $4$,~etc.  Thus:
\Graphic{1}

If we wish to \emph{add} $2$~to~$5$, we begin at~$5$, count $2$~units
\emph{forwards}, and arrive at~$7$, the sum required.  If we wish
to \emph{subtract} $2$~from~$5$, we begin at~$5$, count $2$~units \emph{backwards},
and arrive at~$3$, the difference required.  If we wish
to subtract $5$~from~$5$, we count $5$~units backwards, and arrive
at~$0$.  If we wish to subtract $5$~from~$2$, we cannot do it,
because when we have counted backwards from~$2$ as far as~$0$,
\emph{the natural series of numbers comes to an end}.
%% -----File: 040.png---Folio 34-------

In order to subtract a greater number from a smaller, it
is necessary to \emph{assume} a new series of numbers, beginning
at zero and extending to the left of zero.  The series to the
left of zero must proceed from zero by \emph{the repetitions of the
unit}, precisely like the natural series to the right of zero;
and the \emph{opposition} between the right-hand series and the
left-hand series must be clearly marked.  This opposition
is indicated by calling every number in the right-hand
series a \Defn{positive number}, and prefixing to it, when written,
the sign~$+$; and by calling every number in the left-hand
series a \Defn{negative number}, and prefixing to it the sign~$-$.
The two series of numbers may be called the \Defn{algebraic series
of numbers}, and written thus:
\Graphic{2}

If, now, we wish to subtract $7$~from~$4$, we begin at~$4$ in
the positive series, count $7$~units in the \emph{negative direction}
(to the left), and arrive at~$-3$ in the negative series; that
is, $4 - 7 = -3$.

The result obtained by subtracting a greater number from
a less, when both are positive, is \emph{always a negative number}.

In general, if $a$~and~$b$ represent any two numbers of the
positive series, the expression $a - b$ will be a positive number
when $a$~is greater than~$b$; will be zero when $a$~is equal
to~$b$; will be a negative number when $a$~is less than~$b$.

In counting from left to right in the algebraic series, numbers
\emph{increase} in magnitude; in counting from right to left,
numbers \emph{decrease} in magnitude.  Thus $-3$,~$-1$, $0$, $+2$,~$+4$,
are arranged in \emph{ascending} order of magnitude.

\Paragraph{58.}  Every algebraic number, as $+4$~or~$-4$, consists of a
\emph{sign} $+$~or~$-$ and the \emph{absolute value} of the number.  The
sign shows whether the number belongs to the positive or
%% -----File: 041.png---Folio 35-------
negative series of numbers; the absolute value shows the
place the number has in the positive or negative series.

When no sign stands before a number, the sign~$+$ is
always understood.  Thus $4$~means the same as~$+4$, $a$~means
the same as~$+a$.  But \emph{the sign~$-$ is never omitted}.

\Paragraph{59.}  Two algebraic numbers which have, one the sign~$+$,
and the other the sign~$-$, are said to have \emph{unlike signs}.

Two algebraic numbers which have the same absolute
values, but unlike signs, always cancel each other when
combined.  Thus $+4 - 4 = 0$; $+a - a = 0$.

\Paragraph{60.  Double Meanings of the Signs $+$~and~$-$.}  The use of
the signs $+$~and~$-$ to indicate addition and subtraction
must be carefully distinguished from the use of the signs $+$~and~$-$
to indicate in which series, the positive or the negative,
a given number belongs.  In the first sense they are
signs of \emph{operations}, and are common to Arithmetic and
Algebra; in the second sense they are signs of \emph{opposition},
and are employed in Algebra alone.

\begin{Remark}[Note.] In Arithmetic, if the things counted are \emph{whole units}, the
numbers which count them are called \Defn{whole numbers}, \Defn{integral numbers},
or \Defn{integers}, where the adjective is transferred from the things counted
to the numbers which count them.  But if the things counted are
only \emph{parts of units}, the numbers which count them are called \Defn{fractional
numbers}, or simply \Defn{fractions}, where again the adjective is transferred
from the things counted to the numbers which count them.

Likewise in Algebra, if the units counted are \emph{negative}, the numbers
which count them are called \Defn{negative numbers}, where the adjective
which defines the nature of the units counted is transferred to the
numbers that count them.

A whole number means a number of whole units, a fractional number
means a number of parts of units, and a negative number means
a number of negative units.
\end{Remark}
%% -----File: 042.png---Folio 36-------

\Paragraph{61.  Addition and Subtraction of Algebraic Numbers.}  An
algebraic number which is to be added or subtracted is
often \DPtypo{inclosed}{enclosed} in a parenthesis, in order that the signs $+$~and~$-$,
which are used to distinguish positive and negative
numbers, may not be confounded with the $+$~and~$-$ signs
that denote the operations of addition and subtraction.
Thus $+4 + (-3)$ expresses the sum, and $+4 - (-3)$ expresses
the difference, of the numbers $+4$~and~$-3$.

\Paragraph{62.  Addition.}  In order to add two algebraic numbers, we
begin at the place in the series which the first number occupies,
and count, \emph{in the direction indicated by the sign of the
second number}, as many units as there are in the absolute
value of the second number.
\Graphic{3}

Thus the sum of $+4 + (+3)$ is found by counting from
$+4$ three units in \emph{the positive direction}; that is, to the
right, and is, therefore,~$+7$.

The sum of $+4 + (-3)$ is found by counting from $+4$
three units in \emph{the negative direction}; that is, to the left, and
is, therefore,~$+1$.

The sum of $-4 + (+3)$ is found by counting from $-4$
three units in the positive direction, and is, therefore,~$-1$.

The sum of $-4 + (-3)$ is found by counting from $-4$
three units in the negative direction, and is, therefore,~$-7$.

\Paragraph{63.  Subtraction.}  In order to subtract one algebraic number
from another, we begin at the place in the series which
the minuend occupies, and count, \emph{in the direction opposite to
that indicated by the sign of the subtrahend}, as many units
as there are in the absolute value of the subtrahend.

Thus the result of subtracting $+3$~from~$+4$ is found by
%% -----File: 043.png---Folio 37-------
counting from $+4$ three units in the \emph{negative direction};
that is, in the direction \emph{opposite to that indicated by the sign~$+$
before~$3$}, and is, therefore,~$+1$.

The result of subtracting $-3$~from~$+4$ is found by counting
from $+4$ three units in the \emph{positive direction}, and is,
therefore,~$+7$.

The result of subtracting $+3$~from~$-4$ is found by counting
from $-4$ three units in the \emph{negative direction}, and is,
therefore,~$-7$.

The result of subtracting $-3$~from~$-4$ is found by counting
from $-4$ three units in the \emph{positive direction}, and is,
therefore,~$-1$.

\Paragraph{64.}  Collecting the results obtained in addition and subtraction,
we have:
\[
\begin{array}{c>{\quad}c}
\textsc{Addition.}        & \textsc{Subtraction.} \\
 +4 + (-3) = +4 - 3 = +1. &  +4 - (+3) = +4 - 3 = +1. \\
 +4 + (+3) = +4 + 3 = +7. &  +4 - (-3) = +4 + 3 = +7. \\
- 4 + (-3) = -4 - 3 = -7. &  -4 - (+3) = -4 - 3 = -7. \\
- 4 + (+3) = -4 + 3 = -1. &  -4 - (-3) = -4 + 3 = -1. \\
\end{array}
\]

\Paragraph{65.}  From these four cases of addition, therefore,

% [**** TN: Book uses commas elsewhere]
\Dictum{To Add Algebraic Numbers}\DPtypo{:}{,}
\begin{Theorem}[I.]  If the numbers have like signs, find the sum of their
absolute values, and prefix the common sign to the result.
\end{Theorem}

\begin{Theorem}[II.]  If the numbers have unlike signs, find the difference
of their absolute values, and prefix the sign of the greater
number to the result.
\end{Theorem}

\begin{Theorem}[III.]  If there are more than two numbers, find the sum
of the positive numbers and the sum of the negative numbers,
%% -----File: 044.png---Folio 38-------
take the difference between the absolute values of these two
sums, and prefix the sign of the greater sum to the result.
\end{Theorem}

\begin{Remark}[Note.] Since the order in which numbers are added is immaterial,
we may add any two of the numbers, and then this sum to
any third number, and so on.
\end{Remark}

\Paragraph{66.}  The result is generally called the \Defn{algebraic sum}, in
distinction from the arithmetical sum; that is, the sum of
the absolute values of the numbers.

\Paragraph{67.}  From the four cases of subtraction in §~64, we see
that \textit{subtracting a positive number is equivalent to adding
an equal negative number, and subtracting a negative number
is equivalent to adding an equal positive number}.


\Dictum{To Subtract One Algebraic Number from Another},
\begin{Theorem}
Change the sign of the subtrahend, and add the subtrahend
to the minuend.
\end{Theorem}

\Paragraph{68.  Examples.}

\Item{1.}  Find the sum of $3a$, $2a$, $a$, $5a$, $7a$.

The sum of the coefficients is $3 + 2 + 1 + 5 + 7 = 18$.

Hence the sum of the numbers is~$18a$.

\Item{2.}  Find the sum of $-5c$, $-c$, $-3c$, $-4c$, $-2c$.

The sum of the coefficients is $-5 - 1 - 3 - 4 - 2 = -15$.

Hence the sum of the numbers is~$-15c$.

\Item{3.}  Find the sum of $8x$, $-9x$, $-x$, $3x$, $4x$, $-12x$, $x$.

The sum of the positive coefficients is $8 + 3 + 4 + 1 = 16$.

The sum of the negative coefficients is $-9 - 1 - 12 = -22$.

The difference between $16$~and~$22$ is~$6$, and the sign of
the greater is negative.

Hence the required sum is~$-6x$.
%% -----File: 045.png---Folio 39-------

\PrintBreak
\Exercise{15.}

Find the sum of:
\begin{multicols}{2}
\Item{1.}  $5c$, $23c$, $c$, $11c$.

\Item{2.}  $4a$, $3a$, $7a$, $10a$.

\Item{3.}  $7x$, $12x$, $11x$, $9x$.

\Item{4.}  $6y$, $8y$, $2y$, $35y$.

\Item{5.}  $-3a$, $-5a$, $-18a$.

\Item{6.}  $-5x$, $-6x$, $-18x$, $-11x$.

\Item{7.}  $-3b$, $-b$, $-9b$, $-4b$.

\Item{8.}  $-z$, $-2z$, $-10z$, $-53z$.

\Item{9.}  $-11m$, $-3m$, $-5m$, $-m$.

\Item{10.}  $5d$, $-d$, $-4d$, $2d$.
\end{multicols}

\Item{11.}  $13n$, $13n$, $-11n$, $-6n$, $-9n$, $n$, $2n$, $-3n$.

\Item{12.}  $5g$, $-3g$, $-6g$, $-4g$, $20g$, $-5g$, $-11g$, $-14g$.

\Item{13.}  $-9a^{2}$, $5a^{2}$, $6a^{2}$, $a^{2}$, $2a^{2}$, $-a^{2}$, $-3a^{2}$.

\Item{14.}  $3x^{3}$, $-2x^{3}$, $-5x^{3}$, $-7x^{3}$, $-x^{3}$, $2x^{3}$, $-10x^{3}$, $-x^{3}$.

\Item{15.}  $4a^{2}b^{2}$, $-a^{2}b^{2}$, $-6a^{2}b^{2}$, $4a^{2}b^{2}$, $-2a^{2}b^{2}$, $a^{2}b^{2}$.

\Item{16.}  $6mn$, $-5mn$, $mn$, $-3mn$, $4mn$.

\Item{17.}  $3xyz$, $-2xyz$, $5xyz$, $-7xyz$, $xyz$.

\Item{18.}  $5a^{3}b^{3}c^{3}$, $-7a^{3}b^{3}c^{3}$, $-3a^{3}b^{3}c^{3}$, $2a^{3}b^{3}c^{3}$.

\Item{19.}  $11abcd$, $-10abcd$, $-9abcd$, $-abcd$.

\Item{20.}  Subtract $-a$ from $-b$, and find the value of the
result if $a = -4$, $b = -5$.

When $a = 4$, $b = -2$, $c = -3$, find the difference in
the values of:

\Item{21.}  $a - b + c$ and $-a + b + c$.

\Item{22.}  $a + (-b) + c$ and $a - (-b) + c$.

\Item{23.}  $-a - (-b) + c$ and $-(-a) + (-b) - c$.

\Item{24.}  $a - b + (-c)$ and $a - (-b) - (-c)$.
%% -----File: 046.png---Folio 40-------


\Section{MULTIPLICATION AND DIVISION OF ALGEBRAIC
NUMBERS}

\Paragraph{69.  Multiplication.}  Multiplication is generally defined
in Arithmetic as the process of finding the result when one
number (the multiplicand) is taken as many times as there
are units in another number (the multiplier).  This definition
fails when the \emph{multiplier is a fraction}.  In multiplying
by a fraction, we divide the multiplicand into as many
equal parts as there are units in the denominator, and take
as many of these parts as there are units in the numerator.

If, for example, we multiply $6$~by~$\frac{2}{3}$, we divide $6$ into
\emph{three} equal parts and take \emph{two} of these parts, obtaining $4$
for the product.  The multiplier,~$\frac{2}{3}$, is~$\frac{2}{3}$ of~$1$, and the
product,~$4$, is~$\frac{2}{3}$ of~$6$, in other words, \emph{the product is obtained
from the multiplicand precisely as the multiplier is obtained
from~$1$}.

This statement is also true when the multiplier is a whole
number.  Thus in $5 × 7 = 35$, the multiplier,~$5$, is equal to
\[
1 + 1 + 1 + 1 + 1,
\]
and the product,~$35$, is equal to
\[
7 + 7 + 7 + 7 + 7.
\]

\Paragraph{70.}  \Dictum{Multiplication may be defined}, therefore,

As the operation of finding from two given numbers,
called \emph{multiplicand} and \emph{multiplier}, a third number called
\emph{product}, which is \emph{formed from the multiplicand as the multiplier
is formed from unity}.

\Paragraph{71.}  According to this definition of multiplication,
\begin{DPalign*}[m]
\lintertext{since}
+3 &= + 1 + 1 + 1, \\
3 × (+8) &= +8 + 8 + 8
\Tag{(1)} \\
&= +24, \displaybreak[1] \\
%% -----File: 047.png---Folio 41-------
\lintertext{and}
3 × (-8) &= -8 - 8 - 8
\Tag{(2)} \\
&= -24. \displaybreak[1] \\
\lintertext{\indent Again, since}
-3 &= -1 - 1 - 1; \\
(-3) × 8 &= -8 - 8 - 8
\Tag{(3)} \\
&= -24, \displaybreak[1] \\
\lintertext{and}
(-3) × (-8) &= -(-8) - (-8) - (-8)
\Tag{(4)} \\
&= +8 + 8 + 8 \\
&= +24.
\end{DPalign*}

\Paragraph{72.}  From these four cases it follows that in finding
the product of two algebraic numbers, if the two numbers
have \emph{like} signs, the product will have the \emph{plus} sign, and if
\emph{unlike} signs, the product will have the \emph{minus} sign.

Hence the \Defn{Law of Signs in Multiplication} is:
\begin{Theorem}
Like signs give~$+$, and unlike signs give~$-$.
\end{Theorem}

If $a$~and~$b$ stand for any two numbers, we have
\begin{align*}
(+a) × (+b) &= +ab, \\
(+a) × (-b) &= -ab, \\
(-a) × (+b) &= -ab, \\
(-a) × (-b) &= +ab.
\end{align*}

\Paragraph{73.  The Index Law in Multiplication.}
\begin{DPalign*}
\lintertext{\indent Since}
a^{2} &= aa, \quad\text{and}\quad a^{3} = aaa, \\
a^{2} × a^{3} &= aa × aaa = aaaaa = a^{5} = a^{2 + 3}; \\
a^{4} × a &= aaaa × a = aaaaa = a^{5} = a^{4 + 1}.
\end{DPalign*}

If $a$~stands for any number, and $m$~and~$n$ for any integers,
\[
a^{m} × a^{n} = a^{m + n}. \EqText{Hence,}
\]
\begin{Theorem}
The index of the product of two powers of the same number
is equal to the sum of the indices of the factors.
\end{Theorem}
%% -----File: 048.png---Folio 42-------

\Paragraph{74.  Examples.}

\Item{1.}  Find the product of $6a^{2}b^{2}$ and $7ab^{2}c^{3}$.

Since the order of the factors is immaterial,
\begin{align*}
6a^{2}b^{3} × 7ab^{2}c^{3}
  &= 6 × 7 × a^{2} × a × b^{3} × b^{2} × c^{3} \\
  &= 42a^{3}b^{5}c^{3}.
\end{align*}

\Item{2.}  Find the product of $-3ab$ and $7ab^{3}$.
\begin{align*}
-3ab × 7ab^{3}
  &= -3 × 7 × a × a × b × b^{3} \\
  &= -21a^{2}b^{4}.
\end{align*}

\Paragraph{75.  To Find the Product of Simple Expressions}, therefore,
\begin{Theorem}
Take the product of the coefficients and the sum of the
indices of the like letters.
\end{Theorem}

\Exercise{16.}

Find the product of:
\begin{multicols}{2}
\Item{1.}  $5a^{2}$ and $6a^{3}$.

\Item{2.}  $8ab$ and $5a^{3}b^{2}$.

\Item{3.}  $9xy$ and $7xy$.

\Item{4.}  $2a^{2}b$ and $a^{3}b^{4}c^{2}$.

\Item{5.}  $3a^{3}b^{7}c^{8}$ and $3a^{4}b^{2}c$.

\Item{6.}  $2a$ and $-5a$.

\Item{7.}  $-3a$ and $-4b$.

\Item{8.}  $-ab$ and $a^{3}b^{2}$.

\Item{9.}  $-2ab^{4}$ and $-5a^{4}bc$.

\Item{10.}  $-2x^{6}y^{3}z$ and $-6xy^{2}z$.
\end{multicols}

\Item{11.}  $3a^{2}b$, $-5ab^{2}$, and $-7a^{4}b^{2}$.

\Item{12.}  $2a^{2}bc^{3}$, $-3a^{3}b^{2}c$, and $-ab^{2}c^{3}$.

\Item{13.}  $2b^{2}c^{2}x^{2}$, $2a^{2}b^{2}c^{3}$, and $-3a^{3}bx^{3}$.

\Item{14.}  $2a^{3}b^{2}c$, $-3a^{2}b^{3}c$, and $-4a^{2}bc^{3}$.

\Item{15.}  $7am^{2}x^{3}$, $3a^{4}m^{2}x^{3}$, and $-2amx$.

\Item{16.}  $-3x^{2}y^{2}z^{2}$, $2x^{2}yz^{3}$, and $-5x^{4}yz$.
%% -----File: 049.png---Folio 43-------

If $a = -2$, $b = 3$, and $c = -1$, find the value of:
\begin{multicols}{2}
\Item{17.}  $2ab^{2} - 3bc^{2} + c$.

\Item{18.}  $4a^{2} - 2b^{2} - c^{2}$.

\Item{19.}  $5a + 2b - 4c^{4}$.

\Item{20.}  $2a^{3} - 3b + 8c^{2}$.

\Item{21.}  $-a + 3b - 2c^{2}$.

\Item{22.}  $-a^{3} - 2b - 10c$.

\Item{23.}  $3a^{3} - 3b^{3} - 3c^{3}$.

\Item{24.}  $2ab^{2} - 3bc^{2} + 2ac$.

\Item{25.}  $3abc + 5a^{2}b^{2} - 2a^{2}b$.

\Item{26.}  $ab^{2}c^{2} + 2abc^{2} + a^{2}b^{2}c^{2}$.

\Item{27.}  $2a^{2}bc + 3abc + a^{2}b^{2}c^{2}$.

\Item{28.}  $6a^{2} + 8a^{2}b^{2} - 5a^{2}bc$.
\end{multicols}

\Paragraph{76.  Division.}  To divide $48$~by~$8$ is to find the number
of times it is necessary to take $8$ to make~$48$.  Here the
\emph{product} and \emph{one factor} are given, and \emph{the other factor} is
required.  We may, therefore, take for the definition of
division\Add{:}

The operation by which when \emph{the product} and \emph{one factor}
are given, \emph{the other factor is found}.

With reference to this operation the product is called
the \Defn{dividend}, the given factor the \Defn{divisor}, and the required
factor the \Defn{quotient}.

\Paragraph{77.  Law of Signs in Division.}
\begin{alignat*}{2}
&\text{Since } (+a) × (+b) = +ab,\quad && \therefore +ab ÷ (+a) = +b. \\
&\text{Since } (+a) × (-b) = -ab,      && \therefore -ab ÷ (+a) = -b. \\
&\text{Since } (-a) × (+b) = -ab,      && \therefore -ab ÷ (-a) = +b. \\
&\text{Since } (-a) × (-b) = +ab,      && \therefore +ab ÷ (-a) = -b.
\end{alignat*}

That is, if the dividend and divisor have like signs, the
quotient has the plus sign; and if they have unlike signs,
the quotient has the minus sign.  Hence, in division,
\begin{Theorem}
Like signs give~$+$, and unlike signs give~$-$.
\end{Theorem}
%% -----File: 050.png---Folio 44-------

\Paragraph{78.  Index Law in Division.}

The dividend contains all the factors of the divisor and
of the quotient, and therefore the quotient contains the
factors of the dividend that are not found in the divisor.

Thus, $\dfrac{abc}{bc} = a$, $\dfrac{aabx}{ab} = ax$, $\dfrac{124abc}{-4ab} = -31c$.

Divide $a^{5}$~by~$a^{2}$, $a^{6}$~by~$a^{4}$, $a^{4}$~by~$a$\DPtypo{,}{.}
\begin{alignat*}{4}
\frac{a^{5}}{a^{2}}
  &= \frac{aaaaa}{aa}    &&= aaa &&= a^{3} &&= a^{5-2}; \\
\frac{a^{6}}{a^{4}}
  &= \frac{aaaaaa}{aaaa} &&= aa  &&= a^{2} &&= a^{6-4}\DPtypo{,}{;} \\
\frac{a^{4}}{a}
  &= \frac{aaaa}{a}      &&= aaa &&= a^{3} &&= a^{4-1}\DPtypo{,}{.}
\end{alignat*}

If $m$~and~$n$ stand for any integers, and $m$~is greater than~$n$\Add{,}
\[
a^{m} - a^{n} = a^{m - n}\Add{.}
\]

\begin{Theorem}
The index of the quotient of two powers of the same letter
is equal to the index of the letter in the dividend diminished
by the index of the letter in the divisor.
\end{Theorem}

\Paragraph{79.  Examples.}

\Item{1.}  Divide $15xy$~by~$5x$\Add{.}
\[
\frac{15xy}{5x} = \frac{3 × 5xy}{5x} = 3y.
\]

Here we cancel the factors $5$~and~$x$, which are common
to the dividend and divisor\Add{.}

\Item{2.}  Divide $-21a^{2}b^{3}$~by~$3ab^{2}$\Add{.}
\[
\frac{-21a^{2}b^{3}}{3ab^{2}} = -7ab.
\]
%% -----File: 051.png---Folio 45-------

\Item{3.}  Divide $54a^{5}b^{3}c$~by~$-6ab^{2}c$.
\[
\frac{54a^{5}b^{3}c}{-6ab^{2}c} = -9a^{4}b.
\]

\Item{4.}  Divide $-45x^{4}y^{5}z^{7}$~by~$-15x^{4}y^{5}z^{5}$.
\[
\frac{-45x^{4}y^{5}z^{7}}{-15x^{4}y^{5}z^{5}} = 3z^{2}.
\]

\Item{5.}  Divide $-15a^{3}b^{2}c^{3}$~by~$-60a^{2}bc^{3}$.
\[
\frac{-15a^{3}b^{2}c^{3}}{-60a^{2}bc^{3}} = \frac{ab}{4}.
\]

\Exercise{17.}

Divide:
\begin{multicols}{2}
\Item{1.}  $x^{3}$~by~$x$.

\Item{2.}  $21x^{5}$~by~$7x^{3}$.

\Item{3.}  $35x^{2}$~by~$-5x^{2}$.

\Item{4.}  $-42x^{2}$~by~$6x^{2}$.

\Item{5.}  $-63x^{5}$~by~$-9x$.

\Item{6.}  $-72x^{3}$~by~$-8x^{2}$.

\Item{7.}  $-32a^{2}b^{2}$~by~$8ab^{2}$.

\Item{8.}  $-16x^{3}y^{3}$~by~$-4xy$.

\Item{9.}  $18x^{2}y$~by~$-2xy$.

\Item{10.}  $-25x^{4}y^{2}$~by~$-5x^{3}y^{2}$.

\Item{11.}  $-51x^{2}y^{3}$~by~$-17x^{2}y$.

\Item{12.}  $-28a^{4}b^{3}$~by~$7a^{3}b$.

\Item{13.}  $-36x^{2}y^{6}$~by~$-3xy^{2}$.

\Item{14.}  $-3x^{4}y^{6}$~by~$-5xy^{3}$.

\Item{15.}  $-12a^{2}b^{3}$~by~$8ab^{3}$.

\Item{16.}  $-abcd$~by~$ac$.

\Item{17.}  $-a^{2}b^{3}c^{4}d^{5}$~by~$-ab^{3}c^{3}d^{3}$.

\Item{18.}  $2x^{2}y^{2}z^{3}$~by~$-3xyz^{3}$.

\Item{19.}  $-5a^{5}b^{3}c^{7}$~by~$-a^{4}b^{2}c^{7}$.

\Item{20.}  $52a^{2}m^{3}n^{4}$~by~$13a^{2}m^{2}n^{3}$.

\Item{21.}  $13xy^{2}z^{4}$~by~$39xyz$.

\Item{22.}  $68xc^{2}d^{3}$~by~$-4xcd^{2}$.

\Item{23.}  $-8m^{5}n^{3}p^{2}$~by~$-4m^{5}np$.

\Item{24.}  $-6pqr^{3}$~by~$-2p^{2}qr$.

\Item{25.}  $26a^{2}g^{2}t^{5}$~by~$-2agt^{4}$.

\Item{26.}  $-a^{4}b^{2}c^{3}$~by~$-a^{5}b^{3}c^{4}$.

\Item{27.}  $-3x^{2}y^{2}z^{2}$~by~$-2x^{3}y^{4}z^{5}$.

\Item{28.}  $-6mnp$~by~$-3m^{2}n^{2}p^{2}$.

\Item{29.}  $-17a^{2}b^{3}c^{4}$~by~$51ab^{5}c^{4}$.

\Item{30.}  $-19mg^{2}t^{3}$~by~$57m2^gt^{4}$.
\end{multicols}
%% -----File: 052.png---Folio 46-------


\Chapter{IV.}{Addition and Subtraction.}

\Section{Integral Compound Expressions.}

\Paragraph{80.}  If an algebraic expression contains only \emph{integral
forms}, that is, contains \emph{no letter in the denominator of
any of its terms}, it is called an \Defn{integral expression}.

Thus, $x^{3} + 7cx^{2} - c^{3} - 5c^{2}x$, is an integral expression.

Integral and fractional expressions are so named on
account of the \emph{form of the expressions}, and with no reference
whatever to the \emph{numerical value} of the expressions
when definite numbers are put in place of the letters.

\Paragraph{81.  Addition of Integral Compound Expressions.}  The addition
of two algebraic expressions can be represented by
connecting the second expression with the first by the sign~$+$.
If there are no like terms in the two expressions, the
operation is \emph{algebraically complete} when the two expressions
are thus connected.

If, for example, it is required to add $m + n - p$ to
$a + b + c$, the result will be $a + b + c + (m + n - p)$; or,
removing the parenthesis (§~37), $a + b + c + m + n - p$.

\Paragraph{82.}  If there are like terms in the expressions, the like
terms can be \emph{collected};  that is, every set of like terms can
be replaced by a single term with a coefficient equal to
the algebraic sum of the coefficients of the like terms.
%% -----File: 053.png---Folio 47-------

\Item{1.}  Add $6x^{2} + 5x + 4$ to $x^{2} - 4x - 5$.
\begin{DPalign*}
\lintertext{\indent The sum}
&= x^{2} - 4x - 5 + (6x^{2} + 5x + 4) \\
&= x^{2} - 4x - 5 + 6x^{2} + 5x + 4
\rintertext{§~37} \\
&= x^{2} + 6x^{2} - 4x + 5x - 5 + 4 \\
&= 7x^{2} + x - 1.
\end{DPalign*}

This process is more conveniently represented by arranging
the terms in columns, so that like terms shall stand in
the same column, as follows:
\[
\begin{array}{r*{2}{cr}}
 x^{2} &-& 4x &-& 5 \\
6x^{2} &+& 5x &+& 4 \\
\hline
7x^{2} &+&  x &-& 1 \\
\end{array}
\]

The coefficient of~$x^{2}$ in the result will be $6 + 1$, or~$7$; the
coefficient of~$x$ will be $-4 + 5$, or~$1$; and the last term is
$-5 + 4$, or~$-1$.

\begin{Remark}[Note.] When the coefficient of a term is~$1$, it is not written, but
understood.
\end{Remark}

\Item{2.}  Add $2c^{3} - 5c^{2}d + 6cd^{2} + d^{3}$; $c^{3} + 6c^{2}d - 5cd^{2} - 2d^{3}$;
and $3c^{3} - c^{2}d - 7cd^{2} - 3d^{3}$.
\[
\begin{array}{r*{3}{cr}}
2c^{3} &-& 5c^{2}d &+& 6cd^{2} &+& d^{3} \\
 c^{3} &+& 6c^{2}d &-& 5cd^{2} &-& 2d^{3} \\
3c^{3} &-&  c^{2}d &-& 7cd^{2} &-& 3d^{3} \\
\hline
6c^{3} & &         &-& 6cd^{2} &-& 4d^{3} \\
\end{array}
\]

The coefficient of~$c^{3}$ in the result will be $2 + 1 + 3$, or~$6$;
the coefficient of~$c^{2}d$ will be $-5 + 6 - 1$, or~$0$; therefore
$c^{2}d$ will not appear in the result; the coefficient of~$cd^{2}$ will
be $6 - 5 - 7$, or~$-6$; and the coefficient of~$d^{3}$ will be
$1 - 2 - 3$, or~$-4$.
%% -----File: 054.png---Folio 48-------

\Exercise{18.}

Find the sum of:

\Item{1.}  $a^{2} - ab + b^{2}$; $a^{2} + ab + b^{2}$.

\Item{2.}  $3a^{2} + 5a-7$; $6a^{2} - 7a + 13$.

\Item{3.}  $x + 2y - 3z$; $-3x + y + 2z$; $2x - 3y + z$.

\Item{4.}  $3x + 2y - z$; $-x + 3y + 2z$; $2x - y + 3z$.

\Item{5.}  $-3a + 2b + c$; $a - 3b + 2c$; $2a + 3b - c$.

\Item{6.}  $-a + 3b + 4c$; $3a - b + 2c$; $2a + 2b - 2c$.

\Item{7.}  $4a^{2} + 3a + 5$; $-2a^{2} + 3a - 8$; $a^{2} - a + 1$.

\Item{8.}  $5ab + 6bc - 7ac$; $3ab - 9bc + 4ac$; $3bc + 6ac$.

\Item{9.}  $x^{3} + x^{2} + x$; $2x^{3} + 3x^{2} - 2x$; $3x^{3} - 4x^{2} + x$.

\Item{10.}  $3y^{2} - x^{2} - 3xy$; $5x^{2} + 6xy - 7y^{2}$; $x^{2} + 2y^{2}$.

\Item{11.}  $2a^{2} - 2ab + 3b^{2}$; $4b^{2} + 5ab - 2a^{2}$; $a^{2} - 3ab - 9b^{2}$.

\Item{12.}  $a^{3} - a^{2} + a - 1$; $a^{2} - 2a + 2$; $3a^{3} + 7a + 1$.

\Item{13.}  $2m^{3} - m^{2} - m$; $4m^{3} + 8m^{2} - 7$; $-3m^{3} + m + 9$.

\Item{14.}  $x^{3} - 3x + 6y$; $x^{2} + 2x - 5y$; $x^{3} - 3x^{2} + 5x$.

\Item{15.}  $6x^{3} - 5x + 1$; $x^{3} + 3x + 4$; $7x^{2} + 2x - 3$.

\Item{16.}  $a^{3} + 3a^{2}b - 3ab^{2}$; $-3a^{2}b - 6ab^{2} - b^{3}$; $3a^{2}b + 4ab^{2}$.

\Item{17.}  $a^{3} - 2a^{2}b - 2ab^{2}$; $a^{2}b - 3ab^{2} - b^{3}$; $3ab^{2} - 2a^{3} - b^{3}$.

\Item{18.}  $7x^{3} - 2x^{2}y + 9xy^{2} + 13y^{3}$; $5x^{2}y - 4xy^{2} - 2x^{3} - 3y^{3}$;
$y^{3} - x^{3} - 3x^{2}y - 5xy^{2}$; $2x^{2}y - 5y^{3} - 2x^{3} - xy^{2}$.

\Item{19.}  Show that $x + y + z = 0$, if $x = a - b - c$,
$y = 2b + 2c - 3a$, and $z = 2a - b - c$.

\Item{20.}  Show that $x + y = 3z$, if $x = 3a^{2} - 6a + 12$,
$y = 9a^{2} + 12a - 21$, and $z = 4a^{2} + 2a - 3$.
%% -----File: 055.png---Folio 49-------

\Paragraph{83.  Subtraction of Integral Compound Expressions.}  The
subtraction of one expression from another, if none of the
terms are alike, can be represented only by connecting the
subtrahend with the minuend by means of the sign~$-$.

If, for example, it is required to subtract $a + b + c$ from
$m + n - p$, the result will be represented by
\[
m + n - p - (a + b + c);
\]
or, removing the parenthesis (§~38),
\[
m + n - p - a - b - c.
\]

If, however, some of the terms in the two expressions are
alike, we can replace two like terms by a single term.

Thus, suppose it is required to subtract $a^{3} + 2a^{2} + 3a - 5$
from $2a^{3} - 3a^{2} + 2a - 1$; the result may be expressed as
follows:
\[
2a^{3} - 3a^{2} + 2a - 1 - (a^{3} + 2a^{2} + 3a - 5);
\]
or, removing the parenthesis (§~38),
\begin{align*}
&2a^{3} - 3a^{2} + 2a - 1 - a^{3} - 2a^{2} - 3a + 5 \\
&\quad= 2a^{3} - a^{3} - 3a^{2} - 2a^{2} + 2a - 3a - 1 + 5 \\
&\quad= a^{3} - 5a^{2} - a + 4.
\end{align*}

This process is more easily performed by writing the subtrahend
below the minuend, \emph{mentally} changing the sign of
each term in the subtrahend, and adding.
\[
\begin{array}{r*{3}{cr}}
2a^{3} &-& 3a^{2} &+& 2a &-& 1 \\
 a^{3} &+& 2a^{2} &+& 3a &-& 5 \\
\hline
 a^{3} &-& 5a^{2} &-&  a &+& 4 \\
\end{array}
\]

By changing the sign of each term in the subtrahend,
the coefficient of~$a^{3}$ will be $2 - 1$, or~$1$; the coefficient of~$a^{2}$
will be $-3 - 2$, or~$-5$; the coefficient of~$a$ will be~$2 - 3$,
or~$-1$; the last term will be $-1 + 5$, or~$4$.
%% -----File: 056.png---Folio 50-------

Again, suppose it is required to subtract $x^{5} - 2ax^{4} -
3a^{2} x^{3} + 4a^{3} x^{2}$ from $4a^{3} x^{2} - 2a^{2} x^{3} - 5ax^{4}$. Here terms
which are alike can be written in columns, as before:
\[
\begin{array}{r*{3}{cr}}
       &-& 5ax^{4} &-& 2a^{2} x^{3} &+& 4a^{3} x^{2} \\
 x^{5} &-& 2ax^{4} &-& 3a^{2} x^{3} &+& 4a^{3} x^{2} \\
\hline
-x^{5} &-& 3ax^{4} &+&  a^{2} x^{3} & &
\end{array}
\]

There is no term of~$x^{5}$ in the minuend, hence the coefficient
of~$x^{5}$ in the result will be~$0 - 1$, or~$-1$; the coefficient of~$ax^{4}$
will be $-5 + 2$, or~$-3$; the coefficient\DPtypo{,}{} of~$a^{2}x^{3}$ will be
$-2 + 3$, or~$+1$; the coefficient of~$a^{3}x^{2}$ will be $-4 + 4$, or~$0$,
and therefore the term~$a^{3}x^{2}$ will not appear in the result.

\Exercise{19.}

Subtract:

\Item{1.}  $a - 2b + 3c$ from $2a - 3b + 4c$.

\Item{2.}  $a - 3b - 5c$ from $3a - 5b + c$.

\Item{3.}  $2x - 4y + 6z$ from $4x - y - 2z$.

\Item{4.}  $5x - 11y - 3z$ from $6x - 7y + 2z$.

\Item{5.}  $ab - ac - bc + bd$ from $ab + ac + bc + bd$.

\Item{6.}  $3ab + 2ac - 3bc + bd$ from $5ab - ac + bc + bd$.

\Item{7.}  $2x^{3} - x^{2} - 5x + 3$ from $3x^{3} + 2x^{2} - 3x - 5$.

\Item{8.}  $7x^{2} - 5x + 1 - a$ from $x^{3} - x + 1 - a$.

\Item{9.}  $7b^{3} + 8c^{3} - 15abc$ from $9b^{3} + 3abc - 7c^{3}$.

\Item{10.}  $x^{4} + x - 5x^{3} + 5$ from $7 - 2x^{2} - 3x^{3} + x^{4}$.

\Item{11.}  $a^{3} + b^{3} + c^{3} - 3abc$ from $3abc + a^{3} - 2b^{3} - 3c^{3}$.

\Item{12.}  $2x^{4} - 5x^{2} + 7x - 3$ from $x^{4} + 2 - 2x^{3} - x^{2}$.

\Item{13.}  $1 - x^{5} - x + x^{4} - x^{3}$ from $x^{4} + 1 + x + x^{2}$.

\Item{14.}  $a^{3} - b^{3} + 3a^{2}b - 3ab^{2}$ from $a^{3} + b^{3} - a^{2}b - ab^{2}$.

\Item{15.}  $a^{2} b - ab^{2} - 3a^{3} b^{3} - b^{4}$ from $b^{4} - 5a^{3} b^{3} - 2ab^{2} + a^{2} b$.

\Item{16.}  $-x^{3} + 7x^{2} y - 2y^{3} + 3xy^{2}$ from $3x^{3} + 5y^{3} - xy^{2} + 4x^{2}y$.
%% -----File: 057.png---Folio 51-------

\Paragraph{84.  Parentheses or Brackets.} We have for positive numbers
(§§~37,~38):
\begin{alignat*}{2}
a + (b + c) &= a + b + c,\qquad & \therefore  a + b + c &= a + (b + c); \\
a + (b - c) &= a + b - c,       & \therefore  a + b - c &= a + (b - c); \\
a - (b + c) &= a - b - c,       & \therefore  a - b - c &= a - (b + c); \\
a - (b - c) &= a - b + c,       & \therefore  a - b + c &= a - (b - c).
\end{alignat*}

That is, a parenthesis preceded by~$+$ may be removed
\emph{without changing the sign of any term within the parenthesis};
and any number of terms may be enclosed within a parenthesis
preceded by the sign~$+$, \emph{without changing the sign
of any term}.

A parenthesis preceded by the sign~$-$ may be removed,
\emph{provided the sign of every term within the parenthesis is
changed}, namely, $+$~to~$-$, and $-$~to~$+$; and any number
of terms may be enclosed within a parenthesis preceded
by the sign~$-$, \emph{provided the sign of every term enclosed is
changed}.

The same laws hold for \emph{negative numbers}.

\Paragraph{85.}  Expressions may occur having a parenthesis within
a parenthesis.  In such cases parentheses of different shapes
are used, and the beginner when he meets with a branch
of a parenthesis~$($, or bracket~$[$, or brace~$\{$, must look carefully
for the other part, whatever may intervene; and all
that is included between the two parts of each parenthesis
must be treated as the sign before it directs, without regard
to other parentheses.  It is best to remove each parenthesis
in succession, \emph{beginning with the innermost}.
\begin{align*}
a - &\bigl\{b - [c - (d - e) + f]\bigr\} \\
  &= a - \bigl\{b - [c - d + e + f]\bigr\} \\
  &= a - \bigl\{b - c + d - e - f\bigr\} \\
  &= a - b + c - d + e + f.
\end{align*}
%% -----File: 058.png---Folio 52-------

\Exercise{20.}

Remove the brackets and collect the like terms:

\Item{1.}  $a - b - (b - c) - a + 2b$.

\Item{2.}  $x - [x - (a - b) + a - y]$.

\Item{3.}  $3x - \bigl\{2y - [-7c - 2x] + y\bigr\}$.

\Item{4.}  $5a - [7 - (2b + 5) - 2a]$.

\Item{5.}  $x - [2x + (3a - 2x) - 5a]$.

\Item{6.}  $x - [15y - (13z + 12x)]$.

\Item{7.}  $2a - b + [4c - (b + 2c)]$.

\Item{8.}  $5a - \bigl\{b + [3c - (2b - c)]\bigr\}$.

\Item{9.}  $7x - \bigl\{5y - [3z - (3x + z)]\bigr\}$.

\Item{10.}  $(a - b + c) - (b - a - c) + (a + b - 2c)$.

\Item{11.}  $3x - [-2y - (2y - 3x) + z] + [x - (y - 2z - x)]$.

\Item{12.}  $x - [2x + (x - 2y) + 2y] - 3x - \bigl\{4x - [(x + 2y) - y]\bigr\}$.

\Item{13.}  $x - [y + z - x - (x + y) - z] + (3 x - \Vinc{2y + z})$.

\begin{Remark}[Note.]
The expression $-\Vinc{2y + z}$ is equivalent to~$-(2y + z)$.
\end{Remark}

Consider \emph{all the factors} that precede $x$,~$y$, and~$z$, respectively,
as the \emph{coefficients} of these letters, and collect in
brackets the coefficients of each of these letters:

\Item{14.}  $ax + by + cz - ay + az - bx
= (a - b)x - (a - b)y + (a + c)z$.

\Item{15.}  $ax + az + by - cz - ay + cx$.

\Item{16.}  $2ax - 3ay - 4by + 5cx - 6bz - 7cz$.

\Item{17.}  $az - bmy + 3 cz - anx - cny + acx$.

\Item{18.}  $mnx - x - mny - y + mnz + z$.
%% -----File: 059.png---Folio 53-------


\Chapter{V.}{Multiplication and Division.}

\Section{Compound Integral Expressions.}

\Paragraph{86.  Multiplication.  Polynomials by Monomials.}

We have for positive numbers (§~39),
\begin{align*}
a(b + c) &= ab + ac, \\
a(b - c) &= ab - ac.
\end{align*}

The same law holds for negative numbers.

\Dictum{To multiply a polynomial by a monomial}, therefore,
\begin{Theorem}
Multiply each term of the polynomial by the monomial,
and add the partial products.
\end{Theorem}

\Item{1.}  Find the product of $ab + ac - bc$ and~$abc$.
\[
\begin{array}{rcr}
ab + ac - bc && \\
         abc && \\
\hline
a^{2}b^{2}c + a^{2}bc^{2} &-& ab^{2}c^{2}
\end{array}
\]

\begin{Remark}[Note.]
We multiply~$ab$, the first term of the multiplicand, by~$abc$,
and work to the right.
\end{Remark}

\Exercise{21.}

Find the product of:
\begin{multicols}{2}
\Item{1.}  $x + 7$ and $x$.

\Item{2.}  $2x - 3y$ and $4x$.

\Item{3.}  $2x - 3y$ and $7y$.

\Item{4.}  $x - 2a$ and $2a$.

\Item{5.}  $-x + 3b$ and $-b$.

\Item{6.}  $2a^{2} - 3ab$ and $-3a$.

\Item{7.}  $2x^{2} + 3xz$ and $5z$.

\Item{8.}  $a^{2} - 5ab$ and $5ab$.
%% -----File: 060.png---Folio 54-------

\Item{9.}  $x^{2} - 3 xy$ and $-y^{2}$.

\Item{10.}  $2 x^{3} - 3x^{2}$ and $2x^{2}$.

\Item{11.}  $x^{2} - 3y^{2}$ and $4y$.

\Item{12.}  $x^{2} - 3 y^{2}$ and $-x^{2}$.

\Item{13.}  $b^{3} - a^{2}b^{2}$ and $-a^{3}$.

\Item{14.}  $-a^{2}b^{2} - a^{3}$ and $-a^{2}$.

\Item{15.}  $2x^{3} - 3x^{2} + x$ and $2x^{2}$.

\Item{16.}  $a^{2} - 5ab - b^{2}$ and $5ab$.
\end{multicols}

\Item{17.}  $a^{3} + 2a^{2}b + 2ab^{2}$ and $a^{2}$.

\Item{18.}  $a^{3} + 2a^{2}b + 2ab^{2}$ and $b^{3}$.

\Item{19.}  $4x^{2} - 6xy - 9y^{2}$ and $2x$.

\Item{20.}  $-x^{2} - 2xy + y^{2}$ and $-y$.

\Item{21.}  $-a^{3} - a^{2}b^{2} - b^{3}$ and $-a^{2}$.

\Item{22.}  $-x^{2} + 2xy - y^{2}$ and $-y^{2}$.

\Item{23.}  $3 a^{2}b^{2} - 4 ab^{3} + a^{3}b$ and $5 a^{2}b^{2}$.

\Item{24.}  $-ax^{2} + 3axy^{2} - ay^{4}$ and $-3ay^{2}$.

\Item{25.}  $x^{12} - x^{10}y^{3} - x^{3}y^{10}$ and $x^{3}y^{2}$.

\Item{26.}  $-2x^{3} + 3x^{2}y^{2} - 2xy^{3}$ and $-2x^{2}y^{3}$.

\Item{27.}  $a^{3}x^{2}y^{5} - a^{2}xy^{4} - ay^{3}$ and $a^{7}x^{3}y^{5}$.

\Item{28.}  $3a^{2}b^{2} - 2ab^{3} + 5a^{3}b$ and $5a^{2}b^{3}$.

\Paragraph{87.  Multiplication.  Polynomials by Polynomials.}

If we have $m + n +p$ to be multiplied by $a + b + c$, we
may substitute~$M$ for the multiplier $a + b + c$.  Then
\[
M(m + n + p) = Mm + Mn + Mp.
\]

If now we substitute $a + b + c$ for~$M$, we shall have
\begin{align*}
&(a + b + c) m + (a + b + c) n + (a + b + c) p \\
&= am + bm + cm + an + bn + cn + ap + bp + cp\DPtypo{.}{} \\
&= am + an + ap + bm + bn + bp + cm + cn + cp.
\end{align*}

\Dictum{To find the product of two polynomials}, therefore,
\begin{Theorem}
Multiply every term of the multiplicand by each term of
the multiplier, and add the partial products.
\end{Theorem}
%% -----File: 061.png---Folio 55-------

\Paragraph{88.}  In multiplying polynomials, it is a convenient
arrangement to write the multiplier under the multiplicand,
and place like terms of the partial products in
columns.

\Item{1.}  Multiply $2x - 3y$ by $5x - 4y$.
\[
\begin{array}{ccrrcr}
2x &-& 3&y && \\
5x &-& 4&y && \\
\cline{1-4}
10x^{2} &-& 15&xy & & \\
        &-&  8&xy &+& 12y^{2} \\
\hline
10x^{2} &-& 23&xy &+& 12y^{2} \\
\end{array}
\]

We multiply~$2x$, the first term of the multiplicand, by~$5x$,
the first term of the multiplier, and obtain~$10x^{2}$, then~$-3y$,
the second term of the multiplicand, by~$5x$, and obtain~$-15xy$.
The first line of partial products is $10x^{2} - 15xy$.
In multiplying by~$-4y$, we obtain for a second line of partial
products $-8xy + 12y^{2}$, which is put one place to the
right, so that the like terms $-15xy$~and~$-8xy$ may stand
in the same column.  We then add the coefficients of the
like terms, and obtain the complete product in its simplest
form.

\Item{2.}  Multiply $2a + 3 - 4a^{2}$ by $3 - 2a^{2} - 3a$.

Arrange both multiplicand and multiplier according to
the \emph{ascending} powers of~$a$.
\[
\begin{array}{r*{4}{cr}}
3 &+& 2a &-& 4a^{2} & & && \\
3 &-& 3a &-& 2a^{2} & & && \\
\cline{1-5}
9 &+& 6a &-&12a^{2} & & && \\
  &-& 9a &-& 6a^{2} &+&12a^{3} & & \\
  & &    &-& 6a^{2} &-& 4a^{3} &+& 8a^{4} \\
\hline
9 &-& 3a &-&24a^{2} &+& 8a^{3} &+& 8a^{4} \\
\end{array}
\]
%% -----File: 062.png---Folio 56-------

\Item{3.}  Multiply $3x + x^{4} - 2x^{2}$ by $x^{3} - 2 - x$.

Arrange according to the \emph{descending} powers of~$x$.
\[
\begin{array}{r*{5}{cr}}
x^{4} &-& 2x^{2} &+& \PadTo[l]{3x^4}{3x} && && && \\
x^{3} &-& \PadTo[c]{2x^2}{x}     &-& \PadTo[l]{3x^4}{2}&& && && \\
\cline{1-5}
x^{7} &-& 2x^{5} &+& 3x^{4} && && && \\
      &-&  x^{5} & &        &+& 2x^{3} &-& 3x^{2} && \\
      & &        &-& 2x^{4} & &        &+& 4x^{2} &-& 6x \\
\hline
x^{7} &-& 3x^{5} &+&  x^{4} &+& 2x^{3} &+&  x^{2} &-& 6x \\
\end{array}
\]

\Item{4.}  Multiply $a^{2} + b^{2} + c^{2} - ab - bc - ac$ by $a + b + c$.

Arrange according to descending powers of~$a$.
\[
\begin{array}{l*{8}{cr}}
a^{2} &-& ab &-& ac &+& b^{2} &-& bc &+& c^{2} \\
a     &+&  b &+&  c \\
\cline{1-11}
a^{3} &-& a^{2}b &-& a^{2}c &+& ab^{2} &-& abc   &+& ac^{2} \\
      &+& a^{2}b & &        &-& ab^{2} &-& abc   & &        &+& b^{3} &-& b^{2}c + bc^{2} \\
      & &        &+& a^{2}c & &        &-& abc   &-& ac^{2} & &       &+& b^{2}c - bc^{2} &+& c^{3} \\
\hline
a^{3} & &        & &        & &        &-&3abc   & &        &+& b^{3} & &                &+& c^{3} \\
\end{array}
\]

\begin{Remark}[Note.]
The pupil should observe that, with a view to bringing
like terms of the partial products in columns, the terms of the multiplicand
and multiplier are arranged in the \emph{same order}.
\end{Remark}

\ScreenBreak
\Exercise{22.}

Find the product of:
\begin{multicols}{2}
\Item{1.}  $x + 7$ and $x + 6$.

\Item{2.}  $x - 7$ and $x + 6$.

\Item{3.}  $x + 7$ and $x - 6$.

\Item{4.}  $x - 7$ and $x - 6$.

\Item{5.}  $x + 8$ and $x - 5$.

\Item{6.}  $2x + 3$ and $2x + 3$.

\Item{7.}  $2x - 3$ and $2x - 3$.

\Item{8.}  $2x + 3$ and $2x - 3$.

\Item{9.}  $3x - 2$ and $2 - 3x$.

\Item{10.}  $5x - 3$ and $4x - 7$.

\Item{11.}  $a - 2b$ and $a + 3b$.

\Item{12.}  $a - 7b$ and $a - 5b$.
%% -----File: 063.png---Folio 57-------

\Item{13.}  $5x - 3y$ and $5x - 3y$.

\Item{14.}  $x - b$ and $x - c$.

\Item{15.}  $2m - p$ and $4m - 3p$.

\Item{16.}  $a + b + c$ and $a - c$.

\Item{17.}  $a^{2} - ab + b^{2}$ and $a^{2} + b^{2}$.

\Item{18.}  $x^{3} - 3x^{2} + 7$ and $x^{2} - 3$.

\Item{19.}  $a^{2} + ab + b^{2}$ and $a - b$.

\Item{20.}  $a^{2} - ab + b^{2}$ and $a + b$.
\end{multicols}

\Item{21.}  $x^{2} + 5x - 10$ and $2x^{2} + 3x - 4$.

\Item{22.}  $3x^{3} - 2x^{2} + x$ and $3x^{2} + 2x  - 2$.

\Item{23.}  $x^{3} + 2x^{2}y + 3xy^{2}$ and $x^{2} - 2xy + y^{2}$.

\Item{24.}  $a^{2} - 3ab - b^{2}$ and $-a^{2} + ab + 2b^{2}$.

\Item{25.}  $3a^{2}b^{2} + 2 ab^{3} - 5a^{3}b$ and $5a^{2}b^{2} - ab^{3} - b^{4}$.

\Item{26.}  $a^{2} - 2ab + b^{2}$ and $a^{2} + 2ab + b^{2}$.

\Item{27.}  $ab + ac + cd$ and $ab - ac + cd$.

\Item{28.}  $3x^{2}y^{2} + xy^{3} - 2x^{3}y$ and $x^{2}y^{2} + xy^{3} - 3y^{4}$.

\Item{29.}  $x^{2} + 2xy - y^{2}$ and $x^{2} - 2xy + y^{2}$.

\Item{30.}  $3x^{2} + xy - y^{2}$ and $x^{2} - 2xy - 3y^{2}$.

\Item{31.}  $a^{2} - 2ab - b^{2}$ and $b^{2} - 2ab - a^{2}$.

\Item{32.}  $a^{2} + b^{2} + c^{2} - ac$ and $a^{2} - b^{2} - c^{2}$.

\Item{33.}  $a^{2} + 4abx - 4a^{2}b^{2}x^{2}$ and $a^{2} - 4abx + 4a^{2}b^{2}x^{2}$.

\Item{34.}  $3 a^{2} - 2abx + b^{2}x^{2}$ and $2a^{2} + 3abx - 2b^{2}x^{2}$.

\Item{35.}  $2x^{3}y + 4x^{2}y^{2} - 8xy^{3}$ and $2xy^{3} - 3x^{2}y^{2} + 5x^{3}y$.

\Paragraph{89.  Division.  Polynomials by Monomials.}
\begin{DPalign*}
\lintertext{\indent Since}
a(b + c - d) &= ab + ac - ad, \\
\therefore \frac{ab + ac - ad}{a}
  &= \frac{ab}{a} + \frac{ac}{a} - \frac{ad}{a} \\
  &= b + c - d.
\end{DPalign*}
%% -----File: 064.png---Folio 58-------

\Dictum{To divide a polynomial by a monomial}, therefore,
\begin{Theorem}
Divide each term of the dividend by the divisor, and add
the partial quotients.
\end{Theorem}

Divide $3a^{4}b^{2}c - 9a^{3}bc^{2} - 6a^{2}c^{3}$ by $3a^{2}c$.
\begin{align*}
\frac{3a^{4}b^{2}c - 9a^{3}bc^{2} - 6a^{2}c^{3}}{3a^{2}c}
  &= \frac{3a^{4}b^{2}c}{3a^{2}c}
   - \frac{9a^{3}bc^{2}}{3a^{2}c}
   - \frac{6a^{2}c^{3}}{3a^{2}c} \\
  &= a^{2}b^{2} - 3abc - 2c^{2}.
\end{align*}

\Exercise{23.}

Divide:
\begin{multicols}{2}
\Item{1.}  $2a^{3} - a^{2}$ by $a$.

\Item{2.}  $42a^{5} - 6a^{2}$ by $6a$.

\Item{3.}  $21x^{4} + 3x^{2}$ by $3x^{2}$.

\Item{4.}  $35m^{4} - 7p^{2}$ by $7$.

\Item{5.}  $27x^{5} - 45x^{4}$ by $9x^{2}$.

\Item{6.}  $24x^{6} - 8x^{3}$ by $-8x^{3}$.

\Item{7.}  $34x^{3} - 51x^{2}$ by $17x$.

\Item{8.}  $5x^{5} - 10x^{3}$ by $-5x^{3}$.

\Item{9.}  $-3a^{2} - 6ac$ by $-3a$.

\Item{10.}  $-5x^{3} + x^{2}y$ by $-x^{2}$.

\Item{11.}  $2a^{5}x^{3} - 2a^{4}x^{2}$ by $2a^{4}x^{2}$.

\Item{12.}  $-x^{2}y - x^{2}y^{2}$ by $-xy$.

\Item{13.}  $9a - 12b + 6c$ by $-3$.

\Item{14.}  $a^{3}b^{2} - a^{2}b^{5} - a^{4}b^{2}$ by $a^{2}b$.

\Item{15.}  $3x^{3} - 6x^{2}y - 9xy^{2}$ by $3x$.

\Item{16.}  $x^{2}y^{2} - x^{3}y - xy^{3}$ by $xy$.

\Item{17.}  $a^{3} - a^{2}b - ab^{2}$ by $-a$.

\Item{18.}  $a^{2}b - ab + ab^{2}$ by $-ab$.

\Item{19.}  $xy - x^{2}y^{2} + x^{3}y^{3}$ by $-xy$.

\Item{20.}  $-x^{6} - 2x^{5} - x^{4}$ by $-x^{4}$.
\end{multicols}

\Item{21.}  $a^{2}x - abx - acx$ by $ax$.

\Item{22.}  $3x^{5}y^{2} - 3x^{4}y^{3} - 3x^{2}y^{4}$ by $3x^{2}y^{2}$.

\Item{23.}  $a^{2}b^{2} - 2ab - 3ab^{3}$ by $ab$.

\Item{24.}  $3a^{3}c^{3} + 3a^{2}c - 3ac^{2}$ by $3ac$.
%% -----File: 065.png---Folio 59-------

\Paragraph{90.  Division.  Polynomials by Polynomials.}
\[
\begin{array}{l*{4}{cr}}
\text{If the divisor (one factor)}
  &=& & & a &+& b &+& c, \\
\text{and the quotient (other factor)}
  &=& & & n &+& p &+& q, \\
\cline{5-9}
  & & & &an &+& bn&+& cn \\
\text{then the dividend (product)}
  &=& \smash{\left\{\threelines\right.}\kern-4pt
      &+&ap &+& bp&+& cp \\
  & & &+&aq &+& bq&+& cq\rlap{.} \\
\end{array}
\]

The first term of the dividend is~$an$; that is, the product
of~$a$, the first term of the divisor, by~$n$, the first term of the
quotient.  The first term~$n$ of the quotient is therefore
found by dividing~$an$, the first term of the dividend, by~$a$,
the first term of the divisor.

If the partial product formed by multiplying the entire
divisor by~$n$ be subtracted from the dividend, the first term
of the remainder~$ap$ is the product of~$a$, the first term of
the divisor, by~$p$, the second term of the quotient; that is,
the second term of the quotient is obtained by dividing the
first term of the remainder by the first term of the divisor.
In like manner, the third term of the quotient is obtained
by dividing the first term of the new remainder by the first
term of the divisor; and so on.

\Dictum{To divide one polynomial by another}, therefore,
\begin{Theorem}
Arrange both the dividend and divisor in ascending or
descending powers of some common letter.

Divide the first term of the dividend by the first term of
the divisor.

Write the result as the first term of the quotient.

Multiply all the terms of the divisor by the first term of
the quotient.

Subtract the product from the dividend.

If there is a remainder, consider it as a new dividend,
and proceed as before.
\end{Theorem}
%% -----File: 066.png---Folio 60-------

\Paragraph{91.}  It is of fundamental importance to arrange the dividend
and divisor \emph{in the same order} with respect to a common
letter, and \emph{to keep this order throughout the operation}.

The beginner should study carefully the processes in the
following examples:

\Item{1.}  Divide $x^{2} + 18x + 77$ by $x + 7$.
\[
\begin{array}{r*{2}{cr}|l}
x^{2} &+& 18x &+& 77 & x + 7 \\
\cline{6-6}
x^{2} &+&  7x & &    & x + 11 \\
\cline{1-5}
      & & 11x &+& \NoBar{77} \\
      & & 11x &+& \NoBar{77} \\
\cline{3-5}
\end{array}
\]

\begin{Remark}[Note.]
The pupil will notice that by this process we have in
effect separated the dividend into two parts, $x^{2} + 7x$ and $11x + 77$,
and divided each part by $x + 7$, and that the complete quotient is the
sum of the partial quotients $x$~and~$11$.  Thus,
\begin{align*}
x^{2} + 18x + 77
  &= x^{2} + 7x + 11x + 77 = (x^{2} + 7x) + (11x + 77). \\
\therefore \frac{x^{2} + 18x + 77}{x + 7}
  &= \frac{x^{2} + 7x}{x + 7} + \frac{11x + 77}{x + 7} = x + 11.
\end{align*}
\end{Remark}

\Item{2.}  Divide $a^{2} - 2ab + b^{2}$ by $a - b$.
\[
\begin{array}{r*{2}{cr}|l}
a^{2} &-& 2ab &+& b^{2} & a - b \\
\cline{6-6}
a^{2} &-&  ab & &       & a - b \\
\cline{1-5}
      &-&  ab &+& \NoBar{b^{2}} \\
      &-&  ab &+& \NoBar{b^{2}} \\
\cline{2-5}
\end{array}
\]

\Item{3.}  Divide $a^{4} - ab^{3} + b^{4} + 2a^{2}b^{2} - a^{3}b$ by $a^{2} + b^{2}$.

Arrange according to the descending powers of~$a$.
\[
\begin{array}{r*{4}{cr}|l}
a^{4} &-& a^{3}b &+& 2a^{2}b^{2} &-& ab^{3} &+& b^{4} & a^{2} + b^{2} \\
\cline{10-10}
a^{4} & &        &+&  a^{2}b^{2} & &        & &       & a^{2} - ab + b^{2} \\
\cline{1-9}
      &-& a^{3}b &+&  a^{2}b^{2} &-& ab^{3} &+& \NoBar{b^{4}} \\
      &-& a^{3}b & &             &-& ab^{3} \\
\cline{2-9}
      & &        &+&  a^{2}b^{2} & &        &+& \NoBar{b^{4}} \\
      & &        &+&  a^{2}b^{2} & &        &+& \NoBar{b^{4}} \\
\cline{4-9}
\end{array}
\]
%% -----File: 067.png---Folio 61-------

\Item{4.}  Divide $10a^{2}b^{2} - 20b^{4} - 17a^{3}b + 6a^{4} + ab^{3}$
by $2a^{2} - 4b^{2} - 3ab$.

Arrange according to descending powers of~$a$.
\[
\begin{array}{r*{4}{cr}|l}
6a^{4} &-&17a^{3}b &+& 10a^{2}b^{2} &+&  ab^{3} &-& 20b^{4} & 2a^{2} - 3ab - 4b^{2} \\
\cline{10-10}
6a^{4} &-& 9a^{3}b &-& 12a^{2}b^{2} & &         & &         & 3a^{2} - 4ab + 5b^{2} \\
\cline{1-9}
       &-& 8a^{3}b &+& 22a^{2}b^{2} &+&  ab^{3} &-& \NoBar{20b^{4}} \\
       &-& 8a^{3}b &+& 12a^{2}b^{2} &+&16ab^{3} \\
\cline{2-9}
       & &         & & 10a^{2}b^{2} &-&15ab^{3} &-& \NoBar{20b^{4}} \\
       & &         & & 10a^{2}b^{2} &-&15ab^{3} &-& \NoBar{20b^{4}} \\
\cline{5-9}
\end{array}
\]

\Item{5.}  Divide $5x^{3} - 3x^{4} - 4x^{2} + 1 + x$ by $1 + 2x - 3x^{2}$.

Arrange according to ascending powers of~$x$.
\[
\begin{array}{r*{4}{cr}|l*{2}{cr}}
1 &+&  x &-& 4x^{2} &+& 5x^{3} &-& 3x^{4} & 1 &+& 2x &-& 3x^{2} \\
\cline{10-14}
1 &+& 2x &-& 3x^{2} & &        & &        & 1 &-&  x &+&  x^{2} \\
\cline{1-9}
  &-&  x &-&  x^{2} &+& 5x^{3} &-& \NoBar{3x^{4}} \\
  &-&  x &-& 2x^{2} &+& 3x^{3} \\
\cline{2-9}
  & &    & &  x^{2} &+& 2x^{3} &-& \NoBar{3x^{4}} \\
  & &    & &  x^{2} &+& 2x^{3} &-& \NoBar{3x^{4}} \\
\end{array}
\]

\Item{6.}  Divide $a^{3} + b^{3} + c^{3} - 3abc$ by $a + b + c$.

Arrange according to descending powers of~$a$.
\[
%[** TN: Re-formatted slightly from the original]
\begin{array}{r*{4}{cr}clcr|l}
a^{3} & &        & &        &-& 3abc & &       &+& b^{3} &+& c^{3} & a + b + c \\
\cline{14-14}
a^{3} &+& a^{2}b &+& a^{2}c & &      & &        & &       & &       &
\smash[b]{\begin{aligned}[t]
  a^{2} &- ab - ac \\
  &+b^{2} - bc + c^{2}
\end{aligned}} \\
\cline{1-13}
      &-& a^{2}b &-& a^{2}c &-& 3abc & &        &+& b^{3} &+& \NoBar{c^{3}} \\
      &-& a^{2}b &-& ab^{2} &-&  abc \\
\cline{2-13}
      &-& a^{2}c &+& ab^{2} &-& 2abc & &        &+& b^{3} &+& \NoBar{c^{3}} \\
      &-& a^{2}c & &        &-&  abc &-& ac^{2} \\
\cline{2-13}
      & &        & & ab^{2} &-&  abc &+& ac^{2} &+& b^{3} &+& \NoBar{c^{3}} \\
      & &        & & ab^{2} & &      & &        &+& b^{3} &+& \NoBar{b^{2}c} \\
\cline{4-13}
      & &        & &        &-&  abc &+& ac^{2} &-& b^{2}c&+& \NoBar{c^{3}} \\
      & &        & &        &-&  abc & &        &-& b^{2}c&-& \NoBar{bc^{2}} \\
\cline{6-13}
      & &        & &        & &      & & ac^{2} &+& bc^{2}&+& \NoBar{c^{3}} \\
      & &        & &        & &      & & ac^{2} &+& bc^{2}&+& \NoBar{c^{3}} \\
\cline{9-13}
\end{array}
\]
%% -----File: 068.png---Folio 62-------

\Exercise{24.}

Divide:
\begin{multicols}{2}
\Item{1.}  $x^{2} + 15x + 56$ by $x + 7$.

\Item{2.}  $x^{2} - 15x + 56$ by $x - 7$.

\Item{3.}  $x^{2} + x-56$ by $x - 7$.

\Item{4.}  $x^{2} - x-56$ by $x + 7$.

\Item{5.}  $2a^{2} + 11a + 5$ by $2a + 1$.

\Item{6.}  $6a^{2} - 7a-3$ by $2a - 3$.

\Item{7.}  $4a^{2} + 23a + 15$ by $4a + 3$.

\Item{8.}  $3a^{2} - 4a-4$ by $2 - a$.

\Item{9.}  $x^{4} + x^{2} + 1$ by $x^{2} + x + 1$.

\Item{10.}  $x^{8} + x^{4} + 1$ by $x^{4} - x^{2} + 1$.

\Item{11.}  $1 - a^{3}b^{3}$ by $1 - ab$.

\Item{12.}  $x^{3} - 8x-3$ by $x - 3$.
\end{multicols}

\Item{13.}  $a^{2} - 2ab + b^{2} - c^{2}$ by $a - b - c$.

\Item{14.}  $a^{2} + 2ab + b^{2} - c^{2}$ by $a + b + c$.

\Item{15.}  $x^{2} - y^{2} + 2yz - z^{2}$ by $x - y + z$.

\Item{16.}  $c^{4} + 2c^{2} - c + 2$ by $c^{2} - c + 1$.

\Item{17.}  $x^{2} - 4y^{2} - 4yz - z^{2}$ by $x + 2y + z$.

Arrange and divide:

\Item{18.}  $x^{3} - 6a^{3} + 11a^{2}x - 6ax^{2}$ by $x^{2} + 6a^{2} - 5ax$.

\Item{19.}  $a^{2} - 4b^{2} - 9c^{2} + 12bc$ by $a - 3c + 2b$.

\Item{20.}  $2a^{3} - 8a + a^{4} + 12-7a^{2}$ by $2 + a^{2} - 3a$.

\Item{21.}  $q^{4} + 6q^{3} + 4 + 12q + 13q^{2}$ by $3q + 2 + q^{2}$.

\Item{22.}  $27a^{3} - 8b^{3}$ by $3a - 2b$.

Find the remainder when:

\Item{23.}  $a^{4} + 9a^{2} + 15-11a - 7a^{3}$ is divided by $a - 5$.

\Item{24.}  $7 - 8c^{2} + 5c^{3} + 8c$ is divided by $5c - 3$.

\Item{25.}  $3 + 11a^{3} + 30a^{4} - 82a^{2} - 5a$ is divided by $3a^{2} - 4 + 2a$.

\Item{26.}  $2x^{3} - 16x + 10-39x^{2} + 17x^{4}$ is divided by $2 - 5x^{2} - 4x$.
%% -----File: 069.png---Folio 63-------

\Exercise[Miscellaneous Examples.]{25.}

\Item{1.}  Add $2a^{2} - 3ac - 3ab$; $2b^{2} + 3ac + a^{2}$; $-a^{2} - 2b^{2} + 3ab$.

\Item{2.}  Subtract $3a^{4} - 2 a^{3}b + 4 a^{2}b^{2}$ from $4b^{4} - 2 ab^{3} + 4 a^{2}b^{2}$.

\Item{3.}  Simplify $x - y - \{z - x - (y - x + z)\}$.

\Item{4.}  Multiply $a^{2} + b^{2} + c^{2} - d^{2}$ by $a^{2} + b^{2} - c^{2} + d^{2}$.

\Item{5.}  Divide $10y^{6} + 2 - 12y^{5}$ by $1 + y^{2} - 2y$.

\Item{6.}  If $a = 1$, $b = 2$, and $c = -3$, find the value of
    $a^{2} + b^{2} + c^{2} + 2ab + 2ac + 2bc$.

\Item{7.}  Simplify $x - (y - z) - \bigl\{4y + [2y - (z - x)]\bigr\}$.

\Item{8.}  Multiply $a^{2} + b^{2} + c^{2} - ab - ac - bc$ by $a + b + c$.

\Item{9.}  Divide $16y^{4} - 21x^{2}y^{2} + 21x^{3}y - 10x^{4}$ by $4y^{2} - 5x^{2} + 3xy$.

\Item{10.}  Add $-2a^{4} + 3a^{3}b - 4a^{2}b^{2}$; $2a^{3}b - 3a^{2}b^{2}$; $7a^{2}b^{2} + 2a^{4} - b^{4}$.

\Item{11.}  From $3x^{3} + 5x - 1$ take the sum of $x - 5 + 5x^{2}$ and
$3 + 4x -3x^{2}$.

\Item{12.}  The minuend is $9c^{2} + 11c - 5$, and the remainder is
$6c^{2} - 13c + 7$.  What is the subtrahend?

\Item{13.}  Find the remainder when $a^{4} + 6b^{4}$ is divided by
$a^{2} + 2ab + 2b^{2}$.

\Item{14.}  Multiply $2 - 5x^{2} - 4x$ by $5 + 2x - 3x^{2}$.

\Item{15.}  Divide $a^{6} + a^{5}x + a^{4}x^{2} - a^{3}x^{3} + x^{6}$ by $a^{2} + ax + x^{2}$.

Bracket the coefficients of the different powers of~$x$:

\Item{16.}  $ax^{3} - cx + bx^{2} - bx^{3} + cx^{2} - x$.

\Item{17.}  $ax^{4} - 2x + bx^{4} - cx - ax^{3} + bx^{3}$.

\Item{18.}  $x^{3} - bx^{2} - cx + bx - cx^{2} + ax^{3}$.
%% -----File: 070.png---Folio 64-------


\Chapter[Special Rules in Multiplication and Division.]
{VI.}{Multiplication and Division.}

\Section{Special Rules.}

\Paragraph{92.  Special Rules of Multiplication.}  Some results of multiplication
are of so great utility in shortening algebraic
work that they should be carefully noticed and remembered.
The following are important:

\Paragraph{93.  Square of the Sum of Two Numbers.}
\begin{align*}
(a + b)^{2}
  &= (a + b)(a + b) \\
  &= a(a + b) + b(a + b) \\
  &= a^{2} + ab + ab + b^{2} \\
  &= a^{2} + 2ab + b^{2}.
\end{align*}

Since $a$~and~$b$ stand for any two numbers, we have
\begin{Theorem}[\textsc{Rule 1.}]  The square of the sum of two numbers is the
sum of their squares plus twice their product.
\end{Theorem}

\Paragraph{94.  Square of the Difference of Two Numbers.}
\begin{align*}
(a - b)^{2}
  &= (a - b) (a - b) \\
  &= a(a - b) - b(a - b) \\
  &= a^{2} - ab - ab + b^{2} \\
  &= a^{2} - 2ab + b^{2}.
\end{align*}

Hence we have
\begin{Theorem}[\textsc{Rule 2.}]  The square of the difference of two numbers is
the sum of their squares minus twice their product.
\end{Theorem}
%% -----File: 071.png---Folio 65-------

\Paragraph{95.  Product of the Sum and Difference of Two Numbers.}
\begin{align*}
(a + b)(a - b)
  &= a(a - b) + b(a - b) \\
  &= a^{2} - ab + ab - b^{2} \\
  &= a^{2} - b^{2}.
\end{align*}

Hence, we have
\begin{Theorem}[\textsc{Rule 3.}]  The product of the sum and difference of two
numbers is the difference of their squares.
\end{Theorem}

If we put $2x$~for~$a$, and $3$~for~$b$, we have
\begin{DPalign*}
\lintertext{\indent Rule 1,}  &(2x + 3)^{2} = 4x^{2} + 12x + 9. \\
\lintertext{\indent Rule 2,}  &(2x - 3)^{2} = 4x^{2} - 12x + 9. \\
\lintertext{\indent Rule 3,}  &(2x + 3)(2x - 3) = 4x^{2} - 9.
\end{DPalign*}

\Exercise{26.}

Write by inspection the value of:
\begin{multicols}{2}
\Item{1.}  $(m + n)^{2}$.

\Item{2.}  $(c - a)^{2}$.

\Item{3.}  $(a + 2c)^{2}$.

\Item{4.}  $(3a - 2b)^{2}$.

\Item{5.}  $(2a + 3b)^{2}$.

\Item{6.}  $(a - 3b)^{2}$.

\Item{7.}  $(2x - y)^{2}$.

\Item{8.}  $(y - 2x)^{2}$.

\Item{9.}  $(a + 5b)^{2}$.

\Item{10.}  $(2a - 5c)^{2}$.

\Item{11.}  $(x + y)(x - y)$.

\Item{12.}  $(4a - b)(4a + b)$.

\Item{13.}  $(2b - 3c)(2b + 3c)$.

\Item{14.}  $(x + 5b)(x + 5b)$.

\Item{15.}  $(y - 2z)(y - 2z)$.

\Item{16.}  $(y + 3z)(y - 3z)$.

\Item{17.}  $(2a - 3b)(2a + 3b)$.

\Item{18.}  $(2a - 3b)(2a - 3b)$.

\Item{19.}  $(2a + 3b)(2a + 3b)$.

\Item{20.}  $(5x + 3a)(5x - 3a)$.
\end{multicols}
%% -----File: 072.png---Folio 66-------

\Paragraph{96.  Product of Two Binomials of the Form $x + a$, $x + b$.}
The product of two binomials which have the form $x + a$,
$x + b$, should be carefully noticed and remembered.
\begin{DPalign*}
\lintertext{\Item{1.}}
(x + 5)(x + 3)
  &= x(x + 3) + 5(x + 3) \\
  &= x^{2} + 3x + 5x + 15 \\
  &= x^{2} + 8x + 15. \displaybreak[1] \\
%
\lintertext{\Item{2.}}
(x - 5)(x - 3)
  &= x(x - 3) - 5(x - 3) \\
  &= x^{2} - 3x - 5x + 15 \\
  &= x^{2} - 8x + 15. \displaybreak[1] \\
%
\lintertext{\Item{3.}}
(x + 5)(x - 3)
  &= x(x - 3) + 5(x - 3) \\
  &= x^{2} - 3x + 5x - 15 \\
  &= x^{2} + 2x - 15. \displaybreak[1] \\
%
\lintertext{\Item{4.}}
(x - 5)(x + 3)
  &= x(x + 3) - 5(x + 3) \\
  &= x^{2} + 3x - 5x - 15 \\
  &= x^{2} - 2x - 15.
\end{DPalign*}

\Item{1.}  Each of these results has three terms.

\Item{2.}  The first term of each result is the product of the first
terms of the binomials.

\Item{3.}  The last term of each result is the product of the
second terms of the binomials.

\Item{4.}  The middle term of each result has for a coefficient
the \emph{algebraic sum} of the second terms of the binomials.

\Paragraph{97.}  The intermediate step given above may be omitted,
and the products written at once by \emph{inspection}.  Thus,

\Item{1.}  Multiply $x + 8$ by $x + 7$.
\begin{align*}
&8 + 7 = 15,\quad 8 × 7 = 56. \\
\therefore\ &(x + 8)(x + 7) = x^{2} + 15x + 56.
\end{align*}
%% -----File: 073.png---Folio 67-------

\Item{2.}  Multiply $x - 8$ by $x - 7$.
\begin{align*}
&(-8) + (-7) = -15,\quad (-8)(-7) = +56. \\
\therefore\ &(x - 8)(x - 7) = x^{2} - 15x + 56.
\end{align*}

\Item{3.}  Multiply $x - 7y$ by $x + 6y$.
\begin{align*}
&-7 + 6 = -1,\quad  (-7) × 6y = -42y^{2}. \\
\therefore\ &(x - 7y)(x + 6y) = x^{2} - xy - 42y^{2}.
\end{align*}

\Item{4.}  Multiply $x + 6y$ by $x - 5y$.
\begin{align*}
&+6 - 5 = 1,\quad 6y × (-5y) = -30y^{2}. \\
\DPtypo{}{\therefore}\ &(x + 6y)(x - 5y) = x^{2} + xy - 30y^{2}.
\end{align*}

\Exercise{27.}

Write by inspection the product of:
\begin{multicols}{2}
\Item{1.}  $(x + 7)(x + 4)$.

\Item{2.}  $(x - 3)(x + 7)$.

\Item{3.}  $(x - 2)(x - 4)$.

\Item{4.}  $(x - 6)(x - 10)$.

\Item{5.}  $(x + 7)(x - 4)$.

\Item{6.}  $(x + a)(x - 2a)$.

\Item{7.}  $(x + 3a)(x - a)$.

\Item{8.}  $(a + 3c)(a + 3c)$.

\Item{9.}  $(a + 2x)(a - 4x)$.

\Item{10.}  $(a - 3b)(a - 4b)$.

\Item{11.}  $(a^{2} - c)(a^{2} + 2c)$.

\Item{12.}  $(x - 17)(x - 3)$.

\Item{13.}  $(x + 6y)(x - 5y)$.

\Item{14.}  $(3 + 2x)(3 - x)$.

\Item{15.}  $(5 + 2x)(1 - 2x)$.

\Item{16.}  $(a - 2b)(a + 3b)$.

\Item{17.}  $(a^{2}b^{2} - x^{2})(a^{2}b^{2} - 5x^{2})$.

\Item{18.}  $(a^{3}b - ab^{3})(a^{3}b + 5ab^{3})$.

\Item{19.}  $(x^{2}y - xy^{2})(x^{2}y - 3xy^{2})$.

\Item{20.}  $(x^{2}y + xy^{2})(x^{2}y + xy^{2})$.

\Item{21.}  $(x + a)(x + b)$.

\Item{22.}  $(x + a)(x - b)$.

\Item{23.}  $(x - a)(x + b)$.

\Item{24.}  $(x - a)(x - b)$.

\Item{25.}  $(x + 2a)(x + 2b)$.

\Item{26.}  $(x - 2a)(x + 2b)$.

\Item{27.}  $(x + 2a)(x - 2b)$.

\Item{28.}  $(x - 2a)(x - 2b)$.

\Item{29.}  $(x - a)(x + 3a)$.

\Item{30.}  $(x - 2a)(x + 3a)$.
\end{multicols}
%% -----File: 074.png---Folio 68-------

\Paragraph{98.  Special Rules of Division.}  Some results in division
are so important in abridging algebraic work that they
should be carefully noticed and remembered.

\Paragraph{99.  Difference of Two Squares.}

Since $(a + b)(a - b) = a^{2} - b^{2}$,
\[
\therefore
\frac{a^{2} - b^{2}}{a + b} = a - b;\quad\text{and}\quad
\frac{a^{2} - b^{2}}{a - b} = a + b. \EqText{Hence\Add{,}}
\]
\begin{Theorem}[\textsc{Rule 1.}]  The difference of the squares of two numbers is
divisible by the sum, and by the difference, of the numbers.
\end{Theorem}

\ScreenBreak
\Exercise{28.}

Write by inspection the quotient of:
\begin{multicols}{3}
\Item{1.}  $\dfrac{x^{2} - 4}{x - 2}$.

\Item{2.}  $\dfrac{x^{2} - 4}{x + 2}$.

\Item{3.}  $\dfrac{a^{2} - 9}{a - 3}$.

\Item{4.}  $\dfrac{a^{2} - 9}{a + 3}$.

\Item{5.}  $\dfrac{c^{2} - 25}{c - 5}$.

\Item{6.}  $\dfrac{c^{2} - 25}{c + 5}$.

\Item{7.}  $\dfrac{49x^{2} - y^{2}}{7x - y}$.

\Item{8.}  $\dfrac{49x^{2} - y^{2}}{7x + y}$.

\Item{9.}  $\dfrac{9b^{2} - 1}{3b - 1}$.

\Item{10.}  $\dfrac{9b^{2} - 1}{3b + 1}$.

\Item{11.}  $\dfrac{16x^{4} - 25a^{2}}{4x^{2} - 5a}$.

\Item{12.}  $\dfrac{16x^{4} - 25a^{2}}{4x^{2} + 5a}$.
\end{multicols}

\begin{multicols}{2}
\Item{13.}  $\dfrac{9x^{2} - 25y^{2}}{3x - 5y}$.

\Item{14.}  $\dfrac{a^{2}-(b - c)^{2}}{a-(b - c)}$.

\Item{15.}  $\dfrac{a^{2}-(b - c)^{2}}{a + (b - c)}$.

\Item{16.}  $\dfrac{a^{2}-(2b - c)^{2}}{a-(2b - c)}$.

\Item{17.}  $\dfrac{(5a - 7b)^{2} - 1}{(5a - 7b) - 1}$.

\Item{18.}  $\dfrac{(5a - 7b)^{2} - 1}{(5a - 7b) + 1}$.

\Item{19.}  $\dfrac{z^{2}-(x - y)^{2}}{z-(x - y)}$.

\Item{20.}  $\dfrac{z^{2}-(x - y)^{2}}{z + (x - y)}$.
%% -----File: 075.png---Folio 69-------

\Item{21.}  $\dfrac{a^{2}-(2b - c)^{2}}{a + (2b - c)}$.

\Item{22.}  $\dfrac{(x + 3y)^{2} - z^{2}}{(x + 3y) - z}$.

\Item{23.}  $\dfrac{(x + 3y)^{2} - z^{2}}{x + 3y + z}$.

\Item{24.}  $\dfrac{(a + 2b)^{2} - 4c^{2}}{(a + 2b) - 2c}$.

\Item{25.}  $\dfrac{(a + 2b)^{2} - 4c^{2}}{(a + 2b) + 2c}$.

\Item{26.}  $\dfrac{1 - (3x - 2y)^{2}}{1 + (3x - 2y)}$.
\end{multicols}

\Paragraph{100.  Difference of Two Cubes.}  By performing the division
we have
\[
\frac{a^{3} - b^{3}}{a - b} = a^{2} + ab + b^{2}. \EqText{Hence,}
\]
\begin{Theorem}[\textsc{Rule 2.}]  The difference of the cubes of two numbers is
divisible by the difference of the numbers, and the quotient
is the sum of the squares of the numbers plus their product.
\end{Theorem}

\Exercise{29.}

Write by inspection the quotient of:
\begin{multicols}{3}
\Item{1.}  $\dfrac{1 - x^{3}}{1 - x}$.

\Item{2.}  $\dfrac{1 - 8a^{3}}{1 - 2a}$.

\Item{3.}  $\dfrac{1 - 27c^{3}}{1 - 2c}$.

\Item{4.}  $\dfrac{8a^{3} - b^{3}}{2a - b}$.

\Item{5.}  $\dfrac{64b^{3} - 27c^{3}}{4b - 3c}$.

\Item{6.}  $\dfrac{27x^{3} - 8y^{3}}{3x - 2y}$.

\Item{7.}  $\dfrac{x^{3}y^{3} - z^{3}}{xy - z}$.

\Item{8.}  $\dfrac{a^{3}b^{3} - 8}{ab - 2}$.

\Item{9.}  $\dfrac{125a^{3} - b^{3}}{5a - b}$.

\Item{10.}  $\dfrac{a^{3} - 8b^{3}}{a - 2b}$.

\Item{11.}  $\dfrac{a^{3} - 64}{a - 4}$.

\Item{12.}  $\dfrac{a^{9} - 27}{a^{3} - 3}$.

\Item{13.}  $\dfrac{a^{12} - x^{6}y^{6}}{a^{4} - x^{2}y^{2}}$.

\Item{14.}  $\dfrac{x^{15} - a^{9}b^{9}}{x^{5} - a^{3}b^{3}}$.

\Item{15.}  $\dfrac{27x^{3}y^{3} - z^{12}}{3xy - z^{4}}$.

\Item{16.}  $\dfrac{x^{3}y^{3}z^{3} - 1}{xyz - 1}$.

\Item{17.}  $\dfrac{8a^{3}b^{3}c^{3} - 27}{2abc - 3}$.

\Item{18.}  $\dfrac{1 - 64x^{3}y^{3}z^{3}}{1 - 4xyz}$.
\end{multicols}
%% -----File: 076.png---Folio 70-------

\Paragraph{101.  Sum of Two Cubes.}  By performing the division,
we find that
\[
\frac{a^{3} + b^{3}}{a + b} = a^{2} - ab + b^{2}. \EqText{Hence,}
\]
\begin{Theorem}[\textsc{Rule 3.}]  The sum of the cubes of two numbers is divisible
by the sum of the numbers, and the quotient is the sum
of the squares of the numbers minus their product.
\end{Theorem}

\Exercise{30.}

Write by inspection the quotient of:
\begin{multicols}{3}
\Item{1.}  $\dfrac{1 + x^{3}}{1 + x}$.

\Item{2.}  $\dfrac{1 + 8a^{3}}{1 + 2a}$.

\Item{3.}  $\dfrac{1 + 27c^{3}}{1 + 3c}$.

\Item{4.}  $\dfrac{8a^{3} + b^{3}}{2a + b}$.

\Item{5.}  $\dfrac{64b^{3} + 27c^{3}}{4b + 3c}$.

\Item{6.}  $\dfrac{27x^{3} + 8y^{3}}{3x + 2y}$.

\Item{7.}  $\dfrac{8x^{3} + 125y^{3}}{2x + 5y}$.

\Item{8.}  $\dfrac{x^{3}y^{3} + z^{3}}{xy + z}$.

\Item{9.}  $\dfrac{a^{3}b^{3} + 8}{ab + 2}$.

\Item{10.}  $\dfrac{125a^{3} + b^{3}}{5a + b}$.

\Item{11.}  $\dfrac{a^{3} + 8b^{3}}{a + 2b}$.

\Item{12.}  $\dfrac{a^{6} + 64}{a^{2} + 4}$.

\Item{13.}  $\dfrac{a^{9} + 27}{a^{3} + 3}$.

\Item{14.}  $\dfrac{8a^{6} + b^{3}}{2a^{2} + b}$.

\Item{15.}  $\dfrac{a^{12} + x^{6}y^{6}}{a^{4} + x^{2}y^{2}}$.

\Item{16.}  $\dfrac{x^{15} + a^{9}b^{9}}{x^{5} + a^{3}b^{3}}$.

\Item{17.}  $\dfrac{27x^{3}y^{3} + z^{12}}{3xy + z^{4}}$.

\Item{18.}  $\dfrac{x^{3}y^{3}z^{3} + 1}{xyz + 1}$.

\Item{19.}  $\dfrac{8a^{3}b^{3}c^{3} + 27}{2abc + 3}$.

\Item{20.}  $\dfrac{1 + 64x^{3}y^{3}z^{3}}{1 + 4xyz}$.

\Item{21.}  $\dfrac{1 + 27a^{6}b^{3}c^{3}}{1 + 3a^{2}bc}$.
\end{multicols}

Find by division the quotient of:
\begin{multicols}{3}
\Item{22.}  $\dfrac{x^{4} - y^{4}}{x - y}$.

\Item{23.}  $\dfrac{x^{4} - y^{4}}{x + y}$.

\Item{24.}  $\dfrac{x^{5} - y^{5}}{x - y}$.

\Item{25.}  $\dfrac{x^{5} + y^{5}}{x + y}$.

\Item{26.}  $\dfrac{x^{6} - y^{6}}{x - y}$.

\Item{27.}  $\dfrac{x^{6} - y^{6}}{x + y}$.
\end{multicols}
%% -----File: 077.png---Folio 71-------


\Chapter{VII.}{Factors.}

\Paragraph{102.  Rational Expressions.}  An expression is \emph{rational} when
none of its terms contain square or other roots.

\Paragraph{103.  Factors of Rational and Integral Expressions.}  By factors
of a given integral number in arithmetic we mean
integral numbers that will divide the given number without
remainder.  Likewise by factors of a rational and integral
expression in algebra we mean rational and integral
expressions that will divide the given expression without
remainder.

\Paragraph{104.  Factors of Monomials.}  The factors of a monomial
may be found by inspection.  Thus, the factors of~$21a^{2}b$
are $3$, $7$, $a$, $a$, and~$b$.

\Paragraph{105.  Factors of Polynomials.}  The form of a polynomial
that can be resolved into factors often suggests the process
of finding the factors.


\Section{Case I.}

\Paragraph{106.  When all the terms have a common factor.}

\Item{1.}  Resolve into factors $3a^{2} - 6ab$.

Since $3a$~is seen to be a factor of each term, we have
\begin{align*}
\frac{3a^{2} - 6ab}{3a} &= \frac{3a^{2}}{3a} - \frac{6ab}{3a} = a - 2b. \\
\therefore\ 3a^{2} - 6ab &= 3a(a - 2b).
\end{align*}

Hence, the required factors are $3a$~and~$a - 2b$.
%% -----File: 078.png---Folio 72-------

\Item{2.}  Resolve into factors $4x^{3} + 12x^{2} - 8x$.

Since $4x$ is seen to be a factor of each term, we have
\begin{align*}
\frac{4x^{3} + 12x^{2} - 8x}{4x}
  &= \frac{4x^{3}}{4x} + \frac{12x^{2}}{4x} - \frac{8x}{4x} \\
  &= x^{2} + 3x - 2. \\
\therefore\
4x^{3} + 12x^{2} - 8x &= 4x(x^{2} + 3x - 2).
\end{align*}

Hence the required factors are $4x$~and~$x^{2} + 3x - 2$.

\Exercise{31.}

Resolve into two factors:
\begin{multicols}{2}
\Item{1.}  $2x^{2} - 4x$.

\Item{2.}  $3a^{3} - 6a$.

\Item{3.}  $5a^{2}b^{2} - 10a^{3}b^{3}$.

\Item{4.}  $3x^{2}y + 4xy^{2}$.

\Item{5.}  $8a^{3}b^{2} + 4a^{2}b^{3}$.

\Item{6.}  $3a^{4} - 12a^{2} - 6a^{3}$.

\Item{7.}  $4x^{2} - 8x^{4} - 12x^{5}$.

\Item{8.}  $5 - 10x^{2}y^{2} + 15x^{2}y$.

\Item{9.}  $7a^{2} + 14a - 21a^{3}$.

\Item{10.}  $3x^{3}y^{3} - 6x^{4}y^{4} - 9x^{2}y^{2}$.
\end{multicols}

\Section{Case II.}

\Paragraph{107.  When the terms can be grouped so as to show a common
factor in each group.}

\Item{1.}  Resolve into factors $ac + ad + bc + bd$.
\begin{align*}
ac + ad + bc + bd
  &= (ac + ad) + (bc + bd)
  \Tag{(1)} \\
  &= a(c + d) + b(c + d)
  \Tag{(2)} \\
  &= (a + b)(c + d).
  \Tag{(3)}
\end{align*}

\begin{Remark}[Note.] The first two terms of $ac + ad + bc + bd$ are seen to
have the common factor~$a$, and the last two terms, the common factor~$b$.
Hence we bracket the first two terms and also the last two
terms.  Then we take out the factor~$a$ from $(ac + ad)$ and $b$~from
$(bc + bd)$, and get equation~(2). Since one factor is seen in~(2) to be
$c + d$, dividing by $c + d$, we obtain the other factor, $a + b$.
\end{Remark}
%% -----File: 079.png---Folio 73-------

\Item{2.}  Find the factors of $ac + ad - bc - bd$.
\begin{align*}
ac + ad - bc - bd
  &= (ac + ad) - (bc + bd) \\
  &= a(c + d) - b(c + d) \\
  &= (a - b)(c + d).
\end{align*}

\begin{Remark}[Note.]  Here the last two terms, $-bc - bd$, being put within a
parenthesis preceded by the sign~$-$, have their signs changed to~$+$.
\end{Remark}

\Item{3.}  Resolve into factors $2x^{3} - 3x^{2} - 4x + 6$.
\begin{align*}
2x^{3} - 3x^{2} - 4x + 6
  &= (2x^{3} - 3x^{2}) - (4x - 6) \\
  &= x^{2}(2x - 3) - 2(2x - 3) \\
  &= (x^{2} - 2)(2x - 3).
\end{align*}

\Item{4.}  Resolve into factors $x^{3} + x^{2} - ax - a$.
\begin{align*}
x^{3} + x^{2} - ax - a
  &= (x^{3} + x^{2}) - (ax + a) \\
  &= x^{2}(x + 1) - a(x + 1) \\
  &= (x^{2} - a)(x + 1).
\end{align*}

\Item{5.}  Resolve into factors $x^{3} + 3ax^{2} + x + 3a$.
\begin{align*}
x^{3} + 3ax^{2} + x + 3a
  &= (x^{3} + 3ax^{2}) + (x + 3a) \\
  &= x^{2}(x + 3a) + 1(x + 3a) \\
  &= (x^{2} + 1)(x + 3a).
\end{align*}

\Exercise{32.}

Resolve into factors:
\begin{multicols}{2}
\Item{1.}  $x^{3} + x^{2} + x + 1$.

\Item{2.}  $x^{3} - x^{2} + x - 1$.

\Item{3.}  $x^{2} + xy + xz + yz$.

\Item{4.}  $ax - bx - ay + by$.

\Item{5.}  $a^{2} - ac + ab - bc$.

\Item{6.}  $x^{2} - bx + 3x - 3b$.

\Item{7.}  $2x^{3} - x^{2} + 4x - 2$.

\Item{8.}  $a^{2} - 3a - ab + 3b$.

\Item{9.}  $6a^{2} + 2ab - 3ac - bc$.

\Item{10.}  $abxy + cxy + abc + c^{2}$.

\Item{11.}  $ax - ay - bx + cy - cx + by$.

\Item{12.} $(a - b)^{2} - 2c(a - b)$.
\end{multicols}
%% -----File: 080.png---Folio 74-------


\Section{Case III.}

\Paragraph{108.  When a binomial is the difference of two squares.}

\Item{1.}  Resolve into factors $x^{2} - y^{2}$.
\begin{DPgather*}
\lintertext{\indent Since,}
(x + y)(x - y) = x^{2} - y^{2},
\end{DPgather*}
the factors of $x^{2} - y^{2}$ are $x + y$ and~$x - y$.

\Dictum{To find the factors of a binomial when it is the difference of
two squares}, therefore,
\begin{Theorem}
Take the square root of the first term and the square root
of the second term.

The sum of these roots will form the first factor;

The difference of these roots will form the second factor.
\end{Theorem}

\Paragraph{109.}  The \Defn{square root} of a \emph{monomial} is one of the \textbf{two equal
factors} of the monomial.

Thus $9x^{8}y^{2} = 3x^{4}y × 3x^{4}y$; and $3x^{4}y$ is the square root
of~$9x^{8}y^{2}$.

The rule for extracting the square root of a \emph{monomial,
when a perfect square}, is as follows:
\begin{Theorem}
Extract the square root of the coefficient, and divide the
index of each letter by~$2$.
\end{Theorem}

\Exercise{33.}

Resolve into factors:
\begin{multicols}{3}
\Item{1.}  $4 - x^{2}$.

\Item{2.}  $9 - x^{2}$.

\Item{3.}  $9a^{2} - x^{2}$.

\Item{4.}  $25 - x^{2}$.

\Item{5.}  $25x^{2} - a^{2}$.

\Item{6.}  $16a^{4} - 121$.

\Item{7.}  $121a^{4} - 16$.

\Item{8.}  $4a^{2}b^{2} - c^{2}d^{2}$.

\Item{9.}  $1 - x^{2}y^{2}$.

\Item{10.}  $81x^{2}y^{2} - 1$.

\Item{11.}  $49a^{2}b^{2} - 4$.

\Item{12.}  $25a^{4}b^{4} - 9$.
\end{multicols}
%% -----File: 081.png---Folio 75-------

\begin{multicols}{2}
\Item{13.}  $9a^{8}b^{6} - 16x^{10}$.

\Item{14.}  $144x^{2}y^{2} - 1$.

\Item{15.}  $100x^{6}y^{2}z^{4} - 1$.

\Item{16.}  $1 - 121a^{4}b^{8}c^{12}$.

\Item{17.}  $25a^{2} - 64x^{6}y^{6}$.

\Item{18.}  $16x^{16}-25y^{18}$.
\end{multicols}

Find, by resolving into factors, the value of:
\begin{multicols}{2}
\Item{19.}  $(375)^{2} - (225)^{2}$.

\Item{20.}  $(579)^{2} - (559)^{2}$.

\Item{21.}  $(873)^{2} - (173)^{2}$.

\Item{22.}  $(101)^{2} - (99)^{2}$.

\Item{23.}  $(7244)^{2} - (7242)^{2}$.

\Item{24.}  $(3781)^{2} - (219)^{2}$.
\end{multicols}

\Paragraph{110.}  If the squares are compound expressions, the same
method may be employed.

\Item{1.}  Resolve into factors $(x + 3y)^{2} - 16a^{2}$.
\begin{Soln}
The square root of the first term is~$x + 3y$.

The square root of the second term is~$4a$.

The sum of these roots is~$x + 3y - 4a$.

The difference of these roots is $x + 3y - 4a$.

Therefore $(x + 3y)^{2} - 16a^{2} = (x + 3y + 4a)(x + 3y - 4a)$.
\end{Soln}

\Item{2.}  Resolve into factors $a^{2} - (3b - 5c)^{2}$.
\begin{Soln}
The square roots of the terms are $a$~and~$(3b - 5c)$.

The sum of these roots is $a + (3b - 5c)$, or $a + 3b - 5c$.

The difference of these roots is $a - (3b - 5c)$, or $a - 3b + 5c$.

Therefore $a^{2} - (3b - 5c)^{2} = (a + 3b - 5c)(a - 3b + 5c)$.
\end{Soln}

\Exercise{34.}

\DPtypo{}{Resolve into factors:}
\begin{multicols}{2}
\Item{1.}  $(x + y)^{2} - z^{2}$.

\Item{2.}  $(x - y)^{2} - z^{2}$.

\Item{3.}  $z^{2} - (x + y)^{2}$.

\Item{4.}  $z^{2} - (x - y)^{2}$.

\Item{5.}  $(x + y)^{2} - 4z^{2}$.

\Item{6.}  $4z^{2} - (x - y)^{2}$.

\Item{7.}  $(a + 2b)^{2} - c^{2}$.

\Item{8.}  $(a - 2b)^{2} - c^{2}$.

\Item{9.}  $c^{2} - (a - 2b)^{2}$.

\Item{10.}  $(2a + 5c)^{2} - 1$.
%% -----File: 082.png---Folio 76-------

\Item{11.}  $1 - (2a - 5c)^{2}$.

\Item{12.}  $(a + 3b)^{2} - 16c^{2}$.

\Item{13.}  $(a - 5b)^{2} - 9c^{2}$.

\Item{14.}  $16c^{2} - (a - 5b)^{2}$.

\Item{15.}  $4a^{2} - (x + y)^{2}$.

\Item{16.}  $b^{2} - (a - 2x)^{2}$.

\Item{17.}  $4z^{2} - (x + 3y)^{2}$.

\Item{18.}  $9 - (3a - 7b)^{2}$.

\Item{19.}  $16a^{2} - (2b + 5c)^{2}$.

\Item{20.}  $25c^{2} - (3a - 2x)^{2}$.

\Item{21.}  $9a^{2} - (3b - 5c)^{2}$.

\Item{22.}  $16y^{2} - (a - 3c)^{2}$.

\Item{23.}  $49m^{2} - (p + 2q)^{2}$.

\Item{24.}  $36n^{2} - (d - 2c)^{2}$.

\Item{25.}  $(x + y)^{2} - (a + b)^{2}$.

\Item{26.}  $(x - y)^{2} - (a - b)^{2}$.

\Item{27.}  $(2x + 3)^{2} - (2a + b)^{2}$.

\Item{28.}  $(b - c)^{2} - (a - 2x)^{2}$.

\Item{29.}  $(3x - y)^{2} - (2a - b)^{2}$.

\Item{30.}  $(x - 3y)^{2} - (a + 2b)^{2}$.

\Item{31.}  $(x + 2y)^{2} - (a + 3b)^{2}$.

\Item{32.}  $(x + y)^{2} - (a - z)^{2}$.
\end{multicols}


\Section{Case IV.}

\Paragraph{111.  When a binomial is the difference of two cubes.}
\begin{DPgather*}
\lintertext{\indent Since}
\frac{a^{3} - b^{3}}{a - b} = a^{2} + ab + b^{2},
\end{DPgather*}
the factors of $a^{3} - b^{3}$ are $a - b$ and $a^{2} + ab + b^{2}$.

In like manner we can resolve into factors any expression
which can be written as the difference of two cubes.

\Paragraph{112.}  The rule for extracting the cube root of a \emph{monomial,
when the monomial is a perfect cube}, is,
\begin{Theorem}
Extract the cube root of the coefficient, and divide the index
of each letter by~$3$.
\end{Theorem}

\ScreenBreak
\Item{1.}  Resolve into factors $8a^{3} - 27b^{6}$.

Since $8a^{3} = (2a)^{3}$, and $27b^{6} = (3b^{2})^{3}$, we can write
$8a^{3} - 27b^{6}$ as $(2a)^{3} - (3b^{2})^{3}$.
%% -----File: 083.png---Folio 77-------
\begin{DPgather*}
\lintertext{\indent Since}
a^{3} - b^{3} = (a - b)(a^{2} + ab + b^{2}),
\end{DPgather*}
we have, by putting $2a$ for~$a$ and $3b^{2}$ for~$b$,
\begin{align*}
(2a)^{3} - (3b^{2})^{3}
  &= (2a - 3b^{2})[(2a)^{2} + 2a × 3b^{2} + (3b^{2})^{2}] \\
  &= (2a - 3b^{2})(4a^{2} + 6ab^{2} + 9b^{4}).
\end{align*}

\Item{2.}  Resolve into factors $64x^{3} - 1$.
\begin{align*}
64x^{3} - 1
  &= (4x)^{3} - 1 \\
  &= (4x - 1)[(4x)^{2} + 4x +1] \\
  &= (4x - 1)(16x^{2} + 4x + 1).
\end{align*}

\Dictum{To find the factors of a binomial when it is the difference of
two cubes}, therefore,
\begin{Theorem}
Take the difference of the cube roots of the terms for one
factor, and the sum of the squares of the cube roots of the
terms plus their product for the other factor.
\end{Theorem}

\ScreenBreak
\Exercise{35.}

Resolve into factors:
\begin{multicols}{3}
\Item{1.}  $8x^{3} - y^{3}$.

\Item{2.}  $x^{3} - 1$.

\Item{3.}  $x^{3}y^{3} - z^{3}$.

\Item{4.}  $x^{3} - 64$.

\Item{5.}  $125a^{3} - b^{3}$.

\Item{6.}  $a^{3} - 343$.

\Item{7.}  $a^{3}b^{3} - 27c^{3}$.

\Item{8.}  $x^{3}y^{3}z^{3} - 8$.

\Item{9.}  $8a^{3}b^{3} - 27y^{6}$.

\Item{10.}  $64x^{3} - y^{9}$.

\Item{11.}  $27a^{3} - 64c^{6}$.

\Item{12.}  $x^{3}y^{3} - 216z^{3}$.

\Item{13.}  $64x^{3} - 729y^{3}$.

\Item{14.}  $27a^{3} - 512c^{3}$.

\Item{15.}  $8x^{6} - 125y^{3}$.

\Item{16.}  $64x^{12} - 27y^{15}$.

\Item{17.}  $216 - 8a^{3}$.

\Item{18.}  $343 - 27y^{3}$.
\end{multicols}


\Section{Case V.}

\Paragraph{113.  When a binomial is the sum of two cubes.}
\begin{DPgather*}
\lintertext{\indent Since}
\frac{a^{3} + b^{3}}{a + b} = a^{2} - ab + b^{2},
\end{DPgather*}
the factors of $a^{3} + b^{3}$ are $a + b$ and $a^{2} - ab + b^{2}$.
%% -----File: 084.png---Folio 78-------

In like manner we can resolve into factors any expression
which can be written as the sum of two cubes.

\Item{1.}  Resolve into factors $8x^{3} + 27y^{3}$.

Since by §~112, $8x^{3} = (2x)^{3}$ and $27y^{3} = (3y)^{3}$, we can
write $8x^{3} + 27y^{3}$ as $(2x)^{3} + (3y)^{3}$.

\begin{DPgather*}
\lintertext{\indent Since}
a^{3} + b^{3} = (a + b)(a^{2} - ab + b^{2}),
\end{DPgather*}
we have, by putting $2x$ for~$a$, and $3y$ for~$b$,
\begin{align*}
(2x)^{3} + (3y)^{3}
  &= (2x + 3y)[(2x)^{2} - 2x × 3y + (3y)^{2}] \\
  &= (2x + 3y)(4x^{2} - 6xy + 9y^{2}).
\end{align*}

\Item{2.}  Resolve into factors $125a^{3} + 64x^{6}$\Add{.}
\begin{gather*}
125a^{3} = (5a)^{3},\quad 64x^{6} = (4x^{2})^{3}; \\
\begin{aligned}
\therefore 125a^{3} + 64x^{6}
  &= (5a + 4x^{2})[(5a)^{2} - 5a × 4x^{2} + (4x^{2})^{2}] \\
  &= (5a + 4x^{2})(25a^{2} - 20ax^{2} + 16x^{4})
\end{aligned}
\end{gather*}

\Dictum{To find the factors of a binomial when it is the sum of two
cubes}, therefore,
\begin{Theorem}
Take the sum of the cube roots of the terms for one factor,
and the sum of the squares of the cube roots of the terms
minus their product for the other factor.
\end{Theorem}

\PrintBreak
\Exercise{36.}

Resolve into factors:
\begin{multicols}{3}
\Item{1.}  $x^{3} + 1$.

\Item{2.}  $8x^{3} + y^{3}$.

\Item{3.}  $x^{3} + 125$.

\Item{4.}  $64a^{3} + 27$.

\Item{5.}  $x^{3}y^{3} + z^{3}$.

\Item{6.}  $a^{3} + 64$.

\Item{7.}  $8a^{6} + b^{3}$.

\Item{8.}  $x^{3} + 343$.

\Item{9.}  $8 + x^{3}y^{3}z^{3}$.

\Item{10.}  $y^{9} + 64x^{3}$.

\Item{11.}  $a^{3}b^{3} + 27x^{3}$.

\Item{12.}  $8y^{3}z^{3} + x^{6}$.

\Item{13.}  $y^{9} + 64x^{6}$.

\Item{14.}  $64a^{12} + x^{15}$.

\Item{15.}  $27x^{15} + 8a^{6}$.

\Item{16.}  $27x^{9} + 512$.

\Item{17.}  $343 + 64x^{3}$.

\Item{18.}  $125 + 27y^{3}$.
\end{multicols}
%% -----File: 085.png---Folio 79-------


\Section{Case VI.}

\Paragraph{114.  When a trinomial is a perfect square.}
\begin{DPgather*}
\lintertext{\indent Since}
(x + y)^{2} = x^{2} + 2xy + y^{2},
\end{DPgather*}
the factors of $x^{2} + 2xy + y^{2}$ are $x + y$ and $x + y$.

\begin{DPgather*}
\lintertext{\indent Since}
(x - y)^{2} = x^{2} - 2xy + y^{2},
\end{DPgather*}
the factors of $x^{2} - 2xy + y^{2}$ are $x - y$ and $x - y$.

Therefore, a trinomial is a perfect square, if its first and
last terms are perfect squares and positive, and its middle
term is twice the product of their square roots.

\Dictum{To find the factors of a trinomial when it is a perfect square},
therefore,
\begin{Theorem}
Extract the square roots of the first and last terms, and
connect these square roots by the sign of the middle term.
\end{Theorem}

Thus, if we wish to find the square root of
\[
16a^{2} - 24ab + 9b^{2},
\]
we take the square roots of $16a^{2}$ and $9b^{2}$, which are $4a$
and~$3b$, respectively, and connect these square roots by
the minus sign, the sign of the middle term.  The square
root is therefore
\[
4a - 3b.
\]

Again, if we wish to find the square root of
\[
25x^{2} + 40xy + 16y^{2},
\]
we take the square roots of $25x^{2}$ and $16y^{2}$ and connect these
roots by the plus sign, the sign of the middle term.  The
square root is therefore
\[
5x + 4y.
\]
%% -----File: 086.png---Folio 80-------

\Exercise{37.}

Resolve into factors:
\begin{multicols}{2}
\Item{1.}  $4x^{2} + 4xy + y^{2}$.

\Item{2.}  $x^{2} + 6xy + 9y^{2}$.

\Item{3.}  $x^{2} + 16x + 64$.

\Item{4.}  $x^{2} + 10ax + 25a^{2}$.

\Item{5.}  $a^{2} - 16a + 64$.

\Item{6.}  $a^{2} - 10ab + 25b^{2}$.

\Item{7.}  $c^{2} - 6cd + 9d^{2}$.

\Item{8.}  $4x^{2} - 4x + 1$.

\Item{9.}  $4a^{2} - 12ab + 9b^{2}$.

\Item{10.}  $9a^{2} - 24ab + 16b^{2}$.

\Item{11.}  $x^{2} + 8xy + 16y^{2}$.

\Item{12.}  $x^{2} - 8xy + 16y^{2}$.

\Item{13.}  $4x^{2} - 20xy + 25y^{2}$.

\Item{14.}  $1 + 20a + 100a^{2}$.

\Item{15.}  $49a^{2} - 28a + 4$.

\Item{16.}  $36a^{2} + 60ab + 25b^{2}$.

\Item{17.}  $81x^{2} - 36bx + 4b^{2}$.

\Item{18.}  $m^{2}n^{2} + 14mnx^{2} + 49x^{2}$.
\end{multicols}


\PrintBreak
\Section{Case VII.}

\Paragraph{115.  When a trinomial has the form $x^{2} + ax + b$.}

Where $a$ is the \emph{algebraic sum} of two numbers, and is
either positive or negative; and $b$~is the \emph{product} of these
two numbers, and is either positive or negative.
\begin{DPgather*}
\lintertext{\indent Since}
(x + 5)(x + 3) = x^{2} + 8x + 15,
\end{DPgather*}
the factors of $x^{2} + 8x  + 15$ are $x + 5$ and $x + 3$.
\begin{DPgather*}
\lintertext{\indent Since}
(x + 5)(x - 3) = x^{2} + 2x- 15,
\end{DPgather*}
the factors of $x^{2} + 2x - 15$ are $(x + 5)$ and $(x - 3)$.

Hence, if a trinomial of the form $x^{2} + ax + b$ is such an
expression that it can be resolved into two binomial factors,
it is obvious that the first term of each factor will be~$x$,
and that the second terms of the factors will be two
numbers whose product is~$b$, the last term of the trinomial,
and whose algebraic sum is~$a$, the coefficient of~$x$ in the
middle term of the trinomial.
%% -----File: 087.png---Folio 81-------

\Item{1.}  Resolve into factors $x^{2} + 11x + 30$.
\begin{Soln}
We are required to find two numbers whose product is~$30$ and
whose sum is~$11$.

Two numbers whose product is $30$ are $1$~and~$30$, $2$~and~$15$, $3$~and~$10$,
$5$~and~$6$, and the sum of the last two numbers is~$11$.  Hence,
\[
x^{2} + 11x + 30 = (x + 5)(x + 6).
\]
%[** TN: Solutions sometimes printed in normal-size type; using smaller type]
\end{Soln}

\Item{2.}  Resolve into factors $x^{2} - 7x + 12$.
\begin{Soln}
We are required to find two numbers whose product is~$12$ and
whose algebraic sum is~$-7$.

Since the product is~$+12$, the two numbers are \emph{both positive} or \emph{both
negative}, and since their sum is~$-7$, they must both be negative.

Two negative numbers whose product is~$12$ are $-12$~and~$-1$, $-6$
and~$-2$, $-4$~and~$-3$, and the sum of the last two numbers is~$-7$.
Hence,
\[
x^{2} - 7x + 12 = (x - 4)(x - 3).
\]
\end{Soln}

\Item{3.}  Resolve into factors $x^{2} + 2x - 24$.
\begin{Soln}
We are required to find two numbers whose product is~$-24$ and
whose algebraic sum is~$2$.

Since the product is~$-24$, one of the numbers is positive and the
other negative, and since their sum is~$+2$, the larger number is
positive.

Two numbers whose product is~$-24$, and the larger number positive,
are $24$~and~$-1$, $12$~and~$-2$, $8$~and~$-3$, $6$~and~$-4$, and the sum
of the last two numbers is~$+2$.  Hence,
\[
x^{2} + 2x - 24 = (x + 6)(x - 4).
\]
\end{Soln}

\Item{4.}  Resolve into factors $x^{2} - 3x - 18$.
\begin{Soln}
Since the product is~$-18$, one of the numbers is positive and the
other negative, and since their sum is~$-3$, the larger number is
negative.

Two numbers whose product is~$-18$, and the larger number negative,
are $-18$~and~$1$, $-9$~and~$2$, $-6$~and~$3$, and the sum of the last
two numbers is~$-3$.  Hence,
\[
x^{2} - 3x - 18 = (x - 6)(x + 3).
\]
\end{Soln}
%% -----File: 088.png---Folio 82-------

\Exercise{38.}

Resolve into factors:
\begin{multicols}{2}
\Item{1.}  $a^{2} + 5a + 6$.

\Item{2.}  $a^{2} - 5a + 6$.

\Item{3.}  $a^{2} + 6a + 5$.

\Item{4.}  $a^{2} - 6a + 5$.

\Item{5.}  $a^{2} + 4a - 5$.

\Item{6.}  $a^{2} - 4a - 5$.

\Item{7.}  $c^{2} - 9c + 18$.

\Item{8.}  $c^{2} + 9c + 18$.

\Item{9.}  $c^{2} + 3c - 18$.

\Item{10.}  $c^{2} - 3c - 18$.

\Item{11.}  $x^{2} + 9x + 14$.

\Item{12.}  $x^{2} - 9x + 14$.

\Item{13.}  $x^{2} - 5x - 14$.

\Item{14.}  $x^{2} - 9x + 20$.

\Item{15.}  $x^{2} - x - 20$.

\Item{16.}  $x^{2} + x - 20$.

\Item{17.}  $x^{2} - 10x + 21$.

\Item{18.}  $x^{2} - 4x - 21$.

\Item{19.}  $x^{2} + 4x - 21$.

\Item{20.}  $x^{2} - 15x + 56$.

\Item{21.}  $x^{2} - x - 56$.

\Item{22.}  $x^{2} - 10x + 9$.

\Item{23.}  $x^{2} + 13x + 30$.

\Item{24.}  $x^{2} + 7x - 30$.

\Item{25.}  $x^{2} - 7x - 30$.

\Item{26.}  $a^{2} + ab - 6b^{2}$.

\Item{27.}  $a^{2} - ab - 6b^{2}$.

\Item{28.}  $a^{2} + 3ab - 4b^{2}$.

\Item{29.}  $a^{2} - 3ab - 4b^{2}$.

\Item{30.}  $a^{2}x^{2} - 2ax - 63$.

\Item{31.}  $a^{2} + 2ax - 63x^{2}$.

\Item{32.}  $a^{2} - 9ab + 20b^{2}$.

\Item{33.}  $x^{2}y^{2} - 19xyz + 48z^{2}$.

\Item{34.}  $a^{2}b^{2} + 15abc + 44c^{2}$.

\Item{35.}  $x^{2} - 13xy + 36y^{2}$.

\Item{36.}  $x^{2} + 19xy + 84y^{2}$.

\Item{37.}  $a^{2}x^{2} - 23axy + 102y^{2}$.

\Item{38.}  $x^{4} - 9x^{2}y^{2} + 20y^{4}$.

\Item{39.}  $a^{4}x^{4} - 24a^{2}x^{2}y^{2} + 143y^{4}$.

\Item{40.}  $a^{6}b^{6} - 23a^{3}b^{3}c^{2} + 132c^{4}$.

\Item{41.}  $a^{2} - 20abc - 96b^{2}c^{2}$.

\Item{42.}  $a^{2} - 4abc - 96b^{2}c^{2}$.

\Item{43.}  $a^{2} - 10abc - 96b^{2}c^{2}$.

\Item{44.}  $a^{2} + 29abc - 96b^{2}c^{2}$.

\Item{45.}  $a^{2} - 46abc - 96b^{2}c^{2}$.

\Item{46.}  $a^{2} + 49abc + 48b^{2}c^{2}$.

\Item{47.}  $x^{2} - 18xyz - 243y^{2}z^{2}$.

\Item{48.}  $x^{2}y^{2} - xyz - 182z^{2}$.
\end{multicols}
%% -----File: 089.png---Folio 83-------

\Exercise{39.}

\Section{Examples For Review.}

Resolve into factors:
\begin{multicols}{2}
\Item{1.}  $a^{3} - 7a$.

\Item{2.}  $3a^{2}b^{2} - 2a^{3}b + 3ab^{3}$.

\Item{3.}  $(a - b)^{2} + (a - b)$.

\Item{4.}  $(a + b)^{2} - 1$.

\Item{5.}  $a^{3} + 8b^{3}$.

\Item{6.}  $(x^{2} - 4y^{2}) + (x - 2y)$.

\Item{7.}  $(a^{3} - b^{3}) + (a - b)$.

\Item{8.}  $a^{2} - 6ab + 9b^{2}$.

\Item{9.}  $x^{2} - x -2$.

\Item{10.}  $x^{2} - 2x - 3$.

\Item{11.}  $x^{2} + 4x - 21$.

\Item{12.}  $a^{2} - 11a - 26$.

\Item{13.}  $ax^{2} + bx^{2} + 3a + 3b$.

\Item{14.}  $x^{2} - 3x- xy + 3y$.

\Item{15.}  $x^{2} - 7x + 12$.

\Item{16.}  $a^{2} + 5ab + 6b^{2}$.

\Item{17.}  $x^{4} + 10x^{2} + 25$.

\Item{18.}  $x^{2} - 18x + 81$.

\Item{19.}  $x^{2} - 21x + 110$.

\Item{20.}  $x^{2} + 19x + 88$.

\Item{21.}  $x^{2} - 19x + 88$.

\Item{22.}  $x^{3} - x^{2} + x - 1$.

\Item{23.}  $9x^{4} - x^{2}$.

\Item{24.}  $1 - (a - b)^{2}$.

\Item{25.}  $(a^{3} + b^{3}) + (a + b)$.

\Item{26.}  $m^{2}x - n^{2}x + m^{2}y - n^{2}y$.

\Item{27.}  $(x - y)^{2} - z^{2}$.

\Item{28.}  $z^{2} - (x - y)^{2}$.

\Item{29.}  $4a^{4} - (3a - 1)^{2}$.

\Item{30.}  $8x^{3} - y^{3}$.

\Item{31.}  $x^{3} - 3x^{2}y$.

\Item{32.}  $x^{3} - 27y^{3}$.

\Item{33.}  $x^{2} + 3x - 40$.

\Item{34.}  $x^{2} + 3xy - 10y^{2}$.

\Item{35.}  $1 - 16x^{2}$.

\Item{36.}  $a^{6} - 9a^{2}b^{4}$.

\Item{37.}  $x^{3} + 3x^{2}y + 2xy^{2}$.

\Item{38.}  $x^{4} + 4x^{3}y + 3x^{2}y^{2}$.

\Item{39.}  $x^{2} - 4xy^{2} + 4y^{4}$.

\Item{40.}  $16x^{4} + 8x^{2} + 1$.

\Item{41.}  $9a^{4} - 4a^{2}c^{2}$.

\Item{42.}  $a^{3}b - a^{2}b^{2} - 2ab^{3}$.

\Item{43.}  $x^{4} - x^{3} + 8x - 8$.

\Item{44.}  $a^{4} - a^{3}x + ay^{3} - xy^{3}$.
\end{multicols}
%% -----File: 090.png---Folio 84-------


\Chapter{VIII.}{Common Factors and Multiples.}

\Paragraph{116.  Common Factors.}  A \Defn{common factor} of two or more
\emph{integral numbers} is an integral number which divides each
of them without a remainder.

\Paragraph{117.}  A \Defn{common factor} of two or more integral and rational
\emph{expressions} is an integral and rational expression which
divides each of them without a remainder.

\begin{Remark}
Thus $5a$~is a common factor of $20a$~and~$25a$, $3x^{2}y^{2}$ is a common
factor of~$12x^{2}y^{2}$ and~$15x^{3}y^{3}$.
\end{Remark}

\Paragraph{118.}  Two \emph{numbers} are said to be \Defn{prime} to each other
when they have no common factor except~$1$.

\Paragraph{119.}  Two \emph{expressions} are said to be \Defn{prime} to each other
when they have no common factor except~$1$.

\Paragraph{120.}  The \Defn{highest common factor} of two or more integral
\emph{numbers} is the greatest number that will divide each of
them without a remainder.

\Paragraph{121.}  The \Defn{highest common factor} of two or more integral
and rational \emph{expressions} is an integral and rational expression
of highest degree that will divide each of them without
a remainder.

\begin{Remark}
Thus $3a^{2}$ is the highest common factor of $3a^{2}$, $6a^{3}$, and~$12a^{4}$,
$5x^{2}y^{2}$ is the highest common factor of $10x^{3}y^{2}$ and~$15x^{2}y^{2}$.
\end{Remark}

For brevity, we use \HCF\ for ``highest common factor.''
%% -----File: 091.png---Folio 85-------

\Paragraph{122.  To Find the Highest Common Factor of Two or More
Algebraic Expressions.}

\Item{1.}  Find the \HCF\ of $42a^{3}b^{2}$ and~$30a^{2}b^{4}$.
\begin{alignat*}{2}
&42a^{3}b^{2} &&= 2 × 3 × 7 × aaa × bb; \\
&30a^{2}b^{4} &&= 2 × 3 × 5 × aa × bbbb. \\
\therefore\ &\text{the \HCF} &&= 2 × 3 × aa × bb,
\quad\text{or}\quad 6a^{2}b^{2}.
\end{alignat*}

\Item{2.}  Find the \HCF\ of $x^{2} - 9y^{2}$ and $x^{2} + 6xy + 9y^{2}$.
\begin{alignat*}{2}
&x^{2} - 9y^{2} &&= (x + 3y)(x - 3y); \\
&x^{2} + 6xy + 9y^{2} &&= (x + 3y)(x + 3y). \\
\therefore\ &\text{the \HCF} &&= (x + 3y).
\end{alignat*}

\Item{3.}  Find the \HCF\ of $4x^{2} - 4x - 80$, $2x^{2} - 18x + 40$.
\begin{alignat*}{3}
&4x^{2} -& 4x &- 80 &&= 4(x^{2} - x-20) \\
&&&                 &&= 4(x - 5)(x + 4); \\
&2x^{2} -& 18x &+ 40 &&= 2(x^{2} - 9x + 20) \\
&&&                  &&= 2(x - 5)(x - 4). \\
\therefore\ &\rlap{\text{the \HCF}} &&&&= 2(x - 5).
\end{alignat*}

\Dictum{To find the \HCF\ of two or more expressions}, therefore,
\begin{Theorem}
Resolve each expression into its simplest factors.

Find the product of all the common factors, taking each
factor the least number of times it occurs in any of the given
expressions.
\end{Theorem}

\begin{Remark}[Note.] The \emph{highest common factor} in Algebra corresponds to the
\emph{greatest common measure}, or \emph{greatest common divisor} in Arithmetic.
We cannot apply the terms \emph{greatest} and \emph{least} to an algebraic expression
in which particular values have not been given to the letters
contained in the expression.  Thus $a$~is \emph{greater} than~$a^{2}$, if $a$~stands
for~$\frac{1}{4}$.
\end{Remark}
%% -----File: 092.png---Folio 86-------

\Exercise{40.}

Find the \HCF\ of:
\begin{multicols}{2}
\Item{1.}  $330$ and $546$.

\Item{2.}  $20x^{3}$ and $15x^{4}$.

\Item{3.}  $42ax^{2}$ and $60a^{2}x$.

\Item{4.}  $35a^{2}b^{2}$ and $49ab^{3}$.

\Item{5.}  $28x^{4}$ and $63y^{4}$.

\Item{6.}  $54a^{2}b^{2}$ and $56a^{3}b^{3}$.

\Item{7.}  $x^{3} + 3x^{2}y$ and $x^{3} + 27y^{3}$.

\Item{8.}  $x^{2} + 3x$ and $x^{2} - 9$.

\Item{9.}  $2ax^{3} + x^{3}$ and $8a^{3} + 1$.

\Item{10.}  $(x + y)^{2}$ and $x^{2} - y^{2}$.

\Item{11.}  $a^{3} + a^{2}x$ and $a^{2} - x^{2}$.

\Item{12.}  $a^{2} - 4b^{2}$ and $a^{2} + 2ab$.
\end{multicols}

\Item{13.}  $x^{2} - 1$ and $x^{2} + 2x - 3$.

\Item{14.}  $x^{2} + 5x + 6$ and $x^{2} + 4x + 3$.

\Item{15.}  $x^{2} - 9x + 18$ and $x^{2} - 10x + 24$.

\Item{16.}  $x^{3} + 1$ and $x^{2} - x + 1$.

\Item{17.}  $x^{2} - 3x + 2$ and $x^{2} - 4x + 3$.

\Item{18.}  $x^{2} - 3xy + 2y^{2}$ and $x^{2} - 2xy + y^{2}$.

\Item{19.}  $x^{2} - 4x - 5$ and $x^{2} - 25$.

\Item{20.}  $(a - b)^{2} - c^{2}$ and $ab - b^{2} - bc$.

\Item{21.}  $x^{2} + xy -2y^{2}$ and $x^{2} + 5xy + 6y^{2}$.

\Item{22.}  $x^{2} + 7xy + 12y^{2}$ and $x^{2} + 3xy - 4y^{2}$.

\Item{23.}  $x^{3} - 8y^{3}$ and $x^{2} + 2xy + 4y^{2}$.

\Item{24.}  $x^{3} - 2x^{2} - x + 2$ and $x^{2} - 4x + 4$.

\Item{25.}  $1 - 5a + 6a^{2}$ and $1 - 7a + 12a^{2}$.

\Item{26.}  $x^{2} - 8xy + 7y^{2}$ and $x^{2} - 3xy - 28y^{2}$.

\Item{27.}  $8a^{3} + b^{3}$ and $4a^{2} + 4ab + b^{2}$.

\Item{28.}  $x^{2} - (y - z)^{2}$ and $(x + y)^{2} - z^{2}$.
%% -----File: 093.png---Folio 87-------

\Paragraph{123.  Common Multiples.}  A \Defn{common multiple} of two or
more integral \emph{numbers} is a number which is exactly divisible
by each, of the numbers.

A \Defn{common multiple} of two or more \emph{expressions} is an expression
which is exactly divisible by each of the expressions.

\Paragraph{124.}  The \Defn{lowest common multiple} of two or more \emph{numbers}
is the least number that is exactly divisible by each of the
given numbers.

The \Defn{lowest common multiple} of two or more \emph{expressions} is
the expression of lowest degree that is exactly divisible by
each of the given expressions.

We use \LCM\ for ``lowest common multiple.''

\Paragraph{To find the lowest common multiple of two or more algebraic
expressions.}

\Item{1.}  Find the \LCM\ of $42a^{3}b^{2}$, $30a^{2}b^{4}$, and~$66ab^{3}$.
\begin{align*}
42a^{3}b^{2} &= 2 × 3 × 7 × a^{3} × b^{2}; \\
30a^{2}b^{4} &= 2 × 3 × 5 × a^{2} × b^{4}; \\
66ab^{3} &= 2 × 3 × 11 × a × b^{3}.
\end{align*}

\begin{Soln}
The \LCM\  must evidently contain each factor the greatest number
of times that it occurs in any expression.
\begin{align*}
\therefore\ \text{\LCM}
  &= 2 × 3 × 7 × 5 × 11a^{3} × b^{4}, \\
  &= 2310a^{3}b^{4}.
\end{align*}
\end{Soln}

\Item{2.}  Find the \LCM\ of $4x^{2} - 4x - 80$ and $2x^{2} - 18x + 40$.
\begin{Soln}
\begin{alignat*}{3}
&4x^{2} -&  4x &- 80 &&= 4(x^{2} - x - 20) = 4(x - 5)(x + 4); \\
&2x^{2} -& 18x &+ 40 &&= 2(x^{2} - 9x + 20) = 2(x - 5)(x - 4). \\
\therefore\ &\rlap{\text{\LCM}}&&& &= 4(x - 5)(x + 4)(x - 4).
\end{alignat*}
\end{Soln}

\Dictum{To find the \LCM\ of two or more expressions}, therefore,
%% -----File: 094.png---Folio 88-------
\begin{Theorem}
Resolve each expression into its simplest factors.

Find the product of all the different factors, taking each
factor the greatest number of times it occurs in any of the
given expressions.
\end{Theorem}

\Exercise{41.}

Find the \LCM\ of:
\begin{multicols}{2}
\Item{1.}  $9xy^{3}$ and $6x^{2}y$.

\Item{2.}  $3abc^{2}$ and $2a^{2}bc^{3}$.

\Item{3.}  $4a^{3}b$ and $10ab^{3}$.

\Item{4.}  $6a^{3}b^{3}$ and $15a^{2}b^{4}$.

\Item{5.}  $21xy^{3}$ and $27x^{3}y^{5}$.

\Item{6.}  $xy^{3}z^{2}$ and $x^{2}y^{2}z^{3}$.

\Item{7.}  $a^{2}$ and $a^{2} + a$.

\Item{8.}  $x^{2}$ and $x^{3} - 3x^{2}$.

\Item{9.}  $x^{2} - 1$ and $x^{2} + x$.

\Item{10.}  $x^{2} - 1$ and $x^{2} - x$.

\Item{11.}  $x^{2} + xy$ and $xy + y^{2}$.

\Item{12.}  $x^{2} + 2x$ and $(x + 2)^{2}$.
\end{multicols}

\Item{13.}  $a^{2} + 4a + 4$ and $a^{2} + 5a + 6$.

\Item{14.}  $c^{2} + c - 20$ and $c^{2} - c - 30$.

\Item{15.}  $b^{2} + b - 42$ and $b^{2} - 11b + 30$.

\Item{16.}  $y^{2} - 10y + 24$ and $y^{2} + y -20$.

\Item{17.}  $z^{2} + 2z - 35$ and $z^{2} - 11z + 30$.

\Item{18.}  $x^{2} - 64; x^{3} - 64$; and $x + 8$.

\Item{19.}  $a^{2} - b^{2}; (a + b)^{2}$; and $(a - b)^{2}$.

\Item{20.}  $4ab(a + b)^{2}$ and $2a^{2}(a^{2} - b^{2})$.

\Item{21.}  $y^{2} + 7y + 12; y^{2} + 6y + 8$; and $y^{2} + 5y +6$.

\Item{22.}  $x^{2} - 1; x^{3} + x^{2} + x + 1$; and $x^{3} - x^{2} + x - 1$.

\Item{23.}  $1 - x^{2}; 1 - x^{3}$; and $1 + x$.

\Item{24.}  $x^{2} + 2xy + y^{2}; x^{2} - y^{2}$; and $x^{2} - 2xy + y^{2}$.

\Item{25.}  $x^{3} - 27; x^{2} + 2x - 15; x^{2} + 5x$.

\Item{26.}  $(a + b)^{2} - c^{2}; (a + b + c)^{2}$; and $a + b - c$.

\Item{27.}  $x^{2} - (a + b)x + ab$ and $x^{2} - (a + c)x + ac$.

\Item{28.}  $(a + b)^{2} - c^{2}$ and $a^{2} + ab + ac$.
%% -----File: 095.png---Folio 89-------


\Chapter{IX.}{Fractions.}

\Paragraph{125.}  An \Defn{algebraic fraction} is the indicated quotient of
two expressions, written in the form~$dfrac{a}{b}$.

The dividend~$a$ is called the \Defn{numerator}, and the divisor~$b$
is called the \Defn{denominator}; and the numerator and denominator
are called the \Defn{terms} of the fraction.

\Paragraph{126.}  The introduction of the same factor into the dividend
and divisor does not alter the value of the quotient,
and the rejection of the same factor from the dividend and
divisor does not alter the value of the quotient.
\begin{DPgather*}
\lintertext{\indent Thus}
\frac{12}{4} = 3;\quad
\frac{2 × 12}{2 × 4} = 3; \frac{12÷2}{4÷2} = 3. \EqText{Hence,}
\end{DPgather*}
\begin{Theorem}
The value of a fraction is not altered if the numerator and
denominator are both multiplied, or both divided, by the
same factor.
\end{Theorem}


\Section{Reduction of Fractions.}

\Paragraph{127.}  To reduce a fraction is to change its \emph{form} without
altering its \emph{value}.


\Section{Case I.}

\Paragraph{128.  To Reduce a Fraction to its Lowest Terms.}

A fraction is in its \emph{lowest terms} when the numerator and
denominator have no common factor.  We have, therefore,
the following rule:
%% -----File: 096.png---Folio 90-------
\begin{Theorem}
Resolve the numerator and denominator into their prime
factors, and cancel all the common factors.
\end{Theorem}

Reduce the following fractions to their lowest terms:

\Item{1.}  $\dfrac{38a^{2}b^{3}c^{4}}{57a^{3}bc^{2}} = \dfrac{2 × 19a^{2}b^{3}c^{4}}{3 × 19a^{3}bc^{2}} = \dfrac{2b^{2}c^{2}}{3a}$.

\Item{2.}  $\dfrac{a^{3} - x^{3}}{a^{2} - x^{2}} = \dfrac{(a - x)(a^{2} + ax + x^{2})}{(a - x)(a + x)} = \dfrac{a^{2} + ax + x^{2}}{a + x}$.

\Item{3.}  $\dfrac{a^{2} + 7a + 10}{a^{2} + 5a + 6} = \dfrac{(a + 5)(a + 2)}{(a + 3)(a + 2)} = \dfrac{a + 5}{a + 3}$.

\Exercise{42.}

Reduce to lowest terms:
\begin{multicols}{3}
\Item{1.}  $\dfrac{2a}{6ab}$.

\Item{2.}  $\dfrac{12m^{2}n}{15mn^{2}}$.

\Item{3.}  $\dfrac{21m^{2}p^{2}}{28mp^{4}}$.

\Item{4.}  $\dfrac{3x^{3}y^{2}z}{6xy^{3}z^{2}}$.

\Item{5.}  $\dfrac{5a^{3}b^{3}c^{3}}{15c^{5}}$.

\Item{6.}  $\dfrac{34x^{3}y^{4}z^{5}}{51x^{2}y^{3}z^{5}}$.

\Item{7.}  $\dfrac{46m^{2}np^{3}}{69mnp^{4}}$.

\Item{8.}  $\dfrac{39a^{2}b^{3}c^{4}}{52a^{5}bc^{3}}$.

\Item{9.}  $\dfrac{58xy^{4}z^{6}}{87xy^{2}z^{2}}$.
\end{multicols}

\begin{multicols}{2}
\Item{10.}  $\dfrac{abx - bx^{2}}{acx - cx^{2}}$.

\Item{11.}  $\dfrac{4a^{2} - 9b^{2}}{4a^{2} + 6ab}$.

\Item{12.}  $\dfrac{3a^{2} + 6a}{a^{2} + 4a + 4}$.

\Item{13.}  $\dfrac{x^{2} + 5x}{x^{2} + 4x-5}$.

\Item{14.}  $\dfrac{xy - 3y^{2}}{x^{3} - 27y^{3}}$.

\Item{15.}  $\dfrac{x^{2} + 5x + 4}{x^{2} - x-20}$.

\Item{16.}  $\dfrac{x^{2} + 2x + 1}{x^{2} - x-2}$.

\Item{17.}  $\dfrac{(a + b)^{2} - c^{2}}{a^{2} + ab-ac}$.

\Item{18.}  $\dfrac{x^{2} + 9x + 20}{x^{2} + 7x + 12}$.

\Item{19.}  $\dfrac{x^{2} - 14x - 15}{x^{2} - 12x - 45}$.
\end{multicols}
%% -----File: 097.png---Folio 91-------


\PrintBreak
\Section{Case II.}

\Paragraph{129.  To Reduce a Fraction to an Integral or Mixed Expression.}

\Item{1.}  Reduce $\dfrac{x^{3} - 1}{x - 1}$ to an integral or mixed expression.

By division, $\dfrac{x^{3} - 1}{x - 1} = x^{2} + x + 1$.

\Item{2.}  Reduce $\dfrac{x^{3} - 1}{x + 1}$ to an integral or mixed expression.

By division, $\dfrac{x^{3} - 1}{x + 1} = x^{2} - x + 1 - \dfrac{2}{x + 1}$.

\ScreenBreak
\Dictum{To reduce a fraction to an integral or mixed expression},
therefore,
\begin{Theorem}
Divide the numerator by the denominator.
\end{Theorem}

\begin{Remark}[Note.] If there is a remainder, this remainder must be written
as the numerator of a fraction of which the divisor is the denominator,
and this fraction with its proper sign must be annexed to the integral
part of the quotient.
\end{Remark}

\Exercise{43.}

Reduce to integral or mixed expressions:
\begin{multicols}{2}
\Item{1.}  $\dfrac{a^{2} - b^{2} + 2}{a - b}$.

\Item{2.}  $\dfrac{a^{2} - b^{2} - 2}{a + b}$.

\Item{3.}  $\dfrac{a^{3} - 2a^{2} + 2a + 1}{a^{2} - a - 1}$.

\Item{4.}  $\dfrac{2x^{2} - 2x + 1}{x + 1}$.

\Item{5.}  $\dfrac{8x^{3}}{2x + 1}$.

\Item{6.}  $\dfrac{5x^{3} + 9x^{2} + 3}{x^{2} + x - 1}$.

\Item{7.}  $\dfrac{a^{3} + a^{2} + 7a - 2}{a^{2} + a + 2}$.

\Item{8.}  $\dfrac{y^{4} + y^{2}x^{2} + x^{4}}{y^{2} + yx + x^{2}}$.

\Item{9.}  $\dfrac{x^{4} - 3x^{3} + x - 1}{x^{2} + x + 1}$.

\Item{10.}  $\dfrac{x^{5} - x^{4} + 1}{x^{2} - x - 1}$.
\end{multicols}
%% -----File: 098.png---Folio 92-------


\ScreenBreak
\Section{Case III.}

\Paragraph{130.  To Reduce a Mixed Expression to a Fraction.}

The process is precisely the same as in Arithmetic.  Hence,
\begin{Theorem}
Multiply the integral expression by the denominator, to
the product add the numerator, and under the result write
the denominator.
\end{Theorem}

Reduce to a fraction $a - b - \dfrac{a^{2} - ab - b^{2}}{a + b}$.
\begin{align*}
a - b - \frac{a^{2} - ab - b^{2}}{a + b}
  &= \frac{(a - b)(a + b) - (a^{2} - ab - b^{2})}{a + b} \\
  &= \frac{a^{2} - b^{2} - a^{2} + ab + b^{2}}{a + b} \\
  &= \frac{ab}{a + b}.
\end{align*}

\begin{Remark}[Note.] The dividing line between the terms of a fraction has the
force of a vinculum affecting the numerator.  If, therefore, a \emph{minus
sign} precedes the dividing line, as in the preceding Example, and
this line is removed, the numerator of the given fraction must be
enclosed in a parenthesis preceded by the minus sign, or the sign of
every term of the numerator must be changed.
\end{Remark}

\ScreenBreak
\Exercise{44.}

Reduce to a fraction:
\begin{multicols}{2}
\Item{1.}  $x - y + \dfrac{2xy}{x - y}$.

\Item{2.}  $x + y - \dfrac{2xy}{x + y}$.

\Item{3.}  $1 - \dfrac{x - y}{x + y}$.

\Item{4.}  $a - x - \dfrac{a^{2} + x^{2}}{a - x}$.

\Item{5.}  $x + 2 - \dfrac{x^{2} - 4}{x - 3}$.

\Item{6.}  $\dfrac{x - 3}{x - 2} - 2x + 1$.

\Item{7.}  $\dfrac{x + 3}{x + 2} + x^{2} - x - 1$.

\Item{8.}  $2a - 1 + \dfrac{3 - 4a}{a - 3}$.

\Item{9.}  $1 - 2a^{2} - \dfrac{a^{2} - a + 2}{a - 1}$.

\Item{10.}  $a^{2} + 2a - 5 - \dfrac{2a - 1}{3a^{2} + 1}$.
\end{multicols}
%% -----File: 099.png---Folio 93-------


\Section{Case IV.}

\Paragraph{131.  To Reduce Fractions to their Lowest Common Denominator.}

The process is the same as in Arithmetic.  Hence:
\begin{Theorem}
Find the lowest common multiple of the denominators;
this will be the required denominator.  Divide this denominator
by the denominator of each fraction.

Multiply the first numerator by the first quotient, the second
numerator by the second quotient, and so on.

The products will be the respective numerators of the
equivalent fractions.
\end{Theorem}

\begin{Remark}[Note.] Every fraction should be in its lowest terms before the
common denominator is found.
\end{Remark}

\Item{1.}  Reduce $\dfrac{3x}{4a^{2}}$, $\dfrac{2y}{3a}$, and $\dfrac{5}{6a^{3}}$ to equivalent fractions
having the lowest common denominator.
\begin{Soln}
The \LCM\ of $4a^{2}$, $3a$, and $6a^{3} = 12a^{3}$.

The respective quotients are $3a$, $4a^{2}$, and~$2$.

The products are $9ax$, $8a^{2}y$, and~$10$.

Hence, the required fractions are
\[
\frac{9ax}{12a^{3}},\quad
\frac{8a^{2}y}{12a^{3}}, \quad\text{and}\quad
\frac{10}{12a^{3}}.
\]
\end{Soln}

\Item{2.}  Express $\dfrac{1}{x^{2} + 5x + 6}$ and $\dfrac{1}{x^{2} + 4x + 3}$ with lowest
common denominator.
\begin{Soln}
The factors of the denominators are $x + 3$, $x + 2$; and $x + 3$, $x + 1$.

Hence the lowest common denominator (\LCD) is $(x + 3)(x + 2)(x + 1)$,
and the required numerators are $x + 1$ and $x + 2$.  Hence the required
fractions are
\[
%[** TN: Equations not displayed in the original]
\frac{x + 1}{(x + 3)(x + 2)(x + 1)} \quad\text{and}\quad
\frac{x + 2}{(x + 3)(x + 2)(x + 1)}.
\]
\end{Soln}
%% -----File: 100.png---Folio 94-------

\ScreenBreak
\Exercise{45.}

Express with lowest common denominator:

\begin{multicols}{2}
\Item{1.}  $\dfrac{x}{x - a}$, $\dfrac{x^{2}}{x^{2} - a^{2}}$.

\Item{2.}  $\dfrac{a}{a + b}$, $\dfrac{a^{2}}{a^{2} - b^{2}}$.

\Item{3.}  $\dfrac{1}{1 + 2a}$, $\dfrac{1}{1 - 4a^{2}}$.

\Item{4.}  $\dfrac{9}{16 - x^{2}}$, $\dfrac{4 - x}{4 + x}$.

\Item{5.}  $\dfrac{a^{2}}{27 - a^{3}}$, $\dfrac{a}{3 - a}$.

\Item{6.}  $\dfrac{1}{x^{2} - 5x + 6}$, $\dfrac{1}{x^{2} - x-6}$.
\end{multicols}


\Section{Addition and Subtraction of Fractions.}

\Paragraph{132.}  The algebraic sum of two or more fractions which
have the same denominator, is a fraction whose numerator
is the algebraic sum of the numerators of the given fractions,
and whose denominator is the common denominator
of the given fractions.  Hence,

\Dictum{To add fractions},
\begin{Theorem}
Reduce the fractions to equivalent fractions having the
same denominator; and write the algebraic sum of the
numerators of these fractions over the common denominator.
\end{Theorem}

\Paragraph{133.  When the denominators are simple expressions.}

\Item{1.}  Simplify $\dfrac{3a - 4b}{4} - \dfrac{2a - b + c}{3} + \dfrac{a - 4c}{12}$.
\begin{Soln}
The $\text{\LCD} = 12$.

The multipliers, that is, the quotients obtained by dividing $12$ by
$4$, $3$, and $12$, are $3$,~$4$, and~$1$.

Hence the sum of the fractions equals
\begin{align*}
&\frac{9a - 12b}{12} - \frac{8a - 4b + 4c}{12} + \frac{a - 4c}{12} \\
&\quad= \frac{9a - 12b - 8a + 4b - 4c + a - 4c}{12} \\
&\quad= \frac{2a - 8b-8c}{12} = \frac{a - 4b - 4c}{6}.
\end{align*}
\end{Soln}
%% -----File: 101.png---Folio 95-------

The preceding work may be arranged as follows:
\begin{Soln}
The $\text{\LCD} = 12$.

The multipliers are $3$,~$4$, and~$1$, respectively.
\begin{gather*}
\begin{array}{l*{5}{cr}cl}
3(3a &-& 4b & &  ) &=& 9a &-&12b & &    &=& \text{1st numerator.} \\
\llap{$-$}
4(2a &-&  b &+& c) &=&-8a &+& 4b &-& 4c &=& \text{2d numerator.} \\
1( a &-&    & &4c) &=&  a & &    &-& 4c &=& \text{3d numerator.} \\
\cline{7-11}
      & &    & &    & & 2a &-& 8b &-& 8c \\
      & &    & &    &\rlap{or}& 2(a &-& 4b &-& 4c) &=& \text{the sum of the numerators.}
\end{array} \\
\therefore\ \text{sum of fractions} = \frac{2(a - 4b - 4c)}{12} = \frac{a - 4b - 4c}{6}.
\end{gather*}
\end{Soln}

\Exercise{46.}

Find the sum of:

\Item{1.}  $\dfrac{x + 1}{2} + \dfrac{x - 3}{5} + \dfrac{x + 5}{10}$.

\Item{2.}  $\dfrac{2x - 1}{3} + \dfrac{x + 5}{4} + \dfrac{x - 4}{6}$.

\Item{3.}  $\dfrac{7x - 1}{6} - \dfrac{3x - 2}{7} + \dfrac{x - 5}{3}$.

\Item{4.}  $\dfrac{3x - 2}{9} - \dfrac{x - 2}{6} + \dfrac{5x + 3}{4}$.

\Item{5.}  $\dfrac{x - 1}{6} - \dfrac{x - 3}{3} + \dfrac{x - 5}{2}$.

\Item{6.}  $\dfrac{x - 2y}{2x} + \dfrac{x + 5y}{4x} - \dfrac{x + 7y}{8x}$.

\Item{7.}  $\dfrac{5x - 11}{3} - \dfrac{2x - 1}{10} - \dfrac{11x - 5}{15}$.

\Item{8.}  $\dfrac{x - 3}{3x} - \dfrac{x^{2} - 6x}{5x^{2}} - \dfrac{7x^{2} - x^{3}}{15x^{3}}$.

\Item{9.}  $\dfrac{ac - b^{2}}{ac} - \dfrac{ab - c^{2}}{ab} + \dfrac{a^{2} - bc}{bc}$.
%% -----File: 102.png---Folio 96-------

\PrintBreak
\Paragraph{134.  When the denominators have compound expressions,
arranged in the same order.}

\Item{1.}  Simplify $\dfrac{a + b}{a - b} - \dfrac{a - b}{a + b} - \dfrac{4ab}{a^{2} - b^{2}}$.
\begin{Soln}
The \LCD\ is $(a - b)(a + b)$.

The multipliers are $a + b$, $a - b$, and~$1$, respectively.
\begin{gather*}
\begin{array}{l*{3}{cr}cl}
 (a + b)(a + b) &=& a^{2} &+& 2ab &+& b^{2} &=& \text{1st numerator.} \\
\llap{$-$}
  (a - b)(a - b) &=&-a^{2} &+& 2ab &-& b^{2} &=& \text{2d numerator.} \\
\llap{$-$}
  1(4ab)         &=&       &-& 4ab & &      &=& \text{3d numerator.} \\
\cline{3-7}
                & & \multicolumn{5}{c}{0}  &=& \text{sum of numerators.}
\end{array} \\
\therefore\ \text{Sum of fractions${} = 0$.}
\end{gather*}
\end{Soln}

\Exercise{47.}

Find the sum of:
\begin{multicols}{2}
\Item{1.}  $\dfrac{1}{x + 3} + \dfrac{1}{x - 2}$.

\Item{2.}  $\dfrac{1}{x + 1} + \dfrac{1}{x - 1}$.

\Item{3.}  $\dfrac{4}{x - 8} - \dfrac{1}{x + 2}$.

\Item{4.}  $\dfrac{a + x}{a - x} - \dfrac{a - x}{a + x}$.

\Item{5.}  $\dfrac{x}{x - a} - \dfrac{x^{2}}{x^{2} - a^{2}}$.

\Item{6.}  $\dfrac{4a^{2} + b^{2}}{4a^{2} - b^{2}} - \dfrac{2a + b}{2a - b}$.

\Item{7.}  $\dfrac{7}{9 - a^{2}} - \dfrac{1}{3 + a} - \dfrac{1}{3 - a}$.

\Item{8.}  $\dfrac{1}{a - b} - \dfrac{1}{a + b} - \dfrac{b}{a^{2} - b^{2}}$.
\end{multicols}
%[** TN: Adjusted layout]

\Item{9.}  $\dfrac{2}{x - 2} - \dfrac{2}{x + 2} + \dfrac{5x}{x^{2} - 4}$.

\Item{10.}  $\dfrac{3 - x}{1 - 3x} - \dfrac{3 + x}{1 + 3x} - \dfrac{15x - 1}{1 - 9x^{2}}$.

\begin{multicols}{2}
\Item{11.}  $\dfrac{1}{a} - \dfrac{1}{a + 3} + \dfrac{3}{a + 1}$.

\Item{12.}  $\dfrac{x}{x - 1} - \dfrac{1} - \dfrac{1}{x + 1}$.
\end{multicols}

\Item{13.}  $\dfrac{x + 1}{x + 2} + \dfrac{x - 2}{x - 3} + \dfrac{2x + 7}{x^{2} - x - 6}$.

\Item{14.}  $\dfrac{1}{x(x - 1)} - \dfrac{2}{x^{2} - 1} + \dfrac{1}{x(x + 1)}$.
%% -----File: 103.png---Folio 97-------

\Exercise{48.}

Find the sum of:

\Item{1.}  $\dfrac{1}{2x + 1} + \dfrac{1}{2x - 1} - \dfrac{4x}{4x^{2} - 1}$.

\Item{2.}  $\dfrac{a^{2} + b^{2}}{a^{2} - b^{2}} + \dfrac{a}{a + b} - \dfrac{b}{a - b}$.

\Item{3.}  $\dfrac{3a}{1 - a^{2}} + \dfrac{2}{1 - a} - \dfrac{2}{1 + a}$.

\Item{4.}  $\dfrac{1}{2x + 5y} - \dfrac{3x}{4x^{2} - 25y^{2}} + \dfrac{1}{2x + 5y}$.

\Item{5.}  $\dfrac{1}{x + 4y} - \dfrac{8y}{x^{2} - 16y^{2}} + \dfrac{1}{x - 4y}$.

\Item{6.}  $\dfrac{3}{2x - 3} - \dfrac{2}{2x + 3} - \dfrac{3}{4x^{2} - 9}$.


\Section{Multiplication and Division of Fractions.}

\Paragraph{135.}  Find the product of $\dfrac{a}{b} × \dfrac{c}{d}$.

Let $\dfrac{a}{b} = x$, and $\dfrac{c}{d} = y$.

Then $a = bx$, and $c = dy$.

The product of these two equations is
\begin{DPalign*}
ac &= bdxy. \\
\lintertext{\indent Divide by~$bd$,}
\frac{ac}{bd} &= xy. \displaybreak[1] \\
\lintertext{\indent But}
\frac{a}{b} × \frac{c}{d} = xy. \displaybreak[1] \\
\lintertext{\indent Therefore}
\frac{a}{b} × \frac{c}{d} = \frac{ac}{bd}.
\end{DPalign*}
%% -----File: 104.png---Folio 98-------

\Dictum{To find the product of two fractions}, therefore,
\begin{Theorem}
Find the product of the numerators for the required
numerator, and the product of the denominators for the
required denominator.
\end{Theorem}

In like manner,
\[
\frac{a}{b} × \frac{c}{d} × \frac{e}{f} = \frac{ac}{bd} × \frac{e}{f} = \frac{ace}{bdf}.
\]

\Paragraph{136.  Reciprocals.}  If the product of two numbers is equal
to~$1$, each of the numbers is called the \Defn{reciprocal} of the
other.

The reciprocal of $\dfrac{a}{b}$ is $\dfrac{b}{a}$, for $\dfrac{b}{a} × \dfrac{a}{b} = \dfrac{ba}{ab} = 1$.

The reciprocal of a fraction, therefore, is the fraction
inverted.

%[** TN: Next line broken in the original]
Since $\dfrac{a}{b} ÷ \dfrac{a}{b} = 1$, and $\dfrac{b}{a} × \dfrac{a}{b} = 1$,
it follows that
\begin{Theorem}
To divide by a fraction is the same as to multiply by its
reciprocal.
\end{Theorem}

\Paragraph{137.}  \Dictum{To Divide by a Fraction}, therefore,
\begin{Theorem}
Invert the divisor and multiply.
\end{Theorem}

\begin{Remark}[Note.] Every mixed expression should first be reduced to a fraction,
and every integral expression should be written as a fraction
having $1$~for the denominator.  Both terms of each fraction should
be expressed in their prime factors, and if a factor is common to a
numerator and denominator, it should be cancelled, as the cancelling
of a common factor \emph{before} the multiplication is evidently equivalent
to cancelling it \emph{after} the multiplication.
\end{Remark}
%% -----File: 105.png---Folio 99-------

\Item{1.}  Find the product of $\dfrac{3a^{2}b}{2x^{2}y} × \dfrac{6xy^{2}}{7ab} × \dfrac{7abc}{9a^{2}by^{2}}$.
\[
\frac{3a^{2}b}{2x^{2}y} × \frac{6xy^{2}}{7ab} × \frac{7abc}{9a^{2}by^{2}}
  = \frac{3 × 6 × 7a^{3}b^{2}cxy^{2}}{2 × 7 × 9a^{3}b^{2}x^{2}y^{3}}
  = \frac{c}{xy}.
\]

\Item{2.}  Find the product of $\dfrac{ab - b^{2}}{a + b} × \dfrac{ab + b^{2}}{a^{2} - b^{2}}$.
\[
\frac{ab - b^{2}}{a + b} × \frac{ab + b^{2}}{a^{2} - b^{2}}
  = \frac{b(a - b)}{(a + b)} × \frac{b(a + b)}{(a - b)(a + b)}
  = \frac{b^{2}}{a + b}.
\]

\Item{3.}  Find quotient of $\dfrac{ab}{(a - b)^{2}} ÷ \dfrac{ac}{a^{2} - b^{2}}$.
\begin{align*}
\frac{ab}{(a - b)^{2}} ÷ \frac{ac}{a^{2} - b^{2}}
  &= \frac{ab}{(a - b)(a - b)} × \frac{(a - b)(a + b)}{ac} \\
  &= \frac{b(a + b)}{c(a - b)}.
\end{align*}

\Item{4.}  Find the result of $\dfrac{1}{x} × \dfrac{x^{2} - 1}{x^{2} - 4x - 5} ÷ \dfrac{x^{2} + 2x - 3}{x^{2} - 25}$.
\begin{align*}
&\frac{1}{x} × \frac{x^{2} - 1}{x^{2} - 4x - 5} ÷ \frac{x^{2} + 2x - 3}{x^{2} - 25} \\
&\qquad= \frac{1}{x} × \frac{x^{2} - 1}{x^{2} - 4x - 5} × \frac{x^{2} - 25}{x^{2} + 2x - 3} \\
&\qquad= \frac{1}{x} × \frac{(x - 1)(x + 1)}{(x - 5)(x + 1)} × \frac{(x - 5)(x + 5)}{(x + 3)(x - 1)} \\
&\qquad= \frac{x + 5}{x(x + 3)}.
\end{align*}
%% -----File: 106.png---Folio 100-------

\Exercise{49.}

Express in the simplest form:
\begin{multicols}{2}
\Item{1.}  $\dfrac{15a^{2}}{7b^{2}} × \dfrac{28ab}{9a^{3}c}$.

\Item{2.}  $\dfrac{3x^{2}y^{2}z^{3}}{4a^{2}b^{2}c^{2}} × \dfrac{8a^{3}b^{2}c^{2}}{9x^{2}yz^{3}}$.

\Item{3.}  $\dfrac{5m^{2}n^{2}p^{4}}{3x^{2}yz^{3}} × \dfrac{21xyz^{2}}{20m^{2}n^{2}p^{2}}$.

\Item{4.}  $\dfrac{16a^{4}b^{2}c^{3}}{21m^{2}x^{3}y^{4}} × \dfrac{3m^{3}x^{3}y^{4}}{8a^{2}b^{2}c^{2}}$.

\Item{5.}  $\dfrac{2a}{bc} × \dfrac{3b}{ac} × \dfrac{5c}{ab}$.

\Item{6.}  $\dfrac{2a^{3}}{3bc} × \dfrac{3b^{3}}{5ac} × \dfrac{5c^{3}}{2ab}$.

\Item{7.}  $\dfrac{5abc^{3}}{3x^{2}} ÷ \dfrac{10ac^{3}}{6bx^{2}}$.

\Item{8.}  $\dfrac{x^{2} - a^{2}}{x^{2} - 4a^{2}} × \dfrac{x + 2a}{x - a}$.

\Item{9.}  $\dfrac{x^{2}y^{2} + 3xy}{4c^{2} - 1} × \dfrac{2c + 1}{xy + 3}$.

\Item{10.}  $\dfrac{a^{2} - 100}{a^{2} - 9} × \dfrac{a - 3}{a - 10}$.

\Item{11.}  $\dfrac{9x^{2} - 4y^{2}}{x^{2} - 4} × \dfrac{x + 2}{3x - 2y}$.

\Item{12.}  $\dfrac{25a^{2} - b^{2}}{16a^{2} - 9b^{2}} ÷ \dfrac{5a - b}{4a - 3b}$.
\end{multicols}
%[** TN: Moved end of two-column layout up two questions]

\Item{13.}  $\dfrac{x^{2} - 49}{(a + b)^{2} - c^{2}} ÷ \dfrac{x + 7}{(a + b) - c}$.

\Item{14.}  $\dfrac{x^{2} + 2x + 1}{x^{2} - 25} ÷ \dfrac{x + 1}{x^{2} + 5x}$.

\Item{15.}  $\dfrac{a^{2} + 3a + 2}{a^{2} + 5a + 6} × \dfrac{a^{2} + 7a + 12}{a^{2} + 9a + 20}$.

\Item{16.}  $\dfrac{y^{2} - y-30}{y^{2} - 36} × \dfrac{y^{2} - y-2}{y^{2} + 3y-10} × \dfrac{y^{2} + 6y}{y^{2} + y}$.

\Item{17.}  $\dfrac{x^{2} - 2x + 1}{x^{2} - y^{2}} × \dfrac{x^{2} + 2xy + y^{2}}{x - 1} ÷ \dfrac{x^{2} - 1}{x^{2} - xy}$.

\Item{18.}  $\dfrac{a^{2} - b^{2}}{a^{2} - 3ab + 2b^{2}} × \dfrac{ab - 2b^{2}}{a^{2} + ab} ÷ \dfrac{(a - b)^{2}}{a(a - b)}$.

\Item{19.}  $\dfrac{(a + b)^{2} - c^{2}}{a^{2} + ab-ac} × \dfrac{a^{2}b^{2}c^{2}}{a^{2} + ab + ac} ÷ \dfrac{b^{2}c^{2}}{abc}$.

\Item{20.}  $\dfrac{x^{2} + 7xy + 10y^{2}}{x^{2} + 6xy + 5y^{2}} × \dfrac{x + 1}{x^{2} + 4x + 4} ÷ \dfrac{1}{x + 2}$.
%% -----File: 107.png---Folio 101-------

\Paragraph{138.  Complex Fractions.}  A complex fraction is one that
has a fraction in the numerator, or in the denominator, or
in both.

The shortest way to reduce to its simplest form a complex
fraction is to multiply both terms of the fraction by
the \LCD\ of the fractions contained in the numerator and
denominator.

\Item{1.}  Simplify $\dfrac{3x}{x - \frac{1}{4}}$.

Multiply both terms by~$4$, and we have
\[
\frac{12x}{4x - 1}.
\]

\Item{2.}  Simplify $\dfrac{\dfrac{a + x}{a - x} - \dfrac{a - x}{a + x}}{\dfrac{a + x}{a - x} + \dfrac{a - x}{a + x}}$.

The \LCD\ of the fractions in the numerator and denominator
is
\[
(a - x)(a + x).
\]

Multiply by $(a - x)(a + x)$, and the result is
\begin{align*}
&\frac{(a + x)^{2} - (a - x)^{2}}{(a + x)^{2} + (a - x)^{2}} \\
&\qquad= \frac{(a^{2} + 2ax + x^{2}) - (a^{2} - 2ax + x^{2})}
               {(a^{2} + 2ax + x^{2}) + (a^{2} - 2ax + x^{2})} \\
&\qquad= \frac{a^{2} + 2ax + x^{2} - a^{2} + 2ax - x^{2}}
               {a^{2} + 2ax + x^{2} + a^{2} - 2ax + x^{2}} \\
&\qquad= \frac{4ax}{2a^{2} + 2x^{2}} \\
&\qquad= \frac{2ax}{a^{2} + x^{2}}.
\end{align*}
%% -----File: 108.png---Folio 102-------

\Exercise{50.}

Reduce to the simplest form:
\begin{multicols}{2}
\Item{1.}  $\dfrac{\dfrac{x}{b} + \dfrac{y}{b}}{\dfrac{z}{b}}$.

\Item{2.}  $\dfrac{x + \dfrac{y}{4}}{x - \dfrac{y}{3}}$.

\Item{3.}  $\dfrac{\dfrac{ab}{7} - 3d}{3c - \dfrac{ab}{d}}$.

\Item{4.}  $\dfrac{1 + \dfrac{1}{x + 1}}{1 - \dfrac{1}{x - 1}}$.

\Item{5.}  $\dfrac{\dfrac{2m + x}{m + x} - 1}{1 - \dfrac{x}{m + x}}$.

\Item{6.}  $\dfrac{\dfrac{x + y}{x^{2} - y^{2}}}{\dfrac{x - y}{x + y}}$.

\Item{7.}  $\dfrac{a + \dfrac{ab}{a - b}}{a - \dfrac{ab}{a + b}}$.

\Item{8.}  $\dfrac{9a^{2} - 64}{a - 1 - \dfrac{a + 4}{4}}$.

\Item{9.}  $\dfrac{\dfrac{1}{x} + \dfrac{1}{y}}{\dfrac{1}{x} - \dfrac{1}{y}}$.

\Item{10.}  $\dfrac{x + 3 + \dfrac{2}{x}}{1 + \dfrac{3}{x} + \dfrac{2}{x^{2}}}$.

\Item{11.}  $\dfrac{\dfrac{1}{x} - \dfrac{2}{x^{2}} + \dfrac{1}{x^{3}}}{\dfrac{(1 - x)^{2}}{x^{2}}}$.

\Item{12.}  $\dfrac{x^{2} - x-6}{1 - \dfrac{4}{x^{2}}}$.

\Item{13.}  $\dfrac{a^{2} - a + \dfrac{a - 1}{a + 1}}{a + \dfrac{1}{a + 1}}$.

\Item{14.}  $\dfrac{\dfrac{4a(a - x)}{a^{2} - x^{2}}}{\dfrac{a - x}{a + x}}$.
\end{multicols}
%% -----File: 109.png---Folio 103-------


\Chapter{X.}{Fractional Equations.}

\Paragraph{139.  To Reduce Equations containing Fractions.}

\Item{1.}  Solve $\dfrac{x}{3} - \dfrac{x - 1}{11} = x - 9$.
\begin{Soln}
Multiply by~$33$, the \LCM\ of the denominators.
\begin{DPalign*}
\lintertext{\indent Then,}
11x - 3x + 3 &= 33x - 297. \\
\lintertext{\indent Transpose,}
11x - 3x - 33x &= - 297 - 3. \\
\lintertext{\indent Combine,}
- 25x &= - 300. \\
\lintertext{\indent Divide by~$-25$,}
x &= 12.
\end{DPalign*}
\end{Soln}

\begin{Remark}[Note.] Since the minus sign precedes the second fraction, in removing
the denominator, the~$+$ (understood) before~$x$, the first term
of the numerator, is changed to~$-$, and the~$-$ before~$1$, the second
term of the numerator, is changed to~$+$.
\end{Remark}

\Dictum{To clear an equation of fractions}, therefore,
\begin{Theorem}
Multiply each term by the \LCM\ of the denominators.
\end{Theorem}

If a fraction is preceded by a \textbf{minus sign}, \emph{the sign of
every term of the numerator must be changed when the
denominator is removed}.

\Item{2.}  Solve $\dfrac{x + 1}{4} - \frac{1}{5}(x - 1) = 1$.
\begin{Soln}
Multiply by~$20$, the \LCD
\begin{DPalign*}
5x + 5 - 4(x - 1) &= 20. \\
5x + 5 - 4x + 4 &= 20. \\
\lintertext{\indent Transpose,}
5x - 4x &= 20 - 5 - 4. \\
\lintertext{\indent Combine,}
x &= 11.
\end{DPalign*}
\end{Soln}
%% -----File: 110.png---Folio 104-------

\Item{3.}  Solve
\[
7x - \frac{(2x - 3)(3x - 5)}{5} = \frac{153}{10} - \frac{(4x - 5)(3x - 1)}{10}.
\]
\begin{Soln}
Multiply by~$10$, the \LCD, and we have
\[
70x - 2(2x - 3)(3x - 5) = 153 - (4x - 5) (3x - 1).
\]

Find the products of $(2x - 3)(3x - 5)$ and $(4x - 5)(3x - 1)$.
\[
70x - 2(6x^{2} - 19x + 15) = 153 - (12x^{2} - 19x + 5).
\]

Remove the parentheses,
\[
70x - 12x^{2} + 38x - 30 = 153 - 12x^{2} + 19x - 5.
\]

Cancel the~$-12x^{2}$ on each side and transpose,
\begin{DPalign*}
70x + 38x - 19x &= 153 + 30 - 5. \\
\lintertext{\indent Combine,}
89x &= 178. \\
\lintertext{\indent Divide by~$89$,}
x &= 2.
\end{DPalign*}
\end{Soln}

\Item{4.}  Solve $\dfrac{2x + 1}{2x - 1} - \dfrac{2x - 1}{2x + 1} = \dfrac{8}{4x^{2} - 1}$.
\begin{Soln}
\begin{DPalign*}
\lintertext{\indent Since}
4x^{2} - 1 &= (2x + 1)(2x - 1), \\
\lintertext{the \LCD}
  &= (2x + 1)(2x - 1).
\end{DPalign*}
Multiply by the \LCD, and we have,
\begin{DPalign*}
4x^{2} + 4x + 1-(4x^{2} - 4x + 1) &= 8. \\
\therefore 4x^{2} + 4x + 1 - 4x^{2} + 4x - 1 &= 8. \\
\therefore 8x &= 8. \\
\therefore x &= 1.
\end{DPalign*}
\end{Soln}

\Item{5.}  Solve $\dfrac{4}{x + 1} - \dfrac{x + 1}{x - 1} + \dfrac{x^{2} - 3}{x^{2} - 1} = 0$.
\begin{Soln}
\begin{DPalign*}
\lintertext{\indent Since}
x^{2} - 1 &= (x + 1)(x - 1) \\
\lintertext{the \LCD}
  &= (x + 1)(x - 1).
\end{DPalign*}

Multiply by the \LCD, $x^{2} - 1$, and we have,
\begin{DPalign*}
4(x - 1) - (x + 1)(x + 1) + x^{2} - 3 &= 0. \\
\therefore 4x - 4 - x^{2} - 2x - 1 + x^{2} - 3 &= 0. \\
\therefore 2x &= 8. \\
\therefore x &= 4.
\end{DPalign*}
\end{Soln}
%% -----File: 111.png---Folio 105-------

%[** TN: Force page break in both print and screen layout]
\newpage
\Exercise{51.}

Solve:
\begin{multicols}{2}
\Item{1.}  $\dfrac{x - 1}{2} = \dfrac{x + 1}{3}$.

\Item{2.}  $\dfrac{3x - 1}{4} = \dfrac{2x + 1}{3}$.

\Item{3.}  $\dfrac{6x - 19}{2} = \dfrac{2x - 11}{3}$.

\Item{4.}  $\dfrac{7x - 40}{8} = \dfrac{9x - 80}{10}$.
\end{multicols}
%[** TN: Move end of two-column layout up six questions]

\Item{5.}  $\dfrac{3x - 116}{4} + \dfrac{180 - 5x}{6} = 0$.

\Item{6.}  $\dfrac{3x - 4}{2} - \dfrac{3x - 1}{16} = \dfrac{6x - 5}{8}$.

\Item{7.}  $\dfrac{x - 1}{8} - \dfrac{x + 1}{18} = 1$.

\Item{8.}  $\dfrac{60 - x}{14} - \dfrac{3x - 5}{7} = \dfrac{3x}{4}$.

\Item{9.}  $\dfrac{3x - 1}{11} - \dfrac{2 - x}{10} = \dfrac{6}{5}$.

\Item{10.}  $\dfrac{4x}{x + 1} - \dfrac{x}{x - 2} = 3$.

\Item{11.}  $\dfrac{2x + 1}{4} - \dfrac{4x - 1}{10} + 1 - \frac{1}{4} = 0$.

\Item{12.}  $\dfrac{x - 1}{5} - \dfrac{43 - 5x}{6} - \dfrac{3x - 1}{8} = 0$.

\Item{13.}  $\dfrac{1}{x + 7} = \dfrac{2}{x + 1} - \dfrac{1}{x + 3}$.

\Item{14.}  $\dfrac{1}{x + 4} + \dfrac{2}{x + 6} - \dfrac{3}{x + 5} = 0$.

\Item{15.}  $\dfrac{4}{x^{2} - 1} + \dfrac{1}{x - 1} + \dfrac{1}{x + 1} = 0$.

\Item{16.}  $\dfrac{3x + 1}{4} - \dfrac{5x - 4}{7} = 12 - 2x - \dfrac{x - 2}{3}$.

\Item{17.}  $\frac{1}{8}(5x + 3) - \frac{1}{3}(3 - 4x) + \frac{1}{6}(9 - 5x) = \frac{1}{2}(31 - x)$.

\Item{18.}  $\frac{1}{15} (34x - 56) - \frac{1}{5}(7x - 3) - \frac{1}{3}(7x - 5) = 0$.
%% -----File: 112.png---Folio 106-------

\Exercise{52.}

Solve:

\Item{1.}  $\frac{2}{3}(x + 1) - \frac{1}{7}(x + 5) = 1$.

\Item{2.}  $\frac{6}{7}(x - 9) - \frac{1}{3}(5 - x) + 3x + 1 = 0$.

\Item{3.}  $\frac{1}{3}(5x - 24) + \frac{1}{7}(x - 2) - 2(x - 1) = 0$.

\Item{4.}  $\dfrac{x + 3}{4} + \dfrac{7x - 2}{5} = \dfrac{5x - 1}{4} + \dfrac{5x + 4}{9}$.

\Item{5.}  $\dfrac{x + 1}{3} - \dfrac{x - 1}{4} = \dfrac{x - 2}{5} - \dfrac{x - 3}{6} + \dfrac{31}{60}$.

\Item{6.}  $\dfrac{(2x - 1)(2 - x)}{2} + x^{2} - \dfrac{1 + 3x}{2} = 0$.

\Item{7.}  $\dfrac{6x - 11}{4} - \dfrac{3 - 4x}{6} = \dfrac{4}{3} - \dfrac{x}{8}$.

\Item{8.}  $\dfrac{x + 6}{4} - \dfrac{16 - 3x}{12} = 4\frac{1}{6}$.

\Item{9.}  $x - \dfrac{x - 2}{3} = \dfrac{x + 23}{4} - \dfrac{10 + x}{5}$.

\Item{10.}  $\dfrac{5x + 3}{x - 1} + \dfrac{2x - 3}{2x - 1} = 6$.

\Item{11.}  $\dfrac{3x}{4x + 1} + 1 = 2 - \dfrac{x}{2(2x - 1)}$.

\Item{12.}  $\dfrac{8x + 7}{5x + 4} - 1 = 1 - \dfrac{2x}{5x + 1}$.

\Item{13.}  $\dfrac{x + 1}{2(x - 1)} - \dfrac{x - 1}{x + 1} = \dfrac{17 - x^{2}}{2(x^{2} - 1)}$.
%% -----File: 113.png---Folio 107-------

\Paragraph{140.}  If the denominators contain both simple and compound
expressions, it is generally best to remove the simple
expressions first, and then the compound expressions.  After
each multiplication the result should be reduced to the
simplest form.

\Item{1.}  Solve $\dfrac{4x + 3}{10} - \dfrac{2x + 3}{5x - 1} = \dfrac{2x - 1}{5}$.
\begin{Soln}
\begin{DPalign*}
\lintertext{\indent Multiply by~$10$,}
4x + 3 - \frac{10(2x + 3)}{5x - 1} &= 4x - 2. \\
\lintertext{\indent Transpose,}
4x + 3 - 4x + 2 &= \frac{10(2x + 3)}{5x - 1}. \\
\lintertext{\indent Combine,}
5 &= \frac{10(2x + 3)}{5x - 1}. \displaybreak[1] \\
\lintertext{\indent Divide by~$5$,}
1 &= \frac{2(2x + 3)}{5x - 1}. \\
\lintertext{\indent Multiply by $5x - 1$,}
5x - 1 &= 4x + 6. \\
\lintertext{\indent Transpose and combine,}
x &= 7.
\end{DPalign*}
\end{Soln}

\Exercise{53.}

Solve:

\Item{1.}  $\dfrac{10x + 13}{18} - \dfrac{x + 2}{x - 3} = \dfrac{5x - 4}{9}$.

\Item{2.}  $\dfrac{6x + 7}{10} - \dfrac{3x + 1}{5} = \dfrac{x - 1}{3x - 4}$.

\Item{3.}  $\dfrac{11x - 12}{14} - \dfrac{11x - 7}{19x + 7} = \dfrac{22x - 36}{28}$.

\Item{4.}  $\dfrac{2x - 1}{5} + \dfrac{2x - 3}{17x - 12} = \dfrac{4x - 3}{10}$.

\Item{5.}  $\dfrac{11x - 13}{7} - \dfrac{13x + 7}{3x + 7} = \dfrac{22x - 75}{14}$.

\Item{6.}  $\dfrac{6x - 13}{2x + 3} + \dfrac{6x + 7}{9} - \dfrac{2x + 4}{3} = 0$.
%% -----File: 114.png---Folio 108-------

\Paragraph{141.  Literal Equations.}  Literal equations are equations
in which some or all of the known numbers are represented
by letters; the numbers regarded as known numbers are
usually represented by the \emph{first} letters of the alphabet.

\Item{1.}  Solve $\dfrac{x + a}{x - b} + \dfrac{x + b}{x - a} = 2$.
\begin{Soln}
Multiply by $(x - a)(x - b)$,
\begin{DPalign*}
\llap{$(x + a)(x - a) + (x + b)(x - b)$} &= 2(x - a)(x - b), \\
\lintertext{or}
x^{2} - a^{2} + x^{2} - b^{2} &= 2x^{2} - 2ax - 2bx + 2ab. \displaybreak[1] \\
\lintertext{\indent Transpose,}
x^{2} + x^{2} - 2x^{2} + 2ax + 2bx &= a^{2} + 2ab + b^{2}\Add{.} \displaybreak[1] \\
\lintertext{\indent Combine,}
2ax + 2bx &= a^{2} + 2ab + b^{2}, \\
\lintertext{or}
2(a + b)x &= a^{2} + 2ab + b^{2}. \displaybreak[1] \\
\lintertext{\indent \rlap{Divide by~$a + b$,}}
2x &= a + b\Add{,} \\
\therefore x &= \frac{a + b}{2}.
\end{DPalign*}
\end{Soln}

\Exercise{54.}

Solve:

\Item{1.}  $a(x - a) = b(x - b)$.

\Item{2.}  $(a + b)x + (a - b)x = a^{2}$.

\Item{3.}  $(a + b)x - (a - b)x = b^{2}$.

\Item{4.}  $(2x - a) + (x - 2a) = 3a$.

\Item{5.}  $(x + a + b) + (x + a - b) = 2b$.

\Item{6.}  $(x - a)(x - b) = x(x + c)$.

\Item{7.}  $x^{2} + b^{2} = (a - x)(a - x)$.

\Item{8.}  $(a + b)(2 - x) = (a - b)(2 + x)$.

\Item{9.}  $(x - a)(2x - a) = 2(x - b)^{2}$.

\Item{10.}  $(a + bx)(c + d) = (a + b)(c + dx)$.

\Item{11.}  $\dfrac{x}{a - b} - \dfrac{3a}{a + b} = \dfrac{bx}{a^{2} - b^{2}}$.
%% -----File: 115.png---Folio 109-------

\PrintBreak
\Paragraph{142.  Problems involving Fractional Equations.}

\Exercise{55.}

Ex.  The sum of the third and fifth parts of a certain
number exceeds two times the difference of the fourth and
sixth parts by~$22$.  Find the number.
\begin{Soln}
\begin{DPalign*}
\lintertext{\indent Let}
x &= \text{the number\Add{.}} \\
\lintertext{\indent Then}
\frac{x}{3} + \frac{x}{5} &= \text{the sum of its third and fifth parts,} \\
\frac{x}{4} - \frac{x}{6} &= \text{the difference of its fourth and sixth parts,} \\
2\left(\frac{x}{4} - \frac{x}{6}\right)
  &= \text{$2$~times the difference of} \\
  &\qquad\text{its fourth and sixth parts,}
\end{DPalign*}
\begin{DPalign*}
\frac{x}{3} + \frac{x}{5} - 2\left(\frac{x}{4} - \frac{x}{6}\right)
  &= \text{the given excess.} \\
\lintertext{\indent But}
22 &= \text{the given excess.} \\
\therefore \frac{x}{3} + \frac{x}{5} - 2\left(\frac{x}{4} - \frac{x}{6}\right)
  &= 22.
\end{DPalign*}

Multiply by~$60$ the \LCD\ of the fractions.
\begin{DPalign*}
20x + 12x - 30x + 20x &= 60 × 22. \\
\lintertext{\indent Combining,}
22x &= 60 × 22\Add{,} \\
\therefore x &= 60.
\end{DPalign*}

The required number, therefore, is~$60$.
\end{Soln}

\Item{1.}  The difference between the fifth and seventh parts of
a certain number is~$2$.  Find the number.

\Item{2.}  One-half of a certain number exceeds the sum of its
fifth and seventh parts by~$11$.  Find the number.

\Item{3.}  The sum of the third and sixth parts of a certain
number exceeds the difference of its sixth and ninth parts
by~$16$.  Find the number.

\Item{4.}  There are two consecutive numbers, $x$~and~$x + 1$, such
that one-half the larger exceeds one-third the smaller
number by~$10$.  Find the numbers.
%% -----File: 116.png---Folio 110-------

\Exercise{56.}

Ex.  The sum of two numbers is~$63$, and if the greater
is divided by the smaller number, the quotient is~$2$ and the
remainder~$3$.  Find the numbers.
\begin{Soln}
\begin{DPalign*}
\lintertext{\indent Let}
x &= \text{the greater number.} \\
\lintertext{\indent Then}
63 - x &= \text{the smaller number.} \\
\text{Since the quotient}
  &= \frac{\text{Dividend} - \text{Remainder}}{\text{Divisor}},
\end{DPalign*}
and since, in this problem, the dividend is~$x$, the remainder is~$3$,
and the divisor is~$63 - x$, we have
\begin{DPalign*}
\frac{x - 3}{63 - x} = 2. \\
\lintertext{\indent Solving,}
x &= 43.
\end{DPalign*}

The two numbers, therefore, are $43$~and~$20$.
\end{Soln}

\Item{1.}  The sum of two numbers is~$100$, and if the greater is
divided by the smaller number, the quotient is~$4$ and the
remainder~$5$.  Find the numbers.

\Item{2.}  The sum of two numbers is~$124$, and if the greater is
divided by the smaller number, the quotient is~$4$ and the
remainder~$4$.  Find the numbers.

\Item{3.}  The difference of two numbers is~$49$, and if the greater
is divided by the smaller, the quotient is~$4$ and the remainder~$4$.
Find the numbers.

\Item{4.}  The difference of two numbers is~$91$, and if the
greater is divided by the smaller, the quotient is~$8$ and the
remainder~$7$.  Find the numbers.

\Item{5.}  Divide $320$ into two parts such that the smaller part
is contained in the larger part $11$~times, with a remainder
of~$20$.
%% -----File: 117.png---Folio 111-------

\Exercise{57.}

Ex.  Eight years ago a boy was one-fourth as old as he
will be one year hence.  How old is he now?
\begin{Soln}
\begin{DPalign*}
\lintertext{\indent Let}
x &= \text{the number of years old he is now.} \\
\lintertext{\indent Then}
x - 8 &= \text{the number of years old he was eight years ago,} \\
\lintertext{and}
x + 1 &= \text{the number of years old he will be one year hence.}
\end{DPalign*}
\begin{DPalign*}
\therefore x - 8 &= \tfrac{1}{4}(x + 1). \\
\lintertext{\indent Solving,}
x &= 11.
\end{DPalign*}

Therefore the boy is $11$~years old.
\end{Soln}

\Item{1.}  A son is one-fourth as old as his father.  In $24$~years
he will be one-half as old.  Find the age of the son.

\Item{2.}  B's~age is one-sixth of A's~age.  In $15$~years B's~age
will be one-third of A's~age.  Find their ages.

\Item{3.}  The sum of the ages of A~and~B is $30$~years, and $5$~years
hence B's~age will be one-third of~A's.  Find their
ages.

\Item{4.}  A father is $35$~years old, and his son is one-fourth of
that age.  In how many years will the son be half as old
as his father?

\Item{5.}  A is $60$~years old, and B's~age is two-thirds of~A's.
How many years ago was B's~age one-fifth of~A's?

\Item{6.}  A son is one-third as old as his father.  Four years
ago he was only one-fourth as old as his father.  What is
the age of each?

\Item{7.}  A is $50$~years old, and B~is half as old as~A\@.  In how
many years will B be two-thirds as old as~A?

\Item{8.}  B~is one-half as old as~A\@.  Ten years ago he was
one-fourth as old as~A\@.  What are their present ages?

\Item{9.}  The sum of the ages of a father and his son is $80$~years.
The son's age increased by $5$~years is one-fourth of
the father's age.  Find their ages.
%% -----File: 118.png---Folio 112-------

\Exercise{58.}

Ex.  A~can do a piece of work in $2$~days, and B~can do
it in $3$~days.  How long will it take both together to do
the work?
\begin{Soln}
\begin{DPalign*}
\lintertext{\indent Let}
x &= \text{the number of days it will take both together.} \\
\lintertext{\indent Then}
\frac{1}{x} &= \text{\emph{the part both together can do in one day},} \\
\tfrac{1}{2} &= \text{the part A can do in one day,} \\
\tfrac{1}{3} &= \text{the part B can do in one day,} \\
\lintertext{and}
\tfrac{1}{2} + \tfrac{1}{3}
  &= \text{\emph{the part both together can do in one day}.}
\end{DPalign*}
\begin{DPalign*}
\therefore \frac{1}{2} + \frac{1}{3} &= \frac{1}{x}. \\
\lintertext{\indent Solving,}
x &= 1\tfrac{1}{5}.
\end{DPalign*}

Therefore they together can do the work in $1\frac{1}{5}$~days.
\end{Soln}

\Item{1.}  A~can do a piece of work in $3$~days, B~in $5$~days, and
C~in $6$~days.  How long will it take them to do it working
together?

\Item{2.}  A~can do a piece of work in $5$~days, B~in $4$~days, and
C~in $3$~days.  How long will it take them together to do
the work?

\Item{3.}  A~can do a piece of work in $2\frac{1}{2}$~days, B~in $3\frac{1}{2}$~days,
and C~in $3\frac{3}{4}$~days.  How long will it take them together to
do the work?

\Item{4.}  A~can do a piece of work in $10$~days, B~in $12$~days;
A~and~B together, with the help of~C, can do the work in
$4$~days.  How long will it take C~alone to do the work?

\Item{5.}  A~and~B together can mow a field in $10$~hours, A~and~C
in $12$~hours, and A~alone in $20$~hours.  In what time
can B~and~C together mow the field?

\Item{6.}  A~and~B together can build a wall in $12$~days, A~and~C
in $15$~days, B~and~C in $20$~days.  In what time can they
build the wall if they all work together?

\begin{Remark}[Hint.]  By working \emph{$2$~days each} they build $\frac{1}{12} + \frac{1}{15} + \frac{1}{20}$ of it.
\end{Remark}
%% -----File: 119.png---Folio 113-------

\Exercise{59.}

Ex.  A cistern can be filled by three pipes in $15$, $20$, and
$30$~hours, respectively.  In what time will it be filled by
all the pipes together?
\begin{Soln}
\begin{DPalign*}
\lintertext{\indent Let}
x &= \text{the number of hours it will take all together.} \\
\lintertext{\indent Then}
\frac{1}{x} &= \text{the part all together can fill in one hour\Add{.}} \\
\llap{$\tfrac{1}{15} + \tfrac{1}{20}$} + \tfrac{1}{30}
  &= \text{the part all together can fill in one hour\Add{.}}
\end{DPalign*}
\begin{DPalign*}
\frac{1}{15} + \frac{1}{20} + \frac{1}{30} &= \frac{1}{x}\Add{.} \\
\lintertext{\indent Solving,}
x &= 6 \tfrac{2}{3}\Add{.}
\end{DPalign*}

Therefore the pipes together can fill it in $6 \frac{2}{3}$~hours.
\end{Soln}

\Item{1.}  A cistern can be filled by three pipes in $16$, $24$, and
$32$~hours, respectively.  In what time will it be filled by
all the pipes together?

\Item{2.}  A tank can be filled by two pipes in $3$~hours and $4$~hours,
respectively, and can be emptied by a third pipe in
$6$~hours.  In what time will the cistern be filled if the pipes
are all running together?

\Item{3.}  A tank can be filled by three pipes in $1$~hour and $40$
minutes, $3$~hours and $20$ minutes, and $5$~hours, respectively.
In what time will the tank be filled if all three pipes are
running together?

\Item{4.}  A cistern can be filled by three pipes in $2\frac{1}{3}$~hours, $3\frac{1}{2}$~hours,
and $4\frac{2}{3}$~hours, respectively.  In what time will the
cistern be filled if all the pipes are running together?

\Item{5.}  A cistern has three pipes.  The first pipe will fill the
cistern in $12$~hours, the second in $20$~hours, and all three
pipes together will fill it in $6$~hours.  How long will it take
the third pipe alone to fill it?
%% -----File: 120.png---Folio 114-------

\Exercise{60.}

Ex.  A courier who travels $6$~miles an hour is followed,
after $2$~hours, by a second courier who travels $7\frac{1}{2}$~miles an
hour.  In how many hours will the second courier overtake
the first?
\begin{Soln}
\begin{DPalign*}
\lintertext{\indent Let}
x &= \text{the number of hours the first travels.} \\
\lintertext{\indent Then}
x - 2 &= \text{the number of hours the second travels,} \\
6x &= \text{the number of miles the first travels,} \\
\lintertext{and}
(x - 2) 7\tfrac{1}{2} &= \text{the number of miles the second travels.}
\displaybreak[1] \\
\intertext{\indent They both travel the same distance.}
\therefore 6x &= (x - 2) 7\tfrac{1}{2}, \\
\lintertext{or}
12x &= 15x - 30. \\
\therefore x &= 10.
\end{DPalign*}

Therefore the second courier will overtake the first in $10 - 2$, or
$8$~hours.
\end{Soln}

\Item{1.}  A sets out from Boston and walks towards Portland
at the rate of $3$~miles an hour.  Three hours afterward B
sets out from the same place and walks in the same direction
at the rate of $4$~miles an hour.  How far from Boston
will B~overtake~A?

\Item{2.}  A courier who goes at the rate of $6\frac{1}{2}$~miles an hour is
followed, after $4$~hours, by another who goes at the rate of
$7\frac{1}{2}$~miles an hour.  In how many hours will the second
overtake the first?

\Item{3.}  A person walks to the top of a mountain at the rate
of two miles an hour, and down the same way at the rate
of $4$~miles an hour.  If he is out $6$~hours, how far is it to
the top of the mountain?

\Item{4.}  In going a certain distance, a train travelling at the
rate of $40$~miles an hour takes $2$~hours less than a train
travelling $30$~miles an hour.  Find the distance.
%% -----File: 121.png---Folio 115-------

\Exercise{61.}

Ex.  A hare takes $4$~leaps to a greyhound's~$3$; but $2$~of
the greyhound's leaps are equivalent to $3$~of the hare's.
The hare has a start of $50$~leaps.  How many leaps must
the greyhound take to catch the hare?
\begin{Soln}
\begin{DPalign*}
\lintertext{\indent Let}
3x &= \text{the number of leaps taken by the greyhound.} \\
\lintertext{\indent Then}
4x &= \text{the number of leaps of the hare in the same time.} \\
%[** TN: Hack to coax spacing]
\lintertext{\indent Also, l\rlap{et}}
a &= \text{the number of feet in one leap of the hare.} \displaybreak[1] \\
\lintertext{\indent Then}
\frac{3a}{2} &= \text{the number of feet in one leap of the hound.}
\displaybreak[1] \\
\lintertext{\indent \rlap{Therefore,}}
&\quad 3x x \frac{3a}{2} \quad\text{or}\quad
\frac{9ax}{2} = \text{the whole distance}.
\end{DPalign*}

As the hare has a start of $50$~leaps, and takes $4x$~leaps more before
she is caught, and as each leap is $a$~feet,
\begin{DPalign*}
(50 + 4x)a &= \text{the whole distance.} \\
\therefore \frac{9ax}{2} &= (50 + 4x)a. \\
\lintertext{\indent Multiply by~$2$,}
9ax &= (100 + 8x)a, \\
\lintertext{\indent Divide by~$a$,}
9x &= 100 + 8x, \\
x &= 100, \\
\therefore 3x &= 300.
\end{DPalign*}

Therefore the greyhound must take $300$~leaps.
\end{Soln}

\Item{1.}  A hound makes $3$~leaps while a rabbit makes~$5$; but
$1$~of the hound's leaps is equivalent to $2$~of the rabbit's.
The rabbit has a start of $120$~leaps.  How many leaps
will the rabbit take before she is caught?

\Item{2.}  A rabbit takes $6$~leaps to a dog's~$5$, and $7$~of the dog's
leaps are equivalent to $9$~of the rabbit's.  The rabbit has
a start of~$60$ of her own leaps.  How many leaps must the
dog take to catch the rabbit?

\Item{3.}  A dog makes $4$~leaps while a rabbit makes~$5$; but $3$~of
the dog's leaps are equivalent to $4$~of the rabbit's.  The
rabbit has a start of $90$~of the \emph{dog's leaps}.  How many
leaps will each take before the rabbit is caught?
%% -----File: 122.png---Folio 116-------

\Exercise{62.}

Ex.  Find the time between $2$ and $3$~o'clock when the
hands of a clock are together.
\begin{Soln}
At 2 o'clock the hour-hand is  10 minute-spaces ahead of the
minute-hand.
\begin{DPalign*}
\lintertext{\indent Let}
x &= \text{\emph{the number of spaces the minute-hand moves over}.} \\
\lintertext{\indent The\rlap{n}}
x - 10 &= \text{the number of spaces the hour-hand moves over.}
\displaybreak[1] \\
\intertext{\indent Now, as the minute-hand moves $12$~times as fast as the hour-hand,}
\llap{$12(x - 10)$} &= \text{\emph{the number of spaces the minute-hand moves over}.}
\end{DPalign*}
\begin{DPalign*}
\therefore 12(x - 10) &= x, \\
\lintertext{and}
11x &= 120. \\
\therefore x &= 10\tfrac{10}{11}.
\end{DPalign*}

Therefore the time is $10\frac{10}{11}$~minutes past $2$~o'clock.
\end{Soln}

\Item{1.}  Find the time between $5$ and $6$~o'clock when the
hands of a clock are together.

\Item{2.}  Find the time between $2$ and $3$~o'clock when the
hands of a clock are at right angles to each other.

\begin{Remark}[Hint.]  In this case the minute-hand is $15$~minutes ahead of the
hour-hand.
\end{Remark}

\Item{3.}  Find the time between $2$ and $3$~o'clock when the
hands of a clock point in opposite directions.

\begin{Remark}[Hint.]  In this case the minute-hand is $30$~minutes ahead of the
hour-hand.
\end{Remark}

\Item{4.}  Find the time between $1$ and $2$~o'clock when the
hands of a clock are at right angles to each other.

\Item{5.}  Find the time between $1$ and $2$~o'clock when the
hands of a clock point in opposite directions.

\Item{6.}  At what time between $7$ and $8$~o'clock are the hands
of a watch together?
%% -----File: 123.png---Folio 117-------

\Exercise{63.}

Ex.  A rectangle has its length $6$~feet more and its width
$5$~feet less than the side of its equivalent square.  Find the
dimensions of the rectangle.
\begin{Soln}
\begin{DPalign*}
\lintertext{\indent Let}
x &= \text{the number of feet in a side of the square.} \\
\lintertext{\indent Then}
x + 6 &= \text{the number of feet in the length of the rectangle,} \\
\lintertext{and}
x - 5 &= \text{the number of feet in the width of the rectangle.}
\intertext{Since the area of a rectangle is equal to the product of the number
of units of length in the length and width of the rectangle,}
\llap{$(x + 6)(x - 5)$} &= \text{the area of the rectangle in square feet,} \\
\lintertext{and}
x × x &= \text{the area of the square in square feet.}
\end{DPalign*}

But these areas are equal.
\begin{DPalign*}
\therefore (x + 6)(x - 5) &= x^{2}. \\
\lintertext{\indent Solving,}
x &= 30.
\end{DPalign*}

Therefore the dimensions of the rectangle are $36$~feet and $25$~feet.
\end{Soln}

\Item{1.}  A rectangle has its length and breadth respectively $7$~feet
longer and $6$~feet shorter than the side of the equivalent
square.  Find its area.

\Item{2.}  The length of a floor exceeds the breadth by $5$~feet.
If each dimension were $1$~foot more, the area of the floor
would be $42$~sq.~ft.\ more.  Find its dimensions.

\Item{3.}  A rectangle whose length is $6$~feet more than its
breadth, would have its area $35$~sq.~ft.\ more, if each dimension
were $1$~foot more.  Find its dimensions.

\Item{4.}  The length of a rectangle exceeds its width by $3$~feet.
If the length is increased by $3$~feet and the width diminished
by $2$~feet, the area will not be altered.  Find its
dimensions.

\Item{5.}  The length of a floor exceeds its width by $10$~feet
If each dimension were $2$~feet more, the area would be $144$~sq.~ft.\
more.  Find its dimensions.
%% -----File: 124.png---Folio 118-------

\Paragraph{143.  Formulas and Rules.}  When the \emph{given} numbers of a
problem are represented by letters, the result obtained from
solving the problem is a general expression which includes
all problems of that class.  Such an expression is called a
\Defn{formula}, and the translation of this formula into words is
called a \Defn{rule}.

\Paragraph{144.}  We will illustrate by examples.

\Item{1.}  The sum of two numbers is~$s$, and their difference~$d$.
Find the numbers.
\begin{Soln}
\begin{DPalign*}
\lintertext{\indent Let}
x &= \text{the smaller number;} \\
\lintertext{then}
x + d &= \text{the larger number.} \displaybreak[1] \\
\lintertext{\indent Hence}
x + x + d &= s, \\
\lintertext{or}
2x &= s - d. \\
\therefore x &= \frac{s - d}{2}, \displaybreak[1] \\
\lintertext{and}
x + d &= \frac{s - d}{2} + d = \frac{s - d + 2d}{2}, \\
  &= \frac{s + d}{2}.
\end{DPalign*}

Therefore the numbers are $\dfrac{s + d}{2}$ and $\dfrac{s - d}{2}$.
\end{Soln}

As these formulas hold true whatever numbers $s$~and~$d$
stand for, we have the general rule for finding two numbers
when their sum and difference are given:
\begin{Theorem}
Add the difference to the sum and take half the result for
the greater number.

Subtract the difference from the sum and take half the
result for the smaller number.
\end{Theorem}

\Item{2.}  If A~can do a piece of work in $a$~days, and B~can
do the same work in $b$~days, in how many days can both
together do it?
%% -----File: 125.png---Folio 119-------
\begin{Soln}
\begin{DPalign*}
\lintertext{\indent Let}
x &= \text{the required number of days\Add{.}} \\
\lintertext{\indent Then,}
\frac{1}{x} &= \text{\emph{the part both together can do in one day}\Add{.}}
\displaybreak[1] \\
\lintertext{\indent Now}
\frac{1}{a} &= \text{the part A can do in one day,} \\
\lintertext{and}
\frac{1}{b} &= \text{the part B can do in one day\Add{,}} \\
\lintertext{therefore}
\frac{1}{a} + \frac{1}{b}
  &= \text{\emph{the part both together can do in one day}}
\end{DPalign*}
\begin{DPalign*}
\frac{1}{a} + \frac{1}{b} &= \frac{1}{x}\Add{.} \\
\lintertext{\indent Whence}
x &= \frac{ab}{a + b}\Add{.}
\end{DPalign*}
\end{Soln}

The translation of this formula gives the following rule
for finding the time required by two agents together to
produce a given result when the time required by each
agent separately is known.
\begin{Theorem}
Divide the product of the numbers which express the units
of time required by each to do the work by the sum of these
numbers, the quotient is the time required by both together.
\end{Theorem}

\Paragraph{145.  Interest Formulas.}  The elements involved in computation
of interest are the \emph{principal}, \emph{rate}, \emph{time}, \emph{interest},
and \emph{amount}.
\begin{DPalign*}
\lintertext{\indent Let}
p &= \text{the principal,} \\
r &= \text{the interest of \$$1$ for $1$~year, at the given rate,} \\
t &= \text{the time expressed in years,} \\
i &= \text{the interest for the given time and rate,} \\
a &= \text{the amount (sum of principal and interest).}
\end{DPalign*}

\Paragraph{146.  Given the Principal, Rate, and Time.  Find the Interest.}

Since $r$ is the interest of~\$$1$ for $1$~year, $pr$~is the interest
of~\$$p$ for $1$~year, and $prt$~is the interest of~\$$p$ for $t$~years
\begin{DPgather*}
i = prt.
\rintertext{(Formula 1.)}
\end{DPgather*}

\begin{Theorem}[\textsc{Rule.}]  Find the product of the principal, rate, and time\Add{.}
\end{Theorem}
%% -----File: 126.png---Folio 120-------

\Paragraph{147.  Given the Interest, Rate, and Time.  Find the Principal.}
\begin{DPalign*}
\lintertext{\indent By formula 1,}
prt &= i. \\
\lintertext{\indent Divide by~$rt$,}
p &= \frac{i}{rt}.
\rintertext{(Formula 2.)}
\end{DPalign*}

\begin{Theorem}[\textsc{Rule.}]  Divide the interest by the product of the rate and
time.
\end{Theorem}

\Paragraph{148.  Given the Amount, Rate, and Time.  Find the Principal.}
\begin{DPalign*}
\lintertext{\indent From formula 1,}
p + prt &= a, \\
\lintertext{or}
p(1 + rt) &= a. \\
\lintertext{\indent Divide by $1 + rt$,}
p &= \frac{a}{1 + rt}.
\rintertext{(Formula 3.)}
\end{DPalign*}

\begin{Theorem}[\textsc{Rule.}]  Divide the amount by $1$~plus the product of the
rate and time.
\end{Theorem}

\Paragraph{149.  Given the Amount, Principal, and Rate.  Find the Time.}
\begin{DPalign*}
\lintertext{\indent From formula 1,}
p + prt &= a. \\
\lintertext{\indent Transpose~$p$,}
prt &= a - p. \\
\lintertext{\indent Divide by~$pr$,}
t &= \frac{a - p}{pr}.
\rintertext{(Formula 4.)}
\end{DPalign*}

\begin{Theorem}[\textsc{Rule.}]  Subtract the principal from the amount, and
divide the result by the product of the principal and rate.
\end{Theorem}

\Paragraph{150.  Given the Amount, Principal, and Time.  Find the Rate.}
\begin{DPalign*}
\lintertext{\indent From formula 1,}
p + prt &= a. \\
\lintertext{\indent Transpose~$p$,}
prt &= a - p. \\
\lintertext{\indent Divide by~$pt$,}
r &= \frac{a - p}{pt}.
\rintertext{(Formula 5.)}
\end{DPalign*}

\begin{Theorem}[\textsc{Rule.}]  Subtract the principal from the amount, and
divide the result by the product of the principal and time.
\end{Theorem}
%% -----File: 127.png---Folio 121-------

\Exercise{64.}

Solve the following examples by the preceding formulas:

\Item{1.}  The sum of two angles is $120°\, 30'\, 30''$ and their difference
$59°\, 30'\, 30''$.  Find the angles.

\Item{2.}  Find the interest of \$$1000$ for $3$~years and $4$~months
at~$4$\%.

\Item{3.}  Find the principal that will amount to \$$2280$ in $3$~years
and $6$~months at~$4$\%.

\Item{4.}  Find the principal that will produce \$$280$ interest in
$2$~years and $4$~months at~$3$\%.

\Item{5.}  Find the principal that will produce \$$270$ interest
in $1$~year and $6$~months at~$6$\%.

\Item{6.}  Find the principal that will amount to \$$590$ in $4$~years
at~$4\frac{1}{2}$\%.

\Item{7.}  Find the rate if the amount of \$$250$ for $4$~years
is~\$$300$.

\Item{8.}  Find the rate if \$$1000$ amounts to \$$2000$ in $16$~years
and $8$~months.

\Item{9.}  Find the time required for the interest on \$$400$ to
be \$$54$ at~$4\frac{1}{2}$\%.

\Item{10.}  Find the time required for \$$160$ to amount to \$$250$
at~$6$\%.

\Item{11.}  How much money must be invested at~$5$\% to yield
an annual income of~\$$1250$?

\Item{12.}  Find the principal that will produce \$$100$ a month
if invested at $6$\%~per annum.

\Item{13.}  Find the rate if the interest on \$$1000$ for $8$~months
is~\$$40$.

\Item{14.}  Find the time for a sum of money on interest at~$5$\%
to double itself.
%% -----File: 128.png---Folio 122-------


\Chapter{XI.}{Simultaneous Equations of the First
Degree.}

\Paragraph{151.}  If we have two unknown numbers and but one relation
between them, we can find an unlimited number of
pairs of values for which the given relation will hold true.
Thus, if $x$~and~$y$ are unknown, and we have given only the
one relation $x + y = 10$, we can \emph{assume} any value for~$x$,
and then from the relation $x + y = 10$ find the corresponding
value of~$y$.  For from $x + y = 10$ we find $y = 10 - x$.
If $x$~stands for~$1$, $y$~stands for~$9$; if $x$~stands for~$2$, $y$~stands
for~$8$; if $x$~stands for~$-2$, $y$~stands for~$12$; and so on without
end.

\Paragraph{152.}  We may, however, have two equations that express
\emph{different} relations between the two unknown numbers.
Such equations are called \Defn{independent equations}.  Thus,
$x + y = 10$ and $x - y = 2$ are independent equations, for
they evidently express \emph{different} relations between $x$~and~$y$.

\Paragraph{153.}  Independent equations involving the \emph{same} unknown
numbers are called \Defn{simultaneous equations}.

If we have two unknown numbers, and two independent
equations involving them, there is but \emph{one} pair of values
which will hold true for both equations.  Thus, if besides
the relation $x + y = 10$, we have also the relation $x - y = 2$,
the only pair of values for which both equations will hold
true is the pair $x = 6$, $y = 4$.

Observe that in this problem $x$~stands for the same number
in \emph{both} equations; so also does~$y$.
%% -----File: 129.png---Folio 123-------

\Paragraph{154.}  Simultaneous equations are solved by combining
the equations so as to obtain a single equation with one
unknown number.

This process is called \Defn{Elimination}.

\Paragraph{155.  Elimination by Addition or Subtraction.}

%[** TN: Omitted large brace for grouping, horizontal bar indicating summation]
\Item{1.} Solve:
\begin{alignat*}{2}
5x &- 3y &&= 20
\Tag{(1)} \\
2x &+ 5y &&= 39
\Tag{(2)}
\end{alignat*}
\begin{Soln}
Multiply (1) by~$5$, and (2) by~$3$,
\begin{DPalign*}
25x - 15y &= 100
\Tag{(3)} \\
 6x + 15y &= 117
\Tag{(4)} \\
\lintertext{\indent Add (3) and (4),}
31x \PadTo{{} + 15y}{} &= 217 \\
\therefore x &= 7.
\intertext{\indent Substitute the value of~$x$ in~(2),}
14 + 5y &= 39. \\
5y &= 25. \\
\therefore y &= 5.
\end{DPalign*}

In this solution $y$~is eliminated by \emph{addition}.
\end{Soln}

\Item{2.}  Solve:
\begin{alignat*}{2}
6x &+ 35y &&= 177
\Tag{(1)} \\
8x &- 21y &&= \Z33
\Tag{(2)}
\end{alignat*}
\begin{Soln}
Multiply (1) by~$4$, and (2) by~$3$,
\begin{DPalign*}
24x + 140y &= 708
\Tag{(3)} \\
24x - \Z63y &= \Z99
\Tag{(4)} \\
\lintertext{\indent Subtract,}
203y &= 609 \\
\therefore y &= 3. \\
\intertext{\indent Substitute the value of~$y$ in (2),}
8x - 63 &= 33. \\
8x &= 96. \\
\therefore x &= 12.
\end{DPalign*}

In this solution $x$~is eliminated by \emph{subtraction}.
\end{Soln}
%% -----File: 130.png---Folio 124-------

\Paragraph{156.}  \Dictum{To eliminate by addition or subtraction}, therefore,
\begin{Theorem}
Multiply the equations by such numbers as will make the
coefficients of one of the unknown numbers equal in the
resulting equations.

Add the resulting equations if these equal coefficients have
unlike signs; subtract one from the other if these equal coefficients
have like signs.
\end{Theorem}

\begin{Remark}[Note.] It is generally best to select the letter to be eliminated
which requires the smallest multipliers to make its coefficients equal;
and the smallest multiplier for each equation is found by dividing
the \LCM\ of the coefficients of this letter by the given coefficient
in that equation.  Thus, in example~2, the \LCM\ of $6$~and~$8$ (the
coefficients of~$x$) is~$24$, and hence the smallest multipliers of the two
equations are $4$~and~$3$, respectively.
\end{Remark}

Sometimes the solution is simplified by first adding the
given equations, or by subtracting one from the other.
\begin{Soln}
\begin{DPalign*}
\lintertext{\indent Ex.}
x + 49y &= 51
\Tag{(1)} \\
49x + \PadTo{49y}{y} &= 99
\Tag{(2)} \displaybreak[1] \\
\lintertext{\indent Add (1) and (2),}
50x + 50y &= 150
\Tag{(3)} \displaybreak[1] \\
\lintertext{\indent Divide (3) by~$50$,}
x + y &= 3.
\Tag{(4)} \displaybreak[1] \\
\lintertext{\indent \rlap{Subtract (4) from (1),}}
48y &= 48. \\
\therefore y &= 1. \displaybreak[1] \\
\lintertext{\indent \rlap{Subtract (4) from (2),}}
48x &= 96. \\
\therefore x &= 2.
\end{DPalign*}
\end{Soln}

\Exercise{65.}

Solve by addition or subtraction:
\begin{multicols}{2}
\Item{1.}  \raisebox{-0.5\baselineskip}{$\left.\begin{alignedat}{2}
 5x &+ 4y &&= 14 \\
17x &- 3y &&= 31
\end{alignedat}\right\}$}

\Item{2.}  \raisebox{-0.5\baselineskip}{$\left.\begin{alignedat}{2}
 3x &- 2y &&= \Z5 \\
 2x &+ 5y &&= 16
\end{alignedat}\right\}$}

\Item{3.}  \raisebox{-0.5\baselineskip}{$\left.\begin{alignedat}{2}
 2x &- 3y &&= \Z7 \\
 5x &+ 2y &&= 27
\end{alignedat}\right\}$}

\Item{4.}  \raisebox{-0.5\baselineskip}{$\left.\begin{alignedat}{2}
 7x &+ 6y &&= 20 \\
 2x &+ 5y &&= \Z9
\end{alignedat}\right\}$}

\Item{5.}  \raisebox{-0.5\baselineskip}{$\left.\begin{alignedat}{2}
  x &+ 5y &&= 11 \\
 3x &+ 2y &&= \Z7
\end{alignedat}\right\}$}

\Item{6.}  \raisebox{-0.5\baselineskip}{$\left.\begin{alignedat}{2}
 3x &- 5y &&= 13 \\
 4x &- 7y &&= 17
\end{alignedat}\right\}$}
%% -----File: 131.png---Folio 125-------

\Item{7.}  \raisebox{-0.5\baselineskip}{$\left.\begin{alignedat}{2}
 8x &- \Z y &&= \Z3 \\
 7x &+ 2y &&= 63
\end{alignedat}\right\}$}

\Item{8.}  \raisebox{-0.5\baselineskip}{$\left.\begin{alignedat}{2}
 5x &- 4y &&= \Z7 \\
 7x &+ 3y &&= 70
\end{alignedat}\right\}$}

\Item{9.}  \raisebox{-0.5\baselineskip}{$\left.\begin{alignedat}{2}
  x &+ 21y &&= \Z2 \\
 2x &+ 27y &&= 19
\end{alignedat}\right\}$}

\Item{10.}  \raisebox{-0.5\baselineskip}{$\left.\begin{alignedat}{2}
 6x &- 13y &&= -1 \\
 5x &- 12y &&= -2
\end{alignedat}\right\}$}

\Item{11.}  \raisebox{-0.5\baselineskip}{$\left.\begin{alignedat}{2}
 7x &+ \Z y &&= 265 \\
 3x &- 5y &&= \Z\Z5
\end{alignedat}\right\}$}

\Item{12.}  \raisebox{-0.5\baselineskip}{$\left.\begin{alignedat}{2}
 2x &+ 3y &&= \Z7 \\
 8x &- 5y &&= 11
\end{alignedat}\right\}$}

\Item{13.}  \raisebox{-0.5\baselineskip}{$\left.\begin{alignedat}{2}
 5x &+ 7y &&= 19 \\
 7x &+ 4y &&= 15
\end{alignedat}\right\}$}

\Item{14.}  \raisebox{-0.5\baselineskip}{$\left.\begin{alignedat}{2}
11x &- 12y &&= \Z9 \\
 4x &+ \Z5y &&= 22
\end{alignedat}\right\}$}

\Item{15.}  \raisebox{-0.5\baselineskip}{$\left.\begin{alignedat}{2}
  x &+ 8y &&= 17 \\
 7x &- 3y &&= \Z1
\end{alignedat}\right\}$}

\Item{16.}  \raisebox{-0.5\baselineskip}{$\left.\begin{alignedat}{2}
 4x &+ 3y &&= 25 \\
 5x &- 4y &&= \Z8
\end{alignedat}\right\}$}
\end{multicols}

Clear of fractions and solve:
\begin{multicols}{2}
\Item{17.}  \raisebox{-0.5\baselineskip}{$\left.\begin{alignedat}{2}
 \frac{2x}{3} &- \frac{5y}{4} &&= 3 \\
 \frac{7x}{4} &- \frac{5y}{3} &&= \frac{43}{3}
\end{alignedat}\right\}$}

\Item{18.}  \raisebox{-0.5\baselineskip}{$\left.\begin{alignedat}{2}
 \frac{7x}{6} &+ \frac{6y}{7} &&= 32 \\
 \frac{5x}{4} &- \frac{2y}{3} &&= 1
\end{alignedat}\right\}$}

\Item{19.}  \raisebox{-0.5\baselineskip}{$\left.\begin{aligned}
  \frac{x + y}{4} - \frac{7x - 5y}{11} &= 3 \\
  \frac{x}{5} - \frac{2y}{7} + 1 &= 0
\end{aligned}\right\}$}

\Item{20.}  \raisebox{-0.5\baselineskip}{$\left.\begin{aligned}
 \frac{ 6x + 7y}{2} &= 22 \\
 \frac{55y - 2x}{5} &= 20
\end{aligned}\right\}$}
\end{multicols}

\Item{21.}  \raisebox{-0.5\baselineskip}{$\left.\begin{alignedat}{2}
  \frac{x + y}{2} &- \frac{x - y}{3} &&= \Z8 \\
  \frac{x + y}{3} &+ \frac{x - y}{4} &&= 11
\end{alignedat}\right\}$}

\Item{22.}  \raisebox{-0.5\baselineskip}{$\left.\begin{aligned}
 \frac{8x - 5y}{7} + \frac{11y - 4x}{5} &= 4 \\
 \frac{17x - 13y}{5} + \frac{2x}{3} &= 7
\end{aligned}\right\}$}
%% -----File: 132.png---Folio 126-------

\Item{23.}  \raisebox{-0.5\baselineskip}{$\left.\begin{alignedat}{2}
 \frac{ 5x - 3y}{3} &+ \frac{7x - 5y}{11} &&= 4 \\
 \frac{15y - 3x}{7} &+ \frac{7y - 3x}{5} &&= 4
   \end{alignedat}\right\}$}

\Item{24.}  \raisebox{-0.5\baselineskip}{$\left.\begin{aligned}
 \frac{2x - 3}{4} - \frac{y - 8}{5} &= \frac{y + 3}{4} \\
 \frac{x - 7}{3} + \frac{4y + 1}{11} &= 3
  \end{aligned}\right\}$}

\Item{25.}  \raisebox{-0.5\baselineskip}{$\left.\begin{alignedat}{2}
  \frac{x - 2y}{6} &- \frac{x + 3y}{4} &&= \frac{3}{2} \\
  \frac{2x - y}{6} &- \frac{3x + y}{4} &&= \frac{5y}{4}
 \end{alignedat}\right\}$}

\Item{26.}  \raisebox{-0.5\baselineskip}{$\left.\begin{alignedat}{2}
  \frac{x}{a + b} &+ \frac{y}{a - b} &&= \frac{1}{a - b} \\
  \frac{x}{a + b} &- \frac{y}{a - b} &&= \frac{1}{a + b}
  \end{alignedat}\right\}$}

\begin{Remark}[Note.]  To find $x$ in problem~26, add the equations; to find~$y$,
subtract one from the other.  Do not clear of fractions.
\end{Remark}

\Paragraph{157.  Problems involving Two Unknown Numbers.}

Ex.  If A gives B \$$10$, B~will have three times as much
money as~A\@. If B gives A \$$10$, A~will have twice as
much money as~B\@. How much has each?
\begin{Soln}
\begin{DPalign*}
\lintertext{\indent Let}
x &= \text{the number of dollars A has,} \\
\lintertext{and}
y &= \text{the number of dollars B has.} \displaybreak[1] \\
\intertext{\indent Then, after A gives B \$$10,$}
x - 10 &= \text{the number of dollars A has,} \\
y + 10 &= \text{the number of dollars B has.} \displaybreak[1] \\
\intertext{\indent Since B's money is now $3$~times A's, we have,}
y + 10 &= 3(x - 10).
\Tag{(1)} \displaybreak[1] \\
%% -----File: 133.png---Folio 127-------
\intertext{\indent If B gives A \$$10$,}
x + 10 &= \text{the number of dollars A has,} \\
y - 10 &= \text{the number of dollars B has.} \displaybreak[1] \\
\intertext{\indent Since A's money is now $2$~times B's, we have}
x + 10 &= 2(y - 10).
\Tag{(2)}
\end{DPalign*}

From the solution of equations (1) and (2), $x = 22$, and $y = 26$.

Therefore A has \$$22$, and B has~\$$26$.
\end{Soln}

\Exercise{66.}

\Item{1.}  If A gives B \$$200$, A~will then have half as much
money as~B; but if B gives A \$$200$, B~will have one-third
as much as A\@.  How much has each?

\Item{2.}  Half the sum of two numbers is~$20$, and three times
their difference is~$18$.  Find the numbers.

\Item{3.}  The sum of two numbers is~$36$, and their difference
is equal to one-eighth of the smaller number increased by~$2$.
Find the numbers.

\Item{4.}  If $4$~yards of velvet and $3$~yards of silk are sold for~\$$33$,
and $5$~yards of velvet and $6$~yards of silk for~\$$48$,
what is the price per yard of the velvet and of the silk?

\Item{5.}  If $7$~bushels of wheat and $10$~of rye are sold for~\$$15$,
and $4$~bushels of wheat and $5$~of rye are sold for~\$$8$, what
is the price per bushel of the wheat and of the rye?

\Item{6.}  If $12$~pounds of tea and $4$~pounds of coffee cost~\$$7$,
and $4$~pounds of tea and $12$~pounds of coffee cost~\$$5$, what is
the price per pound of tea and of coffee?

\Item{7.}  Six horses and $7$~cows can be bought for~\$$1000$,
and $11$~horses and $13$~cows for~\$$1844$.  Find the value of
a horse and of a cow.
%% -----File: 134.png---Folio 128-------

\Exercise{67.}

Ex.  A certain fraction becomes equal to~$\frac{1}{2}$ if $2$~is added
to its numerator, and equal to~$\frac{1}{3}$ if $3$~is added to its denominator.
Find the fraction.
\begin{Soln}
\begin{DPalign*}
\lintertext{\indent Let}
\frac{x}{y} &= \text{the required fraction.} \\
\lintertext{\indent Then}
\frac{x + 2}{y} &= \frac{1}{2}, \\
\lintertext{and}
\frac{x}{y + 3} &= \frac{1}{3}.
\end{DPalign*}

The solution of these equations gives $7$ for~$x$, and $18$ for~$y$.

Therefore the required fraction is~$\frac{7}{18}$.
\end{Soln}

\Item{1.}  If the numerator of a certain fraction is increased
by~$2$ and its denominator diminished by~$2$, its value will
be~$1$.  If the numerator is increased by the denominator
and the denominator is diminished by~$5$, its value will be~$5$.
Find the fraction.

\Item{2.}  If $1$~is added to the denominator of a fraction, its
value will be~$\frac{1}{2}$.  If $2$~is added to its numerator, its value
will be~$\frac{3}{5}$.  Find the fraction.

\Item{3.}  If $1$~is added to the numerator of a fraction, its value
will be~$\frac{1}{5}$.  If $1$~is added to its denominator, its value will
be~$\frac{1}{7}$.  Find the fraction.

\Item{4.}  If the numerator of a fraction is doubled and its
denominator diminished by~$1$, its value will be~$\frac{1}{2}$.  If its
denominator is doubled and its numerator increased by~$1$,
its value will be~$\frac{1}{7}$.  Find the fraction.

\Item{5.}  In a certain proper fraction the difference between
the numerator and the denominator is~$15$.  If the numerator
is multiplied by~$4$ and the denominator increased by~$6$,
its value will be~$1$.  Find the fraction.
%% -----File: 135.png---Folio 129-------

\Exercise{68.}

The expression $64$ means $60 + 4$, that is, $10$~\emph{times} $6 + 4$,
and has for its \emph{digits} $6$~and~$4$.  If the digits were unknown
and represented by $x$~and~$y$, the number would be represented
by~$10x + y$.

Ex.  The sum of the two digits of a number is~$10$, and if
$18$~is added to the number, the digits will be reversed.
Find the number.
\begin{Soln}
\begin{DPalign*}[m]
\lintertext{\indent Let}
x &= \text{the tens' digit,} \\
\lintertext{and}
y &= \text{the units' digit.} \\
\lintertext{\indent Then}
10x + y &= \text{the number.} \displaybreak[1] \\
\lintertext{\indent Hence}
x + y &= 10,
\Tag{(1)} \\
\lintertext{and}
10x + y + 18 &= 10y + x.
\Tag{(2)} \displaybreak[1] \\
\lintertext{\indent From (2),}
9x - 9y &= -18, \\
\lintertext{or}
x - y &= -2.
\Tag{(3)} \displaybreak[1] \\
\lintertext{\indent \rlap{Add (1) and (3),}}
2x &= 8, \\
\lintertext{and therefore}
x &= 4. \\
\lintertext{\indent \rlap{Subtract (3) from (1),}}
2y &= 12, \\
\lintertext{and therefore}
y &= 6.
\end{DPalign*}

Therefore the number is~$46$.
\end{Soln}

\Item{1.}  The sum of the two digits of a number is~$9$, and if $9$
is added to the number, the digits will be reversed.  Find
the number.

\Item{2.}  A certain number of two digits is equal to eight times
the sum of its digits, and if $45$~is subtracted from the
number, the digits will be reversed.  Find the number.

\Item{3.}  The sum of a certain number of two digits and the
number formed by reversing the digits is~$132$, and the
difference of these numbers is~$18$.  Find the numbers.

\Item{4.}  The sum of the two digits of a number is~$9$, and if
the number is divided by the sum of its digits, the quotient
is~$6$.  Find the number.
%% -----File: 136.png---Folio 130-------

\Exercise{69.}

Ex.  A sum of money, at simple interest, amounted to
\$$2480$ in $4$~years, and to \$$2600$ in $5$~years.  Find the sum
and the rate of interest.
\begin{Soln}
\begin{DPalign*}
\lintertext{\indent Let}
x &= \text{the number of dollars in the principal,} \\
\lintertext{and}
y &= \text{the rate of interest.}
\end{DPalign*}

The interest for one year is $\dfrac{y}{100}$ of the principal; that is, $\dfrac{xy}{100}$.
For $4$~years the interest is~$\dfrac{4xy}{100}$, and for $5$~years~$\dfrac{5xy}{100}$.  The amount
is principal $+$~interest,
\begin{DPalign*}
\lintertext{or}
x + \frac{4xy}{100} &= 2480. \\
x + \frac{5xy}{100} &= 2600. \\
\lintertext{\indent Hence}
100x + 4xy &= 248,000.
\Tag{(1)} \\
100x + 5xy &= 260,000.
\Tag{(2)} \displaybreak[1] \\
\intertext{\indent Divide (1) by~$4$ and (2) by~$5$, and we have}
25x + xy &= 62,000 \\
20x + xy &= 52,000 \displaybreak[1] \\
\lintertext{\indent Subtract,}
5x \PadTo{{}+ xy}{} &= 10,000. \displaybreak[1] \\
\lintertext{\indent Therefore}
\PadTo{5x}{x} \PadTo{{}+ xy}{} &= \PadTo{10,000}{2000}.
\end{DPalign*}

Substitute the value of~$x$ in~(1), $y = 6$.

Therefore the sum is \$$2000$, and the rate~$6$%.
\end{Soln}

\Item{1.}  A sum of money, at simple interest, amounted in $5$~years
to~\$$3000$, and in $6$~years to~\$$3100$.  Find the sum
and the rate of interest.

\Item{2.}  A sum of money, at simple interest, amounted in $10$~months
to~\$$1680$, and in $18$~months to~\$$1744$.  Find the
sum and the rate of interest.

\Item{3.}  A man has \$$10,000$ invested, a part at~$4$\%, and the
remainder at~$5$\%.  The annual income from his $4$\%~investment
is \$$40$~more than from his $5$\%~investment.  Find the
sum invested at~$4$\% and at~$5$\%.
%% -----File: 137.png---Folio 131-------

\Exercise{70.}

\Section{Miscellaneous Examples.}

\Item{1.}  Half the sum of two numbers is~$20$; and $5$~times
their difference is~$20$.  Find the numbers.

\Item{2.}  A certain number when divided by a second number
gives $7$ for a quotient and $4$ for a remainder.  If three
times the first number is divided by twice the second
number, the quotient is~$11$ and the remainder~$4$.  Find
the numbers.

\Item{3.}  A fraction becomes $\frac{4}{5}$ in value by the addition of~$2$ to
its numerator and $3$~to its denominator.  If $2$~is subtracted
from its numerator and $1$~from its denominator, the value
of the fraction is~$\frac{3}{4}$.  Find the fraction.

\Item{4.}  A farmer sold $50$~bushels of wheat and $30$~of barley
for $74$~dollars; and at the same prices he sold $30$~bushels
of wheat and $50$~bushels of barley for $70$~dollars.  What
was the price of the wheat and of the barley per bushel?

\Item{5.}  If A gave \$$10$ to~B, he would then have three times
as much money as~B; but if B gave \$$5$ to~A, A~would
have four times as much as~B\@. How much has each?

\Item{6.}  A and B have together \$$100$.  If A~were to spend
one-half of his money, and B~one-third of his, they would
then have only \$$55$ between them.  How much money
has each?

\Item{7.}  A fruit-dealer sold $6$~lemons and $3$~oranges for $21$~cents,
and $3$~lemons and $8$~oranges for $30$~cents.  What
was the price of each?

\Item{8.}  If A gives me $10$~apples, he will have just twice as
many as~B\@.  If he gives the $10$~apples to~$B$ instead of to
me, A~and~B will each have the same number.  How
many apples has each?
%% -----File: 138.png---Folio 132-------


\Chapter{XII.}{Quadratic Equations.}

\Paragraph{158.}  An equation which contains the \emph{square} of the
unknown number, but no higher power, is called a \Defn{quadratic
equation}.

\Paragraph{159.}  A quadratic equation which involves but one unknown
number as~$x$, can contain only:

1.  Terms involving the square of~$x$.

2.  Terms involving the first power of~$x$.

3.  Terms which do not involve~$x$.

Collecting similar terms, every quadratic equation can be
made to assume the form
\[
ax^{2} + bx + c = 0,
\]
where $a$,~$b$, and~$c$ are known numbers, and $x$~the unknown
number.

If $a$,~$b$,~$c$ are numbers expressed by figures, the equation
is a \Defn{numerical quadratic}.  If $a$,~$b$,~$c$ are numbers represented
wholly or in part by letters, the equation is a \Defn{literal quadratic}.

\Paragraph{160.}  In the equation $ax^{2} + bx + c = 0$, $a$,~$b$, and~$c$ are
called the \Defn{coefficients} of the equation.  The third term,~$c$, is
called the \Defn{constant term}.

If the first power of~$x$ is wanting, the equation is a \Defn{pure
quadratic}; in this case $b = 0$.

If the first power of~$x$ is present, the equation is an
\Defn{affected} or \Defn{complete quadratic}.
%% -----File: 139.png---Folio 133-------


\Section{Pure Quadratic Equations.}

\Paragraph{161.  Examples.}

\Item{1.}  Solve the equation $5x^{2} - 48 = 2x^{2}$.
\begin{Soln}
\begin{DPalign*}
\lintertext{\indent We have}
5x^{2} - 48 &= 2x^{2}. \\
\lintertext{\indent Collect the terms,}
3x^{2} &= 48. \displaybreak[1] \\
\lintertext{\indent Divide by~$3$,}
x^{2} &= 16. \\
\lintertext{\indent Extract the square root,}
x &= ±4.
\end{DPalign*}

The sign~$±$ before the~$4$, read \emph{plus or minus}, shows that the root
is either $+$~or~$-$. For $(+4) × (+4) = 16$, and $(-4) × (-4) = 16$

The square root of any number is positive or negative.  Hitherto
we have given only the positive value.  In this chapter we shall
give both values.  This sign~$\surd$, called the \Defn{radical sign}, is used to
indicate that a root is to be extracted.  Thus $\sqrt{4}$~means the square
root of~$4$ is required. $\sqrt[3]{4}$~means the third root of~$4$ is required; the
small figure placed over the radical sign is called the \Defn{index} of the
root, and shows the root required.
\end{Soln}

\Item{2.}  Solve the equation $3x^{2} - 15 = 0$.
\begin{Soln}
\begin{DPalign*}
\lintertext{\indent We have}
3x^{2} &= 15, \\
\lintertext{or}
x^{2} &= 5. \\
\lintertext{\indent Extract the square root,}
x &= ±\sqrt{5}.
\end{DPalign*}

The roots cannot be found exactly, since the square root of~$5$ cannot
be found exactly; it can, however, be determined approximately
to any required degree of accuracy; for example, the roots lie between
$2.23606$ and $2.23607$; and between $-2.23606$ and~$-2.23607$.
\end{Soln}

\Item{3.}  Solve the equation $3x^{2} + 15 = 0$.
\begin{Soln}
\begin{DPalign*}
\lintertext{\indent We have}
3x^{2} &= -15, \\
\lintertext{or}
x^{2} &= -5. \\
\lintertext{\indent Extract the square root,}
x &= ±\sqrt{-5}.
\end{DPalign*}

There is no square root of a negative number, since the square of
any number, positive or negative, is positive; $(-5) × (-5) = +25$.

The square root of~$-5$ differs from the square root of~$+5$ in that
the latter can be found as accurately as we please, while the former
cannot be found at all.
\end{Soln}
%% -----File: 140.png---Folio 134-------

\Paragraph{162.}  A root which can be found exactly is called an \Defn{exact}
or \Defn{rational} root.  Such roots are either whole numbers or
fractions.

A root which is indicated but can be found only approximately
is called a \Defn{surd}.  Such roots involve the roots of
imperfect powers.

Rational and surd roots are together called \Defn{real} roots.

A root which is indicated but cannot be found, either
exactly or approximately, is called an \Defn{imaginary} root.  Such
roots involve the even roots of negative numbers.

\Exercise{71.}

Solve:
\begin{multicols}{2}
\Item{1.}  $5x^{2} - 2 = 3x^{2} + 6$.

\Item{2.}  $3x^{2} + 1 = 2x^{2} + 10$.

\Item{3.}  $4x^{2} - 50 = x^{2} + 25$.

\Item{4.}  $(x - 6)(x + 6) = 28$.
\end{multicols}

\Item{5.}  $(x - 5)(x + 5) = 24$.

\Item{6.}  $3(x^{2} - 11) + 2(x^{2} - 5) = 82$.

\Item{7.}  $11(x^{2} + 5) + 6(3 - x^{2}) = 198$.

\Item{8.}  $5x^{2} + 3 - 2(17 - x^{2}) = 32$.

\Item{9.}  $4(x + 1) - 4(x - 1) = x^{2} - 1$.

\Item{10.}  $86 - 52x = 2(8 - x)(2 - 3x)$.

\Item{11.}  Find two numbers that are to each other as $3$~to~$4$,
and the difference of whose squares is~$112$.

\begin{Remark}[Hint.]  Let $3x$~stand for the smaller and $4x$~for the larger number.
\end{Remark}

\Item{12.}  A boy bought a number of oranges for $36$~cents.
The price of an orange was to the number bought as $1$~to~$4$.
How many oranges did he buy, and how many cents
did each orange cost?

\Item{13.}  A certain street contains $144$~square rods, and the
length is $16$~times the width.  Find the width.
%% -----File: 141.png---Folio 135-------

\Item{14.}  Find the number of rods in the length, and in the
width of a rectangular field containing $3\frac{3}{5}$~acres, if the
length is $4$~times the width.


\Section{Affected Quadratic Equations.}

\Paragraph{163.}  Since
\[
(x + b)^{2} = x^{2} + 2bx + b^{2}, \quad\text{and}\quad
(x - b)^{2} = x^{2} - 2bx + b^{2},
\]
it is evident that the expression $x^{2} + 2bx$ or $x^{2} - 2bx$ lacks
only the \emph{third term},~$b^{2}$, of being a perfect square.

This third term is the square of half the coefficient of~$x$.

Every affected quadratic may be made to assume the
form $x^{2} + 2bx = c$ or $x^{2} - 2bx = c$, by dividing the equation
through by the coefficient of~$x^{2}$.

To \textbf{solve} such an equation:

The first step is to add to both members \emph{the square of
half the coefficient of~$x$}.  This is called \emph{completing the square}.

The second step is to \emph{extract the square root} of each member
of the resulting equation.

The third step is to \emph{reduce} the two resulting simple
equations.

\Item{1.}  Solve the equation $x^{2} - 8x = 20$.
\begin{Soln}
\begin{DPalign*}
\lintertext{\indent We have}
x^{2} - 8x &= 20. \\
\lintertext{\indent Complete the square,}
x^{2} - 8x + 16 &= 36. \\
\lintertext{\indent Extract the square root,}
x - 4 &= ±6. \displaybreak[1] \\
\lintertext{\indent Reduce, using the upper \rlap{sign,}}
x &= 4 + 6 = 10, \\
\lintertext{or using the lower sign,}
x &= 4 - 6 = -2.
\end{DPalign*}

The roots are $10$ and~$-2$.

Verify by putting these numbers for~$x$ in the given equation.
\[
\begin{array}{rcl<{\qquad}|>{\qquad}rcl}
             x &=& 10, &                x &=& -2, \\
10^{2} - 8(10) &=& 20, & (-2)^{2} - 8(-2) &=& 20, \\
      100 - 80 &=& 20. &           4 + 16 &=& 20. \\
\end{array}
\]
\end{Soln}
%% -----File: 142.png---Folio 136-------

\Item{2.}  Solve the equation $\dfrac{x + 1}{x - 1} = \dfrac{4x - 3}{x + 9}$.
\begin{Soln}
\begin{DPalign*}
\lintertext{\indent Free from fractions,}
(x + 1)(x + 9) &= (x - 1)(4x - 3). \\
\lintertext{\indent Therefore,}
-3x^{2} + 17x &= -6.
\end{DPalign*}

Since the square root of a negative number cannot be taken, the
coefficient of~$x^{2}$ must be changed to~$+$.
\begin{DPalign*}
\lintertext{\indent Divide by~$-3$,}
x^{2} - \tfrac{17}{3}x &= 2. \displaybreak[1] \\
\intertext{\indent
Half the coefficient of~$x$ is $\frac{1}{2}$~of $-\frac{17}{3} = -\frac{17}{6}$, and the square of~$-\frac{17}{6}$
is~$\frac{289}{36}$.  Add the square of~$-\frac{17}{6}$ to both sides, and we have}
x^{2} - \frac{17x}{3} + \left(\frac{17}{6}\right)^{2} &= 2 + \frac{289}{36}. \displaybreak[1] \\
\lintertext{\indent Now}
2 + \frac{289}{36} = \frac{72}{36} + \frac{289}{36} &= \frac{361}{36}, \\
\lintertext{therefore,}
x^{2} - \tfrac{17}{3}x + \left(\frac{17}{6}\right)^{2} &= \frac{361}{36}. \displaybreak[1] \\
\lintertext{\indent \rlap{Extract the root,}}
x - \frac{17}{6} &= ±\frac{19}{6}. \displaybreak[1] \\
\lintertext{\indent Reduce,}
x - \frac{17}{6} &= ±\frac{19}{6}. \\
\therefore x &= \frac{17}{6} + \frac{19}{6} = \frac{36}{6} = 6, \\
\lintertext{or}
x &= \frac{17}{6} - \frac{19}{6} = -\frac{2}{6} = -\frac{1}{3}.
\end{DPalign*}

The roots are $6$~and~$-\dfrac{1}{3}$.

Verify by putting these numbers for~$x$ in the original equation.
\[
\begin{array}[t]{rcl<{\qquad}|}
x &=& 6. \\
\dfrac{6 + 1}{6 - 1} &=& \dfrac{24 - 3}{6 + 9}. \\
\dfrac{7}{5} &=& \dfrac{21}{15} \\
\dfrac{7}{5} &=& \dfrac{7}{5} \\
\end{array}
\begin{array}[t]{>{\qquad}rcl}
x &=& -\dfrac{1}{3} \\
\dfrac{-\dfrac{1}{3} + 1}{-\dfrac{1}{3} - 1}
  &=& \dfrac{-\dfrac{4}{3} - 3}{-\dfrac{1}{3} + 9}. \\
-\dfrac{2}{4} &=& -\dfrac{13}{26}.
\end{array}
\]
\end{Soln}
%% -----File: 143.png---Folio 137-------

\PrintBreak
\Exercise{72.}

Solve:
\begin{multicols}{2}
\Item{1.}  $x^{2} - 12x + 27 = 0$.

\Item{2.}  $x^{2} - 6x + 8 = 0$.

\Item{3.}  $x^{2} - 4 = 4x - 7$.

\Item{4.}  $5x^{2} - 4x-1 = 0$.

\Item{5.}  $4x - 3 = 2x - x^{2}$.

\Item{6.}  $9x^{2} - 24x + 16 = 0$.

\Item{7.}  $6x^{2} - 5x-1 = 0$.

\Item{8.}  $4x + 3 = x^{2} + 2x$.

\Item{9.}  $16x^{2} - 16x + 3 = 0$.

\Item{10.}  $3x^{2} - 10x + 3 = 0$.

\Item{11.}  $x^{2} - 14x - 51 = 0$.

\Item{12.}  $34x - x^{2} - 225 = 0$.

\Item{13.}  $x^{2} + x - 20 = 0$.

\Item{14.}  $x^{2} - x - 12 = 0$.

\Item{15.}  $2x^{2} - 12x = - 10$.

\Item{16.}  $3x^{2} + 12x - 36 = 0$.

\Item{17.}  $(2x - 1)^{2} + 9 = 6(2x - 1)$.

\Item{18.}  $6(9x^{2} - x) = 55(x^{2} - 1)$.

\Item{19.}  $32 - 3x^{2} - 10x = 0$.

\Item{20.}  $9x^{2} - 6x - 143 = 0$.

\Item{21.}  $\dfrac{x}{x - 1} - \dfrac{x - 1}{x} = \dfrac{3}{2}$.

\Item{22.}  $\dfrac{1}{x - 2} + \dfrac{2}{x + 2} = \dfrac{5}{6}$.

\Item{23.}  $\dfrac{5x + 7}{x - 1} = 3x + 11$.

\Item{24.}  $\dfrac{7}{x + 4} - \dfrac{1}{4 - x} = \dfrac{2}{3}$.

\Item{25.}  $\dfrac{2}{x + 3} + \dfrac{x + 3}{2} = \dfrac{10}{3}$.

\Item{26.}  $\dfrac{2x}{x + 2} + \dfrac{x + 2}{2x} = 2$.

\Item{27.}  $\dfrac{3(x - 1)}{x + 1} - \dfrac{2(x + 1)}{x - 1} = 5$.

\Item{28.}  $\dfrac{2x + 5}{2x - 5} = \dfrac{7x - 5}{2x}$.

\Item{29.}  $\dfrac{3x - 1}{4x + 7} = \dfrac{x + 1}{x + 7}$.

\Item{30.}  $\dfrac{2x - 1}{x + 3} = \dfrac{x + 3}{2x - 1}$.

\Item{31.}  $\dfrac{x + 4}{x - 4} - \dfrac{x + 2}{x - 3} = 1$.

\Item{32.}  $\dfrac{4}{x - 1} - \dfrac{5}{x + 2} = \dfrac{1}{2}$.

\Item{33.}  $\dfrac{2}{x - 1} = \dfrac{3}{x - 2} + \dfrac{2}{x - 4}$.

\Item{34.}  $\dfrac{5}{x - 2} - \dfrac{3}{x - 1} = \dfrac{1}{2}$.

\Item{35.}  $\dfrac{x}{7 - x} + \dfrac{7 - x}{x} = \dfrac{29}{10}$.

\Item{36.}  $\dfrac{2x - 1}{x - 1} + \dfrac{1}{6} = \dfrac{2x - 3}{x - 2}$.
\end{multicols}
%% -----File: 144.png---Folio 138-------

\ScreenBreak
\Paragraph{164.  Problems involving Quadratics.}  Problems which involve
quadratic equations apparently have two solutions,
since a quadratic equation has two roots.

When both roots of the quadratic equation are positive
integers, they will, in general, both be admissible solutions.
Fractional and negative roots will in some problems give
admissible solutions; in other problems they will not give
admissible solutions.

The reason that every root of the equation will not
always satisfy the conditions of the problem is that the
problem may have certain restrictions, expressed or implied,
that cannot be expressed in the equation.

No difficulty will be found in selecting the result which
belongs to the particular problem we are solving.  Sometimes,
by a change in the statement of the problem, we
may form a new problem which corresponds to the result
that was inapplicable to the original problem.

Here as in simple equations x stands for an unknown
\emph{number}.

\Item{1.}  The sum of the squares of two consecutive numbers is~$41$.
Find the numbers.
\begin{Soln}
\begin{DPalign*}
\lintertext{\indent Let}
x &= \text{one number,} \\
\lintertext{and}
x + 1 &= \text{the other.} \\
\lintertext{\indent Then}
x^{2} + (x + 1)^{2} &= \text{the sum of the squares;} \\
\lintertext{but}
41 &= \text{the sum of the squares.}
\end{DPalign*}
\begin{align*}
\therefore x^{2} + (x + 1)^{2} &= 41. \\
x^{2} + x^{2} + 2x + 1 &= 41. \\
2x^{2} + 2x &= 40. \\
x^{2} + x &= 20.
\end{align*}

The solution of this equation gives $x = 4$, or~$-5$.

The positive root~$4$ gives for the numbers $4$~and~$5$.
\end{Soln}

The negative root~$-5$ is inapplicable to the problem, as
\emph{consecutive numbers} are understood to be integers which
follow each other in the common scale: $1$,~$2$, $3$, $4\dots$.
%% -----File: 145.png---Folio 139-------

\Item{2.}  In a certain nest seven times the number of birds in
the nest is equal to twice the square of the number increased
by~$3$.  Find the number.
\begin{Soln}
\begin{DPalign*}
\lintertext{\indent Let}
x &= \text{the number of birds\Add{.}} \\
\lintertext{\indent Then}
7x &= \text{$7$~times the number} \\
\lintertext{and}
2x^{2} + 3 &= \text{twice the square of the number plus~$3$.}
\intertext{\indent As these two expressions are equal we have}
2x^{2} + 3 &= 7x\Add{.}
\end{DPalign*}

The solution of this \emph{equation} gives $x = 3$ or $x = \frac{1}{2}$.

The value $\frac{1}{2}$ is not applicable to the \emph{problem} for the number of
birds must be a whole number.
\end{Soln}

\Item{3.}  A cistern has two pipes.  By one of them it can be
filled $6$~hours sooner than by the other, and by both
together in $4$~hours.  Find the time it will take each pipe
alone to fill it.
\begin{Soln}
\begin{DPalign*}
\lintertext{\indent Let}
x &= \text{the number of hours it takes the smaller pipe.} \\
\lintertext{\indent Then}
x - 6 &= \text{the number of hours it takes the larger pipe.}
\end{DPalign*}
\begin{DPalign*}
\lintertext{\indent Therefore}
\frac{1}{x} + \frac{1}{x - 6} &= \text{the part both can fill in one hour.} \\
\lintertext{\indent But}
\frac{1}{4} &= \text{the part both can fill in one hour.} \\
\frac{1}{x} + \frac{1}{x - 6} &= \frac{1}{4}
\end{DPalign*}

The solution of this \emph{equation} gives $x = 12$ or $x = 2$.

The value~$2$ is not applicable to the \emph{problem}.

Therefore it takes one pipe $12$~hr.\ and the other $6$~hr.
\end{Soln}

\Item{4.}  A rug is $1$~yard longer than it is broad.  The number
of sq.~yds.\ in the rug is~$12$.  Find its length and breadth.
\begin{Soln}
\begin{DPalign*}
\lintertext{\indent Let}
x &= \text{the number of yards in the breadth\Add{.}} \\
\lintertext{\indent Then}
x + 1 &= \text{the number of yards in the length} \\
\lintertext{and}
x(x + 1) &= \text{the number of sq yds in the rug\Add{.}} \\
\lintertext{\indent Hence}
x(x + 1) &= 12\Add{.}
\end{DPalign*}

The solution of this equation gives $x = 3$ or $x = -4$.

The dimensions are therefore $3$~yards and $4$~yards.
\end{Soln}
%% -----File: 146.png---Folio 140-------

\Exercise{73.}

\Item{1.}  Find two numbers whose sum is~$11$, and whose
product is~$30$.

\Item{2.}  Find two numbers whose difference is~$10$, and the
sum of whose squares is~$250$.

\Item{3.}  A man is five times as old as his son, and the square
of the son's age diminished by the father's age is~$24$.  Find
their ages.

\Item{4.}  A number increased by its square is equal to nine
times the next higher number.  Find the number.

\Item{5.}  The square of the sum of any two consecutive numbers
lacks~$1$ of being twice the sum of the squares of the
numbers.  Show that this statement is true.

\Item{6.}  The length of a rectangular court exceeds its breadth
by $2$~rods.  If the length and breadth were each increased
by $3$~rods, the area of the court would be $80$~square rods.
Find the dimensions of the court.

\Item{7.}  The area of a certain square will be doubled, if its
dimensions are increased by $6$~feet and $4$~feet respectively.
Find its dimensions.

\Item{8.}  The perimeter of a rectangular floor is $76$~feet and
the area of the floor is $360$~square feet.  Find the dimensions
of the floor.

\Item{9.}  The length of a rectangular court exceeds its breadth
by $2$~rods, and its area is $120$~square rods.  Find the
dimensions of the court.

\Item{10.}  The combined ages of a father and son amount to
$64$~years.  Twice the father's age exceeds the square of the
son's age by $8$~years.  Find their respective ages.
%% -----File: 147.png---Folio 141-------

\Exercise{74.}

Ex.  A boat sails $30$~miles at a uniform rate.  If the
rate had been $1$~mile an hour more, the time of the sailing
would have been $1$~hour less.  Find the rate of the sailing.
\begin{Soln}
\begin{DPalign*}
\lintertext{\indent Let}
x &= \text{the rate in miles per hour.} \\
\lintertext{\indent Then}
\frac{30}{x} &= \text{the number of hours.}
\end{DPalign*}

On the other supposition the rate would have been $x + 1$~miles
an hour and the time~$\dfrac{30}{x + 1}$.
\begin{DPalign*}
\lintertext{\indent Hence}
\frac{30}{x} - \frac{30}{x + 1} &= \text{the difference in hours for the sailing.} \\
\lintertext{\indent But}
1 &= \text{the difference in hours for the sailing.} \\
\frac{30}{x} - \frac{30}{x + 1} &= 1\Add{.}
\end{DPalign*}

The solution of this equation gives $x = 5$, or $x = -6$.

Therefore, the rate of the sailing is $5$~miles an hour.
\end{Soln}

\Item{1.}  A boat sails $30$~miles at a uniform rate.  If the rate
had been $1$~mile an hour less, the time of the sailing would
have been $1$~hour more.  Find the rate of the sailing.

\Item{2.}  A laborer built $35$~rods of stone wall.  If he had
built $2$~rods less each day, it would have taken him $2$~days
longer.  How many rods did he build a day on the
average?

\Item{3.}  A man bought flour for~\$$30$.  Had he bought $1$~barrel
more for the same sum, the flour would have cost
him \$$1$~less per barrel.  How many barrels did he buy?

\Item{4.}  A man bought some knives for~\$$6$.  Had he bought
$2$~less for the same money, he would have paid $25$~cents
more for each knife.  How many knives did he buy?

\Item{5.}  What number exceeds its square root by~$30$?

\begin{Remark}[Hint.]  Let $x^{2}$ denote the number.
\end{Remark}
%% -----File: 148.png---Folio 142-------


\Chapter{XIII.}{Arithmetical Progression.}

\Paragraph{165.}  A series of numbers is said to form an \Defn{Arithmetical
Progression} if the difference between any term and the preceding
term is the same throughout the series.

\begin{Remark}
Thus $a$, $b$, $c$, $d$, etc., are in arithmetical progression if $b - a$, $c - b$,
$d - c$,~etc., are all equal.
\end{Remark}

\Paragraph{166.}  This difference is called the \Defn{common difference} of the
progression, and is represented by~$d$.  If $d$~is positive, the
progression is an \emph{increasing} series; if $d$~is negative, the progression
is a \emph{decreasing} series.

What is the common difference in each of the following
series?
\[
\begin{array}{r*{4}{>{\quad}r}}
 1, & \Neg4, &  7, & 10, & \dots \\
 5, &     7, &  9, & 11, & \dots \\
10, &     9, &  8, &  7, & \dots \\
 7, &     3, & -1, & -5, & \dots \\
\end{array}
\]

\Paragraph{167.}  If the first term of an arithmetical progression is
represented by~$a$ and the common difference by~$d$, then
\begin{alignat*}{2}
&\text{the \emph{second} term will be } && a + d,  \\
&\text{the \emph{third} term will be }  && a + 2d, \\
&\text{the \emph{fourth} term will be } && a + 3d,
\end{alignat*}
and so on, the coefficient of~$d$ in each term being always
less by~$1$ than the \emph{number of the term}.

Hence the $n$th~term will be $a + (n - 1)d$.

If we represent the $n$th~term by~$l$, we have
\[
l = a + (n - 1)d.
\Tag{Formula (1)}
\]
%% -----File: 149.png---Folio 143-------

\Paragraph{168.}  We can, therefore, find any term of an arithmetical
progression if the first term and common difference are
given, or if any \emph{two} terms are given.

\Item{1.}  Find the 10th~term of an arithmetical progression if
the 1st~term is~$3$ and the common difference is~$4$.
\begin{Soln}
By formula~(1), the 10th~term is~$3 + (10 - 1)4$, or~$39$.
\end{Soln}

\Item{2.}  If the 8th~term of an arithmetical progression is~$25$,
and the 23d~term~$70$, find the series.
\begin{Soln}
\begin{DPalign*}
\lintertext{\indent By formula (1),}
&\text{the 23d term is $a + 22d$,} \\
\lintertext{and}
&\text{the 8th term is $a + 7d$.}
\end{DPalign*}
\begin{DPalign*}
\lintertext{\indent Therefore,}
a + 22d &= 70 \\
\lintertext{and}
a + \Z7d &= 25\Add{.} \displaybreak[1] \\
\lintertext{\indent Subtract,}
15d &= 45 \\
\lintertext{and}
d &= 3, \\
\lintertext{whence}
a &= 4.
\end{DPalign*}

The series is therefore $4$, $7$, $10$, $13$,~etc.
\end{Soln}

\Paragraph{169. Arithmetical Mean.}  If three numbers are in arithmetical
progression, the middle number is called the arithmetical
mean of the other two numbers.

If $a$, $A$, $b$ are in arithmetical progression, $A$~is the arithmetical
mean of $a$~and~$b$.  Hence, by the definition of an
arithmetical series,
\begin{DPalign*}
A - a &= b - A, \\
\lintertext{whence}
A &= \frac{a + b}{2}.
\rintertext{Formula (2)}
\end{DPalign*}
\begin{Theorem}
Hence, the arithmetical mean of any two numbers is found
by taking half their sum.
\end{Theorem}

\Paragraph{170.}  Sometimes it is required to insert several arithmetical
means between two numbers.
%% -----File: 150.png---Folio 144-------

If $m = {}$the number of means, and $n = {}$the whole number
of terms, then $m + 2 = n$.  If $m + 2$ is substituted for~$n$
in formula~(1),
\begin{DPalign*}
l &= a + (n - 1)d, \\
\lintertext{the result is}
l &= a + (m + 1)d. \displaybreak[1] \\
\lintertext{\indent By transposing~$a$,}
l - a &= (m + 1) d. \\
\therefore \frac{l - a}{m + 1} &= d.
\Tag{\llap{Formula (3)}}
\end{DPalign*}
\begin{Remark}
Thus, if it be required to insert six means between $3$~and~$17$,
the value of~$d$ is found to be $\dfrac{17 - 3}{6 + 1} = 2$; and the series will be $3$,~$5$,
$7$, $9$, $11$, $13$, $15$,~$17$.
\end{Remark}

\Exercise{75.}

\Item{1.}  Find the 25th~term in the series $3$, $6$, $9$,~$\dots$.

\Item{2.}  Find the 13th~term in the series $50$, $49$, $48$,~$\dots$.

\Item{3.}  Find the 15th~term in the series $\frac{1}{7}$, $\frac{3}{7}$, $\frac{5}{7}$,~$\dots$.

\Item{4.}  Find the 19th~term in the series $\frac{1}{4}$, $-\frac{1}{4}$, $-\frac{3}{4}$,~$\dots$.

\Item{5.}  Find the 10th~term in an arithmetical progression
whose 1st~term is~$5$ and 3d~term~$9$.

\Item{6.}  Find the 11th~term in an arithmetical progression
whose 1st~term is~$10$ and whose 6th~term is~$5$.

\Item{7.}  If the 3d~term of an arithmetical progression is~$20$
and the 13th~term is~$100$, what is the 20th~term?

\Item{8.}  Which term of the series $5$, $7$, $9$, $11$,~$\dots$, is~$43$?

\Item{9.}  Which term of the series $\frac{4}{3}$, $\frac{3}{2}$, $\frac{5}{3}$,~$\dots$, is~$18$?

\Item{10.}  What is the arithmetical mean of $20$~and~$32$?

\Item{11.}  What is the arithmetical mean of $a + b$ and $a - b$?

\Item{12.}  Insert $8$ arithmetical means between $20$~and~$29$.
%% -----File: 151.png---Folio 145-------

\Paragraph{171.  To Find the Sum of Any Number of Terms of an Arithmetical
Series.}

If $l$ denote the last term, $a$~the first term, $n$~the number
of terms, $d$~the common difference, and $s$~the sum of the
terms, it is evident that the series beginning with the first
term will be $a$, $a + d$, $a + 2d$,~etc., and beginning with the
last term will be $l$, $l - d$, $l - 2d$,~etc.  Therefore,
\begin{gather*}
\begin{array}{r*{12}{c}}
s &=& a &+& (a + d) &+& (a + 2d) &+& \Add{\dots} &+&  (l - d) &+& l, \rlap{\quad\text{or}} \\
s &=& l &+& (l - d) &+& (l - 2d) &+& \Add{\dots} &+&  (a + d) &+& a \\
\hline
2s &=& (a + l) &+& (a + l) &+& (a + l) &+& \Add{\dots} &+&  (a + l) &+& (a + l) \\
2s &=& \multicolumn{11}{l}{\text{$(a + l)$ taken as many times as there are \emph{terms},}}
\end{array} \displaybreak[1] \\
\begin{aligned}[b]
2s &= n(a + l), \\
\text{and } s &= \frac{n}{2}(a + l)\Add{.}
\end{aligned}
\Tag{Formula (4)}
\end{gather*}

Putting for $l$~its value $a + (n - 1)d$, in formula~(4), we
have
\begin{align*}
s &= \frac{n}{2}\bigl\{a + a + (n - 1)d\bigr\} \\
  &= \frac{n}{2}\bigl\{2a + (n - 1)d\bigr\}
\Tag{Formula (5)}
\end{align*}

\Item{1.}  Find the sum of the first $16$~terms of the series $5$, $7$,
$9$,~$11$\DPtypo{,}{.}
\begin{Soln}
\begin{DPalign*}
\lintertext{\indent Here}
a &= 5,\quad d = 2,\quad n = 16\Add{.} \\
\intertext{\indent Putting these values in formula~(5) we have}
s &= \tfrac{16}{2}(10 + 15 × 2) \\
  &= 320
\end{DPalign*}
\end{Soln}

\Item{2.}  Show that the sum of any number of odd numbers,
beginning with~$1$, is a square number.
\begin{Soln}
The series of odd numbers is $1$, $3$, $5$, $7$,~$\dots$.
%% -----File: 152.png---Folio 146-------
\begin{DPalign*}
\lintertext{\indent Here}
a &= 1 \quad\text{and}\quad d = 2\Add{.}
\intertext{\indent Putting these values in formula (5) we have}
s &= \frac{n}{2} \bigl\{2 + (n - 1)2\bigr\} \\
  &= \frac{n}{2} × 2n \\
  &= n^{2}\Add{.}
\end{DPalign*}

Therefore the sum of the first $5$~odd numbers is~$5^{2}$ or~$25$\Add{,} of the first
$8$ odd numbers is~$8^{2}$ or~$64$, and so on.
\end{Soln}

\Item{3.}  The sum of $20$~terms of an arithmetical progression
is~$420$, and the first term is~$2$.  Find the common difference.
\begin{Soln}
\begin{DPalign*}
\lintertext{\indent Here}
s &= 420,\quad n = 20, \quad\text{and}\quad a = 2\Add{.} \\
\intertext{\indent Putting these values in formula~(5), we have}
420 &= \tfrac{20}{2}(4 + 19d) \\
    &= 40 + 190d \\
190d &= 380 \\
\Add{\therefore} d &= 2\Add{.}
\end{DPalign*}

Therefore the common difference is~$2$.
\end{Soln}

\Exercise{76.}

\Item{1.}  Find the sum of $3$, $5$, $7$, $\dots$, to $20$~terms.

\Item{2.}  Find the sum of $14$, $14\frac{1}{2}$, $15$, $\dots$, to $12$~terms.

\Item{3.}  Find the sum of $\frac{7}{6}$, $1$, $\frac{5}{6}$, $\dots$, to $10$~terms.

\Item{4.}  Find the sum of $-7$, $-5$, $-3$, $\dots$, to $16$~terms.

\Item{5.}  Find the sum of $12$, $9$, $6$, $\dots$, to $21$~terms.

\Item{6.}  Find the sum of $-10\frac{1}{2}$, $-9$, $-7\frac{1}{2}$, $\dots$, to $25$~terms.

\Item{7.}  The sum of three numbers in arithmetical progression
is~$9$, and the sum of their squares is~$35$.  Find the numbers.

\begin{Remark}[Hint.]  Let $x - y$, $x$, $x + y$, stand for the numbers.
\end{Remark}
%% -----File: 153.png---Folio 147-------

\Item{8.}  A common clock strikes the hours from $1$ to~$12$.
How many times does it strike every $24$~hours?

\Item{9.}  The Greenwich clock strikes the hours from $1$ to~$24$.
How many times does it strike in $24$~hours?

\Item{10.}  In a potato race each man picked up $50$~potatoes
placed in line a yard apart, and the first potato one yard
from the basket, picking up one potato at a time and bringing
it to the basket.  How many yards did each man run,
the start being made from the basket?

\Item{11.}  A heavy body falling from a height falls $16.1$~feet
the first second, and in each succeeding second $32.2$~feet
more than in the second next preceding.  How far will a
body fall in $19$~seconds?

\Item{12.}  A stone dropped from a bridge reached the water in
just $3$~seconds.  Find the height of the bridge.  (See Ex.~11.)

\Item{13.}  The arithmetical mean between two numbers is~$13$,
and the mean between the double of the first and the triple
of the second is~$33\frac{1}{2}$.  Find the numbers.

\Item{14.}  Find three numbers of an arithmetical series whose
sum shall be~$27$, and the sum of the first and second shall
be~$frac{4}{5}$ of the sum of the second and third.

\Item{15.}  A travels uniformly $20$~miles a day; B~travels $8$~miles
the first day, $12$~the second, and so on, in arithmetical
progression.  If they start Monday morning from the same
place and travel in the same direction, how far apart will
they be Saturday night?

\Item{16.}  The sum of three terms of an arithmetical progression
is~$36$, and the square of the mean exceeds the product of
the other two terms by~$49$.  Find the numbers.
%% -----File: 154.png---Folio 148-------


\Chapter{XIV.}{Geometrical Progression.}

\Paragraph{172.}  A series of numbers is said to be in \Defn{Geometrical Progression}
when the quotient of any term divided by the
preceding term is the same throughout the series.

\begin{Remark}
Thus $a$, $b$, $c$, $d$, etc., are in geometrical progression if $\dfrac{b}{a} = \dfrac{c}{b} = \dfrac{d}{c}$,~etc.
\end{Remark}

\Paragraph{173.}  This quotient is called the \Defn{common ratio}, and is represented
by~$r$.

State the common ratio of the following series:
\[
\begin{array}{r*{4}{>{\quad}r}}
1, &  3, & \Neg9, &27, & \dots \\
2, &  4, & 8, &16, & \dots \\
16,&  8, & 4, &2, & \dots \\
\frac{2}{3}, &  1, & \frac{3}{2}, & \frac{9}{4}, & \dots \\
4, & -2, & 1, & -\frac{1}{2}, & \dots \\
\end{array}
\]

\Paragraph{174.}  If the first term of a geometrical progression is represented
by~$a$, and the common ratio by~$r$, then
\begin{align*}
&\text{the \emph{second} term will be~$ar$,} \\
&\text{the \emph{third} term will be~$ar^{2}$,} \\
&\text{the \emph{fourth} term will be~$ar^{3}$,}
\end{align*}
and so on, the index of~$r$ being always less by~$1$ than the
\emph{number of the term in the series}.

Hence the $n$th~term will be~$ar^{n - 1}$.
%% -----File: 155.png---Folio 149-------

If we denote the $n$th~term by~$l$, we have
\[
l = ar^{n - 1}.
\Tag{Formula (1)}
\]

\Paragraph{175.}  If the first term and common ratio are given, or if
any \emph{two terms} are given, we can find the series.

\Item{1.}  Find the 5th~term of a geometrical progression if the
first is~$3$ and the common ratio~$2$.
\begin{Soln}
In formula~(1), put $5$~for~$n$, $3$~for~$a$, and $2$~for~$r$.
\begin{DPalign*}
\lintertext{\indent Then}
l &= 3 × 2^{4} = 48.
\end{DPalign*}

Therefore the 5th~term is~$48$.
\end{Soln}

\Item{2.}  Find the geometrical series if the 5th~term is~$48$ and
the 7th~term is~$192$.
\begin{Soln}
The 5th and 7th~terms are $ar^{4}$~and~$ar^{6}$, respectively.
\begin{DPalign*}
\lintertext{\indent Whence}
ar^{4} &= 48,
\Tag{(1)} \\
\lintertext{and}
ar^{6} &= 192.
\Tag{(2)} \\
\lintertext{\indent Divide (2) by~(1),}
r^{2} &= 4. \\
\therefore r &= ±2. \\
\lintertext{\indent From~(1),}
a = \tfrac{48}{16} &= 3.
\end{DPalign*}

Therefore the series is $3$, $±6$, $12$, $±24$,~$\dots$.
\end{Soln}

\Paragraph{176.  Geometrical Mean.}  If three numbers are in geometrical
progression, the middle number is called the \emph{geometrical
mean} of the other two numbers.  Hence, if
$a$,~$G$,~$b$ are in geometrical progression, $G$~is the geometrical
mean of $a$~and~$b$.

By the definition of a geometrical progression,
\begin{DPalign*}
\frac{G}{a} &= \frac{b}{G}. \\
\therefore G^{2} &= ab, \\
\lintertext{and}
G &= ± \sqrt{ab}.
\Tag{Formula (2)}
\end{DPalign*}
\begin{Theorem}[Hence], the geometrical mean of any two numbers is the
square root of their product.
\end{Theorem}
%% -----File: 156.png---Folio 150-------

\PrintBreak
\Paragraph{177.  To Find the Sum of Any Number of Terms of a Geometrical
Progression.}

If $l$~denote the last term, $a$~the first term, $n$~the number
of terms, $r$~the common ratio, and $s$~the sum of the $n$~terms,
then
\begin{DPalign*}
s &= a + ar + ar^{2} + ar^{3} + \dots+ ar^{n - 1}. \\
\lintertext{\indent Multiply by~$r$,}
rs &= ar + ar^{2} + ar^{3} + \dots + ar^{n - 1} + ar^{n}.
\end{DPalign*}

Therefore, by subtracting the first equation from the
second,
\begin{DPalign*}
rs - s &= ar^{n} - a, \\
\lintertext{or}
(r - 1)s &= a(r^{n} - 1). \\
\therefore s &= \frac{a(r^{n} - 1)}{r - 1}.
\Tag{Formula (3)}
\end{DPalign*}

\Paragraph{178.}  When $r$~is $< 1$, this formula will be more convenient
if written
\[
s = \frac{a(1 - r^{n})}{1 - r}.
\]

\Item{1.}  Find the sum of $8$~terms of the series
\[
1,\quad 2,\quad 4,\quad \dots\Add{.}
\]
\begin{DPalign*}
\lintertext{\indent Here}
a &= 1,\quad r = 2,\quad n = 8. \\
\lintertext{\indent From formula~(3),}
s &= 1(2^{8} - 1) = 255.
\end{DPalign*}

\Item{2.}  Find the sum of $6$~terms of the series
\[
2,\quad 3,\quad \tfrac{9}{2},\quad \dots\Add{.}
\]
\begin{DPalign*}
\lintertext{\indent Here}
a &= 2,\quad r = \tfrac{3}{2},\quad n = 6. \\
\lintertext{\indent From formula~(3),}
s &= \frac{2\bigl\{(\frac{3}{2})^{6} - 1\bigr\}}{\frac{3}{2} - 1} \\
  &= \frac{2\bigl\{\frac{729}{64} - 1\bigr\}}{\frac{1}{2}} \\
  &= \frac{4\{729 - 64\}}{64} \\
  &= 41\tfrac{9}{16}.
\end{DPalign*}
%% -----File: 157.png---Folio 151-------

\Exercise{77.}

\Item{1.}  Find the 5th~term of $3$, $9$, $27$\Add{,}~$\dots$.

\Item{2.}  Find the 7th~term of $3$, $6$, $12$\Add{,}~$\dots$.

\Item{3.}  Find the 8th~term of $6$, $3$, $\frac{3}{2}$\Add{,}~$\dots$.

\Item{4.}  Find the 9th~term of $1$, $-2$, $4$\Add{,}~$\dots$.

\Item{5.}  Find the geometrical mean between $2$~and~$8$.

\Item{6.}  Find the common ratio if the 1st~and 3d~terms are
$2$~and~$32$.

Find the sum of the series:

\Item{7.}  $3$, $9$, $27$, $\dots$ to $6$~terms.

\Item{8.}  $3$, $6$, $12$, $\dots$ to $8$~terms.

\Item{9.}  $6$, $3$, $\frac{3}{2}$, $\dots$ to $7$~terms.

\Item{10.}  $8$, $4$, $2$, $\dots$ to $8$~terms.

\Item{11.}  $64$, $32$, $16$, $\dots$ to $9$~terms.

\Item{12.}  $64$, $-32$, $16$, $\dots$ to $5$~terms.

\Item{13.}  $\frac{1}{2}$, $\frac{1}{3}$, $\frac{2}{9}$, $\dots$ to $4$~terms.

\Item{14.}  If a blacksmith uses seven nails in putting a shoe on
a horse's foot, and receives $1$~cent for the first nail, $2$~cents
for the second nail, and so on, what does he receive for
putting on the shoe?

\Item{15.}  If a boy receives $2$~cents for his first day's work,
$4$~cents for his second day, $8$~cents for the third day, and
so on for $12$~days, what will his wages amount to?

\Item{16.}  If the population of a city is $10,000$, and increases
$10$\%~a year for four years, what will be its population at
the end of the four years? (Here $l = ar^{4}$.)
%% -----File: 158.png---Folio 152-------


\Chapter{XV.}{Square and Cube Roots.}

\Section{Square Roots of Compound Expressions.}

\Paragraph{179.}  Since the square of $a + b$ is $a^{2} + 2ab + b^{2}$, the square
root of $a^{2} + 2ab + b^{2}$ is $a + b$.

It is required to find a method of extracting the root
$a + b$ when $a^{2} + 2ab + b^{2}$ is given.
\begin{Soln}
Ex.  The first term,~$a$, of the root is obviously the square root of
the first term,~$a^{2}$, in the expression.
\[
\begin{array}{r*{3}{cr}l}
   & &   a^{2} &+& 2ab &+& \TbBar{b^{2}} & a + b \\
\cline{8-8}
   & &   a^{2} \\
\cline{3-3}
2a &+& \TbBar{b} & & 2ab &+& b^{2} \\
   & & \TbBar{ } & & 2ab &+& b^{2} \\
\cline{4-7}
\end{array}
\]

If the $a^{2}$ is subtracted from the given
expression, the remainder is $2ab + b^{2}$\Add{.}
Therefore the second term,~$b$, of the root
is obtained when the first term of this
remainder is divided by~$2a$; that is, by
\emph{double the part of the root already found}.  Also, since
\[
2ab + b^{2} = (2a + b)b,
\]
the divisor is \emph{completed by adding to the trial-divisor the new term of
the root}\Add{.}
\end{Soln}

\ScreenBreak
Ex.  Find the square root of $25x^{2} - 20x^{3}y + 4x^{4}y^{2}$.
\begin{Soln}
\[
\begin{array}{rcccrll}
&     && & 25x^{2} & \TbBar{-20 x^{3}y + 4x^{4}y^{2}} & 5x - 2x^{2}y \\
\cline{7-7}
\text{Here $a^{2}$}
       &=& (5x)^{2} &=& 25x^{2} \\
\cline{5-5}
2a + b &=& 10x &-& \TbBar{2x^{2}y} & -20x^{3}y + 4x^{4}y^{2} \\
       & &     & & \TbBar{ }       & -20x^{3}y + 4x^{4}y^{2} \\
\cline{6-6}
\end{array}
\]

The expression is \emph{arranged} according to the ascending powers of~$x$\Add{.}

The square root of the first term is~$5x$, hence $5x$~is the first term
of the root.  $(5x)^{2}$ or $25x^{2}$ is subtracted, and the remainder is
\[
-20x^{3}y + 4x^{4}y^{2}.
\]

The second term of the root, $-2x^{2}y$, is obtained by dividing
$-20x^{3}y$ by~$10x$, the double of~$5x$, and this new term of the root is
also annexed to the divisor,~$10x$, to complete the divisor.
\end{Soln}
%% -----File: 159.png---Folio 153-------

\Paragraph{180.}  The same method will apply to longer expressions,
if care be taken to obtain the \emph{trial-divisor} at each stage of
the process, \emph{by doubling the part of the root already found},
and to obtain the \emph{complete divisor by annexing the new term
of the root to the trial-divisor}.

\ScreenBreak
Ex.  Find the square root of
\[
1 + 10x^{2} + 25x^{4} + 16x^{6} - 24x^{5} - 20x^{3} - 4x.
\]
\begin{Soln}
\[
%[** TN: Spacing hack]
\qquad\makebox[0pt][c]{$\begin{array}{r*{2}{cr}lll}
16x^{6} &-& 24x^{5} &+& 25x^{4} &-20x^{3} + 10x^{2} &\TbBar{- 4x + 1} & 4x^{3} - 3x^{2} + 2x - 1 \\
\cline{8-8}
16x^{6} \\
\cline{1-1}
\TbBar{\llap{$8x^{3} - 3x^{2}$}}
        &-& 24x^{5} &+& 25x^{4} \\
\TbBar{}&-& 24x^{5} &+&  9x^{4} \\
\cline{2-5}
\multicolumn{4}{r|}{8x^{3} - 6x^{2} + 2x} &
                        16x^{4} &-20x^{3} +  10x^{2} \\
      & &  & \TbBar{ }& 16x^{4} &-12x^{3} + \Z4x^{2} \\
\cline{5-6}
\multicolumn{5}{r|}{8x^{3} - 6x^{2} + 4x - 1}
                                 &-\Z8x^{3} + \Z6x^{2} & -4x + 1 \\
       & &  &          & \TbBar{}&-\Z8x^{3} + \Z6x^{2} & -4x + 1 \\
\cline{6-7}
\end{array}$}
\]

The expression is arranged according to the descending powers of~$x$.

It will be noticed that each successive trial-divisor may be obtained
by taking the preceding complete divisor with its \emph{last term doubled}.
\end{Soln}

\Exercise{78.}

Find the square root of:
\Item{1.}  $a^{2} + 2ab + 2ac + b^{2} + 2bc + c^{2}$.

\Item{2.}  $x^{4} + 2x^{3} + 3x^{2} + 2x + 1$.

\Item{3.}  $x^{4} - 4x^{3}y + 6x^{2}y^{2} - 4xy^{3} + y^{4}$.

\Item{4.}  $4a^{4} - 12a^{3}b + 29a^{2}b^{2} - 30ab^{3} + 25b^{4}$.

\Item{5.}  $16x^{6} + 24x^{5}y + 9x^{4}y^{2} - 16x^{3}y^{3} - 12x^{2}y^{4} + 4y^{6}$.

\Item{6.}  $4x^{6} - 4x^{4}y^{2} + 12x^{3}y^{3} + x^{2}y^{4} - 6xy^{5} + 9y^{6}$.

\Paragraph{181.  Arithmetical Square Roots.}  In the general method
of extracting the square root of a number expressed by
figures, the first step is to mark off the figures into \emph{groups}.
%% -----File: 160.png---Folio 154-------

Since $1 = 1^{2}$, $100 = 10^{2}$, $10,000 = 100^{2}$, and so on, it is
evident that the square root of a number between $1$~and
$100$ lies between $1$~and~$10$; of a number between $100$ and
$10,000$ lies between $10$~and~$100$.  In other words, the
square root of a number expressed by \emph{one} or \emph{two} figures is
a number of \emph{one} figure, of a number expressed by \emph{three} or
\emph{four} figures is a number of \emph{two} figures, and so on.

If, therefore, an integral square number is divided into
groups of two figures each, from the right to the left, the
number of figures in the root will be equal to the number
of groups of figures.  The last group to the left may have
only one figure.

Ex.  Find the square root of~$3249$.
%[** TN: Tabulated calculation inset in the original]
\begin{Soln}
\[
\begin{array}{rc@{}ll}
    &3&249&(57 \\
    &2&5 \\
\cline{2-3}
107\,\rlap{)}&&749 \\
    &&749 \\
\cline{3-3}
\end{array}
\]

In this case, $a$~in the typical form $a^{2} + 2ab + b^{2}$
represents $5$~tens, that is,~$50$, and $b$~represents~$7$\Add{.} The
$25$ subtracted is really~$2500$, that is,~$a^{2}$, and the complete
divisor $2a + b$ is $2 × 50 + 7 = 107$.
\end{Soln}

\Paragraph{182.}  The same method will apply to numbers of more
than two groups of figures by considering $a$~in the typical
form to represent at each step \emph{the part of the root already
found}.

It must be observed that \emph{$a$~represents so many tens with
respect to the next figure of the root}.

Ex.  Find the square root of~$94,249$.
\begin{Soln}
\[
\begin{array}{r*{3}{c@{\,}}l}
     &9& 42&49&(307 \\
     &9 \\
\cline{2-4}
607\rlap{)} & & \multicolumn{2}{l}{4249} \\
     & & \multicolumn{2}{l}{4249} \\
\cline{3-4}
\end{array}
\]
\end{Soln}

\begin{Remark}[Note.] Since the first trial divisor,~$60$, is not contained in~$42$, we
put a zero in the root, and bring down the next group,~$49$.
\end{Remark}
%% -----File: 161.png---Folio 155-------

\Paragraph{183.}  If the square root of a number has decimal places,
the number itself will have \emph{twice} as many.  Thus, if $0.21$
is the square root of some number, this number will be
$(0.21)^{2} = 0.21 × 0.21 = 0.0441$, and if $0.111$ be the root,
the number will be $(0.111)^{2} = 0.111 x 0.111 = 0.012321$.

Therefore, the number of \emph{decimal} places in every square
decimal will be \emph{even}, and the number of decimal places in
the root will be \emph{half} as many as in the given number itself.

Hence, if a given number contain a decimal, we divide
it into groups of two figures each, by beginning at the
decimal point and marking toward the left for the integral
number, and toward the right for the decimal.  We must
have the last group on the right of the decimal point contain
\emph{two} figures, annexing a cipher when necessary.

Ex.  Find the square roots of $41.2164$ and $965.9664$.
\begin{Soln}
\[
\begin{array}{r@{}r@{}l@{\,}ll}
            & 41&.21&64& (6.42 \\
            & 36 \\
\cline{2-3}
124\rlap{)} &  5&21 \\
            &  4&96 \\
\cline{3-4}
\multicolumn{2}{r}{1282\rlap{)}}&\multicolumn{2}{l}{2564} \\
            &   &\multicolumn{2}{l}{2564} \\
\cline{3-4}
\end{array}\qquad\qquad\qquad
\begin{array}{rl@{}l@{\,}ll}
    9&65&.96& 64&(31.08 \\
    9 \\
\cline{2-2}
  61\rlap{)}&65 \\
     &61 \\
\cline{2-3}
6208\rlap{)}&49&664 \\
     &49&664 \\
\cline{2-3}
\end{array}
\]
\end{Soln}

\Paragraph{184.}  If a number contain an \emph{odd} number of decimal
places, or if any number give a \emph{remainder} when as many
figures in the root have been obtained as the given number
has groups, then its exact square root cannot be found.  We
may, however, approximate to its exact root as near as we
please by annexing ciphers and continuing the operation.

The square root of a common fraction whose denominator
is not a perfect square can be found approximately by
reducing the fraction to a decimal and then extracting the
root; or by reducing the fraction to an equivalent fraction
whose denominator is a perfect square, and extracting the
square root of both terms of the fraction.
%% -----File: 162.png---Folio 156-------

%[** Force page break in both print and screen layout]
\newpage
\Item{1.}  Find the square roots of $3$ and $357.357$.
\begin{Soln}
\[
\begin{array}{rcccl}
   &3.&\multicolumn{3}{l}{(\rlap{$1.732\dots$}} \\
   &1 \\
\cline{2-3}
27\rlap{)}
   &2&00 \\
   &1&89 \\
\cline{3-4}
\multicolumn{2}{r}{343\rlap{)}}
     &11&00 \\
    &&10&29 \\
\cline{4-5}
\multicolumn{3}{r}{3462\rlap{)}}
        &71&00 \\
&&      &69&24 \\
\cline{4-5}
\end{array}\qquad\qquad\qquad
\begin{array}{rcclll}
   &3&57.&35&70&\rlap{$(18.903\dots$} \\
   &1 \\
\cline{2-3}
28\rlap{)}
   &2&57 \\
   &2&24 \\
\cline{3-4}
\multicolumn{2}{r}{369\rlap{)}}
     &33&35 \\
&    &33&21 \\
\cline{4-6}
\multicolumn{3}{r}{37803\rlap{)}}
        &14&70&00 \\
&&      &11&34&09 \\
\cline{4-6}
\end{array}
\]
\end{Soln}

\Item{2.}  Find the square root of~$\frac{5}{8}$.
\begin{Soln}
\begin{DPalign*}
\lintertext{\indent Since}
\frac{5}{8} &= 0.625, \\
\lintertext{the square root of}
\frac{5}{8} &= \sqrt{0.625} \\
            &= 0.79057. \displaybreak[1] \\
\lintertext{\indent Or,}
\frac{5}{8} &= \frac{10}{16}, \\
\lintertext{and the square root of}
\frac{5}{8} &= \frac{\sqrt{10}}{\sqrt{16}} = \tfrac{1}{4}\sqrt{10} \\
  &= \tfrac{1}{4}(3.16227) \\
  &= 0.79057.
\end{DPalign*}
\end{Soln}

\ScreenBreak
\Exercise{79.}

Find the square root of:
\begin{multicols}{3}
\Item{1.}  $324$.

\Item{2.}  $441$.

\Item{3.}  $529$.

\Item{4.}  $961$.

\Item{5.}  $10.24$.

\Item{6.}  $53.29$.

\Item{7.}  $53,824$.

\Item{8.}  $616,225$.

\Item{9.}  $1,500,625$.

\Item{10.}  $346,921$.

\Item{11.}  $31,371,201$.

\Item{12.}  $1,522,756$.
\end{multicols}

\PrintBreak
Find to four decimal places the square root of:
\begin{multicols}{5}
\Item{13.}  $2$.

\Item{14.}  $3$.

\Item{15.}  $5$.

\Item{16.}  $6$.

\Item{17.}  $0.5$.

\Item{18.}  $0.9$.

\Item{19.}  $\frac{2}{3}$.

\Item{20.}  $\frac{3}{4}$.

\Item{21.}  $\frac{4}{5}$.

\Item{22.}  $\frac{5}{8}$.
\end{multicols}
%% -----File: 163.png---Folio 157-------


\Section{Cube Roots of Compound Expressions.}

\Paragraph{185.}  Since the cube of $a + b$ is $a^{3} + 3a^{2}b + 3ab^{2} + b^{3}$,
the cube root of a$^{3} + 3a^{2}b + 3ab^{2} + b^{3}$ is $a + b$.

It is required to devise a method for extracting the cube
root $a + b$ when $a^{3} + 3a^{2}b + 3ab^{2} + b^{3}$ is given.

\Item{1.}  Find the cube root of $a^{3} + 3a^{2}b + 3ab^{2} + b^{3}$.
\begin{Soln}
\[
\begin{array}{rl<{\quad}rcll}
       &              & a^{3} &+& \TbBar{3a^{2}b + 3ab^{2} + b^{3}} & a + b \\
\cline{6-6}
3a^{2} &              & a^{3} \\
\cline{3-3}
       &+ 3ab + b^{2} &\TbBar{}& & 3a^{2}b + 3ab^{2} + b^{3} \\
\cline{1-3}
3a^{2} &+ 3ab + b^{2} &\TbBar{}& & 3a^{2}b + 3ab^{2} + b^{3} \\
\cline{4-5}
\end{array}
\]

The first term $a$ of the root is obviously the cube root of the first
term~$a^{3}$ of the given expression.

If $a^{3}$ be subtracted, the remainder is $3a^{2}b + 3ab^{2} + b^{3}$; therefore,
the second term~$b$ of the root is obtained by dividing the first term of
this remainder by \emph{three times the square of~$a$}.

Also, since $3a^{2}b + 3ab^{2} + b^{3} = (3a^{2} + 3ab + b^{2})b$, the \emph{complete
divisor} is obtained by adding $3ab + b^{2}$ to the \emph{trial divisor~$3a^{2}$}.
\end{Soln}

\Item{2.}  Find the cube root of $8x^{3} + 36x^{2}y + 54xy^{2} + 27y^{3}$.
\begin{Soln}
\[
\begin{array}{rlrlcll}
             &         &        & 8x^{3} &+& \TbBar{36x^{2}y + 54xy^{2} + 27y^{3}} & 2x + 3y \\
\cline{7-7}
             & 12x^{2} &        & 8x^{3} \\
\cline{4-4}
(6x + 3y)3y =&         &  18xy +& \TbBar{9y^{2}} & & 36x^{2}y + 54xy^{2} + 27y^{3} \\
\cline{2-4}
             & 12x^{2} &+ 18xy +& \TbBar{9y^{2}} & & 36x^{2}y + 54xy^{2} + 27y^{3} \\
\cline{5-6}
\end{array}
\]

The cube root of the first term is~$2x$, and this is therefore the first
term of the root.  $8x^{3}$,~the cube of~$2x$, is subtracted.

The second term of the root,~$3y$, is obtained by dividing $36x^{2}y$ by
$3(2x)^{2} = 12x^{2}$, which corresponds to $3a^{2}$ in the typical form, and the
divisor is completed by annexing to~$12x^{2}$ the expression
\[
\bigl\{3(2x) + 3y\bigr\}3y = 18xy + 9y^{2},
\]
which corresponds to $3ab + b^{2}$ in the typical form.
\end{Soln}
%% -----File: 164.png---Folio 158-------

\Paragraph{186.}  The same method may be applied to longer expressions
by considering~$a$ in the typical form $3a^{2} + 3ab + b^{2}$
to represent at each stage of the process \emph{the part of the root
already found}.  Thus, if the part of the root already found
is $x + y$, then $3a^{2}$~of the typical form will be represented
by $3(x + y)^{2}$; and if the third term of the root be~$+z$, \DPtypo{the}{then}
$3ab + b^{2}$ will be represented by $3(x + y)z + z^{2}$. So that
the complete divisor, $3a^{2} + 3ab + b^{2}$, will be represented
by $3(x + y)^{2} + 3(x + y)z + z^{2}$.

Ex.  Find the cube root of $x^{6} - 3x^{5} + 5x^{3} - 3x - 1$.
\begin{Soln}
%[** TN: Special spacing]
\ifthenelse{\not\boolean{ForPrinting}}{\footnotesize}{}
\[
\ifthenelse{\not\boolean{ForPrinting}}{\quad}{}
\begin{array}{*{2}{r}@{}*{3}{r}*{4}{l}}
                       &       &   & \TbBar{}&  x^{2}& \multicolumn{3}{@{}l}{{} - x - 1} \\
\cline{5-6}
                       &       &           & &  x^{6}& {} - 3x^{5} &\multicolumn{2}{@{}l}{\phantom{{}+ 3x^{4}} + 5x^{3} - 3x - 1} \\
                       & 3x^{4}&           & &  x^{6} \\
\cline{5-5}
     (3x^{2} - x)(-x) =&       &{} - 3x^{3}&+&  \TbBar{x^{2}}&{} - 3x^{5}&\phantom{{}+ 3x^{4}} + 5x^{3} \\
\cline{3-5}
                       & 3x^{4}&{} - 3x^{3}&+&  \TbBar{x^{2}}&{} - 3x^{5}&{} + 3x^{4} - \Z x^{3} \\
\cline{6-7}
                       &       &           & &       &\TbBar{}           &{} - 3x^{4} + 6x^{3} - 3x - 1 \\
     3(x^{2} - x)^{2} =& 3x^{4}&{} - 6x^{3}&+& 3x^{2}&\TbBar{} \\
\llap{$(3x^{2} - 3x - 1)$}(-1) =&       &           &-& 3x^{2}& \multicolumn{1}{l|}{{}+ 3x + 1} \\
\cline{2-6}
                       & 3x^{4}&{} - 6x^{3}& &       & \multicolumn{1}{l|}{{}+ 3x + 1}&{} - 3x^{4} + 6x^{3} - 3x - 1 \\
\cline{7-7}
\end{array}
\]
\end{Soln}

\begin{Remark}[Note.] The root is placed \emph{above} the given expression because
there is no room for it on the page at the right of the expression.

The first term of the root,~$x^{2}$, is obtained by taking the cube root
of the first term of the given expression; and the first trial-divisor,~$3x^{4}$,
is obtained by taking three times the square of this term.

The next term of the root is found by dividing~$-3x^{5}$, the first term
of the remainder after~$x^{6}$ is subtracted, by~$3x^{4}$, and the first complete
divisor, $3x^{4} - 3x^{3} + x^{2}$, is found by annexing to the trial divisor
$(3x^{2} - x)(-x)$, which expression corresponds to $(3a + b)b$ in the
typical form.

\emph{The part of the root already found}~($a$) is now represented by $x^{2} - x$,
therefore $3a^{2}$~is represented by $3(x^{2} - x)^{2} = 3x^{4} - 6x^{3} + 3x^{2}$, the
second trial divisor\Add{,} and $(3a + b)b$ by $(3x^{2} - 3x - 1)(-1)$, since $b$~in
this case is found to be~$-1$, therefore, in the second complete
divisor, $3a^{2} + (3a + b)b$ is represented by
\[
(3x^{4} - 6x^{3} + 3x^{2}) + (3x^{2} - 3x - 1)(-1) = 3x^{4} - 6x^{3} + 3x + 1.
\]
\end{Remark}
%% -----File: 165.png---Folio 159-------

\Exercise{80.}

Find the cube root of

\Item{1.}  $x^{3} + 3x^{2} y + 3xy^{2} + y^{3}$\Add{.}

\Item{2.}  $8x^{3} - 12x^{2} + 6x - 1$.

\Item{3.}  $8x^{3} - 36x^{2} y + 54xy^{2} - 27y^{3}$.

\Item{4.}  $64a^{3} - 144a^{2} x + 108ax^{2} - 27x^{3}$\Add{.}

\Item{5.}  $1 + 3x + 6x^{2} + 7x^{3} + 6x^{4} + 3x^{5} + x^{6}$.

\Item{6.}  $x^{6} - 3x^{5} + 6x^{4} - 7x^{3} + 6x^{2} - 3x + 1$.

\ScreenBreak
\Paragraph{187.  Arithmetical Cube Roots.}  In extracting the cube root
of a number expressed by figures, the first step is to mark
it off into groups\Add{.}

Since $1 = 1^{3}$, $1000 = 10^{3}$, $1,000,000 = 100^{3}$, and so on, it
follows that the cube root of any number between $1$~and
$1000$, that is, of any number which has \emph{one}, \emph{two}, or \emph{three}
figures, is a number of \emph{one} figure, and that the cube root
of any number between $1000$ and $1,000,000$, that is, of any
number which has \emph{four}, \emph{five}, or \emph{six} figures, is a number of
\emph{two} figures, and so on.

If, therefore, an integral cube number be divided into
groups of three figures each, from right to left, the number
of figures in the root will be equal to the number of groups\Add{.}
The last group to the left may consist of one, two, or three
figures.

\Paragraph{188.}  If the cube root of a number have decimal places,
the number itself will have \emph{three times} as many.  Thus, if
$0.11$ be the cube root of a number\Add{,} the number is $0.11 × 0.11
x 0.11 = 0.001331$.  Hence\Add{,} if a given number contain a
decimal, we divide the figures of the number into groups
of three figures each, by beginning at the decimal point
and marking toward the left for the integral number, and
%% -----File: 166.png---Folio 160-------
toward the right for the decimal.  We must be careful to
have the last group on the right of the decimal point contain
\emph{three} figures, annexing ciphers when necessary.

\ScreenBreak
Extract the cube root of~$42875$.
\begin{Soln}
\[
\begin{array}{rrrlll}
        &               &              & 42&875&(35 \\
        &               &       a^{3} =& 27 \\
\cline{4-5}
3a^{2} =&   3 × 30^{2} =& \TbBar{2700} & 15&875 \\
   3ab =& 3 × (30 × 5) =&  \TbBar{450} \\
 b^{2} =&        5^{2} =&   \TbBar{25} \\
\cline{3-3}
        &               & \TbBar{3175} & 15&875 \\
\cline{4-5}
\end{array}
\]

Since $42875$ has two groups, the root will have two figures.

The first group, $42$, contains the cube of the tens of the root.

The greatest cube in~$42$ is~$27$, and the cube root of~$27$ is~$3$.  Hence
$3$~is the tens' figure of the root.

We subtract $27$ from~$42$, and bring down the next group,~$875$.
Since $a$~is $3$~tens or~$30$, $3a^{2} = 3 × 30^{2}$, or~$2700$.  This trial-divisor is
contained $5$~times in~$15875$.  The trial-divisor is completed by adding
$3ab + b^{2}$; that is, $450 + 25$, to the trial-divisor.
\end{Soln}

\Paragraph{189.}  The same method will apply to numbers of more
than two groups of figures, by considering in each case~$a$,
the part of the root already found, as so many tens with
respect to the next figure of the root.

\PrintStretch{12pt}
Extract the cube root of~$57512456$.
\begin{Soln}
\[
\begin{array}{*{3}{r}*{4}{l}}
        &                &                &  57&512&456&(386 \\
        &                &         a^{3} =&  27 \\
\cline{4-5}
3a^{2} =&    3 × 30^{2} =&   \TbBar{2700} &  30&512 \\
   3ab =&  3 × (30 × 8) =&    \TbBar{720} \\
 b^{2} =&         8^{2} =&     \TbBar{64} \\
\cline{3-3}
        &                &   \TbBar{3484} &  27&872 \\
\cline{4-6}
        &               &        \TbBar{} & \Z2&640&456 \\
3a^{2} =&   3 × 380^{2} =& \TbBar{433200} \\
   3ab =& 3 × (380 × 6) =&   \TbBar{6840} \\
 b^{2} =&         6^{2} =&     \TbBar{36} \\
\cline{3-3}
        &               &  \TbBar{440076} & \Z2&640&456 \\
\cline{4-6}
\end{array}
\]
\end{Soln}
%% -----File: 167.png---Folio 161-------

Extract the cube root of~$187.149248$.
\begin{Soln}
\[
\begin{array}{*{3}{r}@{}*{4}{r}l}
        &                &             & \multicolumn{2}{r}{187\rlap{.}}&149&248&(5.72 \\
        &                &      a^{3} =& \multicolumn{2}{r}{125} \\
\cline{4-6}
3a^{2} =&    3 × 50^{2} =& \TbBar{7500}&&        62&149 \\
   3ab =&  3 × (50 × 7) =& \TbBar{1050} \\
 b^{2} =&         7^{2} =&   \TbBar{49} \\
\cline{3-3}
        &                & \TbBar{8599}&&        60&193 \\
\cline{4-7}
        &                &             &  \TbBar{}&  1&956&248 \\
3a^{2} =&   3 × 570^{2} =&         9747&\TbBar{00} \\
   3ab =& 3 × (570 × 2) =&           34&\TbBar{20} \\
 b^{2} =&         2^{2} =&             & \TbBar{4} \\
\cline{3-4}
        &                &         9781&\TbBar{24}&  1&956&248 \\
\cline{5-7}
\end{array}
\]

It will be seen from the groups of figures that the root will have
one integral and two decimal places.
\end{Soln}

\Paragraph{190.}  If the given number is not a perfect cube, ciphers
may be annexed, and a value of the root may be found as
near to the \emph{true} value as we please.

Extract the cube root of~$1250.6894$.
\begin{Soln}
\[
\begin{array}{rrrcclll}
        &                 &              & &1&\multicolumn{3}{l}{\rlap{250.689\,400\,(10.77}} \\
        &                 &        a^{3} &=&1 \\
\cline{4-6}
3a^{2} =& \PadTo{3 × (1070 × 7)}{3 × 10^{2}} =&   \TbBar{\Z\Z300}& &&  250 & \phantom{999} & \phantom{999} \\
\end{array}
\]
Since $300$ is not contained in~$250$, the next figure of the root will
be~$0$.
\[
\begin{array}{rrr@{}ccrll}
3a^{2} =&    3 × 100^{2} =& \TbBar{30000}& &\Z&250&689 \\
   3ab =&  3 × (100 × 7) =&  \TbBar{2100} \\
 b^{2} =&          7^{2} =&    \TbBar{49} \\
\cline{3-3}
        &                 & \TbBar{32149}& &&  225&043 \\
\cline{4-8}
        &                 &      &\TbBar{}&&  25&646&400 \\
3a^{2} =&   3 × 1070^{2} =& 34347&\TbBar{00} \\
   3ab =& 3 × (1070 × 7) =&   224&\TbBar{70} \\
 b^{2} =&          7^{2} =&      &\TbBar{49} \\
\cline{3-4}
        &                 & 34572&\TbBar{19}&&  24&200&533 \\
\cline{5-8}
        &                 &      &          &&   1&445&867 \\
\end{array}
\]
\end{Soln}
%% -----File: 168.png---Folio 162-------

\Paragraph{191.}  Notice that if $a$~denotes the first term, and $b$~the
second term of the root, the \emph{first complete divisor} is
\[
3a^{2} + 3ab + b^{2},
\]
and the \emph{second trial-divisor} is $3(a + b)^{2}$, that is,
\[
3a^{2} + 6ab + 3b^{2}.
\]

This expression may be obtained by adding to the preceding
complete divisor, $3a^{2} + 3ab + b^{2}$, \emph{its second term and
twice its third term}.  Thus:
\[
\begin{array}{rr}
3a^{2} + 3ab + & b^{2} \\
         3ab + &2b^{2} \\
\cline{1-2}
3a^{2} + 6ab + &3b^{2} \\
\end{array}
\]

This method of obtaining \emph{trial-divisors} is of great importance
for shortening numerical work, as may be seen in the
following example:

Ex.  Extract the cube root of~$5$ to five places of decimals.
\begin{Soln}
\[
\begin{array}{rrr@{}r@{}rcrrl}
        &               &    &    &&5\rlap{.}&\multicolumn{3}{l}{000\,(1.70997} \\
        &               && \multicolumn{3}{r}{a^{3} = 1} \\
\cline{4-7}
3a^{2} =&   3 × 10^{2} =& 300&\TbBar{}&&4&000 \\
   3ab =&    3(10 × 7) =& 210&\TbBar{} \\
 b^{2} =&        7^{2} =&  49&\TbBar{}\\
\cline{3-3}
        &               & 559&\TbBar{\BB}&&3&913 \\
\cline{5-9}
        &               & 259&  &  &\TbBar{}& 87& 000& 000 \\
\cline{3-5}
3a^{2} =& 3 × 1700^{2} =& 867&00&00&\TbBar{} \\
   3ab =&  3(1700 × 9) =&   4&59&00&\TbBar{} \\
 b^{2} =&        9^{2} =&    &  &81&\TbBar{} \\
\cline{3-5}
        &               & 871&59&81&\TbBar{\BB}& 78& 443& 829 \\
\cline{7-9}
        &               &   4&59&81&\TbBar{}&  8& 556& 1710 \\
\cline{3-5}
3a^{2} =& 3 × 1709^{2} =& 876&20&43&\TbBar{}&  7& 885& 8387 \\
\cline{7-9}
        &               &    &  &  &\TbBar{}&   & 670& 33230 \\
        &               &    &  &  &\TbBar{}&   & 613& 34301 \\
\cline{7-9}
\end{array}
\]
%% -----File: 169.png---Folio 163-------

After the first two figures of the root are found, the next trial-divisor
is obtained by bringing down~$259$, the sum of the $210$ and $49$
obtained in completing the preceding divisor, then \emph{adding the three
lines connected by the brace}, and annexing two ciphers to the result.

This trial divisor is~$86,700$, and if we add $3ab + b^{2}$ to complete the
divisor, when $b = 1$, the complete divisor will be $86,700 + 511 = 87,211$,
and this is larger than the dividend~$87,000$.  We therefore put $0$ for
the next figure of the root.  We then bring down another group
and annex two more ciphers to the trial-divisor.

The last two figures of the root are found by division.  The rule
in such cases is, that two less than the number of figures already
obtained may be found without error by division, the divisor being
three times the square of the part of the root already found.
\end{Soln}

\Paragraph{192.}  The cube root of a common fraction whose denominator
is not a perfect cube can be found approximately by
reducing the fraction to a decimal, and then extracting
the root.

\Exercise{81.}

Find the cube root of:
\begin{multicols}{2}
\Item{1.}  $46,656$.

\Item{2.}  $42,875$.

\Item{3.}  $91,125$.

\Item{4.}  $274,625$.

\Item{5.}  $110,592$.

\Item{6.}  $258,474,853$.

\Item{7.}  $109,215,352$.

\Item{8.}  $259,694,072$.

\Item{9.}  $127,263,527$.

\Item{10.}  $385,828,352$.

\Item{11.}  $1879.080904$.

\Item{12.}  $1838.265625$.
\end{multicols}

Find to four decimal places the cube root of:
\begin{multicols}{4}
\Item{13.}  $0.01$.

\Item{14.}  $0.05$.

\Item{15.}  $0.2$.

\Item{16.}  $4$.

\Item{17.}  $10$.

\Item{18.}  $87$.

\Item{19.}  $2.5$.

\Item{20.}  $2.05$.

\Item{21.}  $3.02$.

\Item{22.}  $\frac{2}{3}$.

\Item{23.}  $\frac{3}{4}$.

\Item{24.}  $\frac{9}{11}$.
\end{multicols}
%% -----File: 170.png---Folio 164-------
%[Blank Page]
%% -----File: 171.png---Folio 165-------


\Answers

\AnsTo[7]{Exercise}{1.} %  Page 10.

\Item{1.}  $14$.

\Item{2.}  $10$.

\Item{3.}  $13$.

\Item{4.}  $11$.

\Item{5.}  $13$.

\Item{6.}  $7$.

\Item{7.}  $9$.

\Item{8.}  $7$.

\Item{9.}  $6$.

\Item{10.}  $2$.

\Item{11.}  $3$.

\Item{12.}  $6$.

\Item{13.}  $2$.

\Item{14.}  $8$.

\Item{15.}  $4$.

\Item{16.}  $3$.

\Item{17.}  $1$.

\Item{18.}  $1$.

\Item{19.}  $3$.

\Item{20.}  $4$.

\Item{21.}  $10$.


\AnsTo[7]{Exercise}{2.} %  Page 12.

\Item{1.}  $91$.

\Item{2.}  $21$.

\Item{3.}  $60$.

\Item{4.}  $24$.

\Item{5.}  $96$.

\Item{6.}  $16$.

\Item{7.}  $36$.

\ResetCols[3]
\Item{8.}  $4a + 4b$.

\Item{9.}  $4a - 4b$.

\Item{10.}  $2a^{2} + 2b^{2}$.

\Item{11.}  $2a^{2} - 2b^{2}$.

\Item{12.}  $3ab + 3c$.

\Item{13.}  $3ab - 3c$.

\Item{14.}  $3c - 3ab$.

\Item{15.}  $ab + ac$.

\Item{16.}  $ab - ac$.

\Item{17.}  $3ab + 3ac$.

\Item{18.}  $3ab - 3ac$.

\Item{19.}  $5ab^{2} + 5ac$.

\Item{20.}  $5ab^{2} - 5ac^{2}$.

\Item{21.}  $5a^{2}b^{2} - 5a^{2}c$.


\AnsTo[5]{Exercise}{3.} %  Page 12.

\Item{1.}  $63$.

\Item{2.}  $280$.

\Item{3.}  $300$.

\Item{4.}  $98$.

\Item{5.}  $81$.

\Item{6.}  $1250$.

\Item{7.}  $105$.

\Item{8.}  $105$.

\Item{9.}  $315$.

\Item{10.}  $35$.

\Item{11.}  $105$.

\Item{12.}  $105$.

\Item{13.}  $0$.

\Item{14.}  $135$.

\Item{15.}  $120$.

\Item{16.}  $0$.

\Item{17.}  $1800$.

\Item{18.}  $540$.

\Item{19.}  $0$.

\Item{20.}  $270$.

\Item{21.}  $540$.


\AnsTo[6]{Exercise}{4.} %  Page 13.

\Item{1.}  $21$.

\Item{2.}  $26$.

\Item{3.}  $72$.

\Item{4.}  $85$.

\Item{5.}  $30$.

\Item{6.}  $17$.

\Item{7.}  $8$.

\Item{8.}  $50$.

\Item{9.}  $24$.

\Item{10.}  $0$.

\Item{11.}  $12$.

\Item{12.}  $100$.

\Item{13.}  $80$.

\Item{14.}  $71$.

\Item{15.}  $139$.

\Item{16.}  $17$.

\Item{17.}  $8$.

\Item{18.}  $5$.

\Item{19.}  $3$.

\Item{20.}  $6$.

\Item{21.}  $5$.

\Item{22.}  $1$.

\Item{23.}  $2$.

\Item{24.}  $2$.


\AnsTo{Exercise}{5.} %  Page 14.

\Item{1.}  $a$~plus~$b$; $a$~minus~$b$; $a$~times~$b$; $a$~divided by~$b$.

\Item{3.}  $a + b$.

\Item{5.}  $a - b$.

\Item{7.}  $x - y$.

\Item{9.}  $4x$; $x^{4}$.

\Item{11.}  $35 - x$.

\Item{12.}  $x - a$.

\Item{14.}  $14 - x$.

\Item{15.}  $a - x$.

\Item{17.}  $xy$.

\Item{18.}  $\dfrac{x}{y}$.
%% -----File: 172.png---Folio 166-------


\AnsTo[2]{Exercise}{6.} %  Page 15.

\Item{2.}  $\dfrac{a}{b}$.

\Item{4.}  $(x - 3)$~yr.; $(x + 7)$~yr.

\Item{6.}  $7(2x - y)$.

\Item{8.}  $x + 1$; $x - 1$.

\Item{9.}  $20 - d$.

\Item{11.}  $x + 8$.

\Item{13.}  $x - 10$.

\Item{14.}  $10$.


\AnsTo[2]{Exercise}{7.} %  Page 16.

\Item{1.}  $(40 - x)$~yr.

\Item{2.}  $(a + y)$~yr.

\Item{3.}  $4$.

\Item{5.}  $ab$.

\Item{6.}  $5x - 3x$.

\Item{8.}  $2x - 3 - (x + 1)$.

\Item{9.}  $40$.

\Item{10.}  $12$.

\Item{11.}  $100a + 25b + 10c$.

\Item{12.}  $100 - x - y$.

\Item{14.}  $xy + c$.


\AnsTo[2]{Exercise}{8.} %  Page 17.

\Item{2.}  $xy - a^{2}$.

\Item{3.}  $\dfrac{ph}{gk}$.

\Item{4.}  $6m^{2} + 5c(d + b - a)$.

\Item{5.}  $5(2n + 1) - 6(c - a + b)$.

\Item{6.}  \$$100 - \text{\$}(a + b + 2c)$.

\Item{8.}  $\dfrac{1}{x}$.

\Item{9.}  $\dfrac{1}{x} + \dfrac{1}{y} + \dfrac{1}{z}$.

\ResetCols[1]
\Item{10.}  $n + (n + 1) + (n + 2)$; $n(n + 1)(n + 2)$.

\ResetCols[4]
\Item{11.}  $\dfrac{36}{x}$.

\Item{12.}  $qd$.

\Item{13.}  $qd + r$.

\Item{14.}  $2x$.

\Item{15.}  $\dfrac{10x}{12}$.

\Item{16.}  $\dfrac{6a}{b}$.

\Item{17.}  $4m$.

\Item{18.}  $x + 17$.

\Item{19.}  $x + y + c$.

\Item{20.}  $36x$.

\Item{21.}  $\dfrac{c}{4}$.

\Item{22.}  $x - 1$, $x$, $x + 1$.

\Item{23.}  $2n + 3$.


\AnsTo[5]{Exercise}{9.} %  Page 24.

\Item{1.}  $4$.

\Item{2.}  $5$.

\Item{3.}  $3$.

\Item{4.}  $1$.

\Item{5.}  $13$.

\Item{6.}  $0$.

\Item{7.}  $5$.

\Item{8.}  $5$.

\Item{9.}  $7$.

\Item{10.}  $4$.

\Item{11.}  $4$.

\Item{12.}  $1$.

\Item{13.}  $25$.

\Item{14.}  $32\frac{1}{2}$.

\Item{15.}  $6$.

\Item{16.}  $10$.

\Item{17.}  $3$.

\Item{18.}  $1$.

\Item{19.}  $7$.

\Item{20.}  $13$.

\Item{21.}  $12$.

\Item{22.}  $14$.

\Item{23.}  $7$.

\Item{24.}  $5$.

\Item{25.}  $0$.

\Item{26.}  $15$.

\Item{27.}  $10\frac{1}{2}$.

\Item{28.}  $4$.

\Item{29.}  $3$.

\Item{30.}  $4$.

\Item{31.}  $\frac{2}{5}$.

\Item{32.}  $6$.

\Item{33.}  $9$.

\Item{34.}  $7$.

\Item{35.}  $5$.

\Item{36.}  $6$.

\Item{37.}  $6$.


\AnsTo[2]{Exercise}{10.} %  Page 27.

\Item{1.}  $30$.

\Item{2.}  Son, $10$~yr.; father, $50$~yr.

\Item{3.}  $78$, $13$.

\Item{4.}  $80$~ft.\ broken off; $10$~ft.\ standing.

\Item{5.}  $23$, $30$.

\Item{6.}  $36$, $48$.

\Item{7.}  $15$, $20$.

\Item{8.}  $20$.

\Item{9.}  $12$.

\Item{10.}  $10$, $40$.

\Item{11.}  $26$, $10$.

\Item{12.}  $25$, $15$.

\Item{13.}  $10$, $20$.

\Item{14.}  $7$, $20$.

\Item{15.}  $12$, $20$.
%% -----File: 173.png---Folio 167-------


\AnsTo{Exercise}{11.} %  Page 29.

\Item{1.}  Cow, \$$42$; horse, \$$168$.

\Item{2.}  $81$.

\Item{3.}  $2$.

\Item{4.}  $30$, $40$.

\Item{5.}  $25$, $26$, $27$.

\Item{6.}  $5$, $6$, $7$, $8$, $9$.

\Item{7.}  A, $30$~yr.; B, $10$~yr.

\Item{8.}  Father, $40$~yr., son, $10$~yr.

\Item{9.}  $40$.

\Item{10.}  $10$.

\Item{11.}  \$$40$.

\Item{12.}  $9$.


\AnsTo[2]{Exercise}{12.} %  Page 30.

\Item{1.}  $15$~men; $30$~women; $45$~children.

\Item{2.}  $50$.

\Item{3.}  $16$.

\Item{4.}  $7$.

\Item{5.}  $35$.

\Item{6.}  $24$.

\Item{7.}  $24$.

\Item{8.}  $20$.

\Item{9.}  $970$; $1074$.


\AnsTo[2]{Exercise}{13.} %  Page 31.

\Item{1.}  $24$.

\Item{2.}  A, \$$60$; B, \$$30$.

\Item{3.}  \$$3$~quarters; $6$~bills.

\Item{4.}  $14$.

\Item{5.}  \$$4$~quarters; $20$~half-dollars.

\Item{6.}  \$$7$~ten-dollar bills; $21$~one-dollar bills.

\Item{7.}  Father, $32$~yr.; son, $8$~yr.

\Item{8.}  $20$.


\AnsTo[2]{Exercise}{14.} %  Page 32.

\Item{1.}  $9$~miles.

\Item{2.}  \$$60$.

\Item{3.}  \$$20$~lb.\ at $65$~cts.;
    $60$~lb.\ at $45$~cts.

\Item{4.}  $40$.

\Item{5.}  $15,000$.

\Item{6.}  $15$~in.; $21$~in.

\Item{7.}  \$$2$~doz.\ at $25$~cts.;
    $5$~doz.\ at $20$~cts.

\Item{8.}  \$$6$~quarters, $18$~ten-cent pieces.


\AnsTo[4]{Exercise}{15.} %  Page 39.

\Item{1.}  $40c$.

\Item{2.}  $24a$.

\Item{3.}  $39x$.

\Item{4.}  $51y$.

\Item{5.}  $-26a$.

\Item{6.}  $-40x$.

\Item{7.}  $-17b$.

\Item{8.}  $-66z$.

\Item{9.}  $-20m$.

\Item{10.}  $2d$.

\Item{11.}  $0$.

\Item{12.}  $-18g$.

\Item{13.}  $a^{2}$.

\Item{14.}  $-21x^{3}$.

\Item{15.}  $0$.

\Item{16.}  $3mn$.

\Item{17.}  $0$.

\Item{18.}  $-3a^{3}b^{3}c^{3}$.

\Item{19.}  $-9abcd$.

\Item{20.}  $1$.

\Item{21.}  $12$.

\Item{22.}  $4$.

\Item{23.}  $-18$.

\Item{24.}  $10$.
%% -----File: 174.png---Folio 168-------


\AnsTo{Exercise}{16.} %  Page 42.

\Item{1.}  $30a^{5}$.

\Item{2.}  $40a^{4}b^{3}$.

\Item{3.}  $63x^{2}y^{2}$.

\Item{4.}  $2a^{5}b^{5}c^{2}$.

\Item{5.}  $9a^{7}b^{9}c^{9}$.

\Item{6.}  $-10a^{2}$.

\Item{7.}  $12ab$.

\Item{8.}  $-a^{4}b^{3}$.

\Item{9.}  $10a^{5}b^{5}c$.

\Item{10.}  $12x^{7}y^{5}z^{2}$.

\Item{11.}  $105a^{7}b^{5}$.

\Item{12.}  $6a^{6}b^{5}c^{7}$.

\Item{13.}  $-12a^{5}b^{5}c^{5}x^{5}$.

\Item{14.}  $24a^{7}b^{6}c^{5}$.

\Item{15.}  $-42a^{6}m^{5}x^{7}$.

\Item{16.}  $30x^{8}y^{4}z^{6}$.

\Item{17.}  $-46$.

\Item{18.}  $-3$.

\Item{19.}  $-8$.

\Item{20.}  $-17$.

\Item{21.}  $9$.

\Item{22.}  $12$.

\Item{23.}  $-102$.

\Item{24.}  $-41$.

\Item{25.}  $174$.

\Item{26.}  $6$.

\Item{27.}  $30$.

\Item{28.}  $372$.


\AnsTo[4]{Exercise}{17.} %  Page 45.

\Item{1.}  $x^{2}$.

\Item{2.}  $3x^{2}$.

\Item{3.}  $-7$.

\Item{4.}  $-7$.

\Item{5.}  $7x^{4}$.

\Item{6.}  $9x$.

\Item{7.}  $-4a$.

\Item{8.}  $4x^{2}y^{2}$.

\Item{9.}  $-9x$.

\Item{10.}  $5x$.

\Item{11.}  $3y^{2}$.

\Item{12.}  $-4ab^{2}$.

\Item{13.}  $12xy^{4}$.

\Item{14.}  $\dfrac{3x^{3}y^{3}}{5}$.

\Item{15.}  $-\dfrac{3a}{2}$.

\Item{16.}  $-bd$.

\Item{17.}  $acd^{2}$.

\Item{18.}  $-\dfrac{2xy}{3}$.

\Item{19.}  $5ab$.

\Item{20.}  $4mn$.

\Item{21.}  $\dfrac{yz^{3}}{3}$.

\Item{22.}  $-17cd$.

\Item{23.}  $2n^{2}p$.

\Item{24.}  $\dfrac{3r^{2}}{p}$.

\Item{25.}  $-13agt$.

\Item{26.}  $\dfrac{1}{abc}$.

\Item{27.}  $\dfrac{3}{2xy^{2}z^{3}}$.

\Item{28.}  $\dfrac{2}{mnp}$.

\Item{29.}  $-\dfrac{a}{3b^{2}}$.

\Item{30.}  $-\dfrac{g}{3mt}$.


\AnsTo[2]{Exercise}{18.} %  Page 48.

\Item{1.}  $2a^{2} + 2b^{2}$.

\Item{2.}  $9a^{2} - 2a + 6$.

\Item{3.}  $0$.

\Item{4.}  $4x + 4y + 4z$.

\Item{5.}  $2b + 2c$.

\Item{6.}  $4a + 4b + 4c$.

\Item{7.}  $3a^{2} + 5a - 2$.

\Item{8.}  $8ab + 3ac$.

\Item{9.}  $6x^{3}$.

\Item{10.}  $5x^{2} + 3xy - 2y^{2}$.

\Item{11.}  $a^{2} - 2b^{2}$.

\Item{12.}  $4a^{3} + 6a + 2$.

\Item{13.}  $3m^{3} + 7m^{2} + 2$.

\Item{14.}  $2x^{3} - 2x^{2} + 4x + y$.

\Item{15.}  $7x^{3} + 7x^{2} + 2$.

\Item{16.}  $a^{3} + 3a^{2}b - 5ab^{2} - b^{3}$.

\Item{17.}  $-a^{3} - a^{2}b - 2ab^{2} - 2b^{3}$.

\Item{18.}  $2x^{3} + 2x^{2}y - xy^{2} + 6y^{3}$.
%% -----File: 175.png---Folio 169-------


\AnsTo[2]{Exercise}{19.} %  Page 50.

\Item{1.}  $a - b + c$.

\Item{2.}  $2a - 2b + 6c$.

\Item{3.}  $2a + 3y -8z$.

\Item{4.}  $x + 4y + 5z$.

\Item{5.}  $2ac + 2bc$.

\Item{6.}  $2ab - 3ac + 4bc$.

\Item{7.}  $x^{3} + 3x^{2} + 2x - 8$.

\Item{8.}  $x^{3} - 7x^{2} + 4x$.

\Item{9.}  $2b^{3} + 18abc - 15c^{3}$.

\Item{10.}  $2 - x - 2x^{2} + 2x^{3}$.

\Item{11.}  $-3b^{3} - 4c^{3} + 6abc$.

\Item{12.}  $-x^{4} - 2x^{3} + 4x^{2} - 7x + 5$.

\Item{13.}  $2x + x^{2} + x^{3} + x^{5}$.

\Item{14.}  $2b^{3} - 4a^{2}b + 2ab^{2}$.

\Item{15.}  $2b^{4} - 2a^{3}b^{3} - ab^{2}$.

\Item{16.}  $4x^{3} - 3x^{2}y - 4xy^{2} + 7y^{3}$.


\AnsTo{Exercise}{20.} %  Page 52.

\Item{1.}  $c$.

\Item{2.}  $y - b$.

\Item{3.}  $x - 3y - 7c$.

\Item{4.}  $7a + 2b-2$.

\Item{5.}  $2a + x$.

\Item{6.}  $13x - 15y + 13z$.

\Item{7.}  $2a - 2b + 2c$.

\Item{8.}  $5a + b - 4c$.

\Item{9.}  $4x - 5y + 2z$.

\Item{10.}  $3a - b$.

\Item{11.}  $2x + 3y + z$.

\Item{12.}  $-8x + y$.

\Item{13.}  $6z - 2y-z$.
\ResetCols

\Item{15.}  $(a + c)x - (a - b)y + (a - c)z$.

\Item{16.}  $(2a + 5c)x - (3a + 4b)y - (6b + 7c)z$.

\Item{17.}  $(ac - an)x - (bm + cn)y + (a + 3c)z$.

\Item{18.}  $(mn - 1)x - (mn + 1)y + (mn + 1)z$.


\AnsTo[2]{Exercise}{21.} %  Page 53.

\Item{1.}  $x^{2} + 7x$.

\Item{2.}  $8x^{2} - 12xy$.

\Item{3.}  $14xy - 21y^{2}$.

\Item{4.}  $2ax - 4a^{2}$.

\Item{5.}  $bx - 3b^{2}$.

\Item{6.}  $-6a^{3} + 9a^{2}b$.

\Item{7.}  $10x^{2}z + 15xz^{2}$.

\Item{8.}  $5a^{3}b - 25a^{2}b^{2}$.

\Item{9.}  $-x^{2}y^{2} + 3xy^{3}$.

\Item{10.}  $4x^{5} - 6x^{4}$.

\Item{11.}  $4x^{2}y - 12y^{3}$.

\Item{12.}  $-x^{4} + 3x^{2}y^{2}$.

\Item{13.}  $-a^{3}b^{3} + a^{5}b^{2}$.

\Item{14.}  $a^{4}b^{2} + a^{5}$.

\Item{15.}  $4x^{5} - 6x^{4} + 2x^{3}$.

\Item{16.}  $5a^{3}b - 25a^{2}b^{2} - 5ab^{3}$.

\Item{17.}  $a^{5} + 2a^{4}b + 2a^{3}b^{2}$.

\Item{18.}  $a^{3}b^{3} + 2a^{2}b^{4} + 2ab^{5}$.

\Item{19.}  $8x^{3} - 12x^{2}y - 18xy^{2}$.

\Item{20.}  $x^{2}y + 2xy^{2} - y^{3}$.

\Item{21.}  $a^{5} + a^{4}b^{2} + a^{2}b^{3}$.

\Item{22.}  $x^{2}y^{2} - 2xy^{3} + y^{4}$.

\Item{23.}  $15a^{4}b^{4} - 20a^{3}b^{5} + 5a^{5}b^{3}$.

\Item{24.}  $3a^{2}x^{2}y^{2} - 9a^{2}xy^{4} + 3a^{2}y^{6}$.

\Item{25.}  $x^{15}y^{2} - x^{13}y^{5} - x^{6}y^{12}$.

\Item{26.}  $4x^{5}y^{3} - 6x^{4}y^{5} + 4x^{3}y^{6}$.

\Item{27.}  $a^{10}x^{5}y^{10} - a^{9}x^{4}y^{9} - a^{8}x^{3}y^{8}$.

\Item{28.}  $15a^{4}b^{5} - 10a^{3}b^{6} + 25a^{5}b^{4}$.
%% -----File: 176.png---Folio 170-------


\AnsTo[2]{Exercise}{22.} %  Page 56.

\Item{1.}  $x^{2} + 13x + 42$.

\Item{2.}  $x^{2} - x - 42$.

\Item{3.}  $x^{2} + x - 42$.

\Item{4.}  $x^{2} - 13x + 42$.

\Item{5.}  $x^{2} + 3x - 40$.

\Item{6.}  $4x^{2} + 12x + 9$.

\Item{7.}  $4x^{2} - 12x + 9$.

\Item{8.}  $4x^{2} - 9$.

\Item{9.}  $-9x^{2} + 12x - 4$.

\Item{10.}  $20x^{2} - 47x + 21$.

\Item{11.}  $a^{2} + ab - 6b^{2}$.

\Item{12.}  $a^{2} - 12ab + 35b^{2}$.

\Item{13.}  $25x^{2} - 30xy + 9y^{2}$.

\Item{14.}  $x^{2} - bx - cx + bc$.

\Item{15.}  $8m^{2} - 10mp + 3p^{2}$.

\Item{16.}  $a^{2} + ab - bc - c^{2}$.

\Item{17.}  $a^{4} - a^{3}b + 2a^{2}b^{2} - ab^{3} + b^{4}$.

\Item{18.}  $x^{5} - 3x^{4} - x^{3} + 16x^{2} - 21$.

\Item{19.}  $a^{3} - b^{3}$.

\Item{20.}  $a^{3} + b^{3}$.

\Item{21.}  $2x^{4} + 13x^{3} - 9x^{2} - 50x + 40$.

\Item{22.}  $9x^{5} - 7x^{3} + 6x^{2} - 2x$.

\Item{23.}  $x^{5} - 4x^{2}y^{3} + 3xy^{4}$.

\Item{24.}  $-a^{4} + 4a^{3}b - 7ab^{3} - 2b^{4}$.
\ResetCols[1]

\Item{25.}  $-25a^{5}b^{3} + 20a^{4}b^{4} + 12a^{3}b^{5} - 5a^{2}b^{6} - 2ab^{7}$.

\Item{26.}  $a^{4} - 2a^{2}b^{2} + b^{4}$.

\Item{27.}  $a^{2}b^{2} + 2abcd - a^{2}c^{2} + c^{2}d^{2}$.

\Item{28.}  $-2x^{5}y^{3} + x^{4}y^{4} + 10x^{3}y^{5} - 8x^{2}y^{6} - 3xy^{7}$.

\Item{29.}  $x^{4} - 4x^{2}y^{2} + 4xy^{3} - y^{4}$.

\Item{30.}  $3x^{4} - 5x^{3}y - 12x^{2}y^{2} - xy^{3} + 3y^{4}$.

\Item{31.}  $-a^{4} + 6a^{2}b^{2} - b^{4}$.

\Item{32.}  $a^{4} - a^{3}c + ab^{2}c - b^{4} - 2b^{2}c^{2} + ac^{3} - c^{4}$.

\Item{33.}  $a^{4} - 16a^{2}b^{2}x^{2} + 32a^{3}b^{3}x^{3} - 16a^{4}b^{4}x^{4}$.

\Item{34.}  $6a^{4} + 5a^{3}bx - 10a^{2}b^{2}x^{2} + 7ab^{3}x^{3} - 2b^{4}x^{4}$.

\Item{35.}  $10x^{6}y^{2} + 14x^{5}y^{3} - 48x^{4}y^{4} + 32x^{3}y^{5} - 16x^{2}y^{6}$.


\AnsTo{Exercise}{23.} %  Page 58.

\Item{1.}  $2a^{2} - a$.

\Item{2.}  $7a^{4} - a$.

\Item{3.}  $7x^{2} + 1$.

\Item{4.}  $5m^{4} - p^{2}$.

\Item{5.}  $3x^{3} - 5x^{2}$.

\Item{6.}  $-3x^{3} + 1$.

\Item{7.}  $2x^{2} - 3x$.

\Item{8.}  $-x^{2} + 2$.

\Item{9.}  $a + 2c$.

\Item{10.}  $5x - y$.

\Item{11.}  $ax - 1$.

\Item{12.}  $x + xy$.

\Item{13.}  $-3a + 4b - 2c$.

\Item{14.}  $ab - b^{4} - a^{2}b$.

\Item{15.}  $x^{2} - 2xy - 3y^{2}$.

\Item{16.}  $xy - x^{2} - y^{2}$.

\Item{17.}  $-a^{2} + ab + b^{2}$.

\Item{18.}  $-a + 1 - b$.

\Item{19.}  $-1 + xy - x^{2}y^{2}$.

\Item{20.}  $x^{2} + 2x + 1$.

\Item{21.}  $a - b - c$.

\Item{22.}  $x^{3} - x^{2}y - y^{2}$.

\Item{23.}  $ab - 2 - 3b^{2}$.

\Item{24.}  $a^{2}c^{2} + a - c$.
%% -----File: 177.png---Folio 171-------


\AnsTo{Exercise}{24.} %  Page 62.

\Item{1.}  $x + 8$.

\Item{2.}  $x - 8$.

\Item{3.}  $x + 8$.

\Item{4.}  $x - 8$.

\Item{5.}  $a + 5$.

\Item{6.}  $3a + 1$.

\Item{7.}  $a + 5$.

\Item{8.}  $-3a - 2$.

\Item{9.}  $x^{2} - x + 1$.

\Item{10.}  $x^{4} + x^{2} + 1$.

\Item{11.}  $1 + ab + a^{2}b^{2}$.

\Item{12.}  $x^{2} + 3x + 1$.

\Item{13.}  $a - b + c$.

\Item{14.}  $a + b - c$.

\Item{15.}  $x + y - z$.

\Item{16.}  $c^{2} + c + 2$.

\Item{17.}  $x - 2y - z$.

\Item{18.}  $x - a$.

\Item{19.}  $a - 2b + 3c$.

\Item{20.}  $a^{2} + 5a + 6$.

\Item{21.}  $q^{2} + 3q + 2$.

\Item{22.}  $9a^{2} + 6ab + 4b^{2}$.

\Item{23.}  $-65$.

\Item{24.}  $10$.

\Item{25.}  $7a - 45$.

\Item{26.}  $2x^{4}$.


\AnsTo[1]{Exercise}{25.} %  Page 63.

\Item{1.}  $2a^{2}$.

\Item{2.}  $-3a^{4} + 2a^{3}b - 2ab^{3} + 4b^{4}$.

\Item{3.}  $x$.

\Item{4.}  $a^{4} + 2a^{2}b^{2} + b^{4} - c^{4} + 2c^{2}d^{2} - d^{4}$.

\Item{5.}  $10y^{4} + 8y^{3} + 6y^{2} + 4y + 2$.

\Item{6.}  $0$.

\Item{7.}  $2z - 7y$.

\Item{8.}  $a^{3} - 3abc + b^{3} + c^{3}$.

\Item{9.}  $4y^{2} - 3xy + 2x^{2}$.

\Item{10.}  $5a^{3}b - b^{4}$.

\Item{11.}  $3x^{3} - 2x^{2} + 1$.

\Item{12.}  $3c^{2} + 24c - 12$.

\Item{13.}  $2b^{4}$.

\Item{14.}  $10 - 16x - 39x^{2} + 2x^{3} + 15x^{4}$.

\Item{15.}  $a^{4} - ax^{3} + x^{4}$.

\Item{16.}  $(a - b)x^{3} + (b + c)x^{2} - (c + 1)x$.

\Item{17.}  $(a + b)x^{4} - (a - b)x^{3} - (c + 2)x$.

\Item{18.}  $(a + 1)x^{3} - (b + c)x^{2} + (b - c)x$.


\AnsTo[2]{Exercise}{26.} %  Page 65.

\Item{1.}  $m^{2} + 2mn + n^{2}$.

\Item{2.}  $c^{2} - 2ac + a^{2}$.

\Item{3.}  $a^{2} + 4ac + 4c^{2}$

\Item{4.}  $9a^{2} - 12ab + 4b^{2}$.

\Item{5.}  $4a^{2} + 12ab + 9b^{2}$.

\Item{6.}  $a^{2} - 6ab + 9b^{2}$.

\Item{7.}  $4x^{2} - 4xy + y^{2}$.

\Item{8.}  $y^{2} - 4xy + 4x^{2}$.

\Item{9.}  $a^{2} + 10ab + 25b^{2}$.

\Item{10.}  $4a2 - 20ac + 25c^{2}$.

\Item{11.}  $x^{2} - y^{2}$.

\Item{12.}  $16a^{2} - b^{2}$.

\Item{13.}  $4b^{2} - 9c^{2}$.

\Item{14.}  $x^{2} + 10bx + 25b^{2}$.

\Item{15.}  $y^{2} - 4yz + 4z^{2}$.

\Item{16.}  $y^{2} - 9z^{2}$.

\Item{17.}  $4a^{2} - 9b^{2}$.

\Item{18.}  $4a^{2} - 12ab + 9b^{2}$.

\Item{19.}  $4a^{2} + 12ab + 9b^{2}$.

\Item{20.}  $25x^{2} - 9a^{2}$.
%% -----File: 178.png---Folio 172-------


\AnsTo[2]{Exercise}{27.} %  Page 67.

\Item{1.}  $x^{2} + 11x + 28$.

\Item{2.}  $x^{2} + 4x - 21$.

\Item{3.}  $x^{2} - 6x + 8$.

\Item{4.}  $x^{2} - 16x + 60$.

\Item{5.}  $x^{2} + 3x - 28$.

\Item{6.}  $x^{2} - ax - 2a^{2}$.

\Item{7.}  $x^{2} + 2ax - 3a^{2}$.

\Item{8.}  $a^{2} + 6ac + 9c^{2}$.

\Item{9.}  $a^{2} - 2ax - 8x^{2}$.

\Item{10.}  $a^{2} - 7ab + 12b^{2}$.

\Item{11.}  $a^{4} + a^{2}c - 2c^{2}$.

\Item{12.}  $x^{2} - 20x + 51$.

\Item{13.}  $x^{2} + xy - 30y^{2}$.

\Item{14.}  $9 + 3x - 2x^{2}$.

\Item{15.}  $5 - 8x - 4x^{2}$.

\Item{16.}  $a^{2} + ab - 6b^{2}$.

\Item{17.}  $a^{4}b^{4} - 6a^{2}b^{2}x^{2} + 5x^{4}$.

\Item{18.}  $a^{6}b^{2} + 4a^{4}b^{4} - 5a^{2}b^{6}$.

\Item{19.}  $x^{4}y^{2} - 4x^{3}y^{3} + 3x^{2}y^{4}$.

\Item{20.}  $x^{4}y^{2} + 2x^{3}y^{3} + x^{2}y^{4}$.

\Item{21.}  $x^{2} + (a + b)x + ab$.

\Item{22.}  $x^{2} + (a - b)x - ab$.

\Item{23.}  $x^{2} - (a - b)x - ab$.

\Item{24.}  $x^{2} - (a + b)x + ab$.

\Item{25.}  $x^{2} + (2a + 2b)x + 4ab$.

\Item{26.}  $x^{2} - (2a - 2b)x - 4ab$.

\Item{27.}  $x^{2} + (2a - 2b)x - 4ab$.

\Item{28.}  $x^{2} - (2a + 2b)x + 4ab$.

\Item{29.}  $x^{2} + 2ax - 3a^{2}$.

\Item{30.}  $x^{2} + ax - 6a^{2}$.


\AnsTo{Exercise}{28.} %  Page 68.

\Item{1.}  $x + 2$.

\Item{2.}  $x - 2$.

\Item{3.}  $a + 3$.

\Item{4.}  $a - 3$.

\Item{5.}  $c + 5$.

\Item{6.}  $c - 5$.

\Item{7.}  $7x + y$.

\Item{8.}  $7x - y$.

\Item{9.}  $3b + 1$.

\Item{10.}  $3b - 1$.

\Item{11.}  $4x^{2} + 5a$.

\Item{12.}  $4x^{2} - 5a$.

\Item{13.}  $3x + 5y$.

\Item{14.}  $a + b - c$.

\Item{15.}  $a - b + c$.

\Item{16.}  $a + 2b - c$.

\Item{17.}  $5a - 7b + 1$.

\Item{18.}  $5a - 7b - 1$.

\Item{19.}  $z + x - y$.

\Item{20.}  $z - x + y$.

\Item{21.}  $a - 2b + c$.

\Item{22.}  $x + 3y + z$.

\Item{23.}  $x + 3y - z$.

\Item{24.}  $a + 2b + 2c$.

\Item{25.}  $a + 2b - 2c$.

\Item{26.}  $1 - 3x + 2y$.


\AnsTo[2]{Exercise}{29.} %  Page 69.

\Item{1.}  $1 + x + x^{2}$.

\Item{2.}  $1 + 2a + 4a^{2}$.

\Item{3.}  $1 + 3c + 9c^{2}$.

\Item{4.}  $4a^{2} + 2ab + b^{2}$.

\Item{5.}  $16b^{2} + 12bc + 9c^{2}$.

\Item{6.}  $9x^{2} + 6xy + 4y^{2}$.

\Item{7.}  $x^{2}y^{2} + xyz + z^{2}$.

\Item{8.}  $a^{2}b^{2} + 2ab + 4$.

\Item{9.}  $25a^{2} + 5ab + b^{2}$.

\Item{10.}  $a^{2} + 2ab + 4b^{2}$.

\Item{11.}  $a^{2} + 4a + 16$.

\Item{12.}  $a^{6} + 3a^{3} + 9$.

\Item{13.}  $a^{8} + a^{4}x^{2}y^{2} + x^{4}y^{4}$.

\Item{14.}  $x^{10} + x^{5}a^{3}b^{3} + a^{6}b^{6}$.

\Item{15.}  $9x^{2}y^{2} + 3xyz^{4} + z^{8}$.

\Item{16.}  $x^{2}y^{2}z^{2} + xyz + 1$.

\Item{17.}  $4a^{2}b^{2}c^{2} - 6abc + 9$.

\Item{18.}  $1 + 4xyz + 16x^{2}y^{2}z^{2}$.
%% -----File: 179.png---Folio 173-------


\AnsTo[2]{Exercise}{30.} %  Page 70.

\Item{1.}  $1 - x + x^{2}$.

\Item{2.}  $1 - 2a + 4a^{2}$.

\Item{3.}  $1 - 3c + 9c^{2}$.

\Item{4.}  $4a^{2} - 2ab + b^{2}$.

\Item{5.}  $16b^{2} - 12bc + 9c^{2}$.

\Item{6.}  $9x^{2} - 6xy + 4y^{2}$.

\Item{7.}  $4x^{2} - 10xy + 25y^{2}$.

\Item{8.}  $x^{2}y^{2} - xyz + z^{2}$.

\Item{9.}  $a^{2}b^{2} - 2ab + 4$.

\Item{10.}  $25a^{2} - 5ab + b^{2}$.

\Item{11.}  $a^{2} - 2ab + 4b^{2}$.

\Item{12.}  $a^{4} - 4a^{2} + 16$.

\Item{13.}  $a^{6} - 3a^{3} + 9$.

\Item{14.}  $4a^{4} - 2a^{2}b + b^{2}$.

\Item{15.}  $a^{8} - a^{4}x^{2}y^{2} + x^{4}y^{4}$.

\Item{16.}  $x^{10} - x^{5}a^{3}b^{3} + a^{6}b^{6}$.

\Item{17.}  $9x^{2}y^{2} - 3xyz^{4} + z^{8}$.

\Item{18.}  $x^{2}y^{2}z^{2} - xyz + 1$.

\Item{19.}  $4a^{2}b^{2}c^{2} - 6abc + 9$.

\Item{20.}  $1 - 4xyz + 16x^{2}y^{2}z^{2}$.

\Item{21.}  $1 - 3a^{2}bc + 9a^{4}b^{2}c^{2}$.

\Item{22.}  $x^{3} + x^{2}y + xy^{2} + y^{3}$.

\Item{23.}  $x^{3} - x^{2}y + xy^{2} - y^{3}$.

\Item{24.}  $x^{4} + x^{3}y + x^{2}y^{2} + xy^{3} + y^{4}$.

\Item{25.}  $x^{4} - x^{3}y + x^{2}y^{2} - xy^{3} + y^{4}$.

\Item{26.}  $x^{5} + x^{4}y + x^{3}y^{2} + x^{2}y^{3} + xy^{4} + y^{5}$.

\Item{27.}  $x^{5} - x^{4}y + x^{3}y^{2} - x^{2}y^{3} + xy^{4} - y^{5}$.


\AnsTo[2]{Exercise}{31.} %  Page 72.

\Item{1.}  $2x(x - 2)$.

\Item{2.}  $3a(a^{2} - 2)$.

\Item{3.}  $5a^{2}b^{2}(1 - 2ab)$.

\Item{4.}  $xy(3x + 4y)$.

\Item{5.}  $4a^{2}b^{2}(2a + b)$.

\Item{6.}  $3a^{2}(a^{2} - 4 - 2a)$.

\Item{7.}  $4x^{2}(1 - 2x^{2} - 3x^{3})$.

\Item{8.}  $5(1 - 2x^{2}y^{2} + 3x^{2}y)$.

\Item{9.}  $7a(a + 2 - 3a^{2})$.

\Item{10.}  $3x^{2}y^{2}(xy - 2x^{2}y^{2} - 3)$.


\AnsTo{Exercise}{32.} %  Page 73.

\Item{1.}  $(x^{2} + 1)(x + 1)$.

\Item{2.}  $(x^{2} + 1)(x - 1)$.

\Item{3.}  $(x + y)(x + z)$.

\Item{4.}  $(x - y)(a - b)$.

\Item{5.}  $(a + b)(a - c)$.

\Item{6.}  $(x + 3)(x - b)$.

\Item{7.}  $(x^{2} + 2)(2x - 1)$.

\Item{8.}  $(a - b)(a - 3)$.

\Item{9.}  $(2a - c)(3a + b)$.

\Item{10.}  $(xy + c)(ab + c)$.

\Item{11.}  $(a - b - c)(x - y)$.

\Item{12.}  $(a - b - 2c)(a - b)$.


\AnsTo[2]{Exercise}{33.} %  Page 74.

\Item{1.}  $(2 + x)(2 - x)$.

\Item{2.}  $(3 + x)(3 - x)$.

\Item{3.}  $(3a + x)(3a - x)$.

\Item{4.}  $(5 + x)(5 - x)$.

\Item{5.}  $(5x + a)(5x - a)$.

\Item{6.}  $(4a^{2} + 11)(4a^{2} - 11)$.

\Item{7.}  $(11a^{2} + 4)(11a^{2} - 4)$.

\Item{8.}  $(2ab + cd)(2ab - cd)$.

\Item{9.}  $(1 + xy)(1 - xy)$.

\Item{10.}  $(9xy + 1)(9xy - 1)$.

\Item{11.}  $(7ab + 2)(7ab - 2)$.

\Item{12.}  $(5a^{2}b^{2} + 3)(5a^{2}b^{2} - 3)$.

\Item{13.}  $(3a^{4}b^{3} + 4x^{5})(3a^{4}b^{3} - 4x^{5})$.

\Item{14.}  $(12xy + 1)(12xy - 1)$.

\Item{15.}  $(10x^{3}yz^{2} + 1)(10x^{3}yz^{2} - 1)$.

\Item{16.}  $(1 + 11a^{2}b^{4}c^{6})(1 - 11a^{2}b^{4}c^{6})$.

\Item{17.}  $(5a + 8x^{3}y^{3})(5a - 8x^{3}y^{3})$.

\Item{18.}  $(4x^{8} + 5y^{9})(4x^{8} - 5y^{9})$.

\Item{19.}  $90,000$.

\Item{20.}  $22,760$.

\Item{21.}  $732,200$.

\Item{22.}  $400$.

\Item{23.}  $28,972$.

\Item{24.}  $14,248,000$.
%% -----File: 180.png---Folio 174-------


\AnsTo[2]{Exercise}{34.} %  Page 75.

\Item{1.}  $(x + y + z)(x + y - z)$.

\Item{2.}  $(x - y + z)(x - y - z)$.

\Item{3.}  $(z + x + y)(z - x - y)$.

\Item{4.}  $(z + x - y)(z - x + y)$.

\Item{5.}  $(x + y + 2z)(x + y - 2z)$.

\Item{6.}  $(2z + x - y)(2z - x + y)$.

\Item{7.}  $(a + 2b + c)(a + 2b - c)$.

\Item{8.}  $(a - 2b + c)(a - 2b - c)$.

\Item{9.}  $(c + a - 2b)(c - a + 2b)$.

\Item{10.}  $(2a + 5c + 1)(2a + 5c - 1)$.

\Item{11.}  $(1 + 2a - 5c)(1 - 2a + 5c)$.

\Item{12.}  $(a + 3b + 4c)(a + 3b - 4c)$.

\Item{13.}  $(a - 5b + 3c)(a - 5b - 3c)$.

\Item{14.}  $(4c + a - 5b)(4c - a + 5b)$.

\Item{15.}  $(2a + x + y)(2a - x - y)$.

\Item{16.}  $(b + a - 2x)(b - a + 2x)$.

\Item{17.}  $(2z + x + 3y)(2z - x - 3y)$.

\Item{18.}  $(3 + 3a - 7b)(3 - 3a + 7b)$.

\Item{19.}  $(4a + 2b + 5c)(4a - 2b - 5c)$.

\Item{20.}  $(5c + 3a - 2x)(4c - 3a + 2x)$.

\Item{21.}  $(3a + 3b -5c)(3a - 3b + 5c)$.

\Item{22.}  $(4y + a - 3c)(4y - a + 3c)$.

\Item{23.}  $(7m + p + 2q)(7m - p - 2q)$.

\Item{24.}  $(6n + d - 2c)(6n - d + 2c)$.

\Item{25.}  $(x + y + a + b)(x + y - a - b)$.

\Item{26.}  $(x - y + a - b)(x - y - a + b)$.

\Item{27.}  $(2x + 3 + 2a + b)(2x + 3 - 2a - b)$.

\Item{28.}  $(b - c + a - 2x)(b - c - a + 2x)$.

\Item{29.}  $(3x - y + 2a - b)(3x - y - 2a + b)$.

\Item{30.}  $(x - 3y + a + 2b)(x - 3y - a - 2b)$.

\Item{31.}  $(x + 2y + a + 3b)(x + 2y - a - 3b)$.

\Item{32.}  $(x + y + a - z)(x + y - a + z)$.


\AnsTo[2]{Exercise}{35.} %  Page 77.

\Item{1.}  $(2x - y)(4x^{2} + 2xy + y^{2})$.

\Item{2.}  $(x - 1)(x^{2} + x + 1)$.

\Item{3.}  $(xy - z)(x^{2}y^{2} + xyz + z^{2})$.

\Item{4.}  $(x - 4)(x^{2} + 4x + 16)$.

\Item{5.}  $(5a - b)(25a^{2} + 5ab + b^{2})$.

\Item{6.}  $(a - 7)(a^{2} + 7a + 49)$.

\Item{7.}  $(ab - 3c)(a^{2}b^{2} + 3abc + 9c^{2})$.

\Item{8.}  $(xyz - 2)(x^{2}y^{2}z^{2} + 2xyz + 4)$.

\Item{9.}  $(2ab - 3y^{2})(4a^{2}b^{2} + 6aby^{2} + 9y^{4})$.

\Item{10.}  $(4x - y^{3})(16x^{2} + 4xy^{3} + y^{6})$.

\Item{11.}  $(3a - 4c^{2})(9a^{2} + 12ac^{2} + 16c^{4})$.

\Item{12.}  $(xy - 6z)(x^{2}y^{2} + 6xyz + 36z^{2})$.

\Item{13.}  $(4x - 9y)(16x^{2} + 36xy + 81y^{2})$.

\Item{14.}  $(3a - 8c)(9a^{2} + 24ac + 64c^{2})$.

\Item{15.}  $(2x^{2} - 5y)(4x^{4} + 10x^{2}y + 25y^{2})$.

\Item{16.}  $(4x^{2} - 3y^{5})(16x^{8} + 12x^{4}y^{5} + 9y^{10})$.

\Item{17.}  $(6 - 2a)(36 + 12a + 4a^{2})$.

\Item{18.}  $(7 - 3y)(49 + 21y + 9y^{2})$.
%% -----File: 181.png---Folio 175-------


\AnsTo[2]{Exercise}{36.} %  Page 78.

\Item{1.}  $(x + 1)(a^{2} - x + 1)$.

\Item{2.}  $(2x + y)(4x^{2} - 2xy + y^{2})$.

\Item{3.}  $(x + 5)(x^{2} - 5x + 25)$.

\Item{4.}  $(4a + 3)(16a^{2} - 12a + 9)$.

\Item{5.}  $(xy + z)(x^{2}y^{2} - xyz + z^{2})$.

\Item{6.}  $(a + 4)(a^{2} - 4a + 16)$.

\Item{7.}  $(2a^{2} + b)(4a^{4} - 2a^{2}b + b^{2})$.

\Item{8.}  $(x + 7)(x^{2} - 7x + 49)$.

\Item{9.}  $(2 + xyz)(4 - 2xyz + x^{2}y^{2}z^{2})$.

\Item{10.}  $(y^{3} + 4x)(y^{6} - 4y^{3}x + 16x^{2})$.

\Item{11.}  $(ab + 3x)(a^{2}b^{2} - 3abx + 9x^{2})$.

\Item{12.}  $(2yz + x^{2})(4y^{2}z^{2} - 2yzx^{2} + x^{4})$.

\Item{13.}  $(y^{3} + 4x^{2})(y^{6} - 4y^{3}x^{2} + 16x^{4})$.

\Item{14.}  $(4a^{4} + x^{5})(16a^{8} - 4a^{4}x^{5} + x^{10})$.

\Item{15.}  $(3x^{5} + 2a^{2})(9x^{10} - 6x^{5}a^{2} + 4a^{4})$.

\Item{16.}  $(3x^{3} + 8)(9x^{6} - 24x^{3} + 64)$.

\Item{17.}  $(7 + 4x)(49 - 28x + 16x^{2})$.

\Item{18.}  $(5 + 3y)(25 - 15y + 9y^{2})$.


\AnsTo{Exercise}{37.} %  Page 80.

\Item{1.}  $(2x + y)(2x + y)$.

\Item{2.}  $(x + 3y)(x + 3y)$.

\Item{3.}  $(x + 8)(x + 8)$.

\Item{4.}  $(x + 5a)(x + 5a)$.

\Item{5.}  $(a - 8)(a - 8)$.

\Item{6.}  $(a - 5b)(a - 5b)$.

\Item{7.}  $(c - 3d)(c - 3d)$.

\Item{8.}  $(2x - 1)(2x - 1)$.

\Item{9.}  $(2a - 3b)(2a - 3b)$.

\Item{10.}  $(3a - 4b)(3a - 4b)$.

\Item{11.}  $(x + 4y)(x + 4y)$.

\Item{12.}  $(x - 4y)(x - 4y)$.

\Item{13.}  $(2x - 5y)(2x - 5y)$.

\Item{14.}  $(1 + 10a)(1 + 10a)$.

\Item{15.}  $(7a - 2)(7a - 2)$.

\Item{16.}  $(6a + 5b)(6a + 5b)$.

\Item{17.}  $(9x - 2b)(9x - 2b)$.

\Item{18.}  $(mn + 7x^{2})(mn + 7x^{2})$.


\AnsTo{Exercise}{38.} %  Page 82.

\Item{1.}  $(a + 2)(a + 3)$.

\Item{2.}  $(a - 2)(a - 3)$.

\Item{3.}  $(a + 1)(a + 5)$.

\Item{4.}  $(a - 1)(a - 5)$.

\Item{5.}  $(a - 1)(a + 5)$.

\Item{6.}  $(a + 1)(a - 5)$.

\Item{7.}  $(c - 3)(c - 6)$.

\Item{8.}  $(c + 3)(c + 6)$.

\Item{9.}  $(c - 3)(c + 6)$.

\Item{10.}  $(c + 3)(c - 6)$.

\Item{11.}  $(x + 2)(x + 7)$.

\Item{12.}  $(x - 2)(x - 7)$.

\Item{13.}  $(x + 2)(x - 7)$.

\Item{14.}  $(x - 4)(x - 5)$.

\Item{15.}  $(x + 4)(x - 5)$.

\Item{16.}  $(x - 4)(x + 5)$.

\Item{17.}  $(x - 3)(x - 7)$.

\Item{18.}  $(x + 3)(x - 7)$.

\Item{19.}  $(x - 3)(x + 7)$.

\Item{20.}  $(x - 7)(x - 8)$.

\Item{21.}  $(x + 7)(x - 8)$.

\Item{22.}  $(x - 1)(x - 9)$.

\Item{23.}  $(x + 3)(x + 10)$.

\Item{24.}  $(x - 3)(x + 10)$.

\Item{25.}  $(x + 3)(x - 10)$.

\Item{26.}  $(a - 2b)(a + 3b)$.

\Item{27.}  $(a + 2b)(a - 3b)$.

\Item{28.}  $(a - b)(a + 4b)$.

\Item{29.}  $(a + b)(a - 4b)$.

\Item{30.}  $(ax + 7)(ax - 9)$.

\Item{31.}  $(a - 7x)(a + 9x)$.

\ResetCols[2]
\Item{32.}  $(a - 4b)(a - 5b)$.

\Item{33.}  $(xy - 3z)(xy - 16z)$.

\Item{34.}  $(ab + 4c)(ab + 11c)$.

\Item{35.}  $(x - 4y)(x - 9y)$.

\Item{36.}  $(x + 7y)(x + 12y)$.

\Item{37.}  $(ax - 6y)(ax - 17y)$.

\Item{38.}  $(x + 2y)(x - 2y)(x^{2} - 5y^{2})$.

\Item{39.}  $(a^{2}x^{2} - 11y^{2})(a^{2}x^{2} - 13y^{2})$.

\Item{40.}  $(a^{3}b^{3} - 11c^{2})(a^{3}b^{3} - 12c^{2})$.

\Item{41.}  $(a + 4bc)(a - 24bc)$.

\Item{42.}  $(a + 8bc)(a - 12bc)$.

\Item{43.}  $(a + 6bc)(a - 16bc)$.

\Item{44.}  $(a - 3bc)(a + 32bc)$.

\Item{45.}  $(a + 2bc)(a - 48bc)$.

\Item{46.}  $(a + bc)(a + 48bc)$.

\Item{47.}  $(x + 9yz)(x - 27yz)$.

\Item{48.}  $(xy + 13z)(xy - 14z)$.
%% -----File: 182.png---Folio 176-------


\AnsTo[2]{Exercise}{39.} %  Page 83.

\Item{1.}  $a(a^{2} - 7)$.

\Item{2.}  $ab(3ab - 2a^{2} + 3b^{2})$.

\Item{3.}  $(a - b + 1)(a - b)$.

\Item{4.}  $(a + b + 1)(a + b - 1)$.

\Item{5.}  $(a + 2b)(a^{2} - 2ab + 4b^{2})$.

\Item{6.}  $(x + 2y + 1)(x - 2y)$.

\Item{7.}  $(a - b)(a^{2} + ab + b^{2} + 1)$.

\Item{8.}  $(a - 3b)(a - 3b)$.

\Item{9.}  $(x + 1)(x - 2)$.

\Item{10.}  $(x + 1)(x - 3)$.

\Item{11.}  $(x - 3)(x + 7)$.

\Item{12.}  $(a + 2)(a - 13)$.

\Item{13.}  $(x^{2} + 3)(a + b)$.

\Item{14.}  $(x - 3)(x - y)$.

\Item{15.}  $(x - 3)(x - 4)$.

\Item{16.}  $(a + 2b)(a + 3b)$.

\Item{17.}  $(x^{2} + 5)(x^{2} + 5)$.

\Item{18.}  $(x - 9)(x - 9)$.

\Item{19.}  $(x - 10)(x - 11)$.

\Item{20.}  $(x + 8)(x + 11)$.

\Item{21.}  $(x - 8)(x + 11)$.

\Item{22.}  $(x^{2} + 1)(x - 1)$.

\Item{23.}  $x^{2}(3x + 1)(3x - 1)$.

\Item{24.}  $(1 + a - b)(1 - a + b)$.

\Item{25.}  $(a + b)(a^{2} - ab + b^{2} + 1)$.

\Item{26.}  $(m + n)(m - n)(x + y)$.

\Item{27.}  $(x - y + z)(x - y - z)$.

\Item{28.}  $(z + x - y)(z - x + y)$.

\Item{29.}  $(2a^{2} + 3a - 1)(2a^{2} - 3a + 1)$.

\Item{30.}  $(2x - y)(4x^{2} + 2xy + y^{2})$.

\Item{31.}  $x^{2}(x - 3y)$.

\Item{32.}  $(x - 3y)(x^{2} + 3xy + 9y^{2})$.

\Item{33.}  $(x - 5)(x + 8)$.

\Item{34.}  $(x - 2y)(x + 5y)$.

\Item{35.}  $(1 + 4x)(1 - 4x)$.

\Item{36.}  $a^{2}(a^{2} + 3b^{2})(a^{2} - 3b^{2})$.

\Item{37.}  $x(x + y)(x + 2y)$.

\Item{38.}  $x^{2}(x + y)(x + 3y)$.

\Item{39.}  $(x - 2y^{2})(x - 2y^{2})$.

\Item{40.}  $(4x^{2} + 1)(4x^{2} + 1)$.

\Item{41.}  $a^{2}(3a + 2c)(3a - 2c)$.

\Item{42.}  $ab(a + b)(a - 2b)$.

\Item{43.}  $(x + 2)(x^{2} - 2x + 4)(x - 1)$.

\Item{44.}  $(a + y)(a^{2} - ay + y^{2})(a - x)$.


\AnsTo{Exercise}{40.} %  Page 86.

\Item{1.}  $6$.

\Item{2.}  $5x^{3}$.

\Item{3.}  $6ax$.

\Item{4.}  $7ab^{2}$.

\Item{5.}  $7$.

\Item{6.}  $2a^{2}b^{2}$.

\Item{7.}  $x + 3y$.

\Item{8.}  $x + 3$.

\Item{9.}  $2a + 1$.

\Item{10.}  $x + y$.

\Item{11.}  $a + x$.

\Item{12.}  $a + 2b$.

\Item{13.}  $x - 1$.

\Item{14.}  $x + 3$.

\Item{15.}  $x - 6$.

\Item{16.}  $x^{2} - x + 1$.

\Item{17.}  $x - 1$.

\Item{18.}  $x - y$.

\Item{19.}  $x - 5$.

\Item{20.}  $a - b - c$.

\Item{21.}  $x + 2y$.

\Item{22.}  $x + 4y$.

\Item{23.}  $x^{2} + 2xy + 4y^{2}$.

\Item{24.}  $x - 2$.

\Item{25.}  $1 - 3a$.

\Item{26.}  $x - 7y$.

\Item{27.}  $2a + b$.

\Item{28.}  $x + y - z$.


\AnsTo[2]{Exercise}{41.} %  Page 88.

\Item{1.}  $18x^{2}y^{3}$.

\Item{2.}  $6a^{2}bc^{3}$.

\Item{3.}  $20a^{3}b^{3}$.

\Item{4.}  $30a^{3}b^{4}$.

\Item{5.}  $189x^{3}y^{5}$.

\Item{6.}  $x^{2}y^{3}z^{3}$.

\Item{7.}  $a^{2}(a + 1)$.

\Item{8.}  $x^{2}(x - 3)$.

\Item{9.}  $x(x + 1)(x - 1)$.

\Item{10.}  $x(x + 1)(x - 1)$.

\Item{11.}  $xy(x + y)$.

\Item{12.}  $x(x + 2)^{2}$.
%% -----File: 183.png---Folio 177-------

\Item{13.}  $(a + 2)(a + 2)(a + 3)$.

\Item{14.}  $(c - 4)(c + 5)(c - 6)$.

\Item{15.}  $(b - 5)(b - 6)(b + 7)$.

\Item{16.}  $(y - 4)(y + 5)(y - 6)$.

\Item{17.}  $(z - 5)(z - 6)(z + 7)$.

\Item{18.}  $(x - 4)(x + 8)(x - 8)(x^{2} + 4x + 16)$.

\Item{19.}  $(a + b)(a + b)(a - b)(a - b)$.

\Item{20.}  $4a^{2}b(a + b)(a + b)(a - b)$.

\Item{21.}  $(y + 2)(y + 3)(y + 4)$.

\Item{22.}  $(x + 1)(x - 1)(x^{2} + 1)$.

\Item{23.}  $(1 + x)(1 - x)(1 + x + x^{2})$.

\Item{24.}  $(x + y)(x + y)(x - y)(x - y)$.

\Item{25.}  $x(x - 3)(x + 5)(x^{2} + 3x + 9)$.

\Item{26.}  $(a + b + c)(a + b + c)(a + b - c)$.

\Item{27.}  $(x - a)(x - b)(x - c)$.

\Item{28.}  $a(a + b + c)(a + b - c)$.


\AnsTo[4]{Exercise}{42.} %  Page 90.

\Item{1.}  $\dfrac{1}{3b}$.

\Item{2.}  $\dfrac{4m}{5n}$.

\Item{3.}  $\dfrac{3m}{4p^{2}}$.

\Item{4.}  $\dfrac{x^{2}}{2yz}$.

\Item{5.}  $\dfrac{a^{3}b^{3}}{3c^{2}}$.

\Item{6.}  $\dfrac{2xy}{3}$.

\Item{7.}  $\dfrac{2m}{3p}$.

\Item{8.}  $\dfrac{3b^{2}c}{4a^{3}}$.

\Item{9.}  $\dfrac{2y^{2}z^{4}}{3}$.

\Item{10.}  $\dfrac{b}{c}$.

\Item{11.}  $\dfrac{2a - 3b}{2a}$.

\Item{12.}  $\dfrac{3a}{a + 2}$.

\Item{13.}  $\dfrac{x}{x - 1}$.

\Item{14.}  $\dfrac{y}{x^{2} + 3xy + 9y^{2}}$.

\Item{15.}  $\dfrac{x + 1}{x - 5}$.

\Item{16.}  $\dfrac{x + 1}{x - 2}$.

\Item{17.}  $\dfrac{a + b + c}{a}$.

\Item{18.}  $\dfrac{x + 5}{x + 3}$.

\Item{19.}  $\dfrac{x + 1}{x + 3}$.


\AnsTo[2]{Exercise}{43.} %  Page 91.

\Item{1.}  $a + b + \dfrac{2}{a - b}$.

\Item{2.}  $a - b - \dfrac{2}{a + b}$.

\Item{3.}  $a - 1 + \dfrac{2a}{a^{2} - a - 1}$.

\Item{4.}  $2x - 4 + \dfrac{5}{x + 1}$.

\Item{5.}  $4x^{2} - 2x + 1 - \dfrac{1}{2x + 1}$.

\Item{6.}  $5x + 4 + \dfrac{x + 7}{x^{2} + x - 1}$.

\Item{7.}  $a + \dfrac{5a - 2}{a^{2} + a + 2}$.

\Item{8.}  $y^{2} - yx + x^{2}$.

\Item{9.}  $x^{2} - 4x + 3 + \dfrac{2x - 4}{x^{2} + x + 1}$.

\Item{10.}  $x^{3} + x + 1 + \dfrac{2x + 2}{x^{2} - x - 1}$.
%% -----File: 184.png---Folio 178-------


\AnsTo[2]{Exercise}{44.} %  Page 92.

\Item{1.}  $\dfrac{x^{2} + y^{2}}{x - y}$.

\Item{2.}  $\dfrac{x^{2} + y^{2}}{x + y}$.

\Item{3.}  $\dfrac{2y}{x + y}$.

\Item{4.}  $-\dfrac{2ax}{a - x}$.

\Item{5.}  $-\dfrac{x + 2}{x - 3}$.

\Item{6.}  $-\dfrac{2x^{2} - 6x + 5}{x - 2}$.

\Item{7.}  $\dfrac{x^{3} + x^{2} - 2x + 1}{x + 2}$.

\Item{8.}  $\dfrac{2a^{2} - 11a + 6}{a - 3}$.

\Item{9.}  $\dfrac{-2a^{3} + a^{2} + 2a-3}{a - 1}$.

\Item{10.}  $\dfrac{3a^{4} + 6a^{3} - 14a^{2} - 4}{3a^{2} + 1}$.


\AnsTo[1]{Exercise}{45.} %  Page 94.

\Item{1.}  $\dfrac{x(x + a)}{(x + a)(x - a)}$, $\dfrac{x^{2}}{(x + a)(x - a)}$.

\Item{2.}  $\dfrac{a(a - b)}{(a + b)(a - b)}$, $\dfrac{a^{2}}{(a + b)(a - b)}$.

\Item{3.}  $\dfrac{1 - 2a}{(1 + 2a)(1 - 2a)}$, $\dfrac{1}{(1 + 2a)(1 - 2a)}$.

\Item{4.}  $\dfrac{9}{(4 + x)(4 - x)}$, $\dfrac{(4 - x)^{2}}{(4 + x)(4 - x)}$.

\Item{5.}  $\dfrac{a^{2}}{(3 - a)(9 + 3a + a^{2})}$, $\dfrac{a(9 + 3a + a^{2})}{(3 - a)(9 + 3a + a^{2})}$.

\Item{6.}  $\dfrac{x + 2}{(x + 2)(x - 2)(x - 3)}$, $\dfrac{x - 2}{(x + 2)(x - 2)(x - 3)}$.


\AnsTo{Exercise}{46.} %  Page 95.

\Item{1.}  $\dfrac{4x + 2}{5}$.

\Item{2.}  $\dfrac{13x + 3}{12}$.

\Item{3.}  $\dfrac{5(9x - 13)}{42}$.

\Item{4.}  $\dfrac{51x + 31}{36}$.

\Item{5.}  $\dfrac{x - 5}{3}$.

\Item{6.}  $\dfrac{5(x - y)}{8x}$.

\Item{7.}  $\dfrac{22x - 97}{30}$.

\Item{8.}  $\dfrac{3x - 4}{15x}$.

\Item{9.}  $\dfrac{a^{3} - b^{3} + c^{3} - abc}{abc}$.


\AnsTo{Exercise}{47.} %  Page 96}.

\Item{1.}  $\dfrac{2x + 1}{(x + 3)(x - 2)}$.

\Item{2.}  $\dfrac{2x}{x^{2} - 1}$.

\Item{3.}  $\dfrac{3x + 16}{(x - 8)(x + 2)}$.

\Item{4.}  $\dfrac{4ax}{a^{3} - x^{2}}$.

\Item{5.}  $\dfrac{ax}{x^{2} - a^{2}}$.

\Item{6.}  $-\dfrac{4ab}{4a^{2} - b^{2}}$.

\Item{7.}  $\dfrac{1}{9 - a^{2}}$.

\Item{8.}  $\dfrac{b}{a^{2} - b^{2}}$.

\Item{9.}  $\dfrac{5x + 8}{x^{2} - 4}$.

\Item{10.}  $\dfrac{1 + x}{1 - 9x^{2}}$.

\Item{11.}  $\dfrac{3(a^{2} + 4a + 1)}{a(a + 1)(a + 3)}$.

\Item{12.}  $\dfrac{2}{x^{2} - 1}$.

\Item{13.}  $\dfrac{2x^{2}}{(x + 2)(x - 3)}$.

\Item{14.}  $0$.
%% -----File: 185.png---Folio 179-------


\AnsTo{Exercise}{48.} %  Page 97.

\Item{1.}  $0$.

\Item{2.}  $\dfrac{2a}{a + b}$.

\Item{3.}  $\dfrac{7a}{1 - a^{2}}$.

\Item{4.}  $\dfrac{x - 10y}{4x^{2} - 25y^{2}}$.

\Item{5.}  $\dfrac{2}{x + 4y}$.

\Item{6.}  $\dfrac{2(x + 6)}{4x^{2} - 9}$.


\AnsTo{Exercise}{49.} %  Page 100.

\Item{1.}  $\dfrac{20}{3bc}$.

\Item{2.}  $\dfrac{2ay}{3}$.

\Item{3.}  $\dfrac{7p^{2}}{4xz}$.

\Item{4.}  $\dfrac{2a^{2}cm}{7}$.

\Item{5.}  $\dfrac{30}{abc}$.

\Item{6.}  $abc$.

\Item{7.}  $b^{2}$.

\Item{8.}  $\dfrac{x + a}{x - 2a}$.

\Item{9.}  $\dfrac{xy}{2c - 1}$.

\Item{10.}  $\dfrac{a + 10}{a + 3}$.

\Item{11.}  $\dfrac{3x + 2y}{x - 2}$.

\Item{12.}  $\dfrac{5a + b}{4a + 3b}$.

\Item{13.}  $\dfrac{x - 7}{a + b + c}$.

\Item{14.}  $\dfrac{x(x + 1)}{x - 5}$.

\Item{15.}  $\dfrac{a + 1}{a + 5}$.

\Item{16.}  $1$.

\Item{17.}  $\dfrac{x(x + y)}{x + 1}$.

\Item{18.}  $\dfrac{b}{a - b}$.

\Item{19.}  $abc$.

\Item{20.}  $\dfrac{(x + 2y)(x + 1)}{(x + y)(x + 2)}$.


\AnsTo{Exercise}{50.} %  Page 102.

\Item{1.}  $\dfrac{x + y}{z}$.

\Item{2.}  $\dfrac{12x + 3y}{12x - 4y}$.

\Item{3.}  $\dfrac{abd - 21d^{2}}{21cd - 7ab}$.

\Item{4.}  $\dfrac{x^{2} + x - 2}{x^{2} - x - 2}$.

\Item{5.}  $1$.

\Item{6.}  $\dfrac{x + y}{x^{2} - 2xy + y^{2}}$.

\Item{7.}  $\dfrac{a + b}{a - b}$.

\Item{8.}  $4(3a + 8)$.

\Item{9.}  $\dfrac{y + x}{y - x}$.

\Item{10.}  $x$.

\Item{11.}  $\dfrac{1}{x}$.

\Item{12.}  $\dfrac{x^{2}(x - 3)}{x - 2}$.

\Item{13.}  $a - 1$.

\Item{14.}  $\dfrac{4a}{a - x}$.


\AnsTo[5]{Exercise}{51.} %  Page 105.

\Item{1.}  $5$.

\Item{2.}  $7$.

\Item{3.}  $2\frac{1}{2}$.

\Item{4.}  $120$.

\Item{5.}  $12$.

\Item{6.}  $2\frac{1}{3}$.

\Item{7.}  $17$.

\Item{8.}  $4$.

\Item{9.}  $4$.

\Item{10.}  $1$.

\Item{11.}  $-16$.

\Item{12.}  $11$.

\Item{13.}  $-4$.

\Item{14.}  $-2$.

\Item{15.}  $-2$.

\Item{16.}  $5$.

\Item{17.}  $9$.

\Item{18.}  $-1$.
%% -----File: 186.png---Folio 180-------


\AnsTo[5]{Exercise}{52.} %  Page 106.

\Item{1.}  $2$.

\Item{2.}  $2$.

\Item{3.}  $-33$.

\Item{4.}  $1$.

\Item{5.}  $\frac{2}{3}$.

\Item{6.}  $1\frac{1}{2}$.

\Item{7.}  $2$.

\Item{8.}  $8$.

\Item{9.}  $5$.

\Item{10.}  $\frac{3}{7}$.

\Item{11.}  $2$.

\Item{12.}  $1$.

\Item{13.}  $3$.


\AnsTo[6]{Exercise}{53.} %  Page 107.

\Item{1.}  $33$.

\Item{2.}  $2$.

\Item{3.}  $3\frac{1}{2}$.

\Item{4.}  $1\frac{5}{37}$.

\Item{5.}  $7$.

\Item{6.}  $3$.


\AnsTo{Exercise}{54.} %  Page 108.

\Item{1.}  $a + b$.

\Item{2.}  $\dfrac{a}{2}$.

\Item{3.}  $\dfrac{b}{2}$.

\Item{4.}  $2a$.

\Item{5.}  $b - a$.

\Item{6.}  $\dfrac{ab}{a + b + c}$.

\Item{7.}  $\dfrac{a^{2} - b^{2}}{2a}$.

\Item{8.}  $\dfrac{2b}{a}$.

\Item{9.}  $\dfrac{2b^{2} - a^{2}}{4b - 3a}$.

\Item{10.}  $1$.

\Item{11.}  $3(a - b)$.


\AnsTo[4]{Exercise}{55.} %  Page 109.

\Item{1.}  $35$.

\Item{2.}  $70$.

\Item{3.}  $36$.

\Item{4.}  $57$, $58$.


\AnsTo[5]{Exercise}{56.} %  Page 110.

\Item{1.}  $81$, $19$.

\Item{2.}  $100$, $24$.

\Item{3.}  $64$, $15$.

\Item{4.}  $103$, $12$.

\Item{5.}  $295$, $25$.


\AnsTo[4]{Exercise}{57.} %  Page 111.

\Item{1.}  $12$~yr.

\Item{2.}  A, $60$~yr.; B, $10$~yr.

\Item{3.}  A, $25$~yr.; B, $5$~yr.

\Item{4.}  $17\frac{1}{2}$~yr.

\Item{5.}  $35$~yr.

\Item{6.}  Son, $12$~yr.; father, $36$~yr.

\Item{7.}  $25$~yr.

\Item{8.}  A, $30$~yr.; B, $15$~yr.

\Item{9.}  Son, $12$~yr.; father, $68$~yr.


\AnsTo[6]{Exercise}{58.} %  Page 112.

\Item{1.}  $1\frac{3}{7}$~dy.

\Item{2.}  $1\frac{13}{47}$~dy.

\Item{3.}  $1\frac{1}{20}$~dy.

\Item{4.}  $15$~dy.

\Item{5.}  $12$~hr.

\Item{6.}  $10$~dy.


\AnsTo[5]{Exercise}{59.} %  Page 113.

\Item{1.}  $7\frac{5}{13}$~hr.

\Item{2.}  $2\frac{2}{5}$~hr.

\Item{3.}  $\frac{10}{11}$~hr.

\Item{4.}  $1\frac{1}{13}$~hr.

\Item{5.}  $30$~hr.


\AnsTo[4]{Exercise}{60.} %  Page 114.

\Item{1.}  $36$~mi.

\Item{2.}  $26$~hr.

\Item{3.}  $8$~mi.

\Item{4.}  $240$~mi.


\AnsTo{Exercise}{61.} %  Page 115.

\Item{1.}  $600$.

\Item{2.}  $700$.

\Item{3.}  Dog, $1440$; rabbit, $1800$.
%% -----File: 187.png---Folio 181-------

\AnsTo[2]{Exercise}{62.} %  Page 116.

\Item{1.}  $27\frac{3}{11}$~min.\ past 5~o'clock.

\Item{2.}  $27\frac{3}{11}$~min.\ past 2~o'clock.

\Item{3.}  $43\frac{7}{11}$~min.\ past 2~o'clock.

\Item{4.}  $21\frac{9}{11}$~min.\ past 1~o'clock.

\Item{5.}  $38\frac{2}{11}$~min.\ past 1~o'clock.

\Item{6.}  $38\frac{2}{11}$~min.\ past 7~o'clock.



\AnsTo{Exercise}{63.} %  Page 117.

\Item{1.}  $1764$~sq.~ft.

\Item{2.}  $18$~ft.\ by $23$~ft.

\Item{3.}  $14$~ft.\ by $20$~ft.

\Item{4.}  $12$~ft.\ by $15$~ft.

\Item{5.}  $30$~ft.\ by $40$~ft.


\AnsTo[4]{Exercise}{64.} %  Page 121.

\Item{1.}  $90°\,0'\,30''$; $30°\,30'$.

\Item{2.}  \$$133\frac{1}{3}$.

\Item{3.}  \$$2000$.

\Item{4.}  \$$4000$.

\Item{5.}  \$$3000$.

\Item{6.}  \$$500$.

\Item{7.}  $5$\%.

\Item{8.}  $6$\%.

\Item{9.}  $3$~yr.

\Item{10.}  $9\frac{3}{8}$~yr.

\Item{11.}  \$$25,000$.

\Item{12.}  \$$20,000$.

\Item{13.}  \$$6$\%.

\Item{14.}  $20$~yr.


\AnsTo{Exercise}{65.} %  Page 124.

\Item{1.}  $x = 2$, $y = 1$.

\Item{2.}  $x = 3$, $y = 2$.

\Item{3.}  $x = 5$, $y = 1$.

\Item{4.}  $x = 2$, $y = 1$.

\Item{5.}  $x = 1$, $y = 2$.

\Item{6.}  $x = 6$, $y = 1$.

\Item{7.}  $x = 3$, $y = 21$.

\Item{8.}  $x = 7$, $y = 7$.

\Item{9.}  $x = 23$, $y = -1$.

\Item{10.}  $x = 2$, $y = 1$.

\Item{11.}  $x = 35$, $y = 20$.

\Item{12.}  $x = 2$, $y = 1$.

\Item{13.}  $x = 1$, $y = 2$.

\Item{14.}  $x = 3$, $y = 2$.

\Item{15.}  $x = 1$, $y = 2$.

\Item{16.}  $x = 4$, $y = 3$.

\Item{17.}  $x = 12$, $y = 4$.

\Item{18.}  $x = 12$, $y = 21$.

\Item{19.}  $x = 5$, $y = 7$.

\Item{20.}  $x = 5$, $y = 2$.

\Item{21.}  $x = 18$, $y = 6$.

\Item{22.}  $x = 3$, $y = 2$.

\Item{23.}  $x = 3$, $y = 2$.

\Item{24.}  $x = 7$, $y = 8$.

\Item{25.}  $x = 8$, $y = -2$.

\Item{26.}  $x = \dfrac{a}{(a - b)}$, $y = \dfrac{b}{(a + b)}$.


\AnsTo[2]{Exercise}{66.} %  Page 127.

\Item{1.}  A, \$$520$; B, \$$440$.

\Item{2.}  $23$~and~$17$.

\Item{3.}  $20$~and~$16$.

\Item{4.}  Velvet, \$$6$; silk, \$$3$.

\Item{5.}  Wheat, \$$1$; rye, \$$\frac{4}{5}$.

\Item{6.}  Tea, \$$\frac{1}{2}$; coffee, \$$\frac{1}{4}$.

\Item{7.}  Horse, \$$92$; cow, \$$64$.
%% -----File: 188.png---Folio 182-------


\AnsTo[5]{Exercise}{67.} %  Page 128.

\Item{1.}  $\frac{3}{7}$.

\Item{2.}  $\frac{13}{25}$.

\Item{3.}  $\frac{3}{20}$.

\Item{4.}  $\frac{5}{21}$.

\Item{5.}  $\frac{7}{22}$.


\AnsTo[4]{Exercise}{68.} %  Page 129.

\Item{1.}  $45$.

\Item{2.}  $72$.

\Item{3.}  $75$~and~$57$.

\Item{4.}  $54$.


\AnsTo{Exercise}{69.} %  Page 130.

\Item{1.}  \$$2500$ at $4$\%.

\Item{2.}  \$$1600$ at $6$\%.

\Item{3.}  \$$6000$ at $4$\%; \$$4000$ at $5$\%.


\AnsTo[4]{Exercise}{70.} %  Page 131.

\Item{1.}  $22$~and~$18$.

\Item{2.}  $60$~and~$8$.

\Item{3.}  $\frac{14}{17}$.

\Item{4.}  Wheat, \$$1$; barley, \$$\frac{4}{5}$.

\Item{5.}  A, \$$235$; B, \$$65$.

\Item{6.}  A, \$$70$; B, \$$30$.

\Item{7.}  Lemon, $2$~cts.; orange, $3$~cts.

\Item{8.}  A, $30$~apples; B, $10$~apples.


\AnsTo[4]{Exercise}{71.} %  Page 134.

\Item{1.}  $±2$.

\Item{2.}  $±3$.

\Item{3.}  $±5$.

\Item{4.}  $±8$.

\Item{5.}  $±7$.

\Item{6.}  $±5$.

\Item{7.}  $±5$.

\Item{8.}  $±3$.

\Item{9.}  $±3$.

\Item{10.}  $±3$.

\Item{11.}  $12$~and~$16$.

\Item{12.}  $12$~oranges at $3$~cts.

\Item{13.}  $3$~rods.

\Item{14.}  Width, $12$~rd.; length, $48$~rd.


\AnsTo{Exercise}{72.} %  Page 137.

\Item{1.}  $9$~or~$3$.

\Item{2.}  $4$~or~$2$.

\Item{3.}  $3$~or~$1$.

\Item{4.}  $1$~or~$-\frac{1}{5}$.

\Item{5.}  $1$~or~$-3$.

\Item{6.}  $\frac{4}{3}$~or~$\frac{4}{3}$.

\Item{7.}  $1$~or~$-\frac{1}{6}$.

\Item{8.}  $3$~or~$-1$.

\Item{9.}  $\frac{3}{4}$~or~$\frac{1}{4}$.

\Item{10.}  $3$~or~$\frac{1}{3}$.

\Item{11.}  $17$~or~$-3$.

\Item{12.}  $25$~or~$9$.

\Item{13.}  $4$~or~$-5$.

\Item{14.}  $4$~or~$-3$.

\Item{15.}  $5$~or~$1$.

\Item{16.}  $2$~or~$-6$.

\Item{17.}  $2$~or~$2$.

\Item{18.}  $5$~or~$-11$.

\Item{19.}  $2$~or~$-5\frac{1}{3}$.

\Item{20.}  $4\frac{1}{3}$~or~$-3\frac{2}{3}$.

\Item{21.}  $2$~or~$\frac{1}{3}$.

\Item{22.}  $4$~or~$-\frac{2}{5}$.

\Item{23.}  $2$~or~$-3$.

\Item{24.}  $10$~or~$2$.

\Item{25.}  $3$~or~$-2\frac{1}{3}$.

\Item{26.}  $2$~or~$2$.

\Item{27.}  $\frac{1}{2}$~or~$-3$.

\Item{28.}  $5$~or~$\frac{1}{2}$.

\Item{29.}  $7$~or~$2$.

\Item{30.}  $4$~or~$-\frac{2}{3}$.

\Item{31.}  $8$~or~$2$.

\Item{32.}  $4$~or~$-7$.

\Item{33.}  $0$~or~$3$.

\Item{34.}  $0$~or~$7$.

\Item{35.}  $5$~or~$2$.

\Item{36.}  $4$~or~$-1$.
%% -----File: 189.png---Folio 183-------


\AnsTo{Exercise}{73.} %  Page 140.

\Item{1.}  $6$~and~$5$.

\Item{2.}  $5$~and~$15$.

\Item{3.}  Son, $8$~yr.; father, $40$~yr.

\Item{4.}  $9$.

\Item{6.}  $5$~rd.\ by $7$~rd.

\Item{7.}  $12$~ft.

\Item{8.}  $20$~ft.\ by $18$~ft.

\Item{9.}  $10$~rd.\ by $12$~rd.

\Item{10.}  Son, $10$~yr.; father, $54$~yr.


\AnsTo{Exercise}{74.} %  Page 141.

\Item{1.}  $6$~miles an hour.

\Item{2.}  $7$.

\Item{3.}  $5$.

\Item{4.}  $8$.

\Item{5.}  $36$.


\AnsTo[4]{Exercise}{75.} %  Page 144.

\Item{1.}  $75$.

\Item{2.}  $38$.

\Item{3.}  $4\frac{1}{7}$.

\Item{4.}  $-8\frac{3}{4}$.

\Item{5.}  $23$.

\Item{6.}  $0$.

\Item{7.}  $156$.

\Item{8.}  $20$th.

\Item{9.}  $101$st.

\Item{10.}  $26$.

\Item{11.}  $a$.

\Item{12.}  $21$, $22$,~etc.


\AnsTo[4]{Exercise}{76.} %  Page 146.

\Item{1.}  $440$.

\Item{2.}  $201$.

\Item{3.}  $4frac{1}{6}$.

\Item{4.}  $128$.

\Item{5.}  $-378$.

\Item{6.}  $187\frac{1}{2}$.

\Item{7.}  $1$, $3$, $5$.

\Item{8.}  $156$.

\Item{9.}  $300$.

\Item{10.}  $2550$~yd.

\Item{11.}  $5812.1$~ft.

\Item{12.}  $144.9$~ft.

\Item{13.}  $11$, $15$.

\Item{14.}  $7$, $9$, $11$.

\Item{15.}  $12$~miles.

\Item{16.}  $5$, $12$, $19$.


\AnsTo[4]{Exercise}{77.} %  Page 151.

\Item{1.}  $243$.

\Item{2.}  $192$.

\Item{3.}  $\frac{3}{64}$.

\Item{4.}  $256$.

\Item{5.}  $±4$.

\Item{6.}  $±4$.

\Item{7.}  $1092$.

\Item{8.}  $765$.

\Item{9.}  $11\frac{29}{32}$.

\Item{10.}  $15\frac{15}{16}$.

\Item{11.}  $127\frac{3}{4}$.

\Item{12.}  $44$.

\Item{13.}  $1\frac{11}{54}$.

\Item{14.}  \$$1.27$.

\Item{15.}  \$$81.90$.

\Item{16.}  $14,641$.


\AnsTo{Exercise}{78.} %  Page 153.

\Item{1.}  $a + b + c$.

\Item{2.}  $x^{2} + x + 1$.

\Item{3.}  $x^{2} - 2xy + y^{2}$.

\Item{4.}  $2a^{2} - 3ab + 5b^{2}$.

\Item{5.}  $4x^{3} + 3x^{2}y - 2y^{3}$.

\Item{6.}  $2x^{3} - xy^{2} + 3y^{3}$.


\AnsTo[4]{Exercise}{79.} %  Page 156.

\Item{1.}  $18$.

\Item{2.}  $21$.

\Item{3.}  $23$.

\Item{4.}  $31$.

\Item{5.}  $3.2$.

\Item{6.}  $7.3$.

\Item{7.}  $232$.

\Item{8.}  $785$.

\Item{9.}  $1225$.

\Item{10.}  $589$.

\Item{11.}  $5601$.

\Item{12.}  $1234$.

\Item{13.}  $1.4142\dots$

\Item{14.}  $1.7320\dots$

\Item{15.}  $2.2360\dots$

\Item{16.}  $2.4494\dots$

\Item{17.}  $0.7071\dots$

\Item{18.}  $0.9486\dots$

\Item{19.}  $0.8164\dots$

\Item{20.}  $0.8660\dots$

\Item{21.}  $0.8944\dots$

\Item{22.}  $0.7905\dots$
%% -----File: 190.png---Folio 184-------


\AnsTo{Exercise}{80.} %  Page 159.

\Item{1.}  $x + y$.

\Item{2.}  $2x - 1$.

\Item{3.}  $2x - 3y$.

\Item{4.}  $4a - 3x$.

\Item{5.}  $1 + x + x^{2}$.

\Item{6.}  $x^{2} - x + 1$.


\AnsTo[4]{Exercise}{81.} %  Page 163.

\Item{1.}  $36$.

\Item{2.}  $35$.

\Item{3.}  $45$.

\Item{4.}  $65$.

\Item{5.}  $48$.

\Item{6.}  $637$.

\Item{7.}  $478$.

\Item{8.}  $638$.

\Item{9.}  $503$.

\Item{10.}  $728$.

\Item{11.}  $12.34$.

\Item{12.}  $12.25$.

\Item{13.}  $0.2154\dots$

\Item{14.}  $0.3684\dots$

\Item{15.}  $0.5848\dots$

\Item{16.}  $1.5874\dots$

\Item{17.}  $2.1544\dots$

\Item{18.}  $4.4310\dots$

\Item{19.}  $1.3572\dots$

\Item{20.}  $1.2703\dots$

\Item{21.}  $1.4454\dots$

\Item{22.}  $0.8735\dots$

\Item{23.}  $0.9085\dots$

\Item{24.}  $0.9352\dots$
%[** TN: Environment opened by \AnsTo]
\end{multicols}


%%%%%%%%%%%%%%%%%%%%%%%%% GUTENBERG LICENSE %%%%%%%%%%%%%%%%%%%%%%%%%%

\BackMatter
\begin{PGtext}
End of the Project Gutenberg EBook of The First Steps in Algebra, by 
G. A. (George Albert) Wentworth

*** END OF THIS PROJECT GUTENBERG EBOOK THE FIRST STEPS IN ALGEBRA ***

***** This file should be named 36670-t.tex or 36670-t.zip *****
This and all associated files of various formats will be found in:
        http://www.gutenberg.org/3/6/6/7/36670/

Produced by Andrew D. Hwang, Peter Vachuska, Chuck Greif
and the Online Distributed Proofreading Team at
http://www.pgdp.net.


Updated editions will replace the previous one--the old editions
will be renamed.

Creating the works from public domain print editions means that no
one owns a United States copyright in these works, so the Foundation
(and you!) can copy and distribute it in the United States without
permission and without paying copyright royalties.  Special rules,
set forth in the General Terms of Use part of this license, apply to
copying and distributing Project Gutenberg-tm electronic works to
protect the PROJECT GUTENBERG-tm concept and trademark.  Project
Gutenberg is a registered trademark, and may not be used if you
charge for the eBooks, unless you receive specific permission.  If you
do not charge anything for copies of this eBook, complying with the
rules is very easy.  You may use this eBook for nearly any purpose
such as creation of derivative works, reports, performances and
research.  They may be modified and printed and given away--you may do
practically ANYTHING with public domain eBooks.  Redistribution is
subject to the trademark license, especially commercial
redistribution.



*** START: FULL LICENSE ***

THE FULL PROJECT GUTENBERG LICENSE
PLEASE READ THIS BEFORE YOU DISTRIBUTE OR USE THIS WORK

To protect the Project Gutenberg-tm mission of promoting the free
distribution of electronic works, by using or distributing this work
(or any other work associated in any way with the phrase "Project
Gutenberg"), you agree to comply with all the terms of the Full Project
Gutenberg-tm License (available with this file or online at
http://gutenberg.net/license).


Section 1.  General Terms of Use and Redistributing Project Gutenberg-tm
electronic works

1.A.  By reading or using any part of this Project Gutenberg-tm
electronic work, you indicate that you have read, understand, agree to
and accept all the terms of this license and intellectual property
(trademark/copyright) agreement.  If you do not agree to abide by all
the terms of this agreement, you must cease using and return or destroy
all copies of Project Gutenberg-tm electronic works in your possession.
If you paid a fee for obtaining a copy of or access to a Project
Gutenberg-tm electronic work and you do not agree to be bound by the
terms of this agreement, you may obtain a refund from the person or
entity to whom you paid the fee as set forth in paragraph 1.E.8.

1.B.  "Project Gutenberg" is a registered trademark.  It may only be
used on or associated in any way with an electronic work by people who
agree to be bound by the terms of this agreement.  There are a few
things that you can do with most Project Gutenberg-tm electronic works
even without complying with the full terms of this agreement.  See
paragraph 1.C below.  There are a lot of things you can do with Project
Gutenberg-tm electronic works if you follow the terms of this agreement
and help preserve free future access to Project Gutenberg-tm electronic
works.  See paragraph 1.E below.

1.C.  The Project Gutenberg Literary Archive Foundation ("the Foundation"
or PGLAF), owns a compilation copyright in the collection of Project
Gutenberg-tm electronic works.  Nearly all the individual works in the
collection are in the public domain in the United States.  If an
individual work is in the public domain in the United States and you are
located in the United States, we do not claim a right to prevent you from
copying, distributing, performing, displaying or creating derivative
works based on the work as long as all references to Project Gutenberg
are removed.  Of course, we hope that you will support the Project
Gutenberg-tm mission of promoting free access to electronic works by
freely sharing Project Gutenberg-tm works in compliance with the terms of
this agreement for keeping the Project Gutenberg-tm name associated with
the work.  You can easily comply with the terms of this agreement by
keeping this work in the same format with its attached full Project
Gutenberg-tm License when you share it without charge with others.

1.D.  The copyright laws of the place where you are located also govern
what you can do with this work.  Copyright laws in most countries are in
a constant state of change.  If you are outside the United States, check
the laws of your country in addition to the terms of this agreement
before downloading, copying, displaying, performing, distributing or
creating derivative works based on this work or any other Project
Gutenberg-tm work.  The Foundation makes no representations concerning
the copyright status of any work in any country outside the United
States.

1.E.  Unless you have removed all references to Project Gutenberg:

1.E.1.  The following sentence, with active links to, or other immediate
access to, the full Project Gutenberg-tm License must appear prominently
whenever any copy of a Project Gutenberg-tm work (any work on which the
phrase "Project Gutenberg" appears, or with which the phrase "Project
Gutenberg" is associated) is accessed, displayed, performed, viewed,
copied or distributed:

This eBook is for the use of anyone anywhere at no cost and with
almost no restrictions whatsoever.  You may copy it, give it away or
re-use it under the terms of the Project Gutenberg License included
with this eBook or online at www.gutenberg.net

1.E.2.  If an individual Project Gutenberg-tm electronic work is derived
from the public domain (does not contain a notice indicating that it is
posted with permission of the copyright holder), the work can be copied
and distributed to anyone in the United States without paying any fees
or charges.  If you are redistributing or providing access to a work
with the phrase "Project Gutenberg" associated with or appearing on the
work, you must comply either with the requirements of paragraphs 1.E.1
through 1.E.7 or obtain permission for the use of the work and the
Project Gutenberg-tm trademark as set forth in paragraphs 1.E.8 or
1.E.9.

1.E.3.  If an individual Project Gutenberg-tm electronic work is posted
with the permission of the copyright holder, your use and distribution
must comply with both paragraphs 1.E.1 through 1.E.7 and any additional
terms imposed by the copyright holder.  Additional terms will be linked
to the Project Gutenberg-tm License for all works posted with the
permission of the copyright holder found at the beginning of this work.

1.E.4.  Do not unlink or detach or remove the full Project Gutenberg-tm
License terms from this work, or any files containing a part of this
work or any other work associated with Project Gutenberg-tm.

1.E.5.  Do not copy, display, perform, distribute or redistribute this
electronic work, or any part of this electronic work, without
prominently displaying the sentence set forth in paragraph 1.E.1 with
active links or immediate access to the full terms of the Project
Gutenberg-tm License.

1.E.6.  You may convert to and distribute this work in any binary,
compressed, marked up, nonproprietary or proprietary form, including any
word processing or hypertext form.  However, if you provide access to or
distribute copies of a Project Gutenberg-tm work in a format other than
"Plain Vanilla ASCII" or other format used in the official version
posted on the official Project Gutenberg-tm web site (www.gutenberg.net),
you must, at no additional cost, fee or expense to the user, provide a
copy, a means of exporting a copy, or a means of obtaining a copy upon
request, of the work in its original "Plain Vanilla ASCII" or other
form.  Any alternate format must include the full Project Gutenberg-tm
License as specified in paragraph 1.E.1.

1.E.7.  Do not charge a fee for access to, viewing, displaying,
performing, copying or distributing any Project Gutenberg-tm works
unless you comply with paragraph 1.E.8 or 1.E.9.

1.E.8.  You may charge a reasonable fee for copies of or providing
access to or distributing Project Gutenberg-tm electronic works provided
that

- You pay a royalty fee of 20% of the gross profits you derive from
     the use of Project Gutenberg-tm works calculated using the method
     you already use to calculate your applicable taxes.  The fee is
     owed to the owner of the Project Gutenberg-tm trademark, but he
     has agreed to donate royalties under this paragraph to the
     Project Gutenberg Literary Archive Foundation.  Royalty payments
     must be paid within 60 days following each date on which you
     prepare (or are legally required to prepare) your periodic tax
     returns.  Royalty payments should be clearly marked as such and
     sent to the Project Gutenberg Literary Archive Foundation at the
     address specified in Section 4, "Information about donations to
     the Project Gutenberg Literary Archive Foundation."

- You provide a full refund of any money paid by a user who notifies
     you in writing (or by e-mail) within 30 days of receipt that s/he
     does not agree to the terms of the full Project Gutenberg-tm
     License.  You must require such a user to return or
     destroy all copies of the works possessed in a physical medium
     and discontinue all use of and all access to other copies of
     Project Gutenberg-tm works.

- You provide, in accordance with paragraph 1.F.3, a full refund of any
     money paid for a work or a replacement copy, if a defect in the
     electronic work is discovered and reported to you within 90 days
     of receipt of the work.

- You comply with all other terms of this agreement for free
     distribution of Project Gutenberg-tm works.

1.E.9.  If you wish to charge a fee or distribute a Project Gutenberg-tm
electronic work or group of works on different terms than are set
forth in this agreement, you must obtain permission in writing from
both the Project Gutenberg Literary Archive Foundation and Michael
Hart, the owner of the Project Gutenberg-tm trademark.  Contact the
Foundation as set forth in Section 3 below.

1.F.

1.F.1.  Project Gutenberg volunteers and employees expend considerable
effort to identify, do copyright research on, transcribe and proofread
public domain works in creating the Project Gutenberg-tm
collection.  Despite these efforts, Project Gutenberg-tm electronic
works, and the medium on which they may be stored, may contain
"Defects," such as, but not limited to, incomplete, inaccurate or
corrupt data, transcription errors, a copyright or other intellectual
property infringement, a defective or damaged disk or other medium, a
computer virus, or computer codes that damage or cannot be read by
your equipment.

1.F.2.  LIMITED WARRANTY, DISCLAIMER OF DAMAGES - Except for the "Right
of Replacement or Refund" described in paragraph 1.F.3, the Project
Gutenberg Literary Archive Foundation, the owner of the Project
Gutenberg-tm trademark, and any other party distributing a Project
Gutenberg-tm electronic work under this agreement, disclaim all
liability to you for damages, costs and expenses, including legal
fees.  YOU AGREE THAT YOU HAVE NO REMEDIES FOR NEGLIGENCE, STRICT
LIABILITY, BREACH OF WARRANTY OR BREACH OF CONTRACT EXCEPT THOSE
PROVIDED IN PARAGRAPH 1.F.3.  YOU AGREE THAT THE FOUNDATION, THE
TRADEMARK OWNER, AND ANY DISTRIBUTOR UNDER THIS AGREEMENT WILL NOT BE
LIABLE TO YOU FOR ACTUAL, DIRECT, INDIRECT, CONSEQUENTIAL, PUNITIVE OR
INCIDENTAL DAMAGES EVEN IF YOU GIVE NOTICE OF THE POSSIBILITY OF SUCH
DAMAGE.

1.F.3.  LIMITED RIGHT OF REPLACEMENT OR REFUND - If you discover a
defect in this electronic work within 90 days of receiving it, you can
receive a refund of the money (if any) you paid for it by sending a
written explanation to the person you received the work from.  If you
received the work on a physical medium, you must return the medium with
your written explanation.  The person or entity that provided you with
the defective work may elect to provide a replacement copy in lieu of a
refund.  If you received the work electronically, the person or entity
providing it to you may choose to give you a second opportunity to
receive the work electronically in lieu of a refund.  If the second copy
is also defective, you may demand a refund in writing without further
opportunities to fix the problem.

1.F.4.  Except for the limited right of replacement or refund set forth
in paragraph 1.F.3, this work is provided to you 'AS-IS' WITH NO OTHER
WARRANTIES OF ANY KIND, EXPRESS OR IMPLIED, INCLUDING BUT NOT LIMITED TO
WARRANTIES OF MERCHANTIBILITY OR FITNESS FOR ANY PURPOSE.

1.F.5.  Some states do not allow disclaimers of certain implied
warranties or the exclusion or limitation of certain types of damages.
If any disclaimer or limitation set forth in this agreement violates the
law of the state applicable to this agreement, the agreement shall be
interpreted to make the maximum disclaimer or limitation permitted by
the applicable state law.  The invalidity or unenforceability of any
provision of this agreement shall not void the remaining provisions.

1.F.6.  INDEMNITY - You agree to indemnify and hold the Foundation, the
trademark owner, any agent or employee of the Foundation, anyone
providing copies of Project Gutenberg-tm electronic works in accordance
with this agreement, and any volunteers associated with the production,
promotion and distribution of Project Gutenberg-tm electronic works,
harmless from all liability, costs and expenses, including legal fees,
that arise directly or indirectly from any of the following which you do
or cause to occur: (a) distribution of this or any Project Gutenberg-tm
work, (b) alteration, modification, or additions or deletions to any
Project Gutenberg-tm work, and (c) any Defect you cause.


Section  2.  Information about the Mission of Project Gutenberg-tm

Project Gutenberg-tm is synonymous with the free distribution of
electronic works in formats readable by the widest variety of computers
including obsolete, old, middle-aged and new computers.  It exists
because of the efforts of hundreds of volunteers and donations from
people in all walks of life.

Volunteers and financial support to provide volunteers with the
assistance they need are critical to reaching Project Gutenberg-tm's
goals and ensuring that the Project Gutenberg-tm collection will
remain freely available for generations to come.  In 2001, the Project
Gutenberg Literary Archive Foundation was created to provide a secure
and permanent future for Project Gutenberg-tm and future generations.
To learn more about the Project Gutenberg Literary Archive Foundation
and how your efforts and donations can help, see Sections 3 and 4
and the Foundation web page at http://www.pglaf.org.


Section 3.  Information about the Project Gutenberg Literary Archive
Foundation

The Project Gutenberg Literary Archive Foundation is a non profit
501(c)(3) educational corporation organized under the laws of the
state of Mississippi and granted tax exempt status by the Internal
Revenue Service.  The Foundation's EIN or federal tax identification
number is 64-6221541.  Its 501(c)(3) letter is posted at
http://pglaf.org/fundraising.  Contributions to the Project Gutenberg
Literary Archive Foundation are tax deductible to the full extent
permitted by U.S. federal laws and your state's laws.

The Foundation's principal office is located at 4557 Melan Dr. S.
Fairbanks, AK, 99712., but its volunteers and employees are scattered
throughout numerous locations.  Its business office is located at
809 North 1500 West, Salt Lake City, UT 84116, (801) 596-1887, email
business@pglaf.org.  Email contact links and up to date contact
information can be found at the Foundation's web site and official
page at http://pglaf.org

For additional contact information:
     Dr. Gregory B. Newby
     Chief Executive and Director
     gbnewby@pglaf.org


Section 4.  Information about Donations to the Project Gutenberg
Literary Archive Foundation

Project Gutenberg-tm depends upon and cannot survive without wide
spread public support and donations to carry out its mission of
increasing the number of public domain and licensed works that can be
freely distributed in machine readable form accessible by the widest
array of equipment including outdated equipment.  Many small donations
($1 to $5,000) are particularly important to maintaining tax exempt
status with the IRS.

The Foundation is committed to complying with the laws regulating
charities and charitable donations in all 50 states of the United
States.  Compliance requirements are not uniform and it takes a
considerable effort, much paperwork and many fees to meet and keep up
with these requirements.  We do not solicit donations in locations
where we have not received written confirmation of compliance.  To
SEND DONATIONS or determine the status of compliance for any
particular state visit http://pglaf.org

While we cannot and do not solicit contributions from states where we
have not met the solicitation requirements, we know of no prohibition
against accepting unsolicited donations from donors in such states who
approach us with offers to donate.

International donations are gratefully accepted, but we cannot make
any statements concerning tax treatment of donations received from
outside the United States.  U.S. laws alone swamp our small staff.

Please check the Project Gutenberg Web pages for current donation
methods and addresses.  Donations are accepted in a number of other
ways including including checks, online payments and credit card
donations.  To donate, please visit: http://pglaf.org/donate


Section 5.  General Information About Project Gutenberg-tm electronic
works.

Professor Michael S. Hart is the originator of the Project Gutenberg-tm
concept of a library of electronic works that could be freely shared
with anyone.  For thirty years, he produced and distributed Project
Gutenberg-tm eBooks with only a loose network of volunteer support.


Project Gutenberg-tm eBooks are often created from several printed
editions, all of which are confirmed as Public Domain in the U.S.
unless a copyright notice is included.  Thus, we do not necessarily
keep eBooks in compliance with any particular paper edition.


Most people start at our Web site which has the main PG search facility:

     http://www.gutenberg.net

This Web site includes information about Project Gutenberg-tm,
including how to make donations to the Project Gutenberg Literary
Archive Foundation, how to help produce our new eBooks, and how to
subscribe to our email newsletter to hear about new eBooks.
\end{PGtext}

% %%%%%%%%%%%%%%%%%%%%%%%%%%%%%%%%%%%%%%%%%%%%%%%%%%%%%%%%%%%%%%%%%%%%%%% %
%                                                                         %
% End of the Project Gutenberg EBook of The First Steps in Algebra, by    %
% G. A. (George Albert) Wentworth                                         %
%                                                                         %
% *** END OF THIS PROJECT GUTENBERG EBOOK THE FIRST STEPS IN ALGEBRA ***  %
%                                                                         %
% ***** This file should be named 36670-t.tex or 36670-t.zip *****        %
% This and all associated files of various formats will be found in:      %
%         http://www.gutenberg.org/3/6/6/7/36670/                         %
%                                                                         %
% %%%%%%%%%%%%%%%%%%%%%%%%%%%%%%%%%%%%%%%%%%%%%%%%%%%%%%%%%%%%%%%%%%%%%%% %

\end{document}
###
@ControlwordReplace = (
  ['\\Answers', 'Answers.'],
  ['\\Preface', 'Preface.'],
  ['\\end{Remark}', ''],
  ['\\end{Soln}', ''],
  ['\\end{Theorem}', ''],
  ['\\HCF', 'H.C.F.'],
  ['\\LCD', 'L.C.D.'],
  ['\\LCM', 'L.C.M.'],
  ['\\end{multicols}', '']
  );

@MathEnvironments = (
  ['\\begin{DPalign*}','\\end{DPalign*}','<DPALIGN>'],
  ['\\begin{DPgather*}','\\end{DPgather*}','<DPGATHER>']
  );

@ControlwordArguments = (
  ['\\BookMark', 1, 0, '', '', 1, 0, '', ''],
  ['\\Chapter', 1, 1, 'Chapter ', '  ', 1, 1, '', ''],
  ['\\Section', 1, 1, '', ''],
  ['\\Exercise', 0, 0, '', '', 1, 1, 'Exercise ', ''],
  ['\\AnsTo', 0, 0, '', '', 1, 1, '', ' ', 1, 1, '', ''],
  ['\\Paragraph', 1, 1, '', ''],
  ['\\Signature', 1, 1, '', '  ', 1, 1, '', ''],
  ['\\begin{Remark}', 0, 1, '', ''],
  ['\\begin{Soln}', 0, 1, '', ''],
  ['\\begin{Theorem}', 0, 1, '', ''],
  ['\\Graphic', 1, 0, '<GRAPHIC>', ''],
  ['\\Item', 0, 0, '', '', 1, 1, '', ''],
  ['\\Ans', 1, 1, '', ' Ans.'],
  ['\\First', 1, 1, '', ''],
  ['\\tb', 0, 0, '', ''],
  ['\\Defn', 1, 1, '', ''],
  ['\\Dictum', 1, 1, '', ''],
  ['\\DPtypo', 1, 0, '', '', 1, 1, '', ''],
  ['\\Add', 1, 1, '', ''],
  ['\\ResetCols', 0, 0, '', ''],
  ['\\begin{multicols}', 1, 0, '', '']
  );
###
This is pdfTeXk, Version 3.141592-1.40.3 (Web2C 7.5.6) (format=pdflatex 2010.5.6)  8 JUL 2011 19:42
entering extended mode
 %&-line parsing enabled.
**36670-t.tex
(./36670-t.tex
LaTeX2e <2005/12/01>
Babel <v3.8h> and hyphenation patterns for english, usenglishmax, dumylang, noh
yphenation, arabic, farsi, croatian, ukrainian, russian, bulgarian, czech, slov
ak, danish, dutch, finnish, basque, french, german, ngerman, ibycus, greek, mon
ogreek, ancientgreek, hungarian, italian, latin, mongolian, norsk, icelandic, i
nterlingua, turkish, coptic, romanian, welsh, serbian, slovenian, estonian, esp
eranto, uppersorbian, indonesian, polish, portuguese, spanish, catalan, galicia
n, swedish, ukenglish, pinyin, loaded.
(/usr/share/texmf-texlive/tex/latex/base/book.cls
Document Class: book 2005/09/16 v1.4f Standard LaTeX document class
(/usr/share/texmf-texlive/tex/latex/base/bk12.clo
File: bk12.clo 2005/09/16 v1.4f Standard LaTeX file (size option)
)
\c@part=\count79
\c@chapter=\count80
\c@section=\count81
\c@subsection=\count82
\c@subsubsection=\count83
\c@paragraph=\count84
\c@subparagraph=\count85
\c@figure=\count86
\c@table=\count87
\abovecaptionskip=\skip41
\belowcaptionskip=\skip42
\bibindent=\dimen102
) (/usr/share/texmf-texlive/tex/latex/base/inputenc.sty
Package: inputenc 2006/05/05 v1.1b Input encoding file
\inpenc@prehook=\toks14
\inpenc@posthook=\toks15
(/usr/share/texmf-texlive/tex/latex/base/latin1.def
File: latin1.def 2006/05/05 v1.1b Input encoding file
)) (/usr/share/texmf-texlive/tex/latex/base/ifthen.sty
Package: ifthen 2001/05/26 v1.1c Standard LaTeX ifthen package (DPC)
) (/usr/share/texmf-texlive/tex/latex/amsmath/amsmath.sty
Package: amsmath 2000/07/18 v2.13 AMS math features
\@mathmargin=\skip43
For additional information on amsmath, use the `?' option.
(/usr/share/texmf-texlive/tex/latex/amsmath/amstext.sty
Package: amstext 2000/06/29 v2.01
(/usr/share/texmf-texlive/tex/latex/amsmath/amsgen.sty
File: amsgen.sty 1999/11/30 v2.0
\@emptytoks=\toks16
\ex@=\dimen103
)) (/usr/share/texmf-texlive/tex/latex/amsmath/amsbsy.sty
Package: amsbsy 1999/11/29 v1.2d
\pmbraise@=\dimen104
) (/usr/share/texmf-texlive/tex/latex/amsmath/amsopn.sty
Package: amsopn 1999/12/14 v2.01 operator names
)
\inf@bad=\count88
LaTeX Info: Redefining \frac on input line 211.
\uproot@=\count89
\leftroot@=\count90
LaTeX Info: Redefining \overline on input line 307.
\classnum@=\count91
\DOTSCASE@=\count92
LaTeX Info: Redefining \ldots on input line 379.
LaTeX Info: Redefining \dots on input line 382.
LaTeX Info: Redefining \cdots on input line 467.
\Mathstrutbox@=\box26
\strutbox@=\box27
\big@size=\dimen105
LaTeX Font Info:    Redeclaring font encoding OML on input line 567.
LaTeX Font Info:    Redeclaring font encoding OMS on input line 568.
\macc@depth=\count93
\c@MaxMatrixCols=\count94
\dotsspace@=\muskip10
\c@parentequation=\count95
\dspbrk@lvl=\count96
\tag@help=\toks17
\row@=\count97
\column@=\count98
\maxfields@=\count99
\andhelp@=\toks18
\eqnshift@=\dimen106
\alignsep@=\dimen107
\tagshift@=\dimen108
\tagwidth@=\dimen109
\totwidth@=\dimen110
\lineht@=\dimen111
\@envbody=\toks19
\multlinegap=\skip44
\multlinetaggap=\skip45
\mathdisplay@stack=\toks20
LaTeX Info: Redefining \[ on input line 2666.
LaTeX Info: Redefining \] on input line 2667.
) (/usr/share/texmf-texlive/tex/latex/amsfonts/amssymb.sty
Package: amssymb 2002/01/22 v2.2d
(/usr/share/texmf-texlive/tex/latex/amsfonts/amsfonts.sty
Package: amsfonts 2001/10/25 v2.2f
\symAMSa=\mathgroup4
\symAMSb=\mathgroup5
LaTeX Font Info:    Overwriting math alphabet `\mathfrak' in version `bold'
(Font)                  U/euf/m/n --> U/euf/b/n on input line 132.
)) (/usr/share/texmf-texlive/tex/latex/base/alltt.sty
Package: alltt 1997/06/16 v2.0g defines alltt environment
) (/usr/share/texmf-texlive/tex/latex/tools/array.sty
Package: array 2005/08/23 v2.4b Tabular extension package (FMi)
\col@sep=\dimen112
\extrarowheight=\dimen113
\NC@list=\toks21
\extratabsurround=\skip46
\backup@length=\skip47
) (/usr/share/texmf-texlive/tex/latex/tools/multicol.sty
Package: multicol 2006/05/18 v1.6g multicolumn formatting (FMi)
\c@tracingmulticols=\count100
\mult@box=\box28
\multicol@leftmargin=\dimen114
\c@unbalance=\count101
\c@collectmore=\count102
\doublecol@number=\count103
\multicoltolerance=\count104
\multicolpretolerance=\count105
\full@width=\dimen115
\page@free=\dimen116
\premulticols=\dimen117
\postmulticols=\dimen118
\multicolsep=\skip48
\multicolbaselineskip=\skip49
\partial@page=\box29
\last@line=\box30
\mult@rightbox=\box31
\mult@grightbox=\box32
\mult@gfirstbox=\box33
\mult@firstbox=\box34
\@tempa=\box35
\@tempa=\box36
\@tempa=\box37
\@tempa=\box38
\@tempa=\box39
\@tempa=\box40
\@tempa=\box41
\@tempa=\box42
\@tempa=\box43
\@tempa=\box44
\@tempa=\box45
\@tempa=\box46
\@tempa=\box47
\@tempa=\box48
\@tempa=\box49
\@tempa=\box50
\@tempa=\box51
\c@columnbadness=\count106
\c@finalcolumnbadness=\count107
\last@try=\dimen119
\multicolovershoot=\dimen120
\multicolundershoot=\dimen121
\mult@nat@firstbox=\box52
\colbreak@box=\box53
) (/usr/share/texmf-texlive/tex/latex/graphics/graphicx.sty
Package: graphicx 1999/02/16 v1.0f Enhanced LaTeX Graphics (DPC,SPQR)
(/usr/share/texmf-texlive/tex/latex/graphics/keyval.sty
Package: keyval 1999/03/16 v1.13 key=value parser (DPC)
\KV@toks@=\toks22
) (/usr/share/texmf-texlive/tex/latex/graphics/graphics.sty
Package: graphics 2006/02/20 v1.0o Standard LaTeX Graphics (DPC,SPQR)
(/usr/share/texmf-texlive/tex/latex/graphics/trig.sty
Package: trig 1999/03/16 v1.09 sin cos tan (DPC)
) (/etc/texmf/tex/latex/config/graphics.cfg
File: graphics.cfg 2007/01/18 v1.5 graphics configuration of teTeX/TeXLive
)
Package graphics Info: Driver file: pdftex.def on input line 90.
(/usr/share/texmf-texlive/tex/latex/pdftex-def/pdftex.def
File: pdftex.def 2007/01/08 v0.04d Graphics/color for pdfTeX
\Gread@gobject=\count108
))
\Gin@req@height=\dimen122
\Gin@req@width=\dimen123
) (/usr/share/texmf-texlive/tex/latex/tools/indentfirst.sty
Package: indentfirst 1995/11/23 v1.03 Indent first paragraph (DPC)
) (/usr/share/texmf-texlive/tex/latex/textcase/textcase.sty
Package: textcase 2004/10/07 v0.07 Text only upper/lower case changing (DPC)
) (/usr/share/texmf-texlive/tex/latex/tools/calc.sty
Package: calc 2005/08/06 v4.2 Infix arithmetic (KKT,FJ)
\calc@Acount=\count109
\calc@Bcount=\count110
\calc@Adimen=\dimen124
\calc@Bdimen=\dimen125
\calc@Askip=\skip50
\calc@Bskip=\skip51
LaTeX Info: Redefining \setlength on input line 75.
LaTeX Info: Redefining \addtolength on input line 76.
\calc@Ccount=\count111
\calc@Cskip=\skip52
) (/usr/share/texmf-texlive/tex/latex/fancyhdr/fancyhdr.sty
\fancy@headwidth=\skip53
\f@ncyO@elh=\skip54
\f@ncyO@erh=\skip55
\f@ncyO@olh=\skip56
\f@ncyO@orh=\skip57
\f@ncyO@elf=\skip58
\f@ncyO@erf=\skip59
\f@ncyO@olf=\skip60
\f@ncyO@orf=\skip61
) (/usr/share/texmf-texlive/tex/latex/geometry/geometry.sty
Package: geometry 2002/07/08 v3.2 Page Geometry
\Gm@cnth=\count112
\Gm@cntv=\count113
\c@Gm@tempcnt=\count114
\Gm@bindingoffset=\dimen126
\Gm@wd@mp=\dimen127
\Gm@odd@mp=\dimen128
\Gm@even@mp=\dimen129
\Gm@dimlist=\toks23
(/usr/share/texmf-texlive/tex/xelatex/xetexconfig/geometry.cfg)) (/usr/share/te
xmf-texlive/tex/latex/hyperref/hyperref.sty
Package: hyperref 2007/02/07 v6.75r Hypertext links for LaTeX
\@linkdim=\dimen130
\Hy@linkcounter=\count115
\Hy@pagecounter=\count116
(/usr/share/texmf-texlive/tex/latex/hyperref/pd1enc.def
File: pd1enc.def 2007/02/07 v6.75r Hyperref: PDFDocEncoding definition (HO)
) (/etc/texmf/tex/latex/config/hyperref.cfg
File: hyperref.cfg 2002/06/06 v1.2 hyperref configuration of TeXLive
) (/usr/share/texmf-texlive/tex/latex/oberdiek/kvoptions.sty
Package: kvoptions 2006/08/22 v2.4 Connects package keyval with LaTeX options (
HO)
)
Package hyperref Info: Option `hyperfootnotes' set `false' on input line 2238.
Package hyperref Info: Option `bookmarks' set `true' on input line 2238.
Package hyperref Info: Option `linktocpage' set `false' on input line 2238.
Package hyperref Info: Option `pdfdisplaydoctitle' set `true' on input line 223
8.
Package hyperref Info: Option `pdfpagelabels' set `true' on input line 2238.
Package hyperref Info: Option `bookmarksopen' set `true' on input line 2238.
Package hyperref Info: Option `colorlinks' set `true' on input line 2238.
Package hyperref Info: Hyper figures OFF on input line 2288.
Package hyperref Info: Link nesting OFF on input line 2293.
Package hyperref Info: Hyper index ON on input line 2296.
Package hyperref Info: Plain pages OFF on input line 2303.
Package hyperref Info: Backreferencing OFF on input line 2308.
Implicit mode ON; LaTeX internals redefined
Package hyperref Info: Bookmarks ON on input line 2444.
(/usr/share/texmf-texlive/tex/latex/ltxmisc/url.sty
\Urlmuskip=\muskip11
Package: url 2005/06/27  ver 3.2  Verb mode for urls, etc.
)
LaTeX Info: Redefining \url on input line 2599.
\Fld@menulength=\count117
\Field@Width=\dimen131
\Fld@charsize=\dimen132
\Choice@toks=\toks24
\Field@toks=\toks25
Package hyperref Info: Hyper figures OFF on input line 3102.
Package hyperref Info: Link nesting OFF on input line 3107.
Package hyperref Info: Hyper index ON on input line 3110.
Package hyperref Info: backreferencing OFF on input line 3117.
Package hyperref Info: Link coloring ON on input line 3120.
\Hy@abspage=\count118
\c@Item=\count119
)
*hyperref using driver hpdftex*
(/usr/share/texmf-texlive/tex/latex/hyperref/hpdftex.def
File: hpdftex.def 2007/02/07 v6.75r Hyperref driver for pdfTeX
\Fld@listcount=\count120
)
\TmpLen=\skip62
\DP@lign@no=\count121
\DP@lignb@dy=\toks26
(./36670-t.aux)
\openout1 = `36670-t.aux'.

LaTeX Font Info:    Checking defaults for OML/cmm/m/it on input line 603.
LaTeX Font Info:    ... okay on input line 603.
LaTeX Font Info:    Checking defaults for T1/cmr/m/n on input line 603.
LaTeX Font Info:    ... okay on input line 603.
LaTeX Font Info:    Checking defaults for OT1/cmr/m/n on input line 603.
LaTeX Font Info:    ... okay on input line 603.
LaTeX Font Info:    Checking defaults for OMS/cmsy/m/n on input line 603.
LaTeX Font Info:    ... okay on input line 603.
LaTeX Font Info:    Checking defaults for OMX/cmex/m/n on input line 603.
LaTeX Font Info:    ... okay on input line 603.
LaTeX Font Info:    Checking defaults for U/cmr/m/n on input line 603.
LaTeX Font Info:    ... okay on input line 603.
LaTeX Font Info:    Checking defaults for PD1/pdf/m/n on input line 603.
LaTeX Font Info:    ... okay on input line 603.
(/usr/share/texmf/tex/context/base/supp-pdf.tex
[Loading MPS to PDF converter (version 2006.09.02).]
\scratchcounter=\count122
\scratchdimen=\dimen133
\scratchbox=\box54
\nofMPsegments=\count123
\nofMParguments=\count124
\everyMPshowfont=\toks27
\MPscratchCnt=\count125
\MPscratchDim=\dimen134
\MPnumerator=\count126
\everyMPtoPDFconversion=\toks28
)
-------------------- Geometry parameters
paper: class default
landscape: --
twocolumn: --
twoside: true
asymmetric: --
h-parts: 9.03374pt, 325.215pt, 9.03375pt
v-parts: 4.15848pt, 495.49379pt, 6.23773pt
hmarginratio: 1:1
vmarginratio: 2:3
lines: --
heightrounded: --
bindingoffset: 0.0pt
truedimen: --
includehead: true
includefoot: true
includemp: --
driver: pdftex
-------------------- Page layout dimensions and switches
\paperwidth  343.28249pt
\paperheight 505.89pt
\textwidth  325.215pt
\textheight 433.62pt
\oddsidemargin  -63.23625pt
\evensidemargin -63.23624pt
\topmargin  -68.11151pt
\headheight 12.0pt
\headsep    19.8738pt
\footskip   30.0pt
\marginparwidth 98.0pt
\marginparsep   7.0pt
\columnsep  10.0pt
\skip\footins  10.8pt plus 4.0pt minus 2.0pt
\hoffset 0.0pt
\voffset 0.0pt
\mag 1000
\@twosidetrue \@mparswitchtrue 
(1in=72.27pt, 1cm=28.45pt)
-----------------------
(/usr/share/texmf-texlive/tex/latex/graphics/color.sty
Package: color 2005/11/14 v1.0j Standard LaTeX Color (DPC)
(/etc/texmf/tex/latex/config/color.cfg
File: color.cfg 2007/01/18 v1.5 color configuration of teTeX/TeXLive
)
Package color Info: Driver file: pdftex.def on input line 130.
)
Package hyperref Info: Link coloring ON on input line 603.
(/usr/share/texmf-texlive/tex/latex/hyperref/nameref.sty
Package: nameref 2006/12/27 v2.28 Cross-referencing by name of section
(/usr/share/texmf-texlive/tex/latex/oberdiek/refcount.sty
Package: refcount 2006/02/20 v3.0 Data extraction from references (HO)
)
\c@section@level=\count127
)
LaTeX Info: Redefining \ref on input line 603.
LaTeX Info: Redefining \pageref on input line 603.
(./36670-t.out) (./36670-t.out)
\@outlinefile=\write3
\openout3 = `36670-t.out'.

LaTeX Font Info:    Try loading font information for U+msa on input line 632.
(/usr/share/texmf-texlive/tex/latex/amsfonts/umsa.fd
File: umsa.fd 2002/01/19 v2.2g AMS font definitions
)
LaTeX Font Info:    Try loading font information for U+msb on input line 632.
(/usr/share/texmf-texlive/tex/latex/amsfonts/umsb.fd
File: umsb.fd 2002/01/19 v2.2g AMS font definitions
) [1

{/var/lib/texmf/fonts/map/pdftex/updmap/pdftex.map}] [2

] [1


] [2] [3


] [4] [5] (./36670-t.toc)
\tf@toc=\write4
\openout4 = `36670-t.toc'.


Underfull \hbox (badness 10000) in paragraph at lines 1--789

 []

[6

] [1




] [2] [3] [4] [5] [6] [7] [8] [9]
Overfull \hbox (1.72365pt too wide) in paragraph at lines 1227--1230
[]\OT1/cmr/m/it/12 If an ex-pres-sion within a paren-the-sis is pre-ceded by th
e sign $\OT1/cmr/m/n/12 +$\OT1/cmr/m/it/12 ,
 []

[10]
Overfull \hbox (1.91821pt too wide) in paragraph at lines 1272--1276
[]\OT1/cmr/m/it/12 If an ex-pres-sion within a paren-the-sis is pre-ceded by th
e sign $\OMS/cmsy/m/n/12 ^^@$\OT1/cmr/m/it/12 ,
 []

[11] [12] [13] [14] [15] [16] [17] [18] [19] [20] [21] [22] [23] [24

] [25]
Underfull \hbox (badness 10000) in paragraph at lines 1954--1961

 []

[26]
LaTeX Font Info:    Try loading font information for OMS+cmr on input line 1964
.
(/usr/share/texmf-texlive/tex/latex/base/omscmr.fd
File: omscmr.fd 1999/05/25 v2.5h Standard LaTeX font definitions
)
LaTeX Font Info:    Font shape `OMS/cmr/m/n' in size <12> not available
(Font)              Font shape `OMS/cmsy/m/n' tried instead on input line 1964.


Underfull \hbox (badness 10000) in paragraph at lines 1966--1973

 []

[27]
Overfull \hbox (2.96944pt too wide) in paragraph at lines 2013--2013
[] 
 []

[28]
Underfull \hbox (badness 10000) in paragraph at lines 2054--2065

 []


Overfull \hbox (10.01392pt too wide) in paragraph at lines 2065--2065
[] 
 []

[29]
Overfull \hbox (36.58296pt too wide) in paragraph at lines 2093--2093
[] 
 []

[30] [31]
Overfull \hbox (11.34335pt too wide) in paragraph at lines 2206--2206
[] 
 []

[32]
Overfull \hbox (11.61705pt too wide) in paragraph at lines 2241--2241
[] 
 []

[33] [34] [35] [36] [37] [38]
Overfull \hbox (1.94836pt too wide) in paragraph at lines 2459--2462
[][] \OT1/cmr/m/n/12 I have $$12\OML/cmm/m/it/12 :\OT1/cmr/m/n/12 75$ in two-do
llar bills and twenty-five cent pieces,
 []

[39] [40] [41] <./images/fig1.pdf, id=408, 361.35pt x 26.0975pt>
File: ./images/fig1.pdf Graphic file (type pdf)
<use ./images/fig1.pdf> [42

 <./images/fig1.pdf>] <./images/fig2.pdf, id=423, 439.6425pt x 24.09pt>
File: ./images/fig2.pdf Graphic file (type pdf)
<use ./images/fig2.pdf> [43 <./images/fig2.pdf>] [44] <./images/fig3.pdf, id=44
1, 427.5975pt x 24.09pt>
File: ./images/fig3.pdf Graphic file (type pdf)
<use ./images/fig3.pdf> [45 <./images/fig3.pdf>] [46] [47] [48] [49] [50]
Overfull \hbox (1.24844pt too wide) in paragraph at lines 2867--2871
\OT1/cmr/m/it/12 mul-ti-pli-cand \OT1/cmr/m/n/12 and \OT1/cmr/m/it/12 mul-ti-pl
ier\OT1/cmr/m/n/12 , a third num-ber called \OT1/cmr/m/it/12 prod-uct\OT1/cmr/m
/n/12 , which
 []


Overfull \hbox (21.76723pt too wide) in paragraph at lines 2872--2894
[]\OT1/cmr/bx/n/12 71. \OT1/cmr/m/n/12 Ac-cord-ing to this def-i-ni-tion of mul
-ti-pli-ca-tion,                             
 []

[51] [52] [53] [54] [55] [56] [57] [58

] [59] [60] [61] [62] [63] [64] [65] [66] [67

] [68] [69] [70] [71] [72] [73] [74] [75] [76] [77] [78] [79] [80] [81] [82

] [83]
Overfull \hbox (42.68674pt too wide) in paragraph at lines 4261--4288
\OT1/cmr/m/n/12 should be care-fully no-ticed and re-mem-bered.                
               
 []

[84] [85] [86] [87] [88]
Overfull \hbox (1.2826pt too wide) in paragraph at lines 4520--4521
[][] $[]$\OT1/cmr/m/n/12 . 
 []

[89]
Overfull \hbox (1.08803pt too wide) in paragraph at lines 4576--4577
[][] $[]$\OT1/cmr/m/n/12 . 
 []

[90] [91] [92

] [93] [94]
Overfull \hbox (2.12788pt too wide) in paragraph at lines 4767--4768
[]\OT1/cmr/bx/n/12 108. When a bi-no-mial is the dif-fer-ence of two squares. 
 []

[95] [96] [97] [98]
Underfull \hbox (badness 10000) in paragraph at lines 4959--4963

 []

[99] [100] [101] [102] [103] [104]
Overfull \hbox (1.29138pt too wide) in paragraph at lines 5268--5270
[]\OT1/cmr/m/n/10.95 Two num-bers whose prod-uct is $30$ are $1$ and $30$, $2$ 
and $15$, $3$ and $10$,
 []

[105] [106] [107] [108] [109] [110

] [111] [112] [113] [114] [115] [116] [117

] [118] [119] [120] [121] [122] [123]
Overfull \hbox (0.88312pt too wide) in paragraph at lines 6065--6066
[][] $[]$\OT1/cmr/m/n/12 , $[]$. 
 []

[124]
Overfull \hbox (3.79523pt too wide) in paragraph at lines 6118--6118
[] 
 []

[125] [126] [127] [128] [129] [130] [131] [132] [133] [134] [135]
Overfull \hbox (12.6967pt too wide) in paragraph at lines 6449--6449
[] 
 []

[136

]
Overfull \hbox (13.73837pt too wide) in paragraph at lines 6508--6508
[] 
 []

[137] [138] [139] [140]
Overfull \hbox (127.43008pt too wide) in paragraph at lines 6644--6644
[] 
 []

[141]
Overfull \hbox (82.71672pt too wide) in paragraph at lines 6685--6685
[] 
 []

[142] [143]
Overfull \hbox (50.27081pt too wide) in paragraph at lines 6733--6733
[] 
 []


Overfull \hbox (2.05074pt too wide) in paragraph at lines 6749--6749
[] 
 []

[144] [145]
Overfull \hbox (60.99086pt too wide) in paragraph at lines 6827--6827
[] 
 []

[146] [147]
Overfull \hbox (41.2357pt too wide) in paragraph at lines 6885--6885
[] 
 []

[148]
Overfull \hbox (33.08694pt too wide) in paragraph at lines 6936--6936
[] 
 []

[149] [150]
Overfull \hbox (30.04593pt too wide) in paragraph at lines 6991--6991
[] 
 []

[151]
Overfull \hbox (57.59785pt too wide) in paragraph at lines 7039--7039
[] 
 []

[152]
Overfull \hbox (24.40863pt too wide) in paragraph at lines 7052--7052
[] 
 []

[153]
Overfull \hbox (54.40073pt too wide) in paragraph at lines 7088--7088
[] 
 []

[154]
Overfull \hbox (55.29538pt too wide) in paragraph at lines 7144--7144
[] 
 []

[155] [156]
Overfull \hbox (25.74243pt too wide) in paragraph at lines 7237--7237
[] 
 []

[157]
Overfull \hbox (22.96907pt too wide) in paragraph at lines 7265--7265
[] 
 []

[158] [159]
Overfull \hbox (2.44133pt too wide) in paragraph at lines 7352--7354
[][] \OT1/cmr/m/n/12 Find the prin-ci-pal that will pro-duce $$280$ in-ter-est 
in $2$ years
 []

[160] [161] [162

] [163] [164] [165] [166] [167] [168] [169] [170] [171] [172]
Overfull \hbox (2.92752pt too wide) in paragraph at lines 7879--7882
[][] \OT1/cmr/m/n/12 A sum of money, at sim-ple in-ter-est, amounted in $10$ mo
nths
 []

[173] [174] [175

]
Overfull \hbox (37.3597pt too wide) in paragraph at lines 7986--7986
[] 
 []


Overfull \hbox (19.65984pt too wide) in paragraph at lines 8009--8009
[] 
 []

[176]
Overfull \hbox (28.17653pt too wide) in paragraph at lines 8026--8026
[] 
 []

[177] [178]
Overfull \hbox (101.83775pt too wide) in paragraph at lines 8136--8136
[] 
 []

[179]
Overfull \hbox (75.6592pt too wide) in paragraph at lines 8158--8158
[] 
 []


Overfull \hbox (81.72191pt too wide) in paragraph at lines 8180--8180
[] 
 []

[180] [181] [182] [183]
Overfull \hbox (14.36713pt too wide) in paragraph at lines 8349--8349
[] 
 []


Overfull \hbox (39.82089pt too wide) in paragraph at lines 8367--8367
[] 
 []

[184]
Overfull \hbox (48.35448pt too wide) in paragraph at lines 8374--8374
[] 
 []


Overfull \hbox (14.62827pt too wide) in paragraph at lines 8395--8395
[] 
 []

[185] [186]
Overfull \hbox (36.34747pt too wide) in paragraph at lines 8465--8465
[] 
 []

[187] [188] [189

] [190]
Overfull \hbox (41.95776pt too wide) in paragraph at lines 8610--8610
[] 
 []

[191] [192] [193] [194] [195] [196] [197

]
Overfull \hbox (5.98813pt too wide) in paragraph at lines 8846--8859
[]\OT1/cmr/m/n/10.95 The 5th and 7th terms are $\OML/cmm/m/it/10.95 ar[]$ \OT1/
cmr/m/n/10.95 and $\OML/cmm/m/it/10.95 ar[]$\OT1/cmr/m/n/10.95 , re-spec-tively
.                 
 []

[198]
Overfull \hbox (1.00482pt too wide) in paragraph at lines 8878--8880
[]\OT1/cmr/m/n/12 Hence\OT1/cmr/m/it/12 , the ge-o-met-ri-cal mean of any two n
um-bers is the square
 []


Overfull \hbox (84.25482pt too wide) in paragraph at lines 8894--8894
[] 
 []

[199]
Overfull \hbox (35.57372pt too wide) in paragraph at lines 8921--8921
[] 
 []


Overfull \hbox (36.34883pt too wide) in paragraph at lines 8935--8935
[] 
 []

[200]
Overfull \hbox (2.11494pt too wide) in paragraph at lines 8950--8952
[][] \OT1/cmr/m/n/12 Find the com-mon ra-tio if the 1st and 3d terms are $2$ an
d $32$. 
 []

[201] [202] [203

] [204] [205] [206] [207] [208] [209] [210] [211] [212] [213] [214] [215] [216]
[217] [218] [219] [220]
Overfull \hbox (2.5148pt too wide) in paragraph at lines 9785--9786
[][]$\OT1/cmr/m/n/10 14$. 
 []


Overfull \hbox (2.5148pt too wide) in paragraph at lines 9787--9788
[][]$\OT1/cmr/m/n/10 10$. 
 []


Overfull \hbox (2.5148pt too wide) in paragraph at lines 9789--9790
[][]$\OT1/cmr/m/n/10 13$. 
 []


Overfull \hbox (2.5148pt too wide) in paragraph at lines 9791--9792
[][]$\OT1/cmr/m/n/10 11$. 
 []


Overfull \hbox (2.5148pt too wide) in paragraph at lines 9793--9794
[][]$\OT1/cmr/m/n/10 13$. 
 []


Overfull \hbox (2.5148pt too wide) in paragraph at lines 9825--9826
[][]$\OT1/cmr/m/n/10 10$. 
 []


Overfull \hbox (2.5148pt too wide) in paragraph at lines 9830--9831
[][]$\OT1/cmr/m/n/10 91$. 
 []


Overfull \hbox (2.5148pt too wide) in paragraph at lines 9832--9833
[][]$\OT1/cmr/m/n/10 21$. 
 []


Overfull \hbox (2.5148pt too wide) in paragraph at lines 9834--9835
[][]$\OT1/cmr/m/n/10 60$. 
 []


Overfull \hbox (2.5148pt too wide) in paragraph at lines 9836--9837
[][]$\OT1/cmr/m/n/10 24$. 
 []


Overfull \hbox (2.5148pt too wide) in paragraph at lines 9838--9839
[][]$\OT1/cmr/m/n/10 96$. 
 []


Overfull \hbox (2.5148pt too wide) in paragraph at lines 9840--9841
[][]$\OT1/cmr/m/n/10 16$. 
 []


Overfull \hbox (2.5148pt too wide) in paragraph at lines 9842--9843
[][]$\OT1/cmr/m/n/10 36$. 
 []

[221


] [222] [223] [224]
Underfull \hbox (badness 10000) in paragraph at lines 10283--10284
[][]\OT1/cmr/m/n/10 $$6$ quar-ters, $18$ ten-cent
 []

[225] [226] [227] [228] [229] [230] [231] [232] [233] [234] [235] [236] [237] [
238]
Underfull \hbox (badness 10000) in paragraph at lines 11388--11389
[][]$\OT1/cmr/m/n/10 (1 + 10\OML/cmm/m/it/10 a\OT1/cmr/m/n/10 )(1 +
 []

[239] [240] [241] [242]
Overfull \hbox (24.36917pt too wide) in paragraph at lines 11738--11739
[][]$[]$\OT1/cmr/m/n/10 . 
 []

[243] [244]
Overfull \hbox (8.9993pt too wide) in paragraph at lines 11831--11832
[][]$[]$\OT1/cmr/m/n/10 . 
 []

[245] [246] [247]
Underfull \hbox (badness 10000) in paragraph at lines 12079--12080
[][]$\OT1/cmr/m/n/10 100$,
 []


Underfull \hbox (badness 10000) in paragraph at lines 12083--12084
[][]$\OT1/cmr/m/n/10 103$,
 []


Underfull \hbox (badness 10000) in paragraph at lines 12085--12086
[][]$\OT1/cmr/m/n/10 295$,
 []


Underfull \hbox (badness 10000) in paragraph at lines 12100--12101
[][]\OT1/cmr/m/n/10 Son,
 []


Underfull \hbox (badness 10000) in paragraph at lines 12106--12107
[][]\OT1/cmr/m/n/10 Son,
 []


Overfull \hbox (8.14185pt too wide) in paragraph at lines 12111--12112
[][]$\OT1/cmr/m/n/10 1[]$ dy. 
 []


Overfull \hbox (12.12798pt too wide) in paragraph at lines 12113--12114
[][]$\OT1/cmr/m/n/10 1[]$ dy. 
 []


Overfull \hbox (12.12798pt too wide) in paragraph at lines 12115--12116
[][]$\OT1/cmr/m/n/10 1[]$ dy. 
 []


Overfull \hbox (6.75574pt too wide) in paragraph at lines 12117--12118
[][]$\OT1/cmr/m/n/10 15$ dy. 
 []


Overfull \hbox (6.22795pt too wide) in paragraph at lines 12119--12120
[][]$\OT1/cmr/m/n/10 12$ hr. 
 []


Overfull \hbox (6.75574pt too wide) in paragraph at lines 12121--12122
[][]$\OT1/cmr/m/n/10 10$ dy. 
 []


Overfull \hbox (0.42635pt too wide) in paragraph at lines 12126--12127
[][]$\OT1/cmr/m/n/10 7[]$ hr. 
 []


Overfull \hbox (0.42635pt too wide) in paragraph at lines 12132--12133
[][]$\OT1/cmr/m/n/10 1[]$ hr. 
 []


Overfull \vbox (1.93811pt too high) has occurred while \output is active []


Overfull \vbox (1.93811pt too high) has occurred while \output is active []

[248]
Underfull \hbox (badness 10000) in paragraph at lines 12188--12189
[][]$\OT1/cmr/m/n/10 90[] 0[] 30[]$;
 []

[249]
Underfull \hbox (badness 10000) in paragraph at lines 12320--12321
[][]\OT1/cmr/m/n/10 $$6000$ at $4$%;
 []


Underfull \hbox (badness 10000) in paragraph at lines 12331--12332
[][]\OT1/cmr/m/n/10 Wheat,
 []


Underfull \hbox (badness 10000) in paragraph at lines 12333--12334
[][]\OT1/cmr/m/n/10 A, $$235$;
 []


Underfull \hbox (badness 10000) in paragraph at lines 12337--12338
[][]\OT1/cmr/m/n/10 Lemon,
 []


Underfull \hbox (badness 10000) in paragraph at lines 12339--12340
\OT1/cmr/m/n/10 ples; B,
 []

[250]
Underfull \hbox (badness 10000) in paragraph at lines 12370--12371
[][]\OT1/cmr/m/n/10 Width,
 []


Underfull \hbox (badness 10000) in paragraph at lines 12370--12371
\OT1/cmr/m/n/10 12$ rd.;
 []


Underfull \hbox (badness 10000) in paragraph at lines 12370--12371
\OT1/cmr/m/n/10 length,
 []

[251] [252] [253] [1



] [2] [3] [4] [5] [6] [7] [8] (./36670-t.aux)

 *File List*
    book.cls    2005/09/16 v1.4f Standard LaTeX document class
    bk12.clo    2005/09/16 v1.4f Standard LaTeX file (size option)
inputenc.sty    2006/05/05 v1.1b Input encoding file
  latin1.def    2006/05/05 v1.1b Input encoding file
  ifthen.sty    2001/05/26 v1.1c Standard LaTeX ifthen package (DPC)
 amsmath.sty    2000/07/18 v2.13 AMS math features
 amstext.sty    2000/06/29 v2.01
  amsgen.sty    1999/11/30 v2.0
  amsbsy.sty    1999/11/29 v1.2d
  amsopn.sty    1999/12/14 v2.01 operator names
 amssymb.sty    2002/01/22 v2.2d
amsfonts.sty    2001/10/25 v2.2f
   alltt.sty    1997/06/16 v2.0g defines alltt environment
   array.sty    2005/08/23 v2.4b Tabular extension package (FMi)
multicol.sty    2006/05/18 v1.6g multicolumn formatting (FMi)
graphicx.sty    1999/02/16 v1.0f Enhanced LaTeX Graphics (DPC,SPQR)
  keyval.sty    1999/03/16 v1.13 key=value parser (DPC)
graphics.sty    2006/02/20 v1.0o Standard LaTeX Graphics (DPC,SPQR)
    trig.sty    1999/03/16 v1.09 sin cos tan (DPC)
graphics.cfg    2007/01/18 v1.5 graphics configuration of teTeX/TeXLive
  pdftex.def    2007/01/08 v0.04d Graphics/color for pdfTeX
indentfirst.sty    1995/11/23 v1.03 Indent first paragraph (DPC)
textcase.sty    2004/10/07 v0.07 Text only upper/lower case changing (DPC)
    calc.sty    2005/08/06 v4.2 Infix arithmetic (KKT,FJ)
fancyhdr.sty    
geometry.sty    2002/07/08 v3.2 Page Geometry
geometry.cfg
hyperref.sty    2007/02/07 v6.75r Hypertext links for LaTeX
  pd1enc.def    2007/02/07 v6.75r Hyperref: PDFDocEncoding definition (HO)
hyperref.cfg    2002/06/06 v1.2 hyperref configuration of TeXLive
kvoptions.sty    2006/08/22 v2.4 Connects package keyval with LaTeX options (HO
)
     url.sty    2005/06/27  ver 3.2  Verb mode for urls, etc.
 hpdftex.def    2007/02/07 v6.75r Hyperref driver for pdfTeX
supp-pdf.tex
   color.sty    2005/11/14 v1.0j Standard LaTeX Color (DPC)
   color.cfg    2007/01/18 v1.5 color configuration of teTeX/TeXLive
 nameref.sty    2006/12/27 v2.28 Cross-referencing by name of section
refcount.sty    2006/02/20 v3.0 Data extraction from references (HO)
 36670-t.out
 36670-t.out
    umsa.fd    2002/01/19 v2.2g AMS font definitions
    umsb.fd    2002/01/19 v2.2g AMS font definitions
  omscmr.fd    1999/05/25 v2.5h Standard LaTeX font definitions
./images/fig1.pdf
./images/fig2.pdf
./images/fig3.pdf
 ***********

 ) 
Here is how much of TeX's memory you used:
 5458 strings out of 94074
 70648 string characters out of 1165154
 142333 words of memory out of 1500000
 7978 multiletter control sequences out of 10000+50000
 19082 words of font info for 72 fonts, out of 1200000 for 2000
 645 hyphenation exceptions out of 8191
 29i,18n,46p,258b,494s stack positions out of 5000i,500n,6000p,200000b,5000s
</usr/share/texmf-texlive/fonts/type1/bluesky/cm/cmbx10.pfb></usr/share/texmf
-texlive/fonts/type1/bluesky/cm/cmbx12.pfb></usr/share/texmf-texlive/fonts/type
1/bluesky/cm/cmcsc10.pfb></usr/share/texmf-texlive/fonts/type1/bluesky/cm/cmex1
0.pfb></usr/share/texmf-texlive/fonts/type1/bluesky/cm/cmmi10.pfb></usr/share/t
exmf-texlive/fonts/type1/bluesky/cm/cmmi12.pfb></usr/share/texmf-texlive/fonts/
type1/bluesky/cm/cmmi8.pfb></usr/share/texmf-texlive/fonts/type1/bluesky/cm/cmr
10.pfb></usr/share/texmf-texlive/fonts/type1/bluesky/cm/cmr12.pfb></usr/share/t
exmf-texlive/fonts/type1/bluesky/cm/cmr17.pfb></usr/share/texmf-texlive/fonts/t
ype1/bluesky/cm/cmr6.pfb></usr/share/texmf-texlive/fonts/type1/bluesky/cm/cmr7.
pfb></usr/share/texmf-texlive/fonts/type1/bluesky/cm/cmr8.pfb></usr/share/texmf
-texlive/fonts/type1/bluesky/cm/cmsy10.pfb></usr/share/texmf-texlive/fonts/type
1/bluesky/cm/cmsy7.pfb></usr/share/texmf-texlive/fonts/type1/bluesky/cm/cmsy8.p
fb></usr/share/texmf-texlive/fonts/type1/bluesky/cm/cmti10.pfb></usr/share/texm
f-texlive/fonts/type1/bluesky/cm/cmti12.pfb></usr/share/texmf-texlive/fonts/typ
e1/bluesky/cm/cmtt8.pfb></usr/share/texmf-texlive/fonts/type1/bluesky/ams/msam1
0.pfb>
Output written on 36670-t.pdf (269 pages, 788109 bytes).
PDF statistics:
 2045 PDF objects out of 2073 (max. 8388607)
 681 named destinations out of 1000 (max. 131072)
 224 words of extra memory for PDF output out of 10000 (max. 10000000)

